% Shared preamble and content for Sudhansh Peddabomma's résumés.
% Each profile (SudhanshPeddabomma-AI.tex, SudhanshPeddabomma-Vision.tex) inputs this file and picks
% which blocks to show and in what order, so every bullet is written once.
% Layout rule: every bullet fills one line, or runs 1.5 to 2 lines. Check with: python3 check_lines.py <pdf>
\documentclass[10pt,letterpaper]{article}

\usepackage[margin=0.55in,top=0.5in,bottom=0.45in]{geometry}
\usepackage[T1]{fontenc}
% Fonts: \resumefont loads the body font; \headfont styles the name and section titles.
% Override either before % Shared preamble and content for Sudhansh Peddabomma's résumés.
% Each profile (SudhanshPeddabomma-AI.tex, SudhanshPeddabomma-Vision.tex) inputs this file and picks
% which blocks to show and in what order, so every bullet is written once.
% Layout rule: every bullet fills one line, or runs 1.5 to 2 lines. Check with: python3 check_lines.py <pdf>
\documentclass[10pt,letterpaper]{article}

\usepackage[margin=0.55in,top=0.5in,bottom=0.45in]{geometry}
\usepackage[T1]{fontenc}
% Fonts: \resumefont loads the body font; \headfont styles the name and section titles.
% Override either before % Shared preamble and content for Sudhansh Peddabomma's résumés.
% Each profile (SudhanshPeddabomma-AI.tex, SudhanshPeddabomma-Vision.tex) inputs this file and picks
% which blocks to show and in what order, so every bullet is written once.
% Layout rule: every bullet fills one line, or runs 1.5 to 2 lines. Check with: python3 check_lines.py <pdf>
\documentclass[10pt,letterpaper]{article}

\usepackage[margin=0.55in,top=0.5in,bottom=0.45in]{geometry}
\usepackage[T1]{fontenc}
% Fonts: \resumefont loads the body font; \headfont styles the name and section titles.
% Override either before % Shared preamble and content for Sudhansh Peddabomma's résumés.
% Each profile (SudhanshPeddabomma-AI.tex, SudhanshPeddabomma-Vision.tex) inputs this file and picks
% which blocks to show and in what order, so every bullet is written once.
% Layout rule: every bullet fills one line, or runs 1.5 to 2 lines. Check with: python3 check_lines.py <pdf>
\documentclass[10pt,letterpaper]{article}

\usepackage[margin=0.55in,top=0.5in,bottom=0.45in]{geometry}
\usepackage[T1]{fontenc}
% Fonts: \resumefont loads the body font; \headfont styles the name and section titles.
% Override either before \input{resume-common} to try another pairing.
% Default: Charter body with Source Sans headings. All-sans alternative:
%   \def\resumefont{\usepackage[default,semibold]{sourcesanspro}}\def\headfont{}
\providecommand{\resumefont}{\usepackage[scaled=0.96]{XCharter}\usepackage[semibold,scale=0.96]{sourcesanspro}}
\providecommand{\headfont}{\sffamily}
\providecommand{\namefont}{\rmfamily\scshape}
\resumefont
\usepackage[final]{microtype}
\usepackage[dvipsnames]{xcolor}
\usepackage{enumitem}
\usepackage{titlesec}
\usepackage{tabularx}
\usepackage{fontawesome5}
\usepackage{etoolbox}
\usepackage[hidelinks]{hyperref}

% The site's green, darkened for print.
\definecolor{accent}{HTML}{3A7A45}
\definecolor{muted}{HTML}{5C675D}

\pagestyle{empty}
\setlength{\parindent}{0pt}
\setlength{\tabcolsep}{0pt}
\providecommand{\leading}{1.1}
\linespread{\leading}

\hypersetup{pdfauthor={Sudhansh Peddabomma}}

% Type scale: name 26pt small caps, section titles 12.5pt caps, entry titles 10.8pt, body 10pt,
% side notes 9pt.
\newcommand{\side}[1]{{\fontsize{9}{11}\selectfont\color{muted}#1}}
\newcommand{\sectitle}[1]{\textls[60]{\MakeUppercase{#1}}}
\titleformat{\section}{\headfont\color{accent}\bfseries\fontsize{12.5}{14}\selectfont}{}{0pt}{\sectitle}[\vspace{1.5pt}{\color{accent}\titlerule[0.8pt]}\global\firstentrytrue]
% The first entry after a section title skips \entrygap, so only \rulegap sits under the rule.
\newif\iffirstentry

% Vertical rhythm, the same idea as the earlier hand-tuned résumés: bullets inside an entry sit
% tight together, entries are separated by a clear gap, and sections by a larger one. A profile
% may override these before \input{resume-common} to fit one page.
\providecommand{\secgap}{12pt}   % above a section title
\providecommand{\rulegap}{5.5pt}   % section rule to first entry
\providecommand{\entrygap}{7pt}  % between entries
\providecommand{\headgap}{1.5pt} % entry heading to its bullets
\providecommand{\itemgap}{0.6pt} % between bullets of one entry
\titlespacing*{\section}{0pt}{\secgap}{\rulegap}

\newcommand{\squarebullet}{\raisebox{0.32ex}{\color{accent}\rule{0.5ex}{0.5ex}}}
\setlist[itemize]{leftmargin=12pt, labelsep=6pt, topsep=\headgap, itemsep=\itemgap, parsep=0pt, partopsep=0pt,
  label=\squarebullet}

% \entry{Organisation or title}{Role or detail (may be empty)}{Dates}; \sub{Role or detail}{Right-hand note}
\newcommand{\entry}[3]{%
  \par\iffirstentry\global\firstentryfalse\else\addvspace{\entrygap}\fi\noindent{\headfont\fontsize{10.8}{13}\selectfont\textbf{#1}}\ifblank{#2}{}{\ {\color{muted}|}\ #2}\hfill\hspace{1em}\side{#3}\par}
\newcommand{\sub}[2]{%
  \noindent\textit{#1}\hfill{\color{muted}\textit{#2}}\par}
\newcommand{\report}[1]{\href{#1}{\color{accent}Report}}
\newcommand{\me}{\textbf{S. Peddabomma}}

% ---------------------------------------------------------------------------------------------------
% Header
\newcommand{\Header}{%
  \begin{center}
    {\namefont\fontsize{26}{28}\selectfont Sudhansh Peddabomma}\\[6pt]
    {\headfont\fontsize{9.5}{11}\selectfont\color{muted}
      \href{mailto:sudhansh6@gmail.com}{\faEnvelope[regular]\ sudhansh6@gmail.com} \quad
      \href{https://sudhansh6.github.io}{\faGlobe\ sudhansh6.github.io} \quad
      \href{https://www.linkedin.com/in/sudhansh-peddabomma-73a208192/}{\faLinkedin\ sudhansh-peddabomma} \quad
      \href{https://github.com/Sudhansh6}{\faGithub\ Sudhansh6}}
  \end{center}\vspace{-4pt}}

% ---------------------------------------------------------------------------------------------------
% Glean bullets, shared by both profiles
\newcommand{\GleanSlides}{\item Led \textbf{AI slide generation} in Glean Assistant, now creating \textbf{180K decks a week} for \textbf{50K weekly active users} at enterprises}
\newcommand{\GleanSlidesPerf}{\item Improved the slides harness, raising reliability to \textbf{97\%}, cutting median latency \textbf{26\%} and cost per deck by \textbf{half}
  in four months}
\newcommand{\GleanSlidesOrch}{\item Architected a \textbf{slide orchestration system} that reuses relevant slides from a company's existing decks and generates
  new ones in its brand and template, so teams can refine decks with each other and Glean in full context, without leaving the product}
\newcommand{\GleanAgents}{\item Moved Glean's \textbf{workflow agents} from one-shot planning to an \textbf{agentic tool-use loop}, greatly improving agentic
  capabilities and instruction following, with \textbf{18\%} higher answer completeness and \textbf{17\%} fewer tokens}
\newcommand{\GleanTeamChat}{\item Engineered the core of \textbf{Team chat}, letting teammates work with Glean in one shared conversation and co-edit its artifacts}
\newcommand{\GleanVideo}{\item Shipped \textbf{video artifacts} end to end, generating narrated videos in chat with streamed playback and admin rollout controls}
\newcommand{\GleanEvals}{\item Created the \textbf{evaluations and metrics} that guide model choices and rollouts across agents, slides and interactive artifacts}

% ---------------------------------------------------------------------------------------------------
% Experience
% \groupbreak separates related groups of bullets inside one long entry (slides, then platform work).
\providecommand{\groupgap}{3.5pt}
\newcommand{\groupbreak}{\vspace{\groupgap}}
\newcommand{\GleanAI}{%
  \entry{Glean}{Software Engineer}{Apr 2025 -- Present}
  \begin{itemize}
    \GleanSlides
    \GleanSlidesOrch
    \GleanSlidesPerf \groupbreak
    \GleanTeamChat
    \GleanVideo
    \GleanAgents
    \GleanEvals
  \end{itemize}}

\newcommand{\GleanVision}{%
  \entry{Glean}{Software Engineer}{Apr 2025 -- Present}
  \begin{itemize}
    \GleanSlides
    \GleanSlidesOrch
    \GleanSlidesPerf \groupbreak
    \GleanVideo
    \GleanEvals
  \end{itemize}}

\newcommand{\Duality}{%
  \entry{Duality AI}{Computer Vision Intern}{Jun 2024 -- Sep 2024}
  \begin{itemize}
    \item Designed automated \textbf{3D reconstruction} techniques for featureless objects, reducing digital-twin generation time
          by \textbf{40\%}
    \item Collaborated with \textbf{Autodesk} to validate Unreal Engine simulations for robotics tasks; leveraging structured
          domain randomization to reduce Sim2Real gap and increase mAP-50 by \textbf{15\%} for object detection and segmentation
  \end{itemize}}

\newcommand{\Microsoft}{%
  \entry{Microsoft}{Data and Applied Scientist Intern}{May 2022 -- Jul 2022}
  \begin{itemize}
    \item Developed a \textbf{decision-tree ranker} to recommend emails without user queries, improving \textbf{Outlook} search capabilities
    \item Proposed hierarchical feature-sets for the ranker, reducing latency for recommendations and improving recall
  \end{itemize}}

% ---------------------------------------------------------------------------------------------------
% Research
\newcommand{\HomeRobots}{%
  \entry{Real-time 3D Perception for Home Robots}{UC San Diego}{Sep 2023 -- Sep 2024}
  \begin{itemize}
    \item Investigated real-time dense \textbf{visual SLAM} methods using NeRFs and Gaussian Splatting for indoor robot navigation
    \item Integrated object segmentation, grasp-pose estimation, and 3D mapping on the Fetch robot via ROS, demonstrating
          a novel \textbf{tabletop object rearrangement} algorithm that reduced cost by \textbf{20\%} compared to the state-of-the-art
  \end{itemize}}

\newcommand{\Tomography}{%
  \entry{3D Tomography with Primal-Dual Neural Networks}{University College London}{May 2021 -- Jul 2023}
  \begin{itemize}
    \item Developed a stochastic neural-network architecture of primal-dual algorithm for online reconstruction of 3D
          volumes from tomographic projections, with \textbf{99.6\,\% structural similarity} in low-dosage conditions
    \item Built a Python library with custom gradient operators for reconstructing volumes in a single pass, reducing
          compute requirements by up to \textbf{5x} over SOTA learning-based approaches for cone vector tomography
  \end{itemize}}

% ---------------------------------------------------------------------------------------------------
% Projects
\newcommand{\OOXML}{%
  \entry{ooxml}{Open-source contributor \textbar\ \href{https://github.com/ChristopherVR/ooxml}{\color{accent}GitHub}}{Sep 2026 -- Present}
  \begin{itemize}
    \item Contributed to an open-source \textbf{TypeScript} engine that renders and edits PowerPoint and Word files in the browser,
          improving editing fidelity, \textbf{real-time collaboration} and load performance across React, Vue and Svelte viewers
  \end{itemize}}

\newcommand{\MirrorAI}{%
  \entry{Mirror AI: Deployable Personas}{Honorable mention in Supabase YC Hackathon}{Oct 2024 -- Dec 2024}
  \begin{itemize}
    \item Designed an \textbf{agentic LLM architecture} with LangGraph to mirror user personalities, creating interactive digital
          personas, and deployed it as a full-stack platform on Supabase and Vercel with one-click publishing
  \end{itemize}}

\newcommand{\LLMReasoning}{%
  \entry{Improving LLM Reasoning for Numerical Problems}{\report{https://sudhansh6.github.io/assets/pdf/MathPrompter++.pdf}}{Sep 2024 -- Dec 2024}
  \begin{itemize}
    \item Enhanced MathPrompter (ACL 2023) with CoT, achieving \textbf{10\% higher accuracy} on Llama 3.1 1B where prior methods failed,
          and reduced hallucination rates significantly by integrating multi-step validation
  \end{itemize}}

\newcommand{\InverseRendering}{%
  \entry{Inverse Rendering with 2D Gaussian Splatting}{\report{https://sudhansh6.github.io/assets/pdf/CSE252D_Report.pdf}}{Mar 2024 -- May 2024}
  \begin{itemize}
    \item Developed a novel \textbf{inverse rendering} framework in CUDA to recover PBR properties of a scene using 2D Gaussian
          Splatting, improving normal map MAE by \textbf{15\%} over the current SOTA methods
  \end{itemize}}

% ---------------------------------------------------------------------------------------------------
% Publications: title on the left, venue on the right, one line each.
\newcommand{\publine}[3]{\par\noindent\href{#1}{#2}\hfill\hspace{1em}\side{#3}\par\addvspace{1.5pt}}
\newcommand{\PubSaturation}{\publine{https://doi.org/10.1016/j.sigpro.2023.109349}{A likelihood based method for compressive signal recovery under Gaussian and saturation noise}{Signal Processing, 2024}}
\newcommand{\PubPermutation}{\publine{https://ieeexplore.ieee.org/document/10888147}{Identification and correction of permutation errors in compressed sensing-based group testing}{ICASSP 2025}}
\newcommand{\PubRearrangement}{\publine{https://arxiv.org/abs/2411.10899}{Planning for tabletop object rearrangement}{ICRA 2025}}

% ---------------------------------------------------------------------------------------------------
% Education, skills, service
% Each profile sets \UCSDCourses, \IITBCourses and \SkillRows before \input{resume-common}.
\providecommand{\UCSDCourses}{}
\providecommand{\IITBCourses}{}
\providecommand{\SkillRows}{}
\newcommand{\coursework}[1]{\ifblank{#1}{}{\par\noindent{\fontsize{9.5}{11.5}\selectfont{\color{muted}Coursework:} #1}\par}}
\newcommand{\Education}{%
  \section{Education}
  \entry{University of California San Diego}{M.S. in Computer Science and Engineering (AI)}{GPA 4.00/4.00 \quad Sep 2023 -- Mar 2025}
  \coursework{\UCSDCourses}
  \entry{Indian Institute of Technology Bombay}{B.Tech.\ with Honors in Computer Science}{CPI 9.66/10 \quad Jul 2019 -- Jul 2023}
  \coursework{\IITBCourses}}

% \skill{Label}{Items} is one row of the skills table.
\newcommand{\skill}[2]{{\headfont\bfseries #1} & #2 \\[1.5pt]}
\newcommand{\cpp}{C\kern-0.04em\raisebox{0.06ex}{\small++}}
\newcommand{\Skills}{%
  \section{Skills}
  \begin{tabularx}{\linewidth}{@{}l@{\hspace{12pt}}>{\raggedright\arraybackslash}X@{}}
    \SkillRows
  \end{tabularx}}
 to try another pairing.
% Default: Charter body with Source Sans headings. All-sans alternative:
%   \def\resumefont{\usepackage[default,semibold]{sourcesanspro}}\def\headfont{}
\providecommand{\resumefont}{\usepackage[scaled=0.96]{XCharter}\usepackage[semibold,scale=0.96]{sourcesanspro}}
\providecommand{\headfont}{\sffamily}
\providecommand{\namefont}{\rmfamily\scshape}
\resumefont
\usepackage[final]{microtype}
\usepackage[dvipsnames]{xcolor}
\usepackage{enumitem}
\usepackage{titlesec}
\usepackage{tabularx}
\usepackage{fontawesome5}
\usepackage{etoolbox}
\usepackage[hidelinks]{hyperref}

% The site's green, darkened for print.
\definecolor{accent}{HTML}{3A7A45}
\definecolor{muted}{HTML}{5C675D}

\pagestyle{empty}
\setlength{\parindent}{0pt}
\setlength{\tabcolsep}{0pt}
\providecommand{\leading}{1.1}
\linespread{\leading}

\hypersetup{pdfauthor={Sudhansh Peddabomma}}

% Type scale: name 26pt small caps, section titles 12.5pt caps, entry titles 10.8pt, body 10pt,
% side notes 9pt.
\newcommand{\side}[1]{{\fontsize{9}{11}\selectfont\color{muted}#1}}
\newcommand{\sectitle}[1]{\textls[60]{\MakeUppercase{#1}}}
\titleformat{\section}{\headfont\color{accent}\bfseries\fontsize{12.5}{14}\selectfont}{}{0pt}{\sectitle}[\vspace{1.5pt}{\color{accent}\titlerule[0.8pt]}\global\firstentrytrue]
% The first entry after a section title skips \entrygap, so only \rulegap sits under the rule.
\newif\iffirstentry

% Vertical rhythm, the same idea as the earlier hand-tuned résumés: bullets inside an entry sit
% tight together, entries are separated by a clear gap, and sections by a larger one. A profile
% may override these before % Shared preamble and content for Sudhansh Peddabomma's résumés.
% Each profile (SudhanshPeddabomma-AI.tex, SudhanshPeddabomma-Vision.tex) inputs this file and picks
% which blocks to show and in what order, so every bullet is written once.
% Layout rule: every bullet fills one line, or runs 1.5 to 2 lines. Check with: python3 check_lines.py <pdf>
\documentclass[10pt,letterpaper]{article}

\usepackage[margin=0.55in,top=0.5in,bottom=0.45in]{geometry}
\usepackage[T1]{fontenc}
% Fonts: \resumefont loads the body font; \headfont styles the name and section titles.
% Override either before \input{resume-common} to try another pairing.
% Default: Charter body with Source Sans headings. All-sans alternative:
%   \def\resumefont{\usepackage[default,semibold]{sourcesanspro}}\def\headfont{}
\providecommand{\resumefont}{\usepackage[scaled=0.96]{XCharter}\usepackage[semibold,scale=0.96]{sourcesanspro}}
\providecommand{\headfont}{\sffamily}
\providecommand{\namefont}{\rmfamily\scshape}
\resumefont
\usepackage[final]{microtype}
\usepackage[dvipsnames]{xcolor}
\usepackage{enumitem}
\usepackage{titlesec}
\usepackage{tabularx}
\usepackage{fontawesome5}
\usepackage{etoolbox}
\usepackage[hidelinks]{hyperref}

% The site's green, darkened for print.
\definecolor{accent}{HTML}{3A7A45}
\definecolor{muted}{HTML}{5C675D}

\pagestyle{empty}
\setlength{\parindent}{0pt}
\setlength{\tabcolsep}{0pt}
\providecommand{\leading}{1.1}
\linespread{\leading}

\hypersetup{pdfauthor={Sudhansh Peddabomma}}

% Type scale: name 26pt small caps, section titles 12.5pt caps, entry titles 10.8pt, body 10pt,
% side notes 9pt.
\newcommand{\side}[1]{{\fontsize{9}{11}\selectfont\color{muted}#1}}
\newcommand{\sectitle}[1]{\textls[60]{\MakeUppercase{#1}}}
\titleformat{\section}{\headfont\color{accent}\bfseries\fontsize{12.5}{14}\selectfont}{}{0pt}{\sectitle}[\vspace{1.5pt}{\color{accent}\titlerule[0.8pt]}\global\firstentrytrue]
% The first entry after a section title skips \entrygap, so only \rulegap sits under the rule.
\newif\iffirstentry

% Vertical rhythm, the same idea as the earlier hand-tuned résumés: bullets inside an entry sit
% tight together, entries are separated by a clear gap, and sections by a larger one. A profile
% may override these before \input{resume-common} to fit one page.
\providecommand{\secgap}{12pt}   % above a section title
\providecommand{\rulegap}{5.5pt}   % section rule to first entry
\providecommand{\entrygap}{7pt}  % between entries
\providecommand{\headgap}{1.5pt} % entry heading to its bullets
\providecommand{\itemgap}{0.6pt} % between bullets of one entry
\titlespacing*{\section}{0pt}{\secgap}{\rulegap}

\newcommand{\squarebullet}{\raisebox{0.32ex}{\color{accent}\rule{0.5ex}{0.5ex}}}
\setlist[itemize]{leftmargin=12pt, labelsep=6pt, topsep=\headgap, itemsep=\itemgap, parsep=0pt, partopsep=0pt,
  label=\squarebullet}

% \entry{Organisation or title}{Role or detail (may be empty)}{Dates}; \sub{Role or detail}{Right-hand note}
\newcommand{\entry}[3]{%
  \par\iffirstentry\global\firstentryfalse\else\addvspace{\entrygap}\fi\noindent{\headfont\fontsize{10.8}{13}\selectfont\textbf{#1}}\ifblank{#2}{}{\ {\color{muted}|}\ #2}\hfill\hspace{1em}\side{#3}\par}
\newcommand{\sub}[2]{%
  \noindent\textit{#1}\hfill{\color{muted}\textit{#2}}\par}
\newcommand{\report}[1]{\href{#1}{\color{accent}Report}}
\newcommand{\me}{\textbf{S. Peddabomma}}

% ---------------------------------------------------------------------------------------------------
% Header
\newcommand{\Header}{%
  \begin{center}
    {\namefont\fontsize{26}{28}\selectfont Sudhansh Peddabomma}\\[6pt]
    {\headfont\fontsize{9.5}{11}\selectfont\color{muted}
      \href{mailto:sudhansh6@gmail.com}{\faEnvelope[regular]\ sudhansh6@gmail.com} \quad
      \href{https://sudhansh6.github.io}{\faGlobe\ sudhansh6.github.io} \quad
      \href{https://www.linkedin.com/in/sudhansh-peddabomma-73a208192/}{\faLinkedin\ sudhansh-peddabomma} \quad
      \href{https://github.com/Sudhansh6}{\faGithub\ Sudhansh6}}
  \end{center}\vspace{-4pt}}

% ---------------------------------------------------------------------------------------------------
% Glean bullets, shared by both profiles
\newcommand{\GleanSlides}{\item Led \textbf{AI slide generation} in Glean Assistant, now creating \textbf{180K decks a week} for \textbf{50K weekly active users} at enterprises}
\newcommand{\GleanSlidesPerf}{\item Improved the slides harness, raising reliability to \textbf{97\%}, cutting median latency \textbf{26\%} and cost per deck by \textbf{half}
  in four months}
\newcommand{\GleanSlidesOrch}{\item Architected a \textbf{slide orchestration system} that reuses relevant slides from a company's existing decks and generates
  new ones in its brand and template, so teams can refine decks with each other and Glean in full context, without leaving the product}
\newcommand{\GleanAgents}{\item Moved Glean's \textbf{workflow agents} from one-shot planning to an \textbf{agentic tool-use loop}, greatly improving agentic
  capabilities and instruction following, with \textbf{18\%} higher answer completeness and \textbf{17\%} fewer tokens}
\newcommand{\GleanTeamChat}{\item Engineered the core of \textbf{Team chat}, letting teammates work with Glean in one shared conversation and co-edit its artifacts}
\newcommand{\GleanVideo}{\item Shipped \textbf{video artifacts} end to end, generating narrated videos in chat with streamed playback and admin rollout controls}
\newcommand{\GleanEvals}{\item Created the \textbf{evaluations and metrics} that guide model choices and rollouts across agents, slides and interactive artifacts}

% ---------------------------------------------------------------------------------------------------
% Experience
% \groupbreak separates related groups of bullets inside one long entry (slides, then platform work).
\providecommand{\groupgap}{3.5pt}
\newcommand{\groupbreak}{\vspace{\groupgap}}
\newcommand{\GleanAI}{%
  \entry{Glean}{Software Engineer}{Apr 2025 -- Present}
  \begin{itemize}
    \GleanSlides
    \GleanSlidesOrch
    \GleanSlidesPerf \groupbreak
    \GleanTeamChat
    \GleanVideo
    \GleanAgents
    \GleanEvals
  \end{itemize}}

\newcommand{\GleanVision}{%
  \entry{Glean}{Software Engineer}{Apr 2025 -- Present}
  \begin{itemize}
    \GleanSlides
    \GleanSlidesOrch
    \GleanSlidesPerf \groupbreak
    \GleanVideo
    \GleanEvals
  \end{itemize}}

\newcommand{\Duality}{%
  \entry{Duality AI}{Computer Vision Intern}{Jun 2024 -- Sep 2024}
  \begin{itemize}
    \item Designed automated \textbf{3D reconstruction} techniques for featureless objects, reducing digital-twin generation time
          by \textbf{40\%}
    \item Collaborated with \textbf{Autodesk} to validate Unreal Engine simulations for robotics tasks; leveraging structured
          domain randomization to reduce Sim2Real gap and increase mAP-50 by \textbf{15\%} for object detection and segmentation
  \end{itemize}}

\newcommand{\Microsoft}{%
  \entry{Microsoft}{Data and Applied Scientist Intern}{May 2022 -- Jul 2022}
  \begin{itemize}
    \item Developed a \textbf{decision-tree ranker} to recommend emails without user queries, improving \textbf{Outlook} search capabilities
    \item Proposed hierarchical feature-sets for the ranker, reducing latency for recommendations and improving recall
  \end{itemize}}

% ---------------------------------------------------------------------------------------------------
% Research
\newcommand{\HomeRobots}{%
  \entry{Real-time 3D Perception for Home Robots}{UC San Diego}{Sep 2023 -- Sep 2024}
  \begin{itemize}
    \item Investigated real-time dense \textbf{visual SLAM} methods using NeRFs and Gaussian Splatting for indoor robot navigation
    \item Integrated object segmentation, grasp-pose estimation, and 3D mapping on the Fetch robot via ROS, demonstrating
          a novel \textbf{tabletop object rearrangement} algorithm that reduced cost by \textbf{20\%} compared to the state-of-the-art
  \end{itemize}}

\newcommand{\Tomography}{%
  \entry{3D Tomography with Primal-Dual Neural Networks}{University College London}{May 2021 -- Jul 2023}
  \begin{itemize}
    \item Developed a stochastic neural-network architecture of primal-dual algorithm for online reconstruction of 3D
          volumes from tomographic projections, with \textbf{99.6\,\% structural similarity} in low-dosage conditions
    \item Built a Python library with custom gradient operators for reconstructing volumes in a single pass, reducing
          compute requirements by up to \textbf{5x} over SOTA learning-based approaches for cone vector tomography
  \end{itemize}}

% ---------------------------------------------------------------------------------------------------
% Projects
\newcommand{\OOXML}{%
  \entry{ooxml}{Open-source contributor \textbar\ \href{https://github.com/ChristopherVR/ooxml}{\color{accent}GitHub}}{Sep 2026 -- Present}
  \begin{itemize}
    \item Contributed to an open-source \textbf{TypeScript} engine that renders and edits PowerPoint and Word files in the browser,
          improving editing fidelity, \textbf{real-time collaboration} and load performance across React, Vue and Svelte viewers
  \end{itemize}}

\newcommand{\MirrorAI}{%
  \entry{Mirror AI: Deployable Personas}{Honorable mention in Supabase YC Hackathon}{Oct 2024 -- Dec 2024}
  \begin{itemize}
    \item Designed an \textbf{agentic LLM architecture} with LangGraph to mirror user personalities, creating interactive digital
          personas, and deployed it as a full-stack platform on Supabase and Vercel with one-click publishing
  \end{itemize}}

\newcommand{\LLMReasoning}{%
  \entry{Improving LLM Reasoning for Numerical Problems}{\report{https://sudhansh6.github.io/assets/pdf/MathPrompter++.pdf}}{Sep 2024 -- Dec 2024}
  \begin{itemize}
    \item Enhanced MathPrompter (ACL 2023) with CoT, achieving \textbf{10\% higher accuracy} on Llama 3.1 1B where prior methods failed,
          and reduced hallucination rates significantly by integrating multi-step validation
  \end{itemize}}

\newcommand{\InverseRendering}{%
  \entry{Inverse Rendering with 2D Gaussian Splatting}{\report{https://sudhansh6.github.io/assets/pdf/CSE252D_Report.pdf}}{Mar 2024 -- May 2024}
  \begin{itemize}
    \item Developed a novel \textbf{inverse rendering} framework in CUDA to recover PBR properties of a scene using 2D Gaussian
          Splatting, improving normal map MAE by \textbf{15\%} over the current SOTA methods
  \end{itemize}}

% ---------------------------------------------------------------------------------------------------
% Publications: title on the left, venue on the right, one line each.
\newcommand{\publine}[3]{\par\noindent\href{#1}{#2}\hfill\hspace{1em}\side{#3}\par\addvspace{1.5pt}}
\newcommand{\PubSaturation}{\publine{https://doi.org/10.1016/j.sigpro.2023.109349}{A likelihood based method for compressive signal recovery under Gaussian and saturation noise}{Signal Processing, 2024}}
\newcommand{\PubPermutation}{\publine{https://ieeexplore.ieee.org/document/10888147}{Identification and correction of permutation errors in compressed sensing-based group testing}{ICASSP 2025}}
\newcommand{\PubRearrangement}{\publine{https://arxiv.org/abs/2411.10899}{Planning for tabletop object rearrangement}{ICRA 2025}}

% ---------------------------------------------------------------------------------------------------
% Education, skills, service
% Each profile sets \UCSDCourses, \IITBCourses and \SkillRows before \input{resume-common}.
\providecommand{\UCSDCourses}{}
\providecommand{\IITBCourses}{}
\providecommand{\SkillRows}{}
\newcommand{\coursework}[1]{\ifblank{#1}{}{\par\noindent{\fontsize{9.5}{11.5}\selectfont{\color{muted}Coursework:} #1}\par}}
\newcommand{\Education}{%
  \section{Education}
  \entry{University of California San Diego}{M.S. in Computer Science and Engineering (AI)}{GPA 4.00/4.00 \quad Sep 2023 -- Mar 2025}
  \coursework{\UCSDCourses}
  \entry{Indian Institute of Technology Bombay}{B.Tech.\ with Honors in Computer Science}{CPI 9.66/10 \quad Jul 2019 -- Jul 2023}
  \coursework{\IITBCourses}}

% \skill{Label}{Items} is one row of the skills table.
\newcommand{\skill}[2]{{\headfont\bfseries #1} & #2 \\[1.5pt]}
\newcommand{\cpp}{C\kern-0.04em\raisebox{0.06ex}{\small++}}
\newcommand{\Skills}{%
  \section{Skills}
  \begin{tabularx}{\linewidth}{@{}l@{\hspace{12pt}}>{\raggedright\arraybackslash}X@{}}
    \SkillRows
  \end{tabularx}}
 to fit one page.
\providecommand{\secgap}{12pt}   % above a section title
\providecommand{\rulegap}{5.5pt}   % section rule to first entry
\providecommand{\entrygap}{7pt}  % between entries
\providecommand{\headgap}{1.5pt} % entry heading to its bullets
\providecommand{\itemgap}{0.6pt} % between bullets of one entry
\titlespacing*{\section}{0pt}{\secgap}{\rulegap}

\newcommand{\squarebullet}{\raisebox{0.32ex}{\color{accent}\rule{0.5ex}{0.5ex}}}
\setlist[itemize]{leftmargin=12pt, labelsep=6pt, topsep=\headgap, itemsep=\itemgap, parsep=0pt, partopsep=0pt,
  label=\squarebullet}

% \entry{Organisation or title}{Role or detail (may be empty)}{Dates}; \sub{Role or detail}{Right-hand note}
\newcommand{\entry}[3]{%
  \par\iffirstentry\global\firstentryfalse\else\addvspace{\entrygap}\fi\noindent{\headfont\fontsize{10.8}{13}\selectfont\textbf{#1}}\ifblank{#2}{}{\ {\color{muted}|}\ #2}\hfill\hspace{1em}\side{#3}\par}
\newcommand{\sub}[2]{%
  \noindent\textit{#1}\hfill{\color{muted}\textit{#2}}\par}
\newcommand{\report}[1]{\href{#1}{\color{accent}Report}}
\newcommand{\me}{\textbf{S. Peddabomma}}

% ---------------------------------------------------------------------------------------------------
% Header
\newcommand{\Header}{%
  \begin{center}
    {\namefont\fontsize{26}{28}\selectfont Sudhansh Peddabomma}\\[6pt]
    {\headfont\fontsize{9.5}{11}\selectfont\color{muted}
      \href{mailto:sudhansh6@gmail.com}{\faEnvelope[regular]\ sudhansh6@gmail.com} \quad
      \href{https://sudhansh6.github.io}{\faGlobe\ sudhansh6.github.io} \quad
      \href{https://www.linkedin.com/in/sudhansh-peddabomma-73a208192/}{\faLinkedin\ sudhansh-peddabomma} \quad
      \href{https://github.com/Sudhansh6}{\faGithub\ Sudhansh6}}
  \end{center}\vspace{-4pt}}

% ---------------------------------------------------------------------------------------------------
% Glean bullets, shared by both profiles
\newcommand{\GleanSlides}{\item Led \textbf{AI slide generation} in Glean Assistant, now creating \textbf{180K decks a week} for \textbf{50K weekly active users} at enterprises}
\newcommand{\GleanSlidesPerf}{\item Improved the slides harness, raising reliability to \textbf{97\%}, cutting median latency \textbf{26\%} and cost per deck by \textbf{half}
  in four months}
\newcommand{\GleanSlidesOrch}{\item Architected a \textbf{slide orchestration system} that reuses relevant slides from a company's existing decks and generates
  new ones in its brand and template, so teams can refine decks with each other and Glean in full context, without leaving the product}
\newcommand{\GleanAgents}{\item Moved Glean's \textbf{workflow agents} from one-shot planning to an \textbf{agentic tool-use loop}, greatly improving agentic
  capabilities and instruction following, with \textbf{18\%} higher answer completeness and \textbf{17\%} fewer tokens}
\newcommand{\GleanTeamChat}{\item Engineered the core of \textbf{Team chat}, letting teammates work with Glean in one shared conversation and co-edit its artifacts}
\newcommand{\GleanVideo}{\item Shipped \textbf{video artifacts} end to end, generating narrated videos in chat with streamed playback and admin rollout controls}
\newcommand{\GleanEvals}{\item Created the \textbf{evaluations and metrics} that guide model choices and rollouts across agents, slides and interactive artifacts}

% ---------------------------------------------------------------------------------------------------
% Experience
% \groupbreak separates related groups of bullets inside one long entry (slides, then platform work).
\providecommand{\groupgap}{3.5pt}
\newcommand{\groupbreak}{\vspace{\groupgap}}
\newcommand{\GleanAI}{%
  \entry{Glean}{Software Engineer}{Apr 2025 -- Present}
  \begin{itemize}
    \GleanSlides
    \GleanSlidesOrch
    \GleanSlidesPerf \groupbreak
    \GleanTeamChat
    \GleanVideo
    \GleanAgents
    \GleanEvals
  \end{itemize}}

\newcommand{\GleanVision}{%
  \entry{Glean}{Software Engineer}{Apr 2025 -- Present}
  \begin{itemize}
    \GleanSlides
    \GleanSlidesOrch
    \GleanSlidesPerf \groupbreak
    \GleanVideo
    \GleanEvals
  \end{itemize}}

\newcommand{\Duality}{%
  \entry{Duality AI}{Computer Vision Intern}{Jun 2024 -- Sep 2024}
  \begin{itemize}
    \item Designed automated \textbf{3D reconstruction} techniques for featureless objects, reducing digital-twin generation time
          by \textbf{40\%}
    \item Collaborated with \textbf{Autodesk} to validate Unreal Engine simulations for robotics tasks; leveraging structured
          domain randomization to reduce Sim2Real gap and increase mAP-50 by \textbf{15\%} for object detection and segmentation
  \end{itemize}}

\newcommand{\Microsoft}{%
  \entry{Microsoft}{Data and Applied Scientist Intern}{May 2022 -- Jul 2022}
  \begin{itemize}
    \item Developed a \textbf{decision-tree ranker} to recommend emails without user queries, improving \textbf{Outlook} search capabilities
    \item Proposed hierarchical feature-sets for the ranker, reducing latency for recommendations and improving recall
  \end{itemize}}

% ---------------------------------------------------------------------------------------------------
% Research
\newcommand{\HomeRobots}{%
  \entry{Real-time 3D Perception for Home Robots}{UC San Diego}{Sep 2023 -- Sep 2024}
  \begin{itemize}
    \item Investigated real-time dense \textbf{visual SLAM} methods using NeRFs and Gaussian Splatting for indoor robot navigation
    \item Integrated object segmentation, grasp-pose estimation, and 3D mapping on the Fetch robot via ROS, demonstrating
          a novel \textbf{tabletop object rearrangement} algorithm that reduced cost by \textbf{20\%} compared to the state-of-the-art
  \end{itemize}}

\newcommand{\Tomography}{%
  \entry{3D Tomography with Primal-Dual Neural Networks}{University College London}{May 2021 -- Jul 2023}
  \begin{itemize}
    \item Developed a stochastic neural-network architecture of primal-dual algorithm for online reconstruction of 3D
          volumes from tomographic projections, with \textbf{99.6\,\% structural similarity} in low-dosage conditions
    \item Built a Python library with custom gradient operators for reconstructing volumes in a single pass, reducing
          compute requirements by up to \textbf{5x} over SOTA learning-based approaches for cone vector tomography
  \end{itemize}}

% ---------------------------------------------------------------------------------------------------
% Projects
\newcommand{\OOXML}{%
  \entry{ooxml}{Open-source contributor \textbar\ \href{https://github.com/ChristopherVR/ooxml}{\color{accent}GitHub}}{Sep 2026 -- Present}
  \begin{itemize}
    \item Contributed to an open-source \textbf{TypeScript} engine that renders and edits PowerPoint and Word files in the browser,
          improving editing fidelity, \textbf{real-time collaboration} and load performance across React, Vue and Svelte viewers
  \end{itemize}}

\newcommand{\MirrorAI}{%
  \entry{Mirror AI: Deployable Personas}{Honorable mention in Supabase YC Hackathon}{Oct 2024 -- Dec 2024}
  \begin{itemize}
    \item Designed an \textbf{agentic LLM architecture} with LangGraph to mirror user personalities, creating interactive digital
          personas, and deployed it as a full-stack platform on Supabase and Vercel with one-click publishing
  \end{itemize}}

\newcommand{\LLMReasoning}{%
  \entry{Improving LLM Reasoning for Numerical Problems}{\report{https://sudhansh6.github.io/assets/pdf/MathPrompter++.pdf}}{Sep 2024 -- Dec 2024}
  \begin{itemize}
    \item Enhanced MathPrompter (ACL 2023) with CoT, achieving \textbf{10\% higher accuracy} on Llama 3.1 1B where prior methods failed,
          and reduced hallucination rates significantly by integrating multi-step validation
  \end{itemize}}

\newcommand{\InverseRendering}{%
  \entry{Inverse Rendering with 2D Gaussian Splatting}{\report{https://sudhansh6.github.io/assets/pdf/CSE252D_Report.pdf}}{Mar 2024 -- May 2024}
  \begin{itemize}
    \item Developed a novel \textbf{inverse rendering} framework in CUDA to recover PBR properties of a scene using 2D Gaussian
          Splatting, improving normal map MAE by \textbf{15\%} over the current SOTA methods
  \end{itemize}}

% ---------------------------------------------------------------------------------------------------
% Publications: title on the left, venue on the right, one line each.
\newcommand{\publine}[3]{\par\noindent\href{#1}{#2}\hfill\hspace{1em}\side{#3}\par\addvspace{1.5pt}}
\newcommand{\PubSaturation}{\publine{https://doi.org/10.1016/j.sigpro.2023.109349}{A likelihood based method for compressive signal recovery under Gaussian and saturation noise}{Signal Processing, 2024}}
\newcommand{\PubPermutation}{\publine{https://ieeexplore.ieee.org/document/10888147}{Identification and correction of permutation errors in compressed sensing-based group testing}{ICASSP 2025}}
\newcommand{\PubRearrangement}{\publine{https://arxiv.org/abs/2411.10899}{Planning for tabletop object rearrangement}{ICRA 2025}}

% ---------------------------------------------------------------------------------------------------
% Education, skills, service
% Each profile sets \UCSDCourses, \IITBCourses and \SkillRows before % Shared preamble and content for Sudhansh Peddabomma's résumés.
% Each profile (SudhanshPeddabomma-AI.tex, SudhanshPeddabomma-Vision.tex) inputs this file and picks
% which blocks to show and in what order, so every bullet is written once.
% Layout rule: every bullet fills one line, or runs 1.5 to 2 lines. Check with: python3 check_lines.py <pdf>
\documentclass[10pt,letterpaper]{article}

\usepackage[margin=0.55in,top=0.5in,bottom=0.45in]{geometry}
\usepackage[T1]{fontenc}
% Fonts: \resumefont loads the body font; \headfont styles the name and section titles.
% Override either before \input{resume-common} to try another pairing.
% Default: Charter body with Source Sans headings. All-sans alternative:
%   \def\resumefont{\usepackage[default,semibold]{sourcesanspro}}\def\headfont{}
\providecommand{\resumefont}{\usepackage[scaled=0.96]{XCharter}\usepackage[semibold,scale=0.96]{sourcesanspro}}
\providecommand{\headfont}{\sffamily}
\providecommand{\namefont}{\rmfamily\scshape}
\resumefont
\usepackage[final]{microtype}
\usepackage[dvipsnames]{xcolor}
\usepackage{enumitem}
\usepackage{titlesec}
\usepackage{tabularx}
\usepackage{fontawesome5}
\usepackage{etoolbox}
\usepackage[hidelinks]{hyperref}

% The site's green, darkened for print.
\definecolor{accent}{HTML}{3A7A45}
\definecolor{muted}{HTML}{5C675D}

\pagestyle{empty}
\setlength{\parindent}{0pt}
\setlength{\tabcolsep}{0pt}
\providecommand{\leading}{1.1}
\linespread{\leading}

\hypersetup{pdfauthor={Sudhansh Peddabomma}}

% Type scale: name 26pt small caps, section titles 12.5pt caps, entry titles 10.8pt, body 10pt,
% side notes 9pt.
\newcommand{\side}[1]{{\fontsize{9}{11}\selectfont\color{muted}#1}}
\newcommand{\sectitle}[1]{\textls[60]{\MakeUppercase{#1}}}
\titleformat{\section}{\headfont\color{accent}\bfseries\fontsize{12.5}{14}\selectfont}{}{0pt}{\sectitle}[\vspace{1.5pt}{\color{accent}\titlerule[0.8pt]}\global\firstentrytrue]
% The first entry after a section title skips \entrygap, so only \rulegap sits under the rule.
\newif\iffirstentry

% Vertical rhythm, the same idea as the earlier hand-tuned résumés: bullets inside an entry sit
% tight together, entries are separated by a clear gap, and sections by a larger one. A profile
% may override these before \input{resume-common} to fit one page.
\providecommand{\secgap}{12pt}   % above a section title
\providecommand{\rulegap}{5.5pt}   % section rule to first entry
\providecommand{\entrygap}{7pt}  % between entries
\providecommand{\headgap}{1.5pt} % entry heading to its bullets
\providecommand{\itemgap}{0.6pt} % between bullets of one entry
\titlespacing*{\section}{0pt}{\secgap}{\rulegap}

\newcommand{\squarebullet}{\raisebox{0.32ex}{\color{accent}\rule{0.5ex}{0.5ex}}}
\setlist[itemize]{leftmargin=12pt, labelsep=6pt, topsep=\headgap, itemsep=\itemgap, parsep=0pt, partopsep=0pt,
  label=\squarebullet}

% \entry{Organisation or title}{Role or detail (may be empty)}{Dates}; \sub{Role or detail}{Right-hand note}
\newcommand{\entry}[3]{%
  \par\iffirstentry\global\firstentryfalse\else\addvspace{\entrygap}\fi\noindent{\headfont\fontsize{10.8}{13}\selectfont\textbf{#1}}\ifblank{#2}{}{\ {\color{muted}|}\ #2}\hfill\hspace{1em}\side{#3}\par}
\newcommand{\sub}[2]{%
  \noindent\textit{#1}\hfill{\color{muted}\textit{#2}}\par}
\newcommand{\report}[1]{\href{#1}{\color{accent}Report}}
\newcommand{\me}{\textbf{S. Peddabomma}}

% ---------------------------------------------------------------------------------------------------
% Header
\newcommand{\Header}{%
  \begin{center}
    {\namefont\fontsize{26}{28}\selectfont Sudhansh Peddabomma}\\[6pt]
    {\headfont\fontsize{9.5}{11}\selectfont\color{muted}
      \href{mailto:sudhansh6@gmail.com}{\faEnvelope[regular]\ sudhansh6@gmail.com} \quad
      \href{https://sudhansh6.github.io}{\faGlobe\ sudhansh6.github.io} \quad
      \href{https://www.linkedin.com/in/sudhansh-peddabomma-73a208192/}{\faLinkedin\ sudhansh-peddabomma} \quad
      \href{https://github.com/Sudhansh6}{\faGithub\ Sudhansh6}}
  \end{center}\vspace{-4pt}}

% ---------------------------------------------------------------------------------------------------
% Glean bullets, shared by both profiles
\newcommand{\GleanSlides}{\item Led \textbf{AI slide generation} in Glean Assistant, now creating \textbf{180K decks a week} for \textbf{50K weekly active users} at enterprises}
\newcommand{\GleanSlidesPerf}{\item Improved the slides harness, raising reliability to \textbf{97\%}, cutting median latency \textbf{26\%} and cost per deck by \textbf{half}
  in four months}
\newcommand{\GleanSlidesOrch}{\item Architected a \textbf{slide orchestration system} that reuses relevant slides from a company's existing decks and generates
  new ones in its brand and template, so teams can refine decks with each other and Glean in full context, without leaving the product}
\newcommand{\GleanAgents}{\item Moved Glean's \textbf{workflow agents} from one-shot planning to an \textbf{agentic tool-use loop}, greatly improving agentic
  capabilities and instruction following, with \textbf{18\%} higher answer completeness and \textbf{17\%} fewer tokens}
\newcommand{\GleanTeamChat}{\item Engineered the core of \textbf{Team chat}, letting teammates work with Glean in one shared conversation and co-edit its artifacts}
\newcommand{\GleanVideo}{\item Shipped \textbf{video artifacts} end to end, generating narrated videos in chat with streamed playback and admin rollout controls}
\newcommand{\GleanEvals}{\item Created the \textbf{evaluations and metrics} that guide model choices and rollouts across agents, slides and interactive artifacts}

% ---------------------------------------------------------------------------------------------------
% Experience
% \groupbreak separates related groups of bullets inside one long entry (slides, then platform work).
\providecommand{\groupgap}{3.5pt}
\newcommand{\groupbreak}{\vspace{\groupgap}}
\newcommand{\GleanAI}{%
  \entry{Glean}{Software Engineer}{Apr 2025 -- Present}
  \begin{itemize}
    \GleanSlides
    \GleanSlidesOrch
    \GleanSlidesPerf \groupbreak
    \GleanTeamChat
    \GleanVideo
    \GleanAgents
    \GleanEvals
  \end{itemize}}

\newcommand{\GleanVision}{%
  \entry{Glean}{Software Engineer}{Apr 2025 -- Present}
  \begin{itemize}
    \GleanSlides
    \GleanSlidesOrch
    \GleanSlidesPerf \groupbreak
    \GleanVideo
    \GleanEvals
  \end{itemize}}

\newcommand{\Duality}{%
  \entry{Duality AI}{Computer Vision Intern}{Jun 2024 -- Sep 2024}
  \begin{itemize}
    \item Designed automated \textbf{3D reconstruction} techniques for featureless objects, reducing digital-twin generation time
          by \textbf{40\%}
    \item Collaborated with \textbf{Autodesk} to validate Unreal Engine simulations for robotics tasks; leveraging structured
          domain randomization to reduce Sim2Real gap and increase mAP-50 by \textbf{15\%} for object detection and segmentation
  \end{itemize}}

\newcommand{\Microsoft}{%
  \entry{Microsoft}{Data and Applied Scientist Intern}{May 2022 -- Jul 2022}
  \begin{itemize}
    \item Developed a \textbf{decision-tree ranker} to recommend emails without user queries, improving \textbf{Outlook} search capabilities
    \item Proposed hierarchical feature-sets for the ranker, reducing latency for recommendations and improving recall
  \end{itemize}}

% ---------------------------------------------------------------------------------------------------
% Research
\newcommand{\HomeRobots}{%
  \entry{Real-time 3D Perception for Home Robots}{UC San Diego}{Sep 2023 -- Sep 2024}
  \begin{itemize}
    \item Investigated real-time dense \textbf{visual SLAM} methods using NeRFs and Gaussian Splatting for indoor robot navigation
    \item Integrated object segmentation, grasp-pose estimation, and 3D mapping on the Fetch robot via ROS, demonstrating
          a novel \textbf{tabletop object rearrangement} algorithm that reduced cost by \textbf{20\%} compared to the state-of-the-art
  \end{itemize}}

\newcommand{\Tomography}{%
  \entry{3D Tomography with Primal-Dual Neural Networks}{University College London}{May 2021 -- Jul 2023}
  \begin{itemize}
    \item Developed a stochastic neural-network architecture of primal-dual algorithm for online reconstruction of 3D
          volumes from tomographic projections, with \textbf{99.6\,\% structural similarity} in low-dosage conditions
    \item Built a Python library with custom gradient operators for reconstructing volumes in a single pass, reducing
          compute requirements by up to \textbf{5x} over SOTA learning-based approaches for cone vector tomography
  \end{itemize}}

% ---------------------------------------------------------------------------------------------------
% Projects
\newcommand{\OOXML}{%
  \entry{ooxml}{Open-source contributor \textbar\ \href{https://github.com/ChristopherVR/ooxml}{\color{accent}GitHub}}{Sep 2026 -- Present}
  \begin{itemize}
    \item Contributed to an open-source \textbf{TypeScript} engine that renders and edits PowerPoint and Word files in the browser,
          improving editing fidelity, \textbf{real-time collaboration} and load performance across React, Vue and Svelte viewers
  \end{itemize}}

\newcommand{\MirrorAI}{%
  \entry{Mirror AI: Deployable Personas}{Honorable mention in Supabase YC Hackathon}{Oct 2024 -- Dec 2024}
  \begin{itemize}
    \item Designed an \textbf{agentic LLM architecture} with LangGraph to mirror user personalities, creating interactive digital
          personas, and deployed it as a full-stack platform on Supabase and Vercel with one-click publishing
  \end{itemize}}

\newcommand{\LLMReasoning}{%
  \entry{Improving LLM Reasoning for Numerical Problems}{\report{https://sudhansh6.github.io/assets/pdf/MathPrompter++.pdf}}{Sep 2024 -- Dec 2024}
  \begin{itemize}
    \item Enhanced MathPrompter (ACL 2023) with CoT, achieving \textbf{10\% higher accuracy} on Llama 3.1 1B where prior methods failed,
          and reduced hallucination rates significantly by integrating multi-step validation
  \end{itemize}}

\newcommand{\InverseRendering}{%
  \entry{Inverse Rendering with 2D Gaussian Splatting}{\report{https://sudhansh6.github.io/assets/pdf/CSE252D_Report.pdf}}{Mar 2024 -- May 2024}
  \begin{itemize}
    \item Developed a novel \textbf{inverse rendering} framework in CUDA to recover PBR properties of a scene using 2D Gaussian
          Splatting, improving normal map MAE by \textbf{15\%} over the current SOTA methods
  \end{itemize}}

% ---------------------------------------------------------------------------------------------------
% Publications: title on the left, venue on the right, one line each.
\newcommand{\publine}[3]{\par\noindent\href{#1}{#2}\hfill\hspace{1em}\side{#3}\par\addvspace{1.5pt}}
\newcommand{\PubSaturation}{\publine{https://doi.org/10.1016/j.sigpro.2023.109349}{A likelihood based method for compressive signal recovery under Gaussian and saturation noise}{Signal Processing, 2024}}
\newcommand{\PubPermutation}{\publine{https://ieeexplore.ieee.org/document/10888147}{Identification and correction of permutation errors in compressed sensing-based group testing}{ICASSP 2025}}
\newcommand{\PubRearrangement}{\publine{https://arxiv.org/abs/2411.10899}{Planning for tabletop object rearrangement}{ICRA 2025}}

% ---------------------------------------------------------------------------------------------------
% Education, skills, service
% Each profile sets \UCSDCourses, \IITBCourses and \SkillRows before \input{resume-common}.
\providecommand{\UCSDCourses}{}
\providecommand{\IITBCourses}{}
\providecommand{\SkillRows}{}
\newcommand{\coursework}[1]{\ifblank{#1}{}{\par\noindent{\fontsize{9.5}{11.5}\selectfont{\color{muted}Coursework:} #1}\par}}
\newcommand{\Education}{%
  \section{Education}
  \entry{University of California San Diego}{M.S. in Computer Science and Engineering (AI)}{GPA 4.00/4.00 \quad Sep 2023 -- Mar 2025}
  \coursework{\UCSDCourses}
  \entry{Indian Institute of Technology Bombay}{B.Tech.\ with Honors in Computer Science}{CPI 9.66/10 \quad Jul 2019 -- Jul 2023}
  \coursework{\IITBCourses}}

% \skill{Label}{Items} is one row of the skills table.
\newcommand{\skill}[2]{{\headfont\bfseries #1} & #2 \\[1.5pt]}
\newcommand{\cpp}{C\kern-0.04em\raisebox{0.06ex}{\small++}}
\newcommand{\Skills}{%
  \section{Skills}
  \begin{tabularx}{\linewidth}{@{}l@{\hspace{12pt}}>{\raggedright\arraybackslash}X@{}}
    \SkillRows
  \end{tabularx}}
.
\providecommand{\UCSDCourses}{}
\providecommand{\IITBCourses}{}
\providecommand{\SkillRows}{}
\newcommand{\coursework}[1]{\ifblank{#1}{}{\par\noindent{\fontsize{9.5}{11.5}\selectfont{\color{muted}Coursework:} #1}\par}}
\newcommand{\Education}{%
  \section{Education}
  \entry{University of California San Diego}{M.S. in Computer Science and Engineering (AI)}{GPA 4.00/4.00 \quad Sep 2023 -- Mar 2025}
  \coursework{\UCSDCourses}
  \entry{Indian Institute of Technology Bombay}{B.Tech.\ with Honors in Computer Science}{CPI 9.66/10 \quad Jul 2019 -- Jul 2023}
  \coursework{\IITBCourses}}

% \skill{Label}{Items} is one row of the skills table.
\newcommand{\skill}[2]{{\headfont\bfseries #1} & #2 \\[1.5pt]}
\newcommand{\cpp}{C\kern-0.04em\raisebox{0.06ex}{\small++}}
\newcommand{\Skills}{%
  \section{Skills}
  \begin{tabularx}{\linewidth}{@{}l@{\hspace{12pt}}>{\raggedright\arraybackslash}X@{}}
    \SkillRows
  \end{tabularx}}
 to try another pairing.
% Default: Charter body with Source Sans headings. All-sans alternative:
%   \def\resumefont{\usepackage[default,semibold]{sourcesanspro}}\def\headfont{}
\providecommand{\resumefont}{\usepackage[scaled=0.96]{XCharter}\usepackage[semibold,scale=0.96]{sourcesanspro}}
\providecommand{\headfont}{\sffamily}
\providecommand{\namefont}{\rmfamily\scshape}
\resumefont
\usepackage[final]{microtype}
\usepackage[dvipsnames]{xcolor}
\usepackage{enumitem}
\usepackage{titlesec}
\usepackage{tabularx}
\usepackage{fontawesome5}
\usepackage{etoolbox}
\usepackage[hidelinks]{hyperref}

% The site's green, darkened for print.
\definecolor{accent}{HTML}{3A7A45}
\definecolor{muted}{HTML}{5C675D}

\pagestyle{empty}
\setlength{\parindent}{0pt}
\setlength{\tabcolsep}{0pt}
\providecommand{\leading}{1.1}
\linespread{\leading}

\hypersetup{pdfauthor={Sudhansh Peddabomma}}

% Type scale: name 26pt small caps, section titles 12.5pt caps, entry titles 10.8pt, body 10pt,
% side notes 9pt.
\newcommand{\side}[1]{{\fontsize{9}{11}\selectfont\color{muted}#1}}
\newcommand{\sectitle}[1]{\textls[60]{\MakeUppercase{#1}}}
\titleformat{\section}{\headfont\color{accent}\bfseries\fontsize{12.5}{14}\selectfont}{}{0pt}{\sectitle}[\vspace{1.5pt}{\color{accent}\titlerule[0.8pt]}\global\firstentrytrue]
% The first entry after a section title skips \entrygap, so only \rulegap sits under the rule.
\newif\iffirstentry

% Vertical rhythm, the same idea as the earlier hand-tuned résumés: bullets inside an entry sit
% tight together, entries are separated by a clear gap, and sections by a larger one. A profile
% may override these before % Shared preamble and content for Sudhansh Peddabomma's résumés.
% Each profile (SudhanshPeddabomma-AI.tex, SudhanshPeddabomma-Vision.tex) inputs this file and picks
% which blocks to show and in what order, so every bullet is written once.
% Layout rule: every bullet fills one line, or runs 1.5 to 2 lines. Check with: python3 check_lines.py <pdf>
\documentclass[10pt,letterpaper]{article}

\usepackage[margin=0.55in,top=0.5in,bottom=0.45in]{geometry}
\usepackage[T1]{fontenc}
% Fonts: \resumefont loads the body font; \headfont styles the name and section titles.
% Override either before % Shared preamble and content for Sudhansh Peddabomma's résumés.
% Each profile (SudhanshPeddabomma-AI.tex, SudhanshPeddabomma-Vision.tex) inputs this file and picks
% which blocks to show and in what order, so every bullet is written once.
% Layout rule: every bullet fills one line, or runs 1.5 to 2 lines. Check with: python3 check_lines.py <pdf>
\documentclass[10pt,letterpaper]{article}

\usepackage[margin=0.55in,top=0.5in,bottom=0.45in]{geometry}
\usepackage[T1]{fontenc}
% Fonts: \resumefont loads the body font; \headfont styles the name and section titles.
% Override either before \input{resume-common} to try another pairing.
% Default: Charter body with Source Sans headings. All-sans alternative:
%   \def\resumefont{\usepackage[default,semibold]{sourcesanspro}}\def\headfont{}
\providecommand{\resumefont}{\usepackage[scaled=0.96]{XCharter}\usepackage[semibold,scale=0.96]{sourcesanspro}}
\providecommand{\headfont}{\sffamily}
\providecommand{\namefont}{\rmfamily\scshape}
\resumefont
\usepackage[final]{microtype}
\usepackage[dvipsnames]{xcolor}
\usepackage{enumitem}
\usepackage{titlesec}
\usepackage{tabularx}
\usepackage{fontawesome5}
\usepackage{etoolbox}
\usepackage[hidelinks]{hyperref}

% The site's green, darkened for print.
\definecolor{accent}{HTML}{3A7A45}
\definecolor{muted}{HTML}{5C675D}

\pagestyle{empty}
\setlength{\parindent}{0pt}
\setlength{\tabcolsep}{0pt}
\providecommand{\leading}{1.1}
\linespread{\leading}

\hypersetup{pdfauthor={Sudhansh Peddabomma}}

% Type scale: name 26pt small caps, section titles 12.5pt caps, entry titles 10.8pt, body 10pt,
% side notes 9pt.
\newcommand{\side}[1]{{\fontsize{9}{11}\selectfont\color{muted}#1}}
\newcommand{\sectitle}[1]{\textls[60]{\MakeUppercase{#1}}}
\titleformat{\section}{\headfont\color{accent}\bfseries\fontsize{12.5}{14}\selectfont}{}{0pt}{\sectitle}[\vspace{1.5pt}{\color{accent}\titlerule[0.8pt]}\global\firstentrytrue]
% The first entry after a section title skips \entrygap, so only \rulegap sits under the rule.
\newif\iffirstentry

% Vertical rhythm, the same idea as the earlier hand-tuned résumés: bullets inside an entry sit
% tight together, entries are separated by a clear gap, and sections by a larger one. A profile
% may override these before \input{resume-common} to fit one page.
\providecommand{\secgap}{12pt}   % above a section title
\providecommand{\rulegap}{5.5pt}   % section rule to first entry
\providecommand{\entrygap}{7pt}  % between entries
\providecommand{\headgap}{1.5pt} % entry heading to its bullets
\providecommand{\itemgap}{0.6pt} % between bullets of one entry
\titlespacing*{\section}{0pt}{\secgap}{\rulegap}

\newcommand{\squarebullet}{\raisebox{0.32ex}{\color{accent}\rule{0.5ex}{0.5ex}}}
\setlist[itemize]{leftmargin=12pt, labelsep=6pt, topsep=\headgap, itemsep=\itemgap, parsep=0pt, partopsep=0pt,
  label=\squarebullet}

% \entry{Organisation or title}{Role or detail (may be empty)}{Dates}; \sub{Role or detail}{Right-hand note}
\newcommand{\entry}[3]{%
  \par\iffirstentry\global\firstentryfalse\else\addvspace{\entrygap}\fi\noindent{\headfont\fontsize{10.8}{13}\selectfont\textbf{#1}}\ifblank{#2}{}{\ {\color{muted}|}\ #2}\hfill\hspace{1em}\side{#3}\par}
\newcommand{\sub}[2]{%
  \noindent\textit{#1}\hfill{\color{muted}\textit{#2}}\par}
\newcommand{\report}[1]{\href{#1}{\color{accent}Report}}
\newcommand{\me}{\textbf{S. Peddabomma}}

% ---------------------------------------------------------------------------------------------------
% Header
\newcommand{\Header}{%
  \begin{center}
    {\namefont\fontsize{26}{28}\selectfont Sudhansh Peddabomma}\\[6pt]
    {\headfont\fontsize{9.5}{11}\selectfont\color{muted}
      \href{mailto:sudhansh6@gmail.com}{\faEnvelope[regular]\ sudhansh6@gmail.com} \quad
      \href{https://sudhansh6.github.io}{\faGlobe\ sudhansh6.github.io} \quad
      \href{https://www.linkedin.com/in/sudhansh-peddabomma-73a208192/}{\faLinkedin\ sudhansh-peddabomma} \quad
      \href{https://github.com/Sudhansh6}{\faGithub\ Sudhansh6}}
  \end{center}\vspace{-4pt}}

% ---------------------------------------------------------------------------------------------------
% Glean bullets, shared by both profiles
\newcommand{\GleanSlides}{\item Led \textbf{AI slide generation} in Glean Assistant, now creating \textbf{180K decks a week} for \textbf{50K weekly active users} at enterprises}
\newcommand{\GleanSlidesPerf}{\item Improved the slides harness, raising reliability to \textbf{97\%}, cutting median latency \textbf{26\%} and cost per deck by \textbf{half}
  in four months}
\newcommand{\GleanSlidesOrch}{\item Architected a \textbf{slide orchestration system} that reuses relevant slides from a company's existing decks and generates
  new ones in its brand and template, so teams can refine decks with each other and Glean in full context, without leaving the product}
\newcommand{\GleanAgents}{\item Moved Glean's \textbf{workflow agents} from one-shot planning to an \textbf{agentic tool-use loop}, greatly improving agentic
  capabilities and instruction following, with \textbf{18\%} higher answer completeness and \textbf{17\%} fewer tokens}
\newcommand{\GleanTeamChat}{\item Engineered the core of \textbf{Team chat}, letting teammates work with Glean in one shared conversation and co-edit its artifacts}
\newcommand{\GleanVideo}{\item Shipped \textbf{video artifacts} end to end, generating narrated videos in chat with streamed playback and admin rollout controls}
\newcommand{\GleanEvals}{\item Created the \textbf{evaluations and metrics} that guide model choices and rollouts across agents, slides and interactive artifacts}

% ---------------------------------------------------------------------------------------------------
% Experience
% \groupbreak separates related groups of bullets inside one long entry (slides, then platform work).
\providecommand{\groupgap}{3.5pt}
\newcommand{\groupbreak}{\vspace{\groupgap}}
\newcommand{\GleanAI}{%
  \entry{Glean}{Software Engineer}{Apr 2025 -- Present}
  \begin{itemize}
    \GleanSlides
    \GleanSlidesOrch
    \GleanSlidesPerf \groupbreak
    \GleanTeamChat
    \GleanVideo
    \GleanAgents
    \GleanEvals
  \end{itemize}}

\newcommand{\GleanVision}{%
  \entry{Glean}{Software Engineer}{Apr 2025 -- Present}
  \begin{itemize}
    \GleanSlides
    \GleanSlidesOrch
    \GleanSlidesPerf \groupbreak
    \GleanVideo
    \GleanEvals
  \end{itemize}}

\newcommand{\Duality}{%
  \entry{Duality AI}{Computer Vision Intern}{Jun 2024 -- Sep 2024}
  \begin{itemize}
    \item Designed automated \textbf{3D reconstruction} techniques for featureless objects, reducing digital-twin generation time
          by \textbf{40\%}
    \item Collaborated with \textbf{Autodesk} to validate Unreal Engine simulations for robotics tasks; leveraging structured
          domain randomization to reduce Sim2Real gap and increase mAP-50 by \textbf{15\%} for object detection and segmentation
  \end{itemize}}

\newcommand{\Microsoft}{%
  \entry{Microsoft}{Data and Applied Scientist Intern}{May 2022 -- Jul 2022}
  \begin{itemize}
    \item Developed a \textbf{decision-tree ranker} to recommend emails without user queries, improving \textbf{Outlook} search capabilities
    \item Proposed hierarchical feature-sets for the ranker, reducing latency for recommendations and improving recall
  \end{itemize}}

% ---------------------------------------------------------------------------------------------------
% Research
\newcommand{\HomeRobots}{%
  \entry{Real-time 3D Perception for Home Robots}{UC San Diego}{Sep 2023 -- Sep 2024}
  \begin{itemize}
    \item Investigated real-time dense \textbf{visual SLAM} methods using NeRFs and Gaussian Splatting for indoor robot navigation
    \item Integrated object segmentation, grasp-pose estimation, and 3D mapping on the Fetch robot via ROS, demonstrating
          a novel \textbf{tabletop object rearrangement} algorithm that reduced cost by \textbf{20\%} compared to the state-of-the-art
  \end{itemize}}

\newcommand{\Tomography}{%
  \entry{3D Tomography with Primal-Dual Neural Networks}{University College London}{May 2021 -- Jul 2023}
  \begin{itemize}
    \item Developed a stochastic neural-network architecture of primal-dual algorithm for online reconstruction of 3D
          volumes from tomographic projections, with \textbf{99.6\,\% structural similarity} in low-dosage conditions
    \item Built a Python library with custom gradient operators for reconstructing volumes in a single pass, reducing
          compute requirements by up to \textbf{5x} over SOTA learning-based approaches for cone vector tomography
  \end{itemize}}

% ---------------------------------------------------------------------------------------------------
% Projects
\newcommand{\OOXML}{%
  \entry{ooxml}{Open-source contributor \textbar\ \href{https://github.com/ChristopherVR/ooxml}{\color{accent}GitHub}}{Sep 2026 -- Present}
  \begin{itemize}
    \item Contributed to an open-source \textbf{TypeScript} engine that renders and edits PowerPoint and Word files in the browser,
          improving editing fidelity, \textbf{real-time collaboration} and load performance across React, Vue and Svelte viewers
  \end{itemize}}

\newcommand{\MirrorAI}{%
  \entry{Mirror AI: Deployable Personas}{Honorable mention in Supabase YC Hackathon}{Oct 2024 -- Dec 2024}
  \begin{itemize}
    \item Designed an \textbf{agentic LLM architecture} with LangGraph to mirror user personalities, creating interactive digital
          personas, and deployed it as a full-stack platform on Supabase and Vercel with one-click publishing
  \end{itemize}}

\newcommand{\LLMReasoning}{%
  \entry{Improving LLM Reasoning for Numerical Problems}{\report{https://sudhansh6.github.io/assets/pdf/MathPrompter++.pdf}}{Sep 2024 -- Dec 2024}
  \begin{itemize}
    \item Enhanced MathPrompter (ACL 2023) with CoT, achieving \textbf{10\% higher accuracy} on Llama 3.1 1B where prior methods failed,
          and reduced hallucination rates significantly by integrating multi-step validation
  \end{itemize}}

\newcommand{\InverseRendering}{%
  \entry{Inverse Rendering with 2D Gaussian Splatting}{\report{https://sudhansh6.github.io/assets/pdf/CSE252D_Report.pdf}}{Mar 2024 -- May 2024}
  \begin{itemize}
    \item Developed a novel \textbf{inverse rendering} framework in CUDA to recover PBR properties of a scene using 2D Gaussian
          Splatting, improving normal map MAE by \textbf{15\%} over the current SOTA methods
  \end{itemize}}

% ---------------------------------------------------------------------------------------------------
% Publications: title on the left, venue on the right, one line each.
\newcommand{\publine}[3]{\par\noindent\href{#1}{#2}\hfill\hspace{1em}\side{#3}\par\addvspace{1.5pt}}
\newcommand{\PubSaturation}{\publine{https://doi.org/10.1016/j.sigpro.2023.109349}{A likelihood based method for compressive signal recovery under Gaussian and saturation noise}{Signal Processing, 2024}}
\newcommand{\PubPermutation}{\publine{https://ieeexplore.ieee.org/document/10888147}{Identification and correction of permutation errors in compressed sensing-based group testing}{ICASSP 2025}}
\newcommand{\PubRearrangement}{\publine{https://arxiv.org/abs/2411.10899}{Planning for tabletop object rearrangement}{ICRA 2025}}

% ---------------------------------------------------------------------------------------------------
% Education, skills, service
% Each profile sets \UCSDCourses, \IITBCourses and \SkillRows before \input{resume-common}.
\providecommand{\UCSDCourses}{}
\providecommand{\IITBCourses}{}
\providecommand{\SkillRows}{}
\newcommand{\coursework}[1]{\ifblank{#1}{}{\par\noindent{\fontsize{9.5}{11.5}\selectfont{\color{muted}Coursework:} #1}\par}}
\newcommand{\Education}{%
  \section{Education}
  \entry{University of California San Diego}{M.S. in Computer Science and Engineering (AI)}{GPA 4.00/4.00 \quad Sep 2023 -- Mar 2025}
  \coursework{\UCSDCourses}
  \entry{Indian Institute of Technology Bombay}{B.Tech.\ with Honors in Computer Science}{CPI 9.66/10 \quad Jul 2019 -- Jul 2023}
  \coursework{\IITBCourses}}

% \skill{Label}{Items} is one row of the skills table.
\newcommand{\skill}[2]{{\headfont\bfseries #1} & #2 \\[1.5pt]}
\newcommand{\cpp}{C\kern-0.04em\raisebox{0.06ex}{\small++}}
\newcommand{\Skills}{%
  \section{Skills}
  \begin{tabularx}{\linewidth}{@{}l@{\hspace{12pt}}>{\raggedright\arraybackslash}X@{}}
    \SkillRows
  \end{tabularx}}
 to try another pairing.
% Default: Charter body with Source Sans headings. All-sans alternative:
%   \def\resumefont{\usepackage[default,semibold]{sourcesanspro}}\def\headfont{}
\providecommand{\resumefont}{\usepackage[scaled=0.96]{XCharter}\usepackage[semibold,scale=0.96]{sourcesanspro}}
\providecommand{\headfont}{\sffamily}
\providecommand{\namefont}{\rmfamily\scshape}
\resumefont
\usepackage[final]{microtype}
\usepackage[dvipsnames]{xcolor}
\usepackage{enumitem}
\usepackage{titlesec}
\usepackage{tabularx}
\usepackage{fontawesome5}
\usepackage{etoolbox}
\usepackage[hidelinks]{hyperref}

% The site's green, darkened for print.
\definecolor{accent}{HTML}{3A7A45}
\definecolor{muted}{HTML}{5C675D}

\pagestyle{empty}
\setlength{\parindent}{0pt}
\setlength{\tabcolsep}{0pt}
\providecommand{\leading}{1.1}
\linespread{\leading}

\hypersetup{pdfauthor={Sudhansh Peddabomma}}

% Type scale: name 26pt small caps, section titles 12.5pt caps, entry titles 10.8pt, body 10pt,
% side notes 9pt.
\newcommand{\side}[1]{{\fontsize{9}{11}\selectfont\color{muted}#1}}
\newcommand{\sectitle}[1]{\textls[60]{\MakeUppercase{#1}}}
\titleformat{\section}{\headfont\color{accent}\bfseries\fontsize{12.5}{14}\selectfont}{}{0pt}{\sectitle}[\vspace{1.5pt}{\color{accent}\titlerule[0.8pt]}\global\firstentrytrue]
% The first entry after a section title skips \entrygap, so only \rulegap sits under the rule.
\newif\iffirstentry

% Vertical rhythm, the same idea as the earlier hand-tuned résumés: bullets inside an entry sit
% tight together, entries are separated by a clear gap, and sections by a larger one. A profile
% may override these before % Shared preamble and content for Sudhansh Peddabomma's résumés.
% Each profile (SudhanshPeddabomma-AI.tex, SudhanshPeddabomma-Vision.tex) inputs this file and picks
% which blocks to show and in what order, so every bullet is written once.
% Layout rule: every bullet fills one line, or runs 1.5 to 2 lines. Check with: python3 check_lines.py <pdf>
\documentclass[10pt,letterpaper]{article}

\usepackage[margin=0.55in,top=0.5in,bottom=0.45in]{geometry}
\usepackage[T1]{fontenc}
% Fonts: \resumefont loads the body font; \headfont styles the name and section titles.
% Override either before \input{resume-common} to try another pairing.
% Default: Charter body with Source Sans headings. All-sans alternative:
%   \def\resumefont{\usepackage[default,semibold]{sourcesanspro}}\def\headfont{}
\providecommand{\resumefont}{\usepackage[scaled=0.96]{XCharter}\usepackage[semibold,scale=0.96]{sourcesanspro}}
\providecommand{\headfont}{\sffamily}
\providecommand{\namefont}{\rmfamily\scshape}
\resumefont
\usepackage[final]{microtype}
\usepackage[dvipsnames]{xcolor}
\usepackage{enumitem}
\usepackage{titlesec}
\usepackage{tabularx}
\usepackage{fontawesome5}
\usepackage{etoolbox}
\usepackage[hidelinks]{hyperref}

% The site's green, darkened for print.
\definecolor{accent}{HTML}{3A7A45}
\definecolor{muted}{HTML}{5C675D}

\pagestyle{empty}
\setlength{\parindent}{0pt}
\setlength{\tabcolsep}{0pt}
\providecommand{\leading}{1.1}
\linespread{\leading}

\hypersetup{pdfauthor={Sudhansh Peddabomma}}

% Type scale: name 26pt small caps, section titles 12.5pt caps, entry titles 10.8pt, body 10pt,
% side notes 9pt.
\newcommand{\side}[1]{{\fontsize{9}{11}\selectfont\color{muted}#1}}
\newcommand{\sectitle}[1]{\textls[60]{\MakeUppercase{#1}}}
\titleformat{\section}{\headfont\color{accent}\bfseries\fontsize{12.5}{14}\selectfont}{}{0pt}{\sectitle}[\vspace{1.5pt}{\color{accent}\titlerule[0.8pt]}\global\firstentrytrue]
% The first entry after a section title skips \entrygap, so only \rulegap sits under the rule.
\newif\iffirstentry

% Vertical rhythm, the same idea as the earlier hand-tuned résumés: bullets inside an entry sit
% tight together, entries are separated by a clear gap, and sections by a larger one. A profile
% may override these before \input{resume-common} to fit one page.
\providecommand{\secgap}{12pt}   % above a section title
\providecommand{\rulegap}{5.5pt}   % section rule to first entry
\providecommand{\entrygap}{7pt}  % between entries
\providecommand{\headgap}{1.5pt} % entry heading to its bullets
\providecommand{\itemgap}{0.6pt} % between bullets of one entry
\titlespacing*{\section}{0pt}{\secgap}{\rulegap}

\newcommand{\squarebullet}{\raisebox{0.32ex}{\color{accent}\rule{0.5ex}{0.5ex}}}
\setlist[itemize]{leftmargin=12pt, labelsep=6pt, topsep=\headgap, itemsep=\itemgap, parsep=0pt, partopsep=0pt,
  label=\squarebullet}

% \entry{Organisation or title}{Role or detail (may be empty)}{Dates}; \sub{Role or detail}{Right-hand note}
\newcommand{\entry}[3]{%
  \par\iffirstentry\global\firstentryfalse\else\addvspace{\entrygap}\fi\noindent{\headfont\fontsize{10.8}{13}\selectfont\textbf{#1}}\ifblank{#2}{}{\ {\color{muted}|}\ #2}\hfill\hspace{1em}\side{#3}\par}
\newcommand{\sub}[2]{%
  \noindent\textit{#1}\hfill{\color{muted}\textit{#2}}\par}
\newcommand{\report}[1]{\href{#1}{\color{accent}Report}}
\newcommand{\me}{\textbf{S. Peddabomma}}

% ---------------------------------------------------------------------------------------------------
% Header
\newcommand{\Header}{%
  \begin{center}
    {\namefont\fontsize{26}{28}\selectfont Sudhansh Peddabomma}\\[6pt]
    {\headfont\fontsize{9.5}{11}\selectfont\color{muted}
      \href{mailto:sudhansh6@gmail.com}{\faEnvelope[regular]\ sudhansh6@gmail.com} \quad
      \href{https://sudhansh6.github.io}{\faGlobe\ sudhansh6.github.io} \quad
      \href{https://www.linkedin.com/in/sudhansh-peddabomma-73a208192/}{\faLinkedin\ sudhansh-peddabomma} \quad
      \href{https://github.com/Sudhansh6}{\faGithub\ Sudhansh6}}
  \end{center}\vspace{-4pt}}

% ---------------------------------------------------------------------------------------------------
% Glean bullets, shared by both profiles
\newcommand{\GleanSlides}{\item Led \textbf{AI slide generation} in Glean Assistant, now creating \textbf{180K decks a week} for \textbf{50K weekly active users} at enterprises}
\newcommand{\GleanSlidesPerf}{\item Improved the slides harness, raising reliability to \textbf{97\%}, cutting median latency \textbf{26\%} and cost per deck by \textbf{half}
  in four months}
\newcommand{\GleanSlidesOrch}{\item Architected a \textbf{slide orchestration system} that reuses relevant slides from a company's existing decks and generates
  new ones in its brand and template, so teams can refine decks with each other and Glean in full context, without leaving the product}
\newcommand{\GleanAgents}{\item Moved Glean's \textbf{workflow agents} from one-shot planning to an \textbf{agentic tool-use loop}, greatly improving agentic
  capabilities and instruction following, with \textbf{18\%} higher answer completeness and \textbf{17\%} fewer tokens}
\newcommand{\GleanTeamChat}{\item Engineered the core of \textbf{Team chat}, letting teammates work with Glean in one shared conversation and co-edit its artifacts}
\newcommand{\GleanVideo}{\item Shipped \textbf{video artifacts} end to end, generating narrated videos in chat with streamed playback and admin rollout controls}
\newcommand{\GleanEvals}{\item Created the \textbf{evaluations and metrics} that guide model choices and rollouts across agents, slides and interactive artifacts}

% ---------------------------------------------------------------------------------------------------
% Experience
% \groupbreak separates related groups of bullets inside one long entry (slides, then platform work).
\providecommand{\groupgap}{3.5pt}
\newcommand{\groupbreak}{\vspace{\groupgap}}
\newcommand{\GleanAI}{%
  \entry{Glean}{Software Engineer}{Apr 2025 -- Present}
  \begin{itemize}
    \GleanSlides
    \GleanSlidesOrch
    \GleanSlidesPerf \groupbreak
    \GleanTeamChat
    \GleanVideo
    \GleanAgents
    \GleanEvals
  \end{itemize}}

\newcommand{\GleanVision}{%
  \entry{Glean}{Software Engineer}{Apr 2025 -- Present}
  \begin{itemize}
    \GleanSlides
    \GleanSlidesOrch
    \GleanSlidesPerf \groupbreak
    \GleanVideo
    \GleanEvals
  \end{itemize}}

\newcommand{\Duality}{%
  \entry{Duality AI}{Computer Vision Intern}{Jun 2024 -- Sep 2024}
  \begin{itemize}
    \item Designed automated \textbf{3D reconstruction} techniques for featureless objects, reducing digital-twin generation time
          by \textbf{40\%}
    \item Collaborated with \textbf{Autodesk} to validate Unreal Engine simulations for robotics tasks; leveraging structured
          domain randomization to reduce Sim2Real gap and increase mAP-50 by \textbf{15\%} for object detection and segmentation
  \end{itemize}}

\newcommand{\Microsoft}{%
  \entry{Microsoft}{Data and Applied Scientist Intern}{May 2022 -- Jul 2022}
  \begin{itemize}
    \item Developed a \textbf{decision-tree ranker} to recommend emails without user queries, improving \textbf{Outlook} search capabilities
    \item Proposed hierarchical feature-sets for the ranker, reducing latency for recommendations and improving recall
  \end{itemize}}

% ---------------------------------------------------------------------------------------------------
% Research
\newcommand{\HomeRobots}{%
  \entry{Real-time 3D Perception for Home Robots}{UC San Diego}{Sep 2023 -- Sep 2024}
  \begin{itemize}
    \item Investigated real-time dense \textbf{visual SLAM} methods using NeRFs and Gaussian Splatting for indoor robot navigation
    \item Integrated object segmentation, grasp-pose estimation, and 3D mapping on the Fetch robot via ROS, demonstrating
          a novel \textbf{tabletop object rearrangement} algorithm that reduced cost by \textbf{20\%} compared to the state-of-the-art
  \end{itemize}}

\newcommand{\Tomography}{%
  \entry{3D Tomography with Primal-Dual Neural Networks}{University College London}{May 2021 -- Jul 2023}
  \begin{itemize}
    \item Developed a stochastic neural-network architecture of primal-dual algorithm for online reconstruction of 3D
          volumes from tomographic projections, with \textbf{99.6\,\% structural similarity} in low-dosage conditions
    \item Built a Python library with custom gradient operators for reconstructing volumes in a single pass, reducing
          compute requirements by up to \textbf{5x} over SOTA learning-based approaches for cone vector tomography
  \end{itemize}}

% ---------------------------------------------------------------------------------------------------
% Projects
\newcommand{\OOXML}{%
  \entry{ooxml}{Open-source contributor \textbar\ \href{https://github.com/ChristopherVR/ooxml}{\color{accent}GitHub}}{Sep 2026 -- Present}
  \begin{itemize}
    \item Contributed to an open-source \textbf{TypeScript} engine that renders and edits PowerPoint and Word files in the browser,
          improving editing fidelity, \textbf{real-time collaboration} and load performance across React, Vue and Svelte viewers
  \end{itemize}}

\newcommand{\MirrorAI}{%
  \entry{Mirror AI: Deployable Personas}{Honorable mention in Supabase YC Hackathon}{Oct 2024 -- Dec 2024}
  \begin{itemize}
    \item Designed an \textbf{agentic LLM architecture} with LangGraph to mirror user personalities, creating interactive digital
          personas, and deployed it as a full-stack platform on Supabase and Vercel with one-click publishing
  \end{itemize}}

\newcommand{\LLMReasoning}{%
  \entry{Improving LLM Reasoning for Numerical Problems}{\report{https://sudhansh6.github.io/assets/pdf/MathPrompter++.pdf}}{Sep 2024 -- Dec 2024}
  \begin{itemize}
    \item Enhanced MathPrompter (ACL 2023) with CoT, achieving \textbf{10\% higher accuracy} on Llama 3.1 1B where prior methods failed,
          and reduced hallucination rates significantly by integrating multi-step validation
  \end{itemize}}

\newcommand{\InverseRendering}{%
  \entry{Inverse Rendering with 2D Gaussian Splatting}{\report{https://sudhansh6.github.io/assets/pdf/CSE252D_Report.pdf}}{Mar 2024 -- May 2024}
  \begin{itemize}
    \item Developed a novel \textbf{inverse rendering} framework in CUDA to recover PBR properties of a scene using 2D Gaussian
          Splatting, improving normal map MAE by \textbf{15\%} over the current SOTA methods
  \end{itemize}}

% ---------------------------------------------------------------------------------------------------
% Publications: title on the left, venue on the right, one line each.
\newcommand{\publine}[3]{\par\noindent\href{#1}{#2}\hfill\hspace{1em}\side{#3}\par\addvspace{1.5pt}}
\newcommand{\PubSaturation}{\publine{https://doi.org/10.1016/j.sigpro.2023.109349}{A likelihood based method for compressive signal recovery under Gaussian and saturation noise}{Signal Processing, 2024}}
\newcommand{\PubPermutation}{\publine{https://ieeexplore.ieee.org/document/10888147}{Identification and correction of permutation errors in compressed sensing-based group testing}{ICASSP 2025}}
\newcommand{\PubRearrangement}{\publine{https://arxiv.org/abs/2411.10899}{Planning for tabletop object rearrangement}{ICRA 2025}}

% ---------------------------------------------------------------------------------------------------
% Education, skills, service
% Each profile sets \UCSDCourses, \IITBCourses and \SkillRows before \input{resume-common}.
\providecommand{\UCSDCourses}{}
\providecommand{\IITBCourses}{}
\providecommand{\SkillRows}{}
\newcommand{\coursework}[1]{\ifblank{#1}{}{\par\noindent{\fontsize{9.5}{11.5}\selectfont{\color{muted}Coursework:} #1}\par}}
\newcommand{\Education}{%
  \section{Education}
  \entry{University of California San Diego}{M.S. in Computer Science and Engineering (AI)}{GPA 4.00/4.00 \quad Sep 2023 -- Mar 2025}
  \coursework{\UCSDCourses}
  \entry{Indian Institute of Technology Bombay}{B.Tech.\ with Honors in Computer Science}{CPI 9.66/10 \quad Jul 2019 -- Jul 2023}
  \coursework{\IITBCourses}}

% \skill{Label}{Items} is one row of the skills table.
\newcommand{\skill}[2]{{\headfont\bfseries #1} & #2 \\[1.5pt]}
\newcommand{\cpp}{C\kern-0.04em\raisebox{0.06ex}{\small++}}
\newcommand{\Skills}{%
  \section{Skills}
  \begin{tabularx}{\linewidth}{@{}l@{\hspace{12pt}}>{\raggedright\arraybackslash}X@{}}
    \SkillRows
  \end{tabularx}}
 to fit one page.
\providecommand{\secgap}{12pt}   % above a section title
\providecommand{\rulegap}{5.5pt}   % section rule to first entry
\providecommand{\entrygap}{7pt}  % between entries
\providecommand{\headgap}{1.5pt} % entry heading to its bullets
\providecommand{\itemgap}{0.6pt} % between bullets of one entry
\titlespacing*{\section}{0pt}{\secgap}{\rulegap}

\newcommand{\squarebullet}{\raisebox{0.32ex}{\color{accent}\rule{0.5ex}{0.5ex}}}
\setlist[itemize]{leftmargin=12pt, labelsep=6pt, topsep=\headgap, itemsep=\itemgap, parsep=0pt, partopsep=0pt,
  label=\squarebullet}

% \entry{Organisation or title}{Role or detail (may be empty)}{Dates}; \sub{Role or detail}{Right-hand note}
\newcommand{\entry}[3]{%
  \par\iffirstentry\global\firstentryfalse\else\addvspace{\entrygap}\fi\noindent{\headfont\fontsize{10.8}{13}\selectfont\textbf{#1}}\ifblank{#2}{}{\ {\color{muted}|}\ #2}\hfill\hspace{1em}\side{#3}\par}
\newcommand{\sub}[2]{%
  \noindent\textit{#1}\hfill{\color{muted}\textit{#2}}\par}
\newcommand{\report}[1]{\href{#1}{\color{accent}Report}}
\newcommand{\me}{\textbf{S. Peddabomma}}

% ---------------------------------------------------------------------------------------------------
% Header
\newcommand{\Header}{%
  \begin{center}
    {\namefont\fontsize{26}{28}\selectfont Sudhansh Peddabomma}\\[6pt]
    {\headfont\fontsize{9.5}{11}\selectfont\color{muted}
      \href{mailto:sudhansh6@gmail.com}{\faEnvelope[regular]\ sudhansh6@gmail.com} \quad
      \href{https://sudhansh6.github.io}{\faGlobe\ sudhansh6.github.io} \quad
      \href{https://www.linkedin.com/in/sudhansh-peddabomma-73a208192/}{\faLinkedin\ sudhansh-peddabomma} \quad
      \href{https://github.com/Sudhansh6}{\faGithub\ Sudhansh6}}
  \end{center}\vspace{-4pt}}

% ---------------------------------------------------------------------------------------------------
% Glean bullets, shared by both profiles
\newcommand{\GleanSlides}{\item Led \textbf{AI slide generation} in Glean Assistant, now creating \textbf{180K decks a week} for \textbf{50K weekly active users} at enterprises}
\newcommand{\GleanSlidesPerf}{\item Improved the slides harness, raising reliability to \textbf{97\%}, cutting median latency \textbf{26\%} and cost per deck by \textbf{half}
  in four months}
\newcommand{\GleanSlidesOrch}{\item Architected a \textbf{slide orchestration system} that reuses relevant slides from a company's existing decks and generates
  new ones in its brand and template, so teams can refine decks with each other and Glean in full context, without leaving the product}
\newcommand{\GleanAgents}{\item Moved Glean's \textbf{workflow agents} from one-shot planning to an \textbf{agentic tool-use loop}, greatly improving agentic
  capabilities and instruction following, with \textbf{18\%} higher answer completeness and \textbf{17\%} fewer tokens}
\newcommand{\GleanTeamChat}{\item Engineered the core of \textbf{Team chat}, letting teammates work with Glean in one shared conversation and co-edit its artifacts}
\newcommand{\GleanVideo}{\item Shipped \textbf{video artifacts} end to end, generating narrated videos in chat with streamed playback and admin rollout controls}
\newcommand{\GleanEvals}{\item Created the \textbf{evaluations and metrics} that guide model choices and rollouts across agents, slides and interactive artifacts}

% ---------------------------------------------------------------------------------------------------
% Experience
% \groupbreak separates related groups of bullets inside one long entry (slides, then platform work).
\providecommand{\groupgap}{3.5pt}
\newcommand{\groupbreak}{\vspace{\groupgap}}
\newcommand{\GleanAI}{%
  \entry{Glean}{Software Engineer}{Apr 2025 -- Present}
  \begin{itemize}
    \GleanSlides
    \GleanSlidesOrch
    \GleanSlidesPerf \groupbreak
    \GleanTeamChat
    \GleanVideo
    \GleanAgents
    \GleanEvals
  \end{itemize}}

\newcommand{\GleanVision}{%
  \entry{Glean}{Software Engineer}{Apr 2025 -- Present}
  \begin{itemize}
    \GleanSlides
    \GleanSlidesOrch
    \GleanSlidesPerf \groupbreak
    \GleanVideo
    \GleanEvals
  \end{itemize}}

\newcommand{\Duality}{%
  \entry{Duality AI}{Computer Vision Intern}{Jun 2024 -- Sep 2024}
  \begin{itemize}
    \item Designed automated \textbf{3D reconstruction} techniques for featureless objects, reducing digital-twin generation time
          by \textbf{40\%}
    \item Collaborated with \textbf{Autodesk} to validate Unreal Engine simulations for robotics tasks; leveraging structured
          domain randomization to reduce Sim2Real gap and increase mAP-50 by \textbf{15\%} for object detection and segmentation
  \end{itemize}}

\newcommand{\Microsoft}{%
  \entry{Microsoft}{Data and Applied Scientist Intern}{May 2022 -- Jul 2022}
  \begin{itemize}
    \item Developed a \textbf{decision-tree ranker} to recommend emails without user queries, improving \textbf{Outlook} search capabilities
    \item Proposed hierarchical feature-sets for the ranker, reducing latency for recommendations and improving recall
  \end{itemize}}

% ---------------------------------------------------------------------------------------------------
% Research
\newcommand{\HomeRobots}{%
  \entry{Real-time 3D Perception for Home Robots}{UC San Diego}{Sep 2023 -- Sep 2024}
  \begin{itemize}
    \item Investigated real-time dense \textbf{visual SLAM} methods using NeRFs and Gaussian Splatting for indoor robot navigation
    \item Integrated object segmentation, grasp-pose estimation, and 3D mapping on the Fetch robot via ROS, demonstrating
          a novel \textbf{tabletop object rearrangement} algorithm that reduced cost by \textbf{20\%} compared to the state-of-the-art
  \end{itemize}}

\newcommand{\Tomography}{%
  \entry{3D Tomography with Primal-Dual Neural Networks}{University College London}{May 2021 -- Jul 2023}
  \begin{itemize}
    \item Developed a stochastic neural-network architecture of primal-dual algorithm for online reconstruction of 3D
          volumes from tomographic projections, with \textbf{99.6\,\% structural similarity} in low-dosage conditions
    \item Built a Python library with custom gradient operators for reconstructing volumes in a single pass, reducing
          compute requirements by up to \textbf{5x} over SOTA learning-based approaches for cone vector tomography
  \end{itemize}}

% ---------------------------------------------------------------------------------------------------
% Projects
\newcommand{\OOXML}{%
  \entry{ooxml}{Open-source contributor \textbar\ \href{https://github.com/ChristopherVR/ooxml}{\color{accent}GitHub}}{Sep 2026 -- Present}
  \begin{itemize}
    \item Contributed to an open-source \textbf{TypeScript} engine that renders and edits PowerPoint and Word files in the browser,
          improving editing fidelity, \textbf{real-time collaboration} and load performance across React, Vue and Svelte viewers
  \end{itemize}}

\newcommand{\MirrorAI}{%
  \entry{Mirror AI: Deployable Personas}{Honorable mention in Supabase YC Hackathon}{Oct 2024 -- Dec 2024}
  \begin{itemize}
    \item Designed an \textbf{agentic LLM architecture} with LangGraph to mirror user personalities, creating interactive digital
          personas, and deployed it as a full-stack platform on Supabase and Vercel with one-click publishing
  \end{itemize}}

\newcommand{\LLMReasoning}{%
  \entry{Improving LLM Reasoning for Numerical Problems}{\report{https://sudhansh6.github.io/assets/pdf/MathPrompter++.pdf}}{Sep 2024 -- Dec 2024}
  \begin{itemize}
    \item Enhanced MathPrompter (ACL 2023) with CoT, achieving \textbf{10\% higher accuracy} on Llama 3.1 1B where prior methods failed,
          and reduced hallucination rates significantly by integrating multi-step validation
  \end{itemize}}

\newcommand{\InverseRendering}{%
  \entry{Inverse Rendering with 2D Gaussian Splatting}{\report{https://sudhansh6.github.io/assets/pdf/CSE252D_Report.pdf}}{Mar 2024 -- May 2024}
  \begin{itemize}
    \item Developed a novel \textbf{inverse rendering} framework in CUDA to recover PBR properties of a scene using 2D Gaussian
          Splatting, improving normal map MAE by \textbf{15\%} over the current SOTA methods
  \end{itemize}}

% ---------------------------------------------------------------------------------------------------
% Publications: title on the left, venue on the right, one line each.
\newcommand{\publine}[3]{\par\noindent\href{#1}{#2}\hfill\hspace{1em}\side{#3}\par\addvspace{1.5pt}}
\newcommand{\PubSaturation}{\publine{https://doi.org/10.1016/j.sigpro.2023.109349}{A likelihood based method for compressive signal recovery under Gaussian and saturation noise}{Signal Processing, 2024}}
\newcommand{\PubPermutation}{\publine{https://ieeexplore.ieee.org/document/10888147}{Identification and correction of permutation errors in compressed sensing-based group testing}{ICASSP 2025}}
\newcommand{\PubRearrangement}{\publine{https://arxiv.org/abs/2411.10899}{Planning for tabletop object rearrangement}{ICRA 2025}}

% ---------------------------------------------------------------------------------------------------
% Education, skills, service
% Each profile sets \UCSDCourses, \IITBCourses and \SkillRows before % Shared preamble and content for Sudhansh Peddabomma's résumés.
% Each profile (SudhanshPeddabomma-AI.tex, SudhanshPeddabomma-Vision.tex) inputs this file and picks
% which blocks to show and in what order, so every bullet is written once.
% Layout rule: every bullet fills one line, or runs 1.5 to 2 lines. Check with: python3 check_lines.py <pdf>
\documentclass[10pt,letterpaper]{article}

\usepackage[margin=0.55in,top=0.5in,bottom=0.45in]{geometry}
\usepackage[T1]{fontenc}
% Fonts: \resumefont loads the body font; \headfont styles the name and section titles.
% Override either before \input{resume-common} to try another pairing.
% Default: Charter body with Source Sans headings. All-sans alternative:
%   \def\resumefont{\usepackage[default,semibold]{sourcesanspro}}\def\headfont{}
\providecommand{\resumefont}{\usepackage[scaled=0.96]{XCharter}\usepackage[semibold,scale=0.96]{sourcesanspro}}
\providecommand{\headfont}{\sffamily}
\providecommand{\namefont}{\rmfamily\scshape}
\resumefont
\usepackage[final]{microtype}
\usepackage[dvipsnames]{xcolor}
\usepackage{enumitem}
\usepackage{titlesec}
\usepackage{tabularx}
\usepackage{fontawesome5}
\usepackage{etoolbox}
\usepackage[hidelinks]{hyperref}

% The site's green, darkened for print.
\definecolor{accent}{HTML}{3A7A45}
\definecolor{muted}{HTML}{5C675D}

\pagestyle{empty}
\setlength{\parindent}{0pt}
\setlength{\tabcolsep}{0pt}
\providecommand{\leading}{1.1}
\linespread{\leading}

\hypersetup{pdfauthor={Sudhansh Peddabomma}}

% Type scale: name 26pt small caps, section titles 12.5pt caps, entry titles 10.8pt, body 10pt,
% side notes 9pt.
\newcommand{\side}[1]{{\fontsize{9}{11}\selectfont\color{muted}#1}}
\newcommand{\sectitle}[1]{\textls[60]{\MakeUppercase{#1}}}
\titleformat{\section}{\headfont\color{accent}\bfseries\fontsize{12.5}{14}\selectfont}{}{0pt}{\sectitle}[\vspace{1.5pt}{\color{accent}\titlerule[0.8pt]}\global\firstentrytrue]
% The first entry after a section title skips \entrygap, so only \rulegap sits under the rule.
\newif\iffirstentry

% Vertical rhythm, the same idea as the earlier hand-tuned résumés: bullets inside an entry sit
% tight together, entries are separated by a clear gap, and sections by a larger one. A profile
% may override these before \input{resume-common} to fit one page.
\providecommand{\secgap}{12pt}   % above a section title
\providecommand{\rulegap}{5.5pt}   % section rule to first entry
\providecommand{\entrygap}{7pt}  % between entries
\providecommand{\headgap}{1.5pt} % entry heading to its bullets
\providecommand{\itemgap}{0.6pt} % between bullets of one entry
\titlespacing*{\section}{0pt}{\secgap}{\rulegap}

\newcommand{\squarebullet}{\raisebox{0.32ex}{\color{accent}\rule{0.5ex}{0.5ex}}}
\setlist[itemize]{leftmargin=12pt, labelsep=6pt, topsep=\headgap, itemsep=\itemgap, parsep=0pt, partopsep=0pt,
  label=\squarebullet}

% \entry{Organisation or title}{Role or detail (may be empty)}{Dates}; \sub{Role or detail}{Right-hand note}
\newcommand{\entry}[3]{%
  \par\iffirstentry\global\firstentryfalse\else\addvspace{\entrygap}\fi\noindent{\headfont\fontsize{10.8}{13}\selectfont\textbf{#1}}\ifblank{#2}{}{\ {\color{muted}|}\ #2}\hfill\hspace{1em}\side{#3}\par}
\newcommand{\sub}[2]{%
  \noindent\textit{#1}\hfill{\color{muted}\textit{#2}}\par}
\newcommand{\report}[1]{\href{#1}{\color{accent}Report}}
\newcommand{\me}{\textbf{S. Peddabomma}}

% ---------------------------------------------------------------------------------------------------
% Header
\newcommand{\Header}{%
  \begin{center}
    {\namefont\fontsize{26}{28}\selectfont Sudhansh Peddabomma}\\[6pt]
    {\headfont\fontsize{9.5}{11}\selectfont\color{muted}
      \href{mailto:sudhansh6@gmail.com}{\faEnvelope[regular]\ sudhansh6@gmail.com} \quad
      \href{https://sudhansh6.github.io}{\faGlobe\ sudhansh6.github.io} \quad
      \href{https://www.linkedin.com/in/sudhansh-peddabomma-73a208192/}{\faLinkedin\ sudhansh-peddabomma} \quad
      \href{https://github.com/Sudhansh6}{\faGithub\ Sudhansh6}}
  \end{center}\vspace{-4pt}}

% ---------------------------------------------------------------------------------------------------
% Glean bullets, shared by both profiles
\newcommand{\GleanSlides}{\item Led \textbf{AI slide generation} in Glean Assistant, now creating \textbf{180K decks a week} for \textbf{50K weekly active users} at enterprises}
\newcommand{\GleanSlidesPerf}{\item Improved the slides harness, raising reliability to \textbf{97\%}, cutting median latency \textbf{26\%} and cost per deck by \textbf{half}
  in four months}
\newcommand{\GleanSlidesOrch}{\item Architected a \textbf{slide orchestration system} that reuses relevant slides from a company's existing decks and generates
  new ones in its brand and template, so teams can refine decks with each other and Glean in full context, without leaving the product}
\newcommand{\GleanAgents}{\item Moved Glean's \textbf{workflow agents} from one-shot planning to an \textbf{agentic tool-use loop}, greatly improving agentic
  capabilities and instruction following, with \textbf{18\%} higher answer completeness and \textbf{17\%} fewer tokens}
\newcommand{\GleanTeamChat}{\item Engineered the core of \textbf{Team chat}, letting teammates work with Glean in one shared conversation and co-edit its artifacts}
\newcommand{\GleanVideo}{\item Shipped \textbf{video artifacts} end to end, generating narrated videos in chat with streamed playback and admin rollout controls}
\newcommand{\GleanEvals}{\item Created the \textbf{evaluations and metrics} that guide model choices and rollouts across agents, slides and interactive artifacts}

% ---------------------------------------------------------------------------------------------------
% Experience
% \groupbreak separates related groups of bullets inside one long entry (slides, then platform work).
\providecommand{\groupgap}{3.5pt}
\newcommand{\groupbreak}{\vspace{\groupgap}}
\newcommand{\GleanAI}{%
  \entry{Glean}{Software Engineer}{Apr 2025 -- Present}
  \begin{itemize}
    \GleanSlides
    \GleanSlidesOrch
    \GleanSlidesPerf \groupbreak
    \GleanTeamChat
    \GleanVideo
    \GleanAgents
    \GleanEvals
  \end{itemize}}

\newcommand{\GleanVision}{%
  \entry{Glean}{Software Engineer}{Apr 2025 -- Present}
  \begin{itemize}
    \GleanSlides
    \GleanSlidesOrch
    \GleanSlidesPerf \groupbreak
    \GleanVideo
    \GleanEvals
  \end{itemize}}

\newcommand{\Duality}{%
  \entry{Duality AI}{Computer Vision Intern}{Jun 2024 -- Sep 2024}
  \begin{itemize}
    \item Designed automated \textbf{3D reconstruction} techniques for featureless objects, reducing digital-twin generation time
          by \textbf{40\%}
    \item Collaborated with \textbf{Autodesk} to validate Unreal Engine simulations for robotics tasks; leveraging structured
          domain randomization to reduce Sim2Real gap and increase mAP-50 by \textbf{15\%} for object detection and segmentation
  \end{itemize}}

\newcommand{\Microsoft}{%
  \entry{Microsoft}{Data and Applied Scientist Intern}{May 2022 -- Jul 2022}
  \begin{itemize}
    \item Developed a \textbf{decision-tree ranker} to recommend emails without user queries, improving \textbf{Outlook} search capabilities
    \item Proposed hierarchical feature-sets for the ranker, reducing latency for recommendations and improving recall
  \end{itemize}}

% ---------------------------------------------------------------------------------------------------
% Research
\newcommand{\HomeRobots}{%
  \entry{Real-time 3D Perception for Home Robots}{UC San Diego}{Sep 2023 -- Sep 2024}
  \begin{itemize}
    \item Investigated real-time dense \textbf{visual SLAM} methods using NeRFs and Gaussian Splatting for indoor robot navigation
    \item Integrated object segmentation, grasp-pose estimation, and 3D mapping on the Fetch robot via ROS, demonstrating
          a novel \textbf{tabletop object rearrangement} algorithm that reduced cost by \textbf{20\%} compared to the state-of-the-art
  \end{itemize}}

\newcommand{\Tomography}{%
  \entry{3D Tomography with Primal-Dual Neural Networks}{University College London}{May 2021 -- Jul 2023}
  \begin{itemize}
    \item Developed a stochastic neural-network architecture of primal-dual algorithm for online reconstruction of 3D
          volumes from tomographic projections, with \textbf{99.6\,\% structural similarity} in low-dosage conditions
    \item Built a Python library with custom gradient operators for reconstructing volumes in a single pass, reducing
          compute requirements by up to \textbf{5x} over SOTA learning-based approaches for cone vector tomography
  \end{itemize}}

% ---------------------------------------------------------------------------------------------------
% Projects
\newcommand{\OOXML}{%
  \entry{ooxml}{Open-source contributor \textbar\ \href{https://github.com/ChristopherVR/ooxml}{\color{accent}GitHub}}{Sep 2026 -- Present}
  \begin{itemize}
    \item Contributed to an open-source \textbf{TypeScript} engine that renders and edits PowerPoint and Word files in the browser,
          improving editing fidelity, \textbf{real-time collaboration} and load performance across React, Vue and Svelte viewers
  \end{itemize}}

\newcommand{\MirrorAI}{%
  \entry{Mirror AI: Deployable Personas}{Honorable mention in Supabase YC Hackathon}{Oct 2024 -- Dec 2024}
  \begin{itemize}
    \item Designed an \textbf{agentic LLM architecture} with LangGraph to mirror user personalities, creating interactive digital
          personas, and deployed it as a full-stack platform on Supabase and Vercel with one-click publishing
  \end{itemize}}

\newcommand{\LLMReasoning}{%
  \entry{Improving LLM Reasoning for Numerical Problems}{\report{https://sudhansh6.github.io/assets/pdf/MathPrompter++.pdf}}{Sep 2024 -- Dec 2024}
  \begin{itemize}
    \item Enhanced MathPrompter (ACL 2023) with CoT, achieving \textbf{10\% higher accuracy} on Llama 3.1 1B where prior methods failed,
          and reduced hallucination rates significantly by integrating multi-step validation
  \end{itemize}}

\newcommand{\InverseRendering}{%
  \entry{Inverse Rendering with 2D Gaussian Splatting}{\report{https://sudhansh6.github.io/assets/pdf/CSE252D_Report.pdf}}{Mar 2024 -- May 2024}
  \begin{itemize}
    \item Developed a novel \textbf{inverse rendering} framework in CUDA to recover PBR properties of a scene using 2D Gaussian
          Splatting, improving normal map MAE by \textbf{15\%} over the current SOTA methods
  \end{itemize}}

% ---------------------------------------------------------------------------------------------------
% Publications: title on the left, venue on the right, one line each.
\newcommand{\publine}[3]{\par\noindent\href{#1}{#2}\hfill\hspace{1em}\side{#3}\par\addvspace{1.5pt}}
\newcommand{\PubSaturation}{\publine{https://doi.org/10.1016/j.sigpro.2023.109349}{A likelihood based method for compressive signal recovery under Gaussian and saturation noise}{Signal Processing, 2024}}
\newcommand{\PubPermutation}{\publine{https://ieeexplore.ieee.org/document/10888147}{Identification and correction of permutation errors in compressed sensing-based group testing}{ICASSP 2025}}
\newcommand{\PubRearrangement}{\publine{https://arxiv.org/abs/2411.10899}{Planning for tabletop object rearrangement}{ICRA 2025}}

% ---------------------------------------------------------------------------------------------------
% Education, skills, service
% Each profile sets \UCSDCourses, \IITBCourses and \SkillRows before \input{resume-common}.
\providecommand{\UCSDCourses}{}
\providecommand{\IITBCourses}{}
\providecommand{\SkillRows}{}
\newcommand{\coursework}[1]{\ifblank{#1}{}{\par\noindent{\fontsize{9.5}{11.5}\selectfont{\color{muted}Coursework:} #1}\par}}
\newcommand{\Education}{%
  \section{Education}
  \entry{University of California San Diego}{M.S. in Computer Science and Engineering (AI)}{GPA 4.00/4.00 \quad Sep 2023 -- Mar 2025}
  \coursework{\UCSDCourses}
  \entry{Indian Institute of Technology Bombay}{B.Tech.\ with Honors in Computer Science}{CPI 9.66/10 \quad Jul 2019 -- Jul 2023}
  \coursework{\IITBCourses}}

% \skill{Label}{Items} is one row of the skills table.
\newcommand{\skill}[2]{{\headfont\bfseries #1} & #2 \\[1.5pt]}
\newcommand{\cpp}{C\kern-0.04em\raisebox{0.06ex}{\small++}}
\newcommand{\Skills}{%
  \section{Skills}
  \begin{tabularx}{\linewidth}{@{}l@{\hspace{12pt}}>{\raggedright\arraybackslash}X@{}}
    \SkillRows
  \end{tabularx}}
.
\providecommand{\UCSDCourses}{}
\providecommand{\IITBCourses}{}
\providecommand{\SkillRows}{}
\newcommand{\coursework}[1]{\ifblank{#1}{}{\par\noindent{\fontsize{9.5}{11.5}\selectfont{\color{muted}Coursework:} #1}\par}}
\newcommand{\Education}{%
  \section{Education}
  \entry{University of California San Diego}{M.S. in Computer Science and Engineering (AI)}{GPA 4.00/4.00 \quad Sep 2023 -- Mar 2025}
  \coursework{\UCSDCourses}
  \entry{Indian Institute of Technology Bombay}{B.Tech.\ with Honors in Computer Science}{CPI 9.66/10 \quad Jul 2019 -- Jul 2023}
  \coursework{\IITBCourses}}

% \skill{Label}{Items} is one row of the skills table.
\newcommand{\skill}[2]{{\headfont\bfseries #1} & #2 \\[1.5pt]}
\newcommand{\cpp}{C\kern-0.04em\raisebox{0.06ex}{\small++}}
\newcommand{\Skills}{%
  \section{Skills}
  \begin{tabularx}{\linewidth}{@{}l@{\hspace{12pt}}>{\raggedright\arraybackslash}X@{}}
    \SkillRows
  \end{tabularx}}
 to fit one page.
\providecommand{\secgap}{12pt}   % above a section title
\providecommand{\rulegap}{5.5pt}   % section rule to first entry
\providecommand{\entrygap}{7pt}  % between entries
\providecommand{\headgap}{1.5pt} % entry heading to its bullets
\providecommand{\itemgap}{0.6pt} % between bullets of one entry
\titlespacing*{\section}{0pt}{\secgap}{\rulegap}

\newcommand{\squarebullet}{\raisebox{0.32ex}{\color{accent}\rule{0.5ex}{0.5ex}}}
\setlist[itemize]{leftmargin=12pt, labelsep=6pt, topsep=\headgap, itemsep=\itemgap, parsep=0pt, partopsep=0pt,
  label=\squarebullet}

% \entry{Organisation or title}{Role or detail (may be empty)}{Dates}; \sub{Role or detail}{Right-hand note}
\newcommand{\entry}[3]{%
  \par\iffirstentry\global\firstentryfalse\else\addvspace{\entrygap}\fi\noindent{\headfont\fontsize{10.8}{13}\selectfont\textbf{#1}}\ifblank{#2}{}{\ {\color{muted}|}\ #2}\hfill\hspace{1em}\side{#3}\par}
\newcommand{\sub}[2]{%
  \noindent\textit{#1}\hfill{\color{muted}\textit{#2}}\par}
\newcommand{\report}[1]{\href{#1}{\color{accent}Report}}
\newcommand{\me}{\textbf{S. Peddabomma}}

% ---------------------------------------------------------------------------------------------------
% Header
\newcommand{\Header}{%
  \begin{center}
    {\namefont\fontsize{26}{28}\selectfont Sudhansh Peddabomma}\\[6pt]
    {\headfont\fontsize{9.5}{11}\selectfont\color{muted}
      \href{mailto:sudhansh6@gmail.com}{\faEnvelope[regular]\ sudhansh6@gmail.com} \quad
      \href{https://sudhansh6.github.io}{\faGlobe\ sudhansh6.github.io} \quad
      \href{https://www.linkedin.com/in/sudhansh-peddabomma-73a208192/}{\faLinkedin\ sudhansh-peddabomma} \quad
      \href{https://github.com/Sudhansh6}{\faGithub\ Sudhansh6}}
  \end{center}\vspace{-4pt}}

% ---------------------------------------------------------------------------------------------------
% Glean bullets, shared by both profiles
\newcommand{\GleanSlides}{\item Led \textbf{AI slide generation} in Glean Assistant, now creating \textbf{180K decks a week} for \textbf{50K weekly active users} at enterprises}
\newcommand{\GleanSlidesPerf}{\item Improved the slides harness, raising reliability to \textbf{97\%}, cutting median latency \textbf{26\%} and cost per deck by \textbf{half}
  in four months}
\newcommand{\GleanSlidesOrch}{\item Architected a \textbf{slide orchestration system} that reuses relevant slides from a company's existing decks and generates
  new ones in its brand and template, so teams can refine decks with each other and Glean in full context, without leaving the product}
\newcommand{\GleanAgents}{\item Moved Glean's \textbf{workflow agents} from one-shot planning to an \textbf{agentic tool-use loop}, greatly improving agentic
  capabilities and instruction following, with \textbf{18\%} higher answer completeness and \textbf{17\%} fewer tokens}
\newcommand{\GleanTeamChat}{\item Engineered the core of \textbf{Team chat}, letting teammates work with Glean in one shared conversation and co-edit its artifacts}
\newcommand{\GleanVideo}{\item Shipped \textbf{video artifacts} end to end, generating narrated videos in chat with streamed playback and admin rollout controls}
\newcommand{\GleanEvals}{\item Created the \textbf{evaluations and metrics} that guide model choices and rollouts across agents, slides and interactive artifacts}

% ---------------------------------------------------------------------------------------------------
% Experience
% \groupbreak separates related groups of bullets inside one long entry (slides, then platform work).
\providecommand{\groupgap}{3.5pt}
\newcommand{\groupbreak}{\vspace{\groupgap}}
\newcommand{\GleanAI}{%
  \entry{Glean}{Software Engineer}{Apr 2025 -- Present}
  \begin{itemize}
    \GleanSlides
    \GleanSlidesOrch
    \GleanSlidesPerf \groupbreak
    \GleanTeamChat
    \GleanVideo
    \GleanAgents
    \GleanEvals
  \end{itemize}}

\newcommand{\GleanVision}{%
  \entry{Glean}{Software Engineer}{Apr 2025 -- Present}
  \begin{itemize}
    \GleanSlides
    \GleanSlidesOrch
    \GleanSlidesPerf \groupbreak
    \GleanVideo
    \GleanEvals
  \end{itemize}}

\newcommand{\Duality}{%
  \entry{Duality AI}{Computer Vision Intern}{Jun 2024 -- Sep 2024}
  \begin{itemize}
    \item Designed automated \textbf{3D reconstruction} techniques for featureless objects, reducing digital-twin generation time
          by \textbf{40\%}
    \item Collaborated with \textbf{Autodesk} to validate Unreal Engine simulations for robotics tasks; leveraging structured
          domain randomization to reduce Sim2Real gap and increase mAP-50 by \textbf{15\%} for object detection and segmentation
  \end{itemize}}

\newcommand{\Microsoft}{%
  \entry{Microsoft}{Data and Applied Scientist Intern}{May 2022 -- Jul 2022}
  \begin{itemize}
    \item Developed a \textbf{decision-tree ranker} to recommend emails without user queries, improving \textbf{Outlook} search capabilities
    \item Proposed hierarchical feature-sets for the ranker, reducing latency for recommendations and improving recall
  \end{itemize}}

% ---------------------------------------------------------------------------------------------------
% Research
\newcommand{\HomeRobots}{%
  \entry{Real-time 3D Perception for Home Robots}{UC San Diego}{Sep 2023 -- Sep 2024}
  \begin{itemize}
    \item Investigated real-time dense \textbf{visual SLAM} methods using NeRFs and Gaussian Splatting for indoor robot navigation
    \item Integrated object segmentation, grasp-pose estimation, and 3D mapping on the Fetch robot via ROS, demonstrating
          a novel \textbf{tabletop object rearrangement} algorithm that reduced cost by \textbf{20\%} compared to the state-of-the-art
  \end{itemize}}

\newcommand{\Tomography}{%
  \entry{3D Tomography with Primal-Dual Neural Networks}{University College London}{May 2021 -- Jul 2023}
  \begin{itemize}
    \item Developed a stochastic neural-network architecture of primal-dual algorithm for online reconstruction of 3D
          volumes from tomographic projections, with \textbf{99.6\,\% structural similarity} in low-dosage conditions
    \item Built a Python library with custom gradient operators for reconstructing volumes in a single pass, reducing
          compute requirements by up to \textbf{5x} over SOTA learning-based approaches for cone vector tomography
  \end{itemize}}

% ---------------------------------------------------------------------------------------------------
% Projects
\newcommand{\OOXML}{%
  \entry{ooxml}{Open-source contributor \textbar\ \href{https://github.com/ChristopherVR/ooxml}{\color{accent}GitHub}}{Sep 2026 -- Present}
  \begin{itemize}
    \item Contributed to an open-source \textbf{TypeScript} engine that renders and edits PowerPoint and Word files in the browser,
          improving editing fidelity, \textbf{real-time collaboration} and load performance across React, Vue and Svelte viewers
  \end{itemize}}

\newcommand{\MirrorAI}{%
  \entry{Mirror AI: Deployable Personas}{Honorable mention in Supabase YC Hackathon}{Oct 2024 -- Dec 2024}
  \begin{itemize}
    \item Designed an \textbf{agentic LLM architecture} with LangGraph to mirror user personalities, creating interactive digital
          personas, and deployed it as a full-stack platform on Supabase and Vercel with one-click publishing
  \end{itemize}}

\newcommand{\LLMReasoning}{%
  \entry{Improving LLM Reasoning for Numerical Problems}{\report{https://sudhansh6.github.io/assets/pdf/MathPrompter++.pdf}}{Sep 2024 -- Dec 2024}
  \begin{itemize}
    \item Enhanced MathPrompter (ACL 2023) with CoT, achieving \textbf{10\% higher accuracy} on Llama 3.1 1B where prior methods failed,
          and reduced hallucination rates significantly by integrating multi-step validation
  \end{itemize}}

\newcommand{\InverseRendering}{%
  \entry{Inverse Rendering with 2D Gaussian Splatting}{\report{https://sudhansh6.github.io/assets/pdf/CSE252D_Report.pdf}}{Mar 2024 -- May 2024}
  \begin{itemize}
    \item Developed a novel \textbf{inverse rendering} framework in CUDA to recover PBR properties of a scene using 2D Gaussian
          Splatting, improving normal map MAE by \textbf{15\%} over the current SOTA methods
  \end{itemize}}

% ---------------------------------------------------------------------------------------------------
% Publications: title on the left, venue on the right, one line each.
\newcommand{\publine}[3]{\par\noindent\href{#1}{#2}\hfill\hspace{1em}\side{#3}\par\addvspace{1.5pt}}
\newcommand{\PubSaturation}{\publine{https://doi.org/10.1016/j.sigpro.2023.109349}{A likelihood based method for compressive signal recovery under Gaussian and saturation noise}{Signal Processing, 2024}}
\newcommand{\PubPermutation}{\publine{https://ieeexplore.ieee.org/document/10888147}{Identification and correction of permutation errors in compressed sensing-based group testing}{ICASSP 2025}}
\newcommand{\PubRearrangement}{\publine{https://arxiv.org/abs/2411.10899}{Planning for tabletop object rearrangement}{ICRA 2025}}

% ---------------------------------------------------------------------------------------------------
% Education, skills, service
% Each profile sets \UCSDCourses, \IITBCourses and \SkillRows before % Shared preamble and content for Sudhansh Peddabomma's résumés.
% Each profile (SudhanshPeddabomma-AI.tex, SudhanshPeddabomma-Vision.tex) inputs this file and picks
% which blocks to show and in what order, so every bullet is written once.
% Layout rule: every bullet fills one line, or runs 1.5 to 2 lines. Check with: python3 check_lines.py <pdf>
\documentclass[10pt,letterpaper]{article}

\usepackage[margin=0.55in,top=0.5in,bottom=0.45in]{geometry}
\usepackage[T1]{fontenc}
% Fonts: \resumefont loads the body font; \headfont styles the name and section titles.
% Override either before % Shared preamble and content for Sudhansh Peddabomma's résumés.
% Each profile (SudhanshPeddabomma-AI.tex, SudhanshPeddabomma-Vision.tex) inputs this file and picks
% which blocks to show and in what order, so every bullet is written once.
% Layout rule: every bullet fills one line, or runs 1.5 to 2 lines. Check with: python3 check_lines.py <pdf>
\documentclass[10pt,letterpaper]{article}

\usepackage[margin=0.55in,top=0.5in,bottom=0.45in]{geometry}
\usepackage[T1]{fontenc}
% Fonts: \resumefont loads the body font; \headfont styles the name and section titles.
% Override either before \input{resume-common} to try another pairing.
% Default: Charter body with Source Sans headings. All-sans alternative:
%   \def\resumefont{\usepackage[default,semibold]{sourcesanspro}}\def\headfont{}
\providecommand{\resumefont}{\usepackage[scaled=0.96]{XCharter}\usepackage[semibold,scale=0.96]{sourcesanspro}}
\providecommand{\headfont}{\sffamily}
\providecommand{\namefont}{\rmfamily\scshape}
\resumefont
\usepackage[final]{microtype}
\usepackage[dvipsnames]{xcolor}
\usepackage{enumitem}
\usepackage{titlesec}
\usepackage{tabularx}
\usepackage{fontawesome5}
\usepackage{etoolbox}
\usepackage[hidelinks]{hyperref}

% The site's green, darkened for print.
\definecolor{accent}{HTML}{3A7A45}
\definecolor{muted}{HTML}{5C675D}

\pagestyle{empty}
\setlength{\parindent}{0pt}
\setlength{\tabcolsep}{0pt}
\providecommand{\leading}{1.1}
\linespread{\leading}

\hypersetup{pdfauthor={Sudhansh Peddabomma}}

% Type scale: name 26pt small caps, section titles 12.5pt caps, entry titles 10.8pt, body 10pt,
% side notes 9pt.
\newcommand{\side}[1]{{\fontsize{9}{11}\selectfont\color{muted}#1}}
\newcommand{\sectitle}[1]{\textls[60]{\MakeUppercase{#1}}}
\titleformat{\section}{\headfont\color{accent}\bfseries\fontsize{12.5}{14}\selectfont}{}{0pt}{\sectitle}[\vspace{1.5pt}{\color{accent}\titlerule[0.8pt]}\global\firstentrytrue]
% The first entry after a section title skips \entrygap, so only \rulegap sits under the rule.
\newif\iffirstentry

% Vertical rhythm, the same idea as the earlier hand-tuned résumés: bullets inside an entry sit
% tight together, entries are separated by a clear gap, and sections by a larger one. A profile
% may override these before \input{resume-common} to fit one page.
\providecommand{\secgap}{12pt}   % above a section title
\providecommand{\rulegap}{5.5pt}   % section rule to first entry
\providecommand{\entrygap}{7pt}  % between entries
\providecommand{\headgap}{1.5pt} % entry heading to its bullets
\providecommand{\itemgap}{0.6pt} % between bullets of one entry
\titlespacing*{\section}{0pt}{\secgap}{\rulegap}

\newcommand{\squarebullet}{\raisebox{0.32ex}{\color{accent}\rule{0.5ex}{0.5ex}}}
\setlist[itemize]{leftmargin=12pt, labelsep=6pt, topsep=\headgap, itemsep=\itemgap, parsep=0pt, partopsep=0pt,
  label=\squarebullet}

% \entry{Organisation or title}{Role or detail (may be empty)}{Dates}; \sub{Role or detail}{Right-hand note}
\newcommand{\entry}[3]{%
  \par\iffirstentry\global\firstentryfalse\else\addvspace{\entrygap}\fi\noindent{\headfont\fontsize{10.8}{13}\selectfont\textbf{#1}}\ifblank{#2}{}{\ {\color{muted}|}\ #2}\hfill\hspace{1em}\side{#3}\par}
\newcommand{\sub}[2]{%
  \noindent\textit{#1}\hfill{\color{muted}\textit{#2}}\par}
\newcommand{\report}[1]{\href{#1}{\color{accent}Report}}
\newcommand{\me}{\textbf{S. Peddabomma}}

% ---------------------------------------------------------------------------------------------------
% Header
\newcommand{\Header}{%
  \begin{center}
    {\namefont\fontsize{26}{28}\selectfont Sudhansh Peddabomma}\\[6pt]
    {\headfont\fontsize{9.5}{11}\selectfont\color{muted}
      \href{mailto:sudhansh6@gmail.com}{\faEnvelope[regular]\ sudhansh6@gmail.com} \quad
      \href{https://sudhansh6.github.io}{\faGlobe\ sudhansh6.github.io} \quad
      \href{https://www.linkedin.com/in/sudhansh-peddabomma-73a208192/}{\faLinkedin\ sudhansh-peddabomma} \quad
      \href{https://github.com/Sudhansh6}{\faGithub\ Sudhansh6}}
  \end{center}\vspace{-4pt}}

% ---------------------------------------------------------------------------------------------------
% Glean bullets, shared by both profiles
\newcommand{\GleanSlides}{\item Led \textbf{AI slide generation} in Glean Assistant, now creating \textbf{180K decks a week} for \textbf{50K weekly active users} at enterprises}
\newcommand{\GleanSlidesPerf}{\item Improved the slides harness, raising reliability to \textbf{97\%}, cutting median latency \textbf{26\%} and cost per deck by \textbf{half}
  in four months}
\newcommand{\GleanSlidesOrch}{\item Architected a \textbf{slide orchestration system} that reuses relevant slides from a company's existing decks and generates
  new ones in its brand and template, so teams can refine decks with each other and Glean in full context, without leaving the product}
\newcommand{\GleanAgents}{\item Moved Glean's \textbf{workflow agents} from one-shot planning to an \textbf{agentic tool-use loop}, greatly improving agentic
  capabilities and instruction following, with \textbf{18\%} higher answer completeness and \textbf{17\%} fewer tokens}
\newcommand{\GleanTeamChat}{\item Engineered the core of \textbf{Team chat}, letting teammates work with Glean in one shared conversation and co-edit its artifacts}
\newcommand{\GleanVideo}{\item Shipped \textbf{video artifacts} end to end, generating narrated videos in chat with streamed playback and admin rollout controls}
\newcommand{\GleanEvals}{\item Created the \textbf{evaluations and metrics} that guide model choices and rollouts across agents, slides and interactive artifacts}

% ---------------------------------------------------------------------------------------------------
% Experience
% \groupbreak separates related groups of bullets inside one long entry (slides, then platform work).
\providecommand{\groupgap}{3.5pt}
\newcommand{\groupbreak}{\vspace{\groupgap}}
\newcommand{\GleanAI}{%
  \entry{Glean}{Software Engineer}{Apr 2025 -- Present}
  \begin{itemize}
    \GleanSlides
    \GleanSlidesOrch
    \GleanSlidesPerf \groupbreak
    \GleanTeamChat
    \GleanVideo
    \GleanAgents
    \GleanEvals
  \end{itemize}}

\newcommand{\GleanVision}{%
  \entry{Glean}{Software Engineer}{Apr 2025 -- Present}
  \begin{itemize}
    \GleanSlides
    \GleanSlidesOrch
    \GleanSlidesPerf \groupbreak
    \GleanVideo
    \GleanEvals
  \end{itemize}}

\newcommand{\Duality}{%
  \entry{Duality AI}{Computer Vision Intern}{Jun 2024 -- Sep 2024}
  \begin{itemize}
    \item Designed automated \textbf{3D reconstruction} techniques for featureless objects, reducing digital-twin generation time
          by \textbf{40\%}
    \item Collaborated with \textbf{Autodesk} to validate Unreal Engine simulations for robotics tasks; leveraging structured
          domain randomization to reduce Sim2Real gap and increase mAP-50 by \textbf{15\%} for object detection and segmentation
  \end{itemize}}

\newcommand{\Microsoft}{%
  \entry{Microsoft}{Data and Applied Scientist Intern}{May 2022 -- Jul 2022}
  \begin{itemize}
    \item Developed a \textbf{decision-tree ranker} to recommend emails without user queries, improving \textbf{Outlook} search capabilities
    \item Proposed hierarchical feature-sets for the ranker, reducing latency for recommendations and improving recall
  \end{itemize}}

% ---------------------------------------------------------------------------------------------------
% Research
\newcommand{\HomeRobots}{%
  \entry{Real-time 3D Perception for Home Robots}{UC San Diego}{Sep 2023 -- Sep 2024}
  \begin{itemize}
    \item Investigated real-time dense \textbf{visual SLAM} methods using NeRFs and Gaussian Splatting for indoor robot navigation
    \item Integrated object segmentation, grasp-pose estimation, and 3D mapping on the Fetch robot via ROS, demonstrating
          a novel \textbf{tabletop object rearrangement} algorithm that reduced cost by \textbf{20\%} compared to the state-of-the-art
  \end{itemize}}

\newcommand{\Tomography}{%
  \entry{3D Tomography with Primal-Dual Neural Networks}{University College London}{May 2021 -- Jul 2023}
  \begin{itemize}
    \item Developed a stochastic neural-network architecture of primal-dual algorithm for online reconstruction of 3D
          volumes from tomographic projections, with \textbf{99.6\,\% structural similarity} in low-dosage conditions
    \item Built a Python library with custom gradient operators for reconstructing volumes in a single pass, reducing
          compute requirements by up to \textbf{5x} over SOTA learning-based approaches for cone vector tomography
  \end{itemize}}

% ---------------------------------------------------------------------------------------------------
% Projects
\newcommand{\OOXML}{%
  \entry{ooxml}{Open-source contributor \textbar\ \href{https://github.com/ChristopherVR/ooxml}{\color{accent}GitHub}}{Sep 2026 -- Present}
  \begin{itemize}
    \item Contributed to an open-source \textbf{TypeScript} engine that renders and edits PowerPoint and Word files in the browser,
          improving editing fidelity, \textbf{real-time collaboration} and load performance across React, Vue and Svelte viewers
  \end{itemize}}

\newcommand{\MirrorAI}{%
  \entry{Mirror AI: Deployable Personas}{Honorable mention in Supabase YC Hackathon}{Oct 2024 -- Dec 2024}
  \begin{itemize}
    \item Designed an \textbf{agentic LLM architecture} with LangGraph to mirror user personalities, creating interactive digital
          personas, and deployed it as a full-stack platform on Supabase and Vercel with one-click publishing
  \end{itemize}}

\newcommand{\LLMReasoning}{%
  \entry{Improving LLM Reasoning for Numerical Problems}{\report{https://sudhansh6.github.io/assets/pdf/MathPrompter++.pdf}}{Sep 2024 -- Dec 2024}
  \begin{itemize}
    \item Enhanced MathPrompter (ACL 2023) with CoT, achieving \textbf{10\% higher accuracy} on Llama 3.1 1B where prior methods failed,
          and reduced hallucination rates significantly by integrating multi-step validation
  \end{itemize}}

\newcommand{\InverseRendering}{%
  \entry{Inverse Rendering with 2D Gaussian Splatting}{\report{https://sudhansh6.github.io/assets/pdf/CSE252D_Report.pdf}}{Mar 2024 -- May 2024}
  \begin{itemize}
    \item Developed a novel \textbf{inverse rendering} framework in CUDA to recover PBR properties of a scene using 2D Gaussian
          Splatting, improving normal map MAE by \textbf{15\%} over the current SOTA methods
  \end{itemize}}

% ---------------------------------------------------------------------------------------------------
% Publications: title on the left, venue on the right, one line each.
\newcommand{\publine}[3]{\par\noindent\href{#1}{#2}\hfill\hspace{1em}\side{#3}\par\addvspace{1.5pt}}
\newcommand{\PubSaturation}{\publine{https://doi.org/10.1016/j.sigpro.2023.109349}{A likelihood based method for compressive signal recovery under Gaussian and saturation noise}{Signal Processing, 2024}}
\newcommand{\PubPermutation}{\publine{https://ieeexplore.ieee.org/document/10888147}{Identification and correction of permutation errors in compressed sensing-based group testing}{ICASSP 2025}}
\newcommand{\PubRearrangement}{\publine{https://arxiv.org/abs/2411.10899}{Planning for tabletop object rearrangement}{ICRA 2025}}

% ---------------------------------------------------------------------------------------------------
% Education, skills, service
% Each profile sets \UCSDCourses, \IITBCourses and \SkillRows before \input{resume-common}.
\providecommand{\UCSDCourses}{}
\providecommand{\IITBCourses}{}
\providecommand{\SkillRows}{}
\newcommand{\coursework}[1]{\ifblank{#1}{}{\par\noindent{\fontsize{9.5}{11.5}\selectfont{\color{muted}Coursework:} #1}\par}}
\newcommand{\Education}{%
  \section{Education}
  \entry{University of California San Diego}{M.S. in Computer Science and Engineering (AI)}{GPA 4.00/4.00 \quad Sep 2023 -- Mar 2025}
  \coursework{\UCSDCourses}
  \entry{Indian Institute of Technology Bombay}{B.Tech.\ with Honors in Computer Science}{CPI 9.66/10 \quad Jul 2019 -- Jul 2023}
  \coursework{\IITBCourses}}

% \skill{Label}{Items} is one row of the skills table.
\newcommand{\skill}[2]{{\headfont\bfseries #1} & #2 \\[1.5pt]}
\newcommand{\cpp}{C\kern-0.04em\raisebox{0.06ex}{\small++}}
\newcommand{\Skills}{%
  \section{Skills}
  \begin{tabularx}{\linewidth}{@{}l@{\hspace{12pt}}>{\raggedright\arraybackslash}X@{}}
    \SkillRows
  \end{tabularx}}
 to try another pairing.
% Default: Charter body with Source Sans headings. All-sans alternative:
%   \def\resumefont{\usepackage[default,semibold]{sourcesanspro}}\def\headfont{}
\providecommand{\resumefont}{\usepackage[scaled=0.96]{XCharter}\usepackage[semibold,scale=0.96]{sourcesanspro}}
\providecommand{\headfont}{\sffamily}
\providecommand{\namefont}{\rmfamily\scshape}
\resumefont
\usepackage[final]{microtype}
\usepackage[dvipsnames]{xcolor}
\usepackage{enumitem}
\usepackage{titlesec}
\usepackage{tabularx}
\usepackage{fontawesome5}
\usepackage{etoolbox}
\usepackage[hidelinks]{hyperref}

% The site's green, darkened for print.
\definecolor{accent}{HTML}{3A7A45}
\definecolor{muted}{HTML}{5C675D}

\pagestyle{empty}
\setlength{\parindent}{0pt}
\setlength{\tabcolsep}{0pt}
\providecommand{\leading}{1.1}
\linespread{\leading}

\hypersetup{pdfauthor={Sudhansh Peddabomma}}

% Type scale: name 26pt small caps, section titles 12.5pt caps, entry titles 10.8pt, body 10pt,
% side notes 9pt.
\newcommand{\side}[1]{{\fontsize{9}{11}\selectfont\color{muted}#1}}
\newcommand{\sectitle}[1]{\textls[60]{\MakeUppercase{#1}}}
\titleformat{\section}{\headfont\color{accent}\bfseries\fontsize{12.5}{14}\selectfont}{}{0pt}{\sectitle}[\vspace{1.5pt}{\color{accent}\titlerule[0.8pt]}\global\firstentrytrue]
% The first entry after a section title skips \entrygap, so only \rulegap sits under the rule.
\newif\iffirstentry

% Vertical rhythm, the same idea as the earlier hand-tuned résumés: bullets inside an entry sit
% tight together, entries are separated by a clear gap, and sections by a larger one. A profile
% may override these before % Shared preamble and content for Sudhansh Peddabomma's résumés.
% Each profile (SudhanshPeddabomma-AI.tex, SudhanshPeddabomma-Vision.tex) inputs this file and picks
% which blocks to show and in what order, so every bullet is written once.
% Layout rule: every bullet fills one line, or runs 1.5 to 2 lines. Check with: python3 check_lines.py <pdf>
\documentclass[10pt,letterpaper]{article}

\usepackage[margin=0.55in,top=0.5in,bottom=0.45in]{geometry}
\usepackage[T1]{fontenc}
% Fonts: \resumefont loads the body font; \headfont styles the name and section titles.
% Override either before \input{resume-common} to try another pairing.
% Default: Charter body with Source Sans headings. All-sans alternative:
%   \def\resumefont{\usepackage[default,semibold]{sourcesanspro}}\def\headfont{}
\providecommand{\resumefont}{\usepackage[scaled=0.96]{XCharter}\usepackage[semibold,scale=0.96]{sourcesanspro}}
\providecommand{\headfont}{\sffamily}
\providecommand{\namefont}{\rmfamily\scshape}
\resumefont
\usepackage[final]{microtype}
\usepackage[dvipsnames]{xcolor}
\usepackage{enumitem}
\usepackage{titlesec}
\usepackage{tabularx}
\usepackage{fontawesome5}
\usepackage{etoolbox}
\usepackage[hidelinks]{hyperref}

% The site's green, darkened for print.
\definecolor{accent}{HTML}{3A7A45}
\definecolor{muted}{HTML}{5C675D}

\pagestyle{empty}
\setlength{\parindent}{0pt}
\setlength{\tabcolsep}{0pt}
\providecommand{\leading}{1.1}
\linespread{\leading}

\hypersetup{pdfauthor={Sudhansh Peddabomma}}

% Type scale: name 26pt small caps, section titles 12.5pt caps, entry titles 10.8pt, body 10pt,
% side notes 9pt.
\newcommand{\side}[1]{{\fontsize{9}{11}\selectfont\color{muted}#1}}
\newcommand{\sectitle}[1]{\textls[60]{\MakeUppercase{#1}}}
\titleformat{\section}{\headfont\color{accent}\bfseries\fontsize{12.5}{14}\selectfont}{}{0pt}{\sectitle}[\vspace{1.5pt}{\color{accent}\titlerule[0.8pt]}\global\firstentrytrue]
% The first entry after a section title skips \entrygap, so only \rulegap sits under the rule.
\newif\iffirstentry

% Vertical rhythm, the same idea as the earlier hand-tuned résumés: bullets inside an entry sit
% tight together, entries are separated by a clear gap, and sections by a larger one. A profile
% may override these before \input{resume-common} to fit one page.
\providecommand{\secgap}{12pt}   % above a section title
\providecommand{\rulegap}{5.5pt}   % section rule to first entry
\providecommand{\entrygap}{7pt}  % between entries
\providecommand{\headgap}{1.5pt} % entry heading to its bullets
\providecommand{\itemgap}{0.6pt} % between bullets of one entry
\titlespacing*{\section}{0pt}{\secgap}{\rulegap}

\newcommand{\squarebullet}{\raisebox{0.32ex}{\color{accent}\rule{0.5ex}{0.5ex}}}
\setlist[itemize]{leftmargin=12pt, labelsep=6pt, topsep=\headgap, itemsep=\itemgap, parsep=0pt, partopsep=0pt,
  label=\squarebullet}

% \entry{Organisation or title}{Role or detail (may be empty)}{Dates}; \sub{Role or detail}{Right-hand note}
\newcommand{\entry}[3]{%
  \par\iffirstentry\global\firstentryfalse\else\addvspace{\entrygap}\fi\noindent{\headfont\fontsize{10.8}{13}\selectfont\textbf{#1}}\ifblank{#2}{}{\ {\color{muted}|}\ #2}\hfill\hspace{1em}\side{#3}\par}
\newcommand{\sub}[2]{%
  \noindent\textit{#1}\hfill{\color{muted}\textit{#2}}\par}
\newcommand{\report}[1]{\href{#1}{\color{accent}Report}}
\newcommand{\me}{\textbf{S. Peddabomma}}

% ---------------------------------------------------------------------------------------------------
% Header
\newcommand{\Header}{%
  \begin{center}
    {\namefont\fontsize{26}{28}\selectfont Sudhansh Peddabomma}\\[6pt]
    {\headfont\fontsize{9.5}{11}\selectfont\color{muted}
      \href{mailto:sudhansh6@gmail.com}{\faEnvelope[regular]\ sudhansh6@gmail.com} \quad
      \href{https://sudhansh6.github.io}{\faGlobe\ sudhansh6.github.io} \quad
      \href{https://www.linkedin.com/in/sudhansh-peddabomma-73a208192/}{\faLinkedin\ sudhansh-peddabomma} \quad
      \href{https://github.com/Sudhansh6}{\faGithub\ Sudhansh6}}
  \end{center}\vspace{-4pt}}

% ---------------------------------------------------------------------------------------------------
% Glean bullets, shared by both profiles
\newcommand{\GleanSlides}{\item Led \textbf{AI slide generation} in Glean Assistant, now creating \textbf{180K decks a week} for \textbf{50K weekly active users} at enterprises}
\newcommand{\GleanSlidesPerf}{\item Improved the slides harness, raising reliability to \textbf{97\%}, cutting median latency \textbf{26\%} and cost per deck by \textbf{half}
  in four months}
\newcommand{\GleanSlidesOrch}{\item Architected a \textbf{slide orchestration system} that reuses relevant slides from a company's existing decks and generates
  new ones in its brand and template, so teams can refine decks with each other and Glean in full context, without leaving the product}
\newcommand{\GleanAgents}{\item Moved Glean's \textbf{workflow agents} from one-shot planning to an \textbf{agentic tool-use loop}, greatly improving agentic
  capabilities and instruction following, with \textbf{18\%} higher answer completeness and \textbf{17\%} fewer tokens}
\newcommand{\GleanTeamChat}{\item Engineered the core of \textbf{Team chat}, letting teammates work with Glean in one shared conversation and co-edit its artifacts}
\newcommand{\GleanVideo}{\item Shipped \textbf{video artifacts} end to end, generating narrated videos in chat with streamed playback and admin rollout controls}
\newcommand{\GleanEvals}{\item Created the \textbf{evaluations and metrics} that guide model choices and rollouts across agents, slides and interactive artifacts}

% ---------------------------------------------------------------------------------------------------
% Experience
% \groupbreak separates related groups of bullets inside one long entry (slides, then platform work).
\providecommand{\groupgap}{3.5pt}
\newcommand{\groupbreak}{\vspace{\groupgap}}
\newcommand{\GleanAI}{%
  \entry{Glean}{Software Engineer}{Apr 2025 -- Present}
  \begin{itemize}
    \GleanSlides
    \GleanSlidesOrch
    \GleanSlidesPerf \groupbreak
    \GleanTeamChat
    \GleanVideo
    \GleanAgents
    \GleanEvals
  \end{itemize}}

\newcommand{\GleanVision}{%
  \entry{Glean}{Software Engineer}{Apr 2025 -- Present}
  \begin{itemize}
    \GleanSlides
    \GleanSlidesOrch
    \GleanSlidesPerf \groupbreak
    \GleanVideo
    \GleanEvals
  \end{itemize}}

\newcommand{\Duality}{%
  \entry{Duality AI}{Computer Vision Intern}{Jun 2024 -- Sep 2024}
  \begin{itemize}
    \item Designed automated \textbf{3D reconstruction} techniques for featureless objects, reducing digital-twin generation time
          by \textbf{40\%}
    \item Collaborated with \textbf{Autodesk} to validate Unreal Engine simulations for robotics tasks; leveraging structured
          domain randomization to reduce Sim2Real gap and increase mAP-50 by \textbf{15\%} for object detection and segmentation
  \end{itemize}}

\newcommand{\Microsoft}{%
  \entry{Microsoft}{Data and Applied Scientist Intern}{May 2022 -- Jul 2022}
  \begin{itemize}
    \item Developed a \textbf{decision-tree ranker} to recommend emails without user queries, improving \textbf{Outlook} search capabilities
    \item Proposed hierarchical feature-sets for the ranker, reducing latency for recommendations and improving recall
  \end{itemize}}

% ---------------------------------------------------------------------------------------------------
% Research
\newcommand{\HomeRobots}{%
  \entry{Real-time 3D Perception for Home Robots}{UC San Diego}{Sep 2023 -- Sep 2024}
  \begin{itemize}
    \item Investigated real-time dense \textbf{visual SLAM} methods using NeRFs and Gaussian Splatting for indoor robot navigation
    \item Integrated object segmentation, grasp-pose estimation, and 3D mapping on the Fetch robot via ROS, demonstrating
          a novel \textbf{tabletop object rearrangement} algorithm that reduced cost by \textbf{20\%} compared to the state-of-the-art
  \end{itemize}}

\newcommand{\Tomography}{%
  \entry{3D Tomography with Primal-Dual Neural Networks}{University College London}{May 2021 -- Jul 2023}
  \begin{itemize}
    \item Developed a stochastic neural-network architecture of primal-dual algorithm for online reconstruction of 3D
          volumes from tomographic projections, with \textbf{99.6\,\% structural similarity} in low-dosage conditions
    \item Built a Python library with custom gradient operators for reconstructing volumes in a single pass, reducing
          compute requirements by up to \textbf{5x} over SOTA learning-based approaches for cone vector tomography
  \end{itemize}}

% ---------------------------------------------------------------------------------------------------
% Projects
\newcommand{\OOXML}{%
  \entry{ooxml}{Open-source contributor \textbar\ \href{https://github.com/ChristopherVR/ooxml}{\color{accent}GitHub}}{Sep 2026 -- Present}
  \begin{itemize}
    \item Contributed to an open-source \textbf{TypeScript} engine that renders and edits PowerPoint and Word files in the browser,
          improving editing fidelity, \textbf{real-time collaboration} and load performance across React, Vue and Svelte viewers
  \end{itemize}}

\newcommand{\MirrorAI}{%
  \entry{Mirror AI: Deployable Personas}{Honorable mention in Supabase YC Hackathon}{Oct 2024 -- Dec 2024}
  \begin{itemize}
    \item Designed an \textbf{agentic LLM architecture} with LangGraph to mirror user personalities, creating interactive digital
          personas, and deployed it as a full-stack platform on Supabase and Vercel with one-click publishing
  \end{itemize}}

\newcommand{\LLMReasoning}{%
  \entry{Improving LLM Reasoning for Numerical Problems}{\report{https://sudhansh6.github.io/assets/pdf/MathPrompter++.pdf}}{Sep 2024 -- Dec 2024}
  \begin{itemize}
    \item Enhanced MathPrompter (ACL 2023) with CoT, achieving \textbf{10\% higher accuracy} on Llama 3.1 1B where prior methods failed,
          and reduced hallucination rates significantly by integrating multi-step validation
  \end{itemize}}

\newcommand{\InverseRendering}{%
  \entry{Inverse Rendering with 2D Gaussian Splatting}{\report{https://sudhansh6.github.io/assets/pdf/CSE252D_Report.pdf}}{Mar 2024 -- May 2024}
  \begin{itemize}
    \item Developed a novel \textbf{inverse rendering} framework in CUDA to recover PBR properties of a scene using 2D Gaussian
          Splatting, improving normal map MAE by \textbf{15\%} over the current SOTA methods
  \end{itemize}}

% ---------------------------------------------------------------------------------------------------
% Publications: title on the left, venue on the right, one line each.
\newcommand{\publine}[3]{\par\noindent\href{#1}{#2}\hfill\hspace{1em}\side{#3}\par\addvspace{1.5pt}}
\newcommand{\PubSaturation}{\publine{https://doi.org/10.1016/j.sigpro.2023.109349}{A likelihood based method for compressive signal recovery under Gaussian and saturation noise}{Signal Processing, 2024}}
\newcommand{\PubPermutation}{\publine{https://ieeexplore.ieee.org/document/10888147}{Identification and correction of permutation errors in compressed sensing-based group testing}{ICASSP 2025}}
\newcommand{\PubRearrangement}{\publine{https://arxiv.org/abs/2411.10899}{Planning for tabletop object rearrangement}{ICRA 2025}}

% ---------------------------------------------------------------------------------------------------
% Education, skills, service
% Each profile sets \UCSDCourses, \IITBCourses and \SkillRows before \input{resume-common}.
\providecommand{\UCSDCourses}{}
\providecommand{\IITBCourses}{}
\providecommand{\SkillRows}{}
\newcommand{\coursework}[1]{\ifblank{#1}{}{\par\noindent{\fontsize{9.5}{11.5}\selectfont{\color{muted}Coursework:} #1}\par}}
\newcommand{\Education}{%
  \section{Education}
  \entry{University of California San Diego}{M.S. in Computer Science and Engineering (AI)}{GPA 4.00/4.00 \quad Sep 2023 -- Mar 2025}
  \coursework{\UCSDCourses}
  \entry{Indian Institute of Technology Bombay}{B.Tech.\ with Honors in Computer Science}{CPI 9.66/10 \quad Jul 2019 -- Jul 2023}
  \coursework{\IITBCourses}}

% \skill{Label}{Items} is one row of the skills table.
\newcommand{\skill}[2]{{\headfont\bfseries #1} & #2 \\[1.5pt]}
\newcommand{\cpp}{C\kern-0.04em\raisebox{0.06ex}{\small++}}
\newcommand{\Skills}{%
  \section{Skills}
  \begin{tabularx}{\linewidth}{@{}l@{\hspace{12pt}}>{\raggedright\arraybackslash}X@{}}
    \SkillRows
  \end{tabularx}}
 to fit one page.
\providecommand{\secgap}{12pt}   % above a section title
\providecommand{\rulegap}{5.5pt}   % section rule to first entry
\providecommand{\entrygap}{7pt}  % between entries
\providecommand{\headgap}{1.5pt} % entry heading to its bullets
\providecommand{\itemgap}{0.6pt} % between bullets of one entry
\titlespacing*{\section}{0pt}{\secgap}{\rulegap}

\newcommand{\squarebullet}{\raisebox{0.32ex}{\color{accent}\rule{0.5ex}{0.5ex}}}
\setlist[itemize]{leftmargin=12pt, labelsep=6pt, topsep=\headgap, itemsep=\itemgap, parsep=0pt, partopsep=0pt,
  label=\squarebullet}

% \entry{Organisation or title}{Role or detail (may be empty)}{Dates}; \sub{Role or detail}{Right-hand note}
\newcommand{\entry}[3]{%
  \par\iffirstentry\global\firstentryfalse\else\addvspace{\entrygap}\fi\noindent{\headfont\fontsize{10.8}{13}\selectfont\textbf{#1}}\ifblank{#2}{}{\ {\color{muted}|}\ #2}\hfill\hspace{1em}\side{#3}\par}
\newcommand{\sub}[2]{%
  \noindent\textit{#1}\hfill{\color{muted}\textit{#2}}\par}
\newcommand{\report}[1]{\href{#1}{\color{accent}Report}}
\newcommand{\me}{\textbf{S. Peddabomma}}

% ---------------------------------------------------------------------------------------------------
% Header
\newcommand{\Header}{%
  \begin{center}
    {\namefont\fontsize{26}{28}\selectfont Sudhansh Peddabomma}\\[6pt]
    {\headfont\fontsize{9.5}{11}\selectfont\color{muted}
      \href{mailto:sudhansh6@gmail.com}{\faEnvelope[regular]\ sudhansh6@gmail.com} \quad
      \href{https://sudhansh6.github.io}{\faGlobe\ sudhansh6.github.io} \quad
      \href{https://www.linkedin.com/in/sudhansh-peddabomma-73a208192/}{\faLinkedin\ sudhansh-peddabomma} \quad
      \href{https://github.com/Sudhansh6}{\faGithub\ Sudhansh6}}
  \end{center}\vspace{-4pt}}

% ---------------------------------------------------------------------------------------------------
% Glean bullets, shared by both profiles
\newcommand{\GleanSlides}{\item Led \textbf{AI slide generation} in Glean Assistant, now creating \textbf{180K decks a week} for \textbf{50K weekly active users} at enterprises}
\newcommand{\GleanSlidesPerf}{\item Improved the slides harness, raising reliability to \textbf{97\%}, cutting median latency \textbf{26\%} and cost per deck by \textbf{half}
  in four months}
\newcommand{\GleanSlidesOrch}{\item Architected a \textbf{slide orchestration system} that reuses relevant slides from a company's existing decks and generates
  new ones in its brand and template, so teams can refine decks with each other and Glean in full context, without leaving the product}
\newcommand{\GleanAgents}{\item Moved Glean's \textbf{workflow agents} from one-shot planning to an \textbf{agentic tool-use loop}, greatly improving agentic
  capabilities and instruction following, with \textbf{18\%} higher answer completeness and \textbf{17\%} fewer tokens}
\newcommand{\GleanTeamChat}{\item Engineered the core of \textbf{Team chat}, letting teammates work with Glean in one shared conversation and co-edit its artifacts}
\newcommand{\GleanVideo}{\item Shipped \textbf{video artifacts} end to end, generating narrated videos in chat with streamed playback and admin rollout controls}
\newcommand{\GleanEvals}{\item Created the \textbf{evaluations and metrics} that guide model choices and rollouts across agents, slides and interactive artifacts}

% ---------------------------------------------------------------------------------------------------
% Experience
% \groupbreak separates related groups of bullets inside one long entry (slides, then platform work).
\providecommand{\groupgap}{3.5pt}
\newcommand{\groupbreak}{\vspace{\groupgap}}
\newcommand{\GleanAI}{%
  \entry{Glean}{Software Engineer}{Apr 2025 -- Present}
  \begin{itemize}
    \GleanSlides
    \GleanSlidesOrch
    \GleanSlidesPerf \groupbreak
    \GleanTeamChat
    \GleanVideo
    \GleanAgents
    \GleanEvals
  \end{itemize}}

\newcommand{\GleanVision}{%
  \entry{Glean}{Software Engineer}{Apr 2025 -- Present}
  \begin{itemize}
    \GleanSlides
    \GleanSlidesOrch
    \GleanSlidesPerf \groupbreak
    \GleanVideo
    \GleanEvals
  \end{itemize}}

\newcommand{\Duality}{%
  \entry{Duality AI}{Computer Vision Intern}{Jun 2024 -- Sep 2024}
  \begin{itemize}
    \item Designed automated \textbf{3D reconstruction} techniques for featureless objects, reducing digital-twin generation time
          by \textbf{40\%}
    \item Collaborated with \textbf{Autodesk} to validate Unreal Engine simulations for robotics tasks; leveraging structured
          domain randomization to reduce Sim2Real gap and increase mAP-50 by \textbf{15\%} for object detection and segmentation
  \end{itemize}}

\newcommand{\Microsoft}{%
  \entry{Microsoft}{Data and Applied Scientist Intern}{May 2022 -- Jul 2022}
  \begin{itemize}
    \item Developed a \textbf{decision-tree ranker} to recommend emails without user queries, improving \textbf{Outlook} search capabilities
    \item Proposed hierarchical feature-sets for the ranker, reducing latency for recommendations and improving recall
  \end{itemize}}

% ---------------------------------------------------------------------------------------------------
% Research
\newcommand{\HomeRobots}{%
  \entry{Real-time 3D Perception for Home Robots}{UC San Diego}{Sep 2023 -- Sep 2024}
  \begin{itemize}
    \item Investigated real-time dense \textbf{visual SLAM} methods using NeRFs and Gaussian Splatting for indoor robot navigation
    \item Integrated object segmentation, grasp-pose estimation, and 3D mapping on the Fetch robot via ROS, demonstrating
          a novel \textbf{tabletop object rearrangement} algorithm that reduced cost by \textbf{20\%} compared to the state-of-the-art
  \end{itemize}}

\newcommand{\Tomography}{%
  \entry{3D Tomography with Primal-Dual Neural Networks}{University College London}{May 2021 -- Jul 2023}
  \begin{itemize}
    \item Developed a stochastic neural-network architecture of primal-dual algorithm for online reconstruction of 3D
          volumes from tomographic projections, with \textbf{99.6\,\% structural similarity} in low-dosage conditions
    \item Built a Python library with custom gradient operators for reconstructing volumes in a single pass, reducing
          compute requirements by up to \textbf{5x} over SOTA learning-based approaches for cone vector tomography
  \end{itemize}}

% ---------------------------------------------------------------------------------------------------
% Projects
\newcommand{\OOXML}{%
  \entry{ooxml}{Open-source contributor \textbar\ \href{https://github.com/ChristopherVR/ooxml}{\color{accent}GitHub}}{Sep 2026 -- Present}
  \begin{itemize}
    \item Contributed to an open-source \textbf{TypeScript} engine that renders and edits PowerPoint and Word files in the browser,
          improving editing fidelity, \textbf{real-time collaboration} and load performance across React, Vue and Svelte viewers
  \end{itemize}}

\newcommand{\MirrorAI}{%
  \entry{Mirror AI: Deployable Personas}{Honorable mention in Supabase YC Hackathon}{Oct 2024 -- Dec 2024}
  \begin{itemize}
    \item Designed an \textbf{agentic LLM architecture} with LangGraph to mirror user personalities, creating interactive digital
          personas, and deployed it as a full-stack platform on Supabase and Vercel with one-click publishing
  \end{itemize}}

\newcommand{\LLMReasoning}{%
  \entry{Improving LLM Reasoning for Numerical Problems}{\report{https://sudhansh6.github.io/assets/pdf/MathPrompter++.pdf}}{Sep 2024 -- Dec 2024}
  \begin{itemize}
    \item Enhanced MathPrompter (ACL 2023) with CoT, achieving \textbf{10\% higher accuracy} on Llama 3.1 1B where prior methods failed,
          and reduced hallucination rates significantly by integrating multi-step validation
  \end{itemize}}

\newcommand{\InverseRendering}{%
  \entry{Inverse Rendering with 2D Gaussian Splatting}{\report{https://sudhansh6.github.io/assets/pdf/CSE252D_Report.pdf}}{Mar 2024 -- May 2024}
  \begin{itemize}
    \item Developed a novel \textbf{inverse rendering} framework in CUDA to recover PBR properties of a scene using 2D Gaussian
          Splatting, improving normal map MAE by \textbf{15\%} over the current SOTA methods
  \end{itemize}}

% ---------------------------------------------------------------------------------------------------
% Publications: title on the left, venue on the right, one line each.
\newcommand{\publine}[3]{\par\noindent\href{#1}{#2}\hfill\hspace{1em}\side{#3}\par\addvspace{1.5pt}}
\newcommand{\PubSaturation}{\publine{https://doi.org/10.1016/j.sigpro.2023.109349}{A likelihood based method for compressive signal recovery under Gaussian and saturation noise}{Signal Processing, 2024}}
\newcommand{\PubPermutation}{\publine{https://ieeexplore.ieee.org/document/10888147}{Identification and correction of permutation errors in compressed sensing-based group testing}{ICASSP 2025}}
\newcommand{\PubRearrangement}{\publine{https://arxiv.org/abs/2411.10899}{Planning for tabletop object rearrangement}{ICRA 2025}}

% ---------------------------------------------------------------------------------------------------
% Education, skills, service
% Each profile sets \UCSDCourses, \IITBCourses and \SkillRows before % Shared preamble and content for Sudhansh Peddabomma's résumés.
% Each profile (SudhanshPeddabomma-AI.tex, SudhanshPeddabomma-Vision.tex) inputs this file and picks
% which blocks to show and in what order, so every bullet is written once.
% Layout rule: every bullet fills one line, or runs 1.5 to 2 lines. Check with: python3 check_lines.py <pdf>
\documentclass[10pt,letterpaper]{article}

\usepackage[margin=0.55in,top=0.5in,bottom=0.45in]{geometry}
\usepackage[T1]{fontenc}
% Fonts: \resumefont loads the body font; \headfont styles the name and section titles.
% Override either before \input{resume-common} to try another pairing.
% Default: Charter body with Source Sans headings. All-sans alternative:
%   \def\resumefont{\usepackage[default,semibold]{sourcesanspro}}\def\headfont{}
\providecommand{\resumefont}{\usepackage[scaled=0.96]{XCharter}\usepackage[semibold,scale=0.96]{sourcesanspro}}
\providecommand{\headfont}{\sffamily}
\providecommand{\namefont}{\rmfamily\scshape}
\resumefont
\usepackage[final]{microtype}
\usepackage[dvipsnames]{xcolor}
\usepackage{enumitem}
\usepackage{titlesec}
\usepackage{tabularx}
\usepackage{fontawesome5}
\usepackage{etoolbox}
\usepackage[hidelinks]{hyperref}

% The site's green, darkened for print.
\definecolor{accent}{HTML}{3A7A45}
\definecolor{muted}{HTML}{5C675D}

\pagestyle{empty}
\setlength{\parindent}{0pt}
\setlength{\tabcolsep}{0pt}
\providecommand{\leading}{1.1}
\linespread{\leading}

\hypersetup{pdfauthor={Sudhansh Peddabomma}}

% Type scale: name 26pt small caps, section titles 12.5pt caps, entry titles 10.8pt, body 10pt,
% side notes 9pt.
\newcommand{\side}[1]{{\fontsize{9}{11}\selectfont\color{muted}#1}}
\newcommand{\sectitle}[1]{\textls[60]{\MakeUppercase{#1}}}
\titleformat{\section}{\headfont\color{accent}\bfseries\fontsize{12.5}{14}\selectfont}{}{0pt}{\sectitle}[\vspace{1.5pt}{\color{accent}\titlerule[0.8pt]}\global\firstentrytrue]
% The first entry after a section title skips \entrygap, so only \rulegap sits under the rule.
\newif\iffirstentry

% Vertical rhythm, the same idea as the earlier hand-tuned résumés: bullets inside an entry sit
% tight together, entries are separated by a clear gap, and sections by a larger one. A profile
% may override these before \input{resume-common} to fit one page.
\providecommand{\secgap}{12pt}   % above a section title
\providecommand{\rulegap}{5.5pt}   % section rule to first entry
\providecommand{\entrygap}{7pt}  % between entries
\providecommand{\headgap}{1.5pt} % entry heading to its bullets
\providecommand{\itemgap}{0.6pt} % between bullets of one entry
\titlespacing*{\section}{0pt}{\secgap}{\rulegap}

\newcommand{\squarebullet}{\raisebox{0.32ex}{\color{accent}\rule{0.5ex}{0.5ex}}}
\setlist[itemize]{leftmargin=12pt, labelsep=6pt, topsep=\headgap, itemsep=\itemgap, parsep=0pt, partopsep=0pt,
  label=\squarebullet}

% \entry{Organisation or title}{Role or detail (may be empty)}{Dates}; \sub{Role or detail}{Right-hand note}
\newcommand{\entry}[3]{%
  \par\iffirstentry\global\firstentryfalse\else\addvspace{\entrygap}\fi\noindent{\headfont\fontsize{10.8}{13}\selectfont\textbf{#1}}\ifblank{#2}{}{\ {\color{muted}|}\ #2}\hfill\hspace{1em}\side{#3}\par}
\newcommand{\sub}[2]{%
  \noindent\textit{#1}\hfill{\color{muted}\textit{#2}}\par}
\newcommand{\report}[1]{\href{#1}{\color{accent}Report}}
\newcommand{\me}{\textbf{S. Peddabomma}}

% ---------------------------------------------------------------------------------------------------
% Header
\newcommand{\Header}{%
  \begin{center}
    {\namefont\fontsize{26}{28}\selectfont Sudhansh Peddabomma}\\[6pt]
    {\headfont\fontsize{9.5}{11}\selectfont\color{muted}
      \href{mailto:sudhansh6@gmail.com}{\faEnvelope[regular]\ sudhansh6@gmail.com} \quad
      \href{https://sudhansh6.github.io}{\faGlobe\ sudhansh6.github.io} \quad
      \href{https://www.linkedin.com/in/sudhansh-peddabomma-73a208192/}{\faLinkedin\ sudhansh-peddabomma} \quad
      \href{https://github.com/Sudhansh6}{\faGithub\ Sudhansh6}}
  \end{center}\vspace{-4pt}}

% ---------------------------------------------------------------------------------------------------
% Glean bullets, shared by both profiles
\newcommand{\GleanSlides}{\item Led \textbf{AI slide generation} in Glean Assistant, now creating \textbf{180K decks a week} for \textbf{50K weekly active users} at enterprises}
\newcommand{\GleanSlidesPerf}{\item Improved the slides harness, raising reliability to \textbf{97\%}, cutting median latency \textbf{26\%} and cost per deck by \textbf{half}
  in four months}
\newcommand{\GleanSlidesOrch}{\item Architected a \textbf{slide orchestration system} that reuses relevant slides from a company's existing decks and generates
  new ones in its brand and template, so teams can refine decks with each other and Glean in full context, without leaving the product}
\newcommand{\GleanAgents}{\item Moved Glean's \textbf{workflow agents} from one-shot planning to an \textbf{agentic tool-use loop}, greatly improving agentic
  capabilities and instruction following, with \textbf{18\%} higher answer completeness and \textbf{17\%} fewer tokens}
\newcommand{\GleanTeamChat}{\item Engineered the core of \textbf{Team chat}, letting teammates work with Glean in one shared conversation and co-edit its artifacts}
\newcommand{\GleanVideo}{\item Shipped \textbf{video artifacts} end to end, generating narrated videos in chat with streamed playback and admin rollout controls}
\newcommand{\GleanEvals}{\item Created the \textbf{evaluations and metrics} that guide model choices and rollouts across agents, slides and interactive artifacts}

% ---------------------------------------------------------------------------------------------------
% Experience
% \groupbreak separates related groups of bullets inside one long entry (slides, then platform work).
\providecommand{\groupgap}{3.5pt}
\newcommand{\groupbreak}{\vspace{\groupgap}}
\newcommand{\GleanAI}{%
  \entry{Glean}{Software Engineer}{Apr 2025 -- Present}
  \begin{itemize}
    \GleanSlides
    \GleanSlidesOrch
    \GleanSlidesPerf \groupbreak
    \GleanTeamChat
    \GleanVideo
    \GleanAgents
    \GleanEvals
  \end{itemize}}

\newcommand{\GleanVision}{%
  \entry{Glean}{Software Engineer}{Apr 2025 -- Present}
  \begin{itemize}
    \GleanSlides
    \GleanSlidesOrch
    \GleanSlidesPerf \groupbreak
    \GleanVideo
    \GleanEvals
  \end{itemize}}

\newcommand{\Duality}{%
  \entry{Duality AI}{Computer Vision Intern}{Jun 2024 -- Sep 2024}
  \begin{itemize}
    \item Designed automated \textbf{3D reconstruction} techniques for featureless objects, reducing digital-twin generation time
          by \textbf{40\%}
    \item Collaborated with \textbf{Autodesk} to validate Unreal Engine simulations for robotics tasks; leveraging structured
          domain randomization to reduce Sim2Real gap and increase mAP-50 by \textbf{15\%} for object detection and segmentation
  \end{itemize}}

\newcommand{\Microsoft}{%
  \entry{Microsoft}{Data and Applied Scientist Intern}{May 2022 -- Jul 2022}
  \begin{itemize}
    \item Developed a \textbf{decision-tree ranker} to recommend emails without user queries, improving \textbf{Outlook} search capabilities
    \item Proposed hierarchical feature-sets for the ranker, reducing latency for recommendations and improving recall
  \end{itemize}}

% ---------------------------------------------------------------------------------------------------
% Research
\newcommand{\HomeRobots}{%
  \entry{Real-time 3D Perception for Home Robots}{UC San Diego}{Sep 2023 -- Sep 2024}
  \begin{itemize}
    \item Investigated real-time dense \textbf{visual SLAM} methods using NeRFs and Gaussian Splatting for indoor robot navigation
    \item Integrated object segmentation, grasp-pose estimation, and 3D mapping on the Fetch robot via ROS, demonstrating
          a novel \textbf{tabletop object rearrangement} algorithm that reduced cost by \textbf{20\%} compared to the state-of-the-art
  \end{itemize}}

\newcommand{\Tomography}{%
  \entry{3D Tomography with Primal-Dual Neural Networks}{University College London}{May 2021 -- Jul 2023}
  \begin{itemize}
    \item Developed a stochastic neural-network architecture of primal-dual algorithm for online reconstruction of 3D
          volumes from tomographic projections, with \textbf{99.6\,\% structural similarity} in low-dosage conditions
    \item Built a Python library with custom gradient operators for reconstructing volumes in a single pass, reducing
          compute requirements by up to \textbf{5x} over SOTA learning-based approaches for cone vector tomography
  \end{itemize}}

% ---------------------------------------------------------------------------------------------------
% Projects
\newcommand{\OOXML}{%
  \entry{ooxml}{Open-source contributor \textbar\ \href{https://github.com/ChristopherVR/ooxml}{\color{accent}GitHub}}{Sep 2026 -- Present}
  \begin{itemize}
    \item Contributed to an open-source \textbf{TypeScript} engine that renders and edits PowerPoint and Word files in the browser,
          improving editing fidelity, \textbf{real-time collaboration} and load performance across React, Vue and Svelte viewers
  \end{itemize}}

\newcommand{\MirrorAI}{%
  \entry{Mirror AI: Deployable Personas}{Honorable mention in Supabase YC Hackathon}{Oct 2024 -- Dec 2024}
  \begin{itemize}
    \item Designed an \textbf{agentic LLM architecture} with LangGraph to mirror user personalities, creating interactive digital
          personas, and deployed it as a full-stack platform on Supabase and Vercel with one-click publishing
  \end{itemize}}

\newcommand{\LLMReasoning}{%
  \entry{Improving LLM Reasoning for Numerical Problems}{\report{https://sudhansh6.github.io/assets/pdf/MathPrompter++.pdf}}{Sep 2024 -- Dec 2024}
  \begin{itemize}
    \item Enhanced MathPrompter (ACL 2023) with CoT, achieving \textbf{10\% higher accuracy} on Llama 3.1 1B where prior methods failed,
          and reduced hallucination rates significantly by integrating multi-step validation
  \end{itemize}}

\newcommand{\InverseRendering}{%
  \entry{Inverse Rendering with 2D Gaussian Splatting}{\report{https://sudhansh6.github.io/assets/pdf/CSE252D_Report.pdf}}{Mar 2024 -- May 2024}
  \begin{itemize}
    \item Developed a novel \textbf{inverse rendering} framework in CUDA to recover PBR properties of a scene using 2D Gaussian
          Splatting, improving normal map MAE by \textbf{15\%} over the current SOTA methods
  \end{itemize}}

% ---------------------------------------------------------------------------------------------------
% Publications: title on the left, venue on the right, one line each.
\newcommand{\publine}[3]{\par\noindent\href{#1}{#2}\hfill\hspace{1em}\side{#3}\par\addvspace{1.5pt}}
\newcommand{\PubSaturation}{\publine{https://doi.org/10.1016/j.sigpro.2023.109349}{A likelihood based method for compressive signal recovery under Gaussian and saturation noise}{Signal Processing, 2024}}
\newcommand{\PubPermutation}{\publine{https://ieeexplore.ieee.org/document/10888147}{Identification and correction of permutation errors in compressed sensing-based group testing}{ICASSP 2025}}
\newcommand{\PubRearrangement}{\publine{https://arxiv.org/abs/2411.10899}{Planning for tabletop object rearrangement}{ICRA 2025}}

% ---------------------------------------------------------------------------------------------------
% Education, skills, service
% Each profile sets \UCSDCourses, \IITBCourses and \SkillRows before \input{resume-common}.
\providecommand{\UCSDCourses}{}
\providecommand{\IITBCourses}{}
\providecommand{\SkillRows}{}
\newcommand{\coursework}[1]{\ifblank{#1}{}{\par\noindent{\fontsize{9.5}{11.5}\selectfont{\color{muted}Coursework:} #1}\par}}
\newcommand{\Education}{%
  \section{Education}
  \entry{University of California San Diego}{M.S. in Computer Science and Engineering (AI)}{GPA 4.00/4.00 \quad Sep 2023 -- Mar 2025}
  \coursework{\UCSDCourses}
  \entry{Indian Institute of Technology Bombay}{B.Tech.\ with Honors in Computer Science}{CPI 9.66/10 \quad Jul 2019 -- Jul 2023}
  \coursework{\IITBCourses}}

% \skill{Label}{Items} is one row of the skills table.
\newcommand{\skill}[2]{{\headfont\bfseries #1} & #2 \\[1.5pt]}
\newcommand{\cpp}{C\kern-0.04em\raisebox{0.06ex}{\small++}}
\newcommand{\Skills}{%
  \section{Skills}
  \begin{tabularx}{\linewidth}{@{}l@{\hspace{12pt}}>{\raggedright\arraybackslash}X@{}}
    \SkillRows
  \end{tabularx}}
.
\providecommand{\UCSDCourses}{}
\providecommand{\IITBCourses}{}
\providecommand{\SkillRows}{}
\newcommand{\coursework}[1]{\ifblank{#1}{}{\par\noindent{\fontsize{9.5}{11.5}\selectfont{\color{muted}Coursework:} #1}\par}}
\newcommand{\Education}{%
  \section{Education}
  \entry{University of California San Diego}{M.S. in Computer Science and Engineering (AI)}{GPA 4.00/4.00 \quad Sep 2023 -- Mar 2025}
  \coursework{\UCSDCourses}
  \entry{Indian Institute of Technology Bombay}{B.Tech.\ with Honors in Computer Science}{CPI 9.66/10 \quad Jul 2019 -- Jul 2023}
  \coursework{\IITBCourses}}

% \skill{Label}{Items} is one row of the skills table.
\newcommand{\skill}[2]{{\headfont\bfseries #1} & #2 \\[1.5pt]}
\newcommand{\cpp}{C\kern-0.04em\raisebox{0.06ex}{\small++}}
\newcommand{\Skills}{%
  \section{Skills}
  \begin{tabularx}{\linewidth}{@{}l@{\hspace{12pt}}>{\raggedright\arraybackslash}X@{}}
    \SkillRows
  \end{tabularx}}
.
\providecommand{\UCSDCourses}{}
\providecommand{\IITBCourses}{}
\providecommand{\SkillRows}{}
\newcommand{\coursework}[1]{\ifblank{#1}{}{\par\noindent{\fontsize{9.5}{11.5}\selectfont{\color{muted}Coursework:} #1}\par}}
\newcommand{\Education}{%
  \section{Education}
  \entry{University of California San Diego}{M.S. in Computer Science and Engineering (AI)}{GPA 4.00/4.00 \quad Sep 2023 -- Mar 2025}
  \coursework{\UCSDCourses}
  \entry{Indian Institute of Technology Bombay}{B.Tech.\ with Honors in Computer Science}{CPI 9.66/10 \quad Jul 2019 -- Jul 2023}
  \coursework{\IITBCourses}}

% \skill{Label}{Items} is one row of the skills table.
\newcommand{\skill}[2]{{\headfont\bfseries #1} & #2 \\[1.5pt]}
\newcommand{\cpp}{C\kern-0.04em\raisebox{0.06ex}{\small++}}
\newcommand{\Skills}{%
  \section{Skills}
  \begin{tabularx}{\linewidth}{@{}l@{\hspace{12pt}}>{\raggedright\arraybackslash}X@{}}
    \SkillRows
  \end{tabularx}}
 to try another pairing.
% Default: Charter body with Source Sans headings. All-sans alternative:
%   \def\resumefont{\usepackage[default,semibold]{sourcesanspro}}\def\headfont{}
\providecommand{\resumefont}{\usepackage[scaled=0.96]{XCharter}\usepackage[semibold,scale=0.96]{sourcesanspro}}
\providecommand{\headfont}{\sffamily}
\providecommand{\namefont}{\rmfamily\scshape}
\resumefont
\usepackage[final]{microtype}
\usepackage[dvipsnames]{xcolor}
\usepackage{enumitem}
\usepackage{titlesec}
\usepackage{tabularx}
\usepackage{fontawesome5}
\usepackage{etoolbox}
\usepackage[hidelinks]{hyperref}

% The site's green, darkened for print.
\definecolor{accent}{HTML}{3A7A45}
\definecolor{muted}{HTML}{5C675D}

\pagestyle{empty}
\setlength{\parindent}{0pt}
\setlength{\tabcolsep}{0pt}
\providecommand{\leading}{1.1}
\linespread{\leading}

\hypersetup{pdfauthor={Sudhansh Peddabomma}}

% Type scale: name 26pt small caps, section titles 12.5pt caps, entry titles 10.8pt, body 10pt,
% side notes 9pt.
\newcommand{\side}[1]{{\fontsize{9}{11}\selectfont\color{muted}#1}}
\newcommand{\sectitle}[1]{\textls[60]{\MakeUppercase{#1}}}
\titleformat{\section}{\headfont\color{accent}\bfseries\fontsize{12.5}{14}\selectfont}{}{0pt}{\sectitle}[\vspace{1.5pt}{\color{accent}\titlerule[0.8pt]}\global\firstentrytrue]
% The first entry after a section title skips \entrygap, so only \rulegap sits under the rule.
\newif\iffirstentry

% Vertical rhythm, the same idea as the earlier hand-tuned résumés: bullets inside an entry sit
% tight together, entries are separated by a clear gap, and sections by a larger one. A profile
% may override these before % Shared preamble and content for Sudhansh Peddabomma's résumés.
% Each profile (SudhanshPeddabomma-AI.tex, SudhanshPeddabomma-Vision.tex) inputs this file and picks
% which blocks to show and in what order, so every bullet is written once.
% Layout rule: every bullet fills one line, or runs 1.5 to 2 lines. Check with: python3 check_lines.py <pdf>
\documentclass[10pt,letterpaper]{article}

\usepackage[margin=0.55in,top=0.5in,bottom=0.45in]{geometry}
\usepackage[T1]{fontenc}
% Fonts: \resumefont loads the body font; \headfont styles the name and section titles.
% Override either before % Shared preamble and content for Sudhansh Peddabomma's résumés.
% Each profile (SudhanshPeddabomma-AI.tex, SudhanshPeddabomma-Vision.tex) inputs this file and picks
% which blocks to show and in what order, so every bullet is written once.
% Layout rule: every bullet fills one line, or runs 1.5 to 2 lines. Check with: python3 check_lines.py <pdf>
\documentclass[10pt,letterpaper]{article}

\usepackage[margin=0.55in,top=0.5in,bottom=0.45in]{geometry}
\usepackage[T1]{fontenc}
% Fonts: \resumefont loads the body font; \headfont styles the name and section titles.
% Override either before % Shared preamble and content for Sudhansh Peddabomma's résumés.
% Each profile (SudhanshPeddabomma-AI.tex, SudhanshPeddabomma-Vision.tex) inputs this file and picks
% which blocks to show and in what order, so every bullet is written once.
% Layout rule: every bullet fills one line, or runs 1.5 to 2 lines. Check with: python3 check_lines.py <pdf>
\documentclass[10pt,letterpaper]{article}

\usepackage[margin=0.55in,top=0.5in,bottom=0.45in]{geometry}
\usepackage[T1]{fontenc}
% Fonts: \resumefont loads the body font; \headfont styles the name and section titles.
% Override either before \input{resume-common} to try another pairing.
% Default: Charter body with Source Sans headings. All-sans alternative:
%   \def\resumefont{\usepackage[default,semibold]{sourcesanspro}}\def\headfont{}
\providecommand{\resumefont}{\usepackage[scaled=0.96]{XCharter}\usepackage[semibold,scale=0.96]{sourcesanspro}}
\providecommand{\headfont}{\sffamily}
\providecommand{\namefont}{\rmfamily\scshape}
\resumefont
\usepackage[final]{microtype}
\usepackage[dvipsnames]{xcolor}
\usepackage{enumitem}
\usepackage{titlesec}
\usepackage{tabularx}
\usepackage{fontawesome5}
\usepackage{etoolbox}
\usepackage[hidelinks]{hyperref}

% The site's green, darkened for print.
\definecolor{accent}{HTML}{3A7A45}
\definecolor{muted}{HTML}{5C675D}

\pagestyle{empty}
\setlength{\parindent}{0pt}
\setlength{\tabcolsep}{0pt}
\providecommand{\leading}{1.1}
\linespread{\leading}

\hypersetup{pdfauthor={Sudhansh Peddabomma}}

% Type scale: name 26pt small caps, section titles 12.5pt caps, entry titles 10.8pt, body 10pt,
% side notes 9pt.
\newcommand{\side}[1]{{\fontsize{9}{11}\selectfont\color{muted}#1}}
\newcommand{\sectitle}[1]{\textls[60]{\MakeUppercase{#1}}}
\titleformat{\section}{\headfont\color{accent}\bfseries\fontsize{12.5}{14}\selectfont}{}{0pt}{\sectitle}[\vspace{1.5pt}{\color{accent}\titlerule[0.8pt]}\global\firstentrytrue]
% The first entry after a section title skips \entrygap, so only \rulegap sits under the rule.
\newif\iffirstentry

% Vertical rhythm, the same idea as the earlier hand-tuned résumés: bullets inside an entry sit
% tight together, entries are separated by a clear gap, and sections by a larger one. A profile
% may override these before \input{resume-common} to fit one page.
\providecommand{\secgap}{12pt}   % above a section title
\providecommand{\rulegap}{5.5pt}   % section rule to first entry
\providecommand{\entrygap}{7pt}  % between entries
\providecommand{\headgap}{1.5pt} % entry heading to its bullets
\providecommand{\itemgap}{0.6pt} % between bullets of one entry
\titlespacing*{\section}{0pt}{\secgap}{\rulegap}

\newcommand{\squarebullet}{\raisebox{0.32ex}{\color{accent}\rule{0.5ex}{0.5ex}}}
\setlist[itemize]{leftmargin=12pt, labelsep=6pt, topsep=\headgap, itemsep=\itemgap, parsep=0pt, partopsep=0pt,
  label=\squarebullet}

% \entry{Organisation or title}{Role or detail (may be empty)}{Dates}; \sub{Role or detail}{Right-hand note}
\newcommand{\entry}[3]{%
  \par\iffirstentry\global\firstentryfalse\else\addvspace{\entrygap}\fi\noindent{\headfont\fontsize{10.8}{13}\selectfont\textbf{#1}}\ifblank{#2}{}{\ {\color{muted}|}\ #2}\hfill\hspace{1em}\side{#3}\par}
\newcommand{\sub}[2]{%
  \noindent\textit{#1}\hfill{\color{muted}\textit{#2}}\par}
\newcommand{\report}[1]{\href{#1}{\color{accent}Report}}
\newcommand{\me}{\textbf{S. Peddabomma}}

% ---------------------------------------------------------------------------------------------------
% Header
\newcommand{\Header}{%
  \begin{center}
    {\namefont\fontsize{26}{28}\selectfont Sudhansh Peddabomma}\\[6pt]
    {\headfont\fontsize{9.5}{11}\selectfont\color{muted}
      \href{mailto:sudhansh6@gmail.com}{\faEnvelope[regular]\ sudhansh6@gmail.com} \quad
      \href{https://sudhansh6.github.io}{\faGlobe\ sudhansh6.github.io} \quad
      \href{https://www.linkedin.com/in/sudhansh-peddabomma-73a208192/}{\faLinkedin\ sudhansh-peddabomma} \quad
      \href{https://github.com/Sudhansh6}{\faGithub\ Sudhansh6}}
  \end{center}\vspace{-4pt}}

% ---------------------------------------------------------------------------------------------------
% Glean bullets, shared by both profiles
\newcommand{\GleanSlides}{\item Led \textbf{AI slide generation} in Glean Assistant, now creating \textbf{180K decks a week} for \textbf{50K weekly active users} at enterprises}
\newcommand{\GleanSlidesPerf}{\item Improved the slides harness, raising reliability to \textbf{97\%}, cutting median latency \textbf{26\%} and cost per deck by \textbf{half}
  in four months}
\newcommand{\GleanSlidesOrch}{\item Architected a \textbf{slide orchestration system} that reuses relevant slides from a company's existing decks and generates
  new ones in its brand and template, so teams can refine decks with each other and Glean in full context, without leaving the product}
\newcommand{\GleanAgents}{\item Moved Glean's \textbf{workflow agents} from one-shot planning to an \textbf{agentic tool-use loop}, greatly improving agentic
  capabilities and instruction following, with \textbf{18\%} higher answer completeness and \textbf{17\%} fewer tokens}
\newcommand{\GleanTeamChat}{\item Engineered the core of \textbf{Team chat}, letting teammates work with Glean in one shared conversation and co-edit its artifacts}
\newcommand{\GleanVideo}{\item Shipped \textbf{video artifacts} end to end, generating narrated videos in chat with streamed playback and admin rollout controls}
\newcommand{\GleanEvals}{\item Created the \textbf{evaluations and metrics} that guide model choices and rollouts across agents, slides and interactive artifacts}

% ---------------------------------------------------------------------------------------------------
% Experience
% \groupbreak separates related groups of bullets inside one long entry (slides, then platform work).
\providecommand{\groupgap}{3.5pt}
\newcommand{\groupbreak}{\vspace{\groupgap}}
\newcommand{\GleanAI}{%
  \entry{Glean}{Software Engineer}{Apr 2025 -- Present}
  \begin{itemize}
    \GleanSlides
    \GleanSlidesOrch
    \GleanSlidesPerf \groupbreak
    \GleanTeamChat
    \GleanVideo
    \GleanAgents
    \GleanEvals
  \end{itemize}}

\newcommand{\GleanVision}{%
  \entry{Glean}{Software Engineer}{Apr 2025 -- Present}
  \begin{itemize}
    \GleanSlides
    \GleanSlidesOrch
    \GleanSlidesPerf \groupbreak
    \GleanVideo
    \GleanEvals
  \end{itemize}}

\newcommand{\Duality}{%
  \entry{Duality AI}{Computer Vision Intern}{Jun 2024 -- Sep 2024}
  \begin{itemize}
    \item Designed automated \textbf{3D reconstruction} techniques for featureless objects, reducing digital-twin generation time
          by \textbf{40\%}
    \item Collaborated with \textbf{Autodesk} to validate Unreal Engine simulations for robotics tasks; leveraging structured
          domain randomization to reduce Sim2Real gap and increase mAP-50 by \textbf{15\%} for object detection and segmentation
  \end{itemize}}

\newcommand{\Microsoft}{%
  \entry{Microsoft}{Data and Applied Scientist Intern}{May 2022 -- Jul 2022}
  \begin{itemize}
    \item Developed a \textbf{decision-tree ranker} to recommend emails without user queries, improving \textbf{Outlook} search capabilities
    \item Proposed hierarchical feature-sets for the ranker, reducing latency for recommendations and improving recall
  \end{itemize}}

% ---------------------------------------------------------------------------------------------------
% Research
\newcommand{\HomeRobots}{%
  \entry{Real-time 3D Perception for Home Robots}{UC San Diego}{Sep 2023 -- Sep 2024}
  \begin{itemize}
    \item Investigated real-time dense \textbf{visual SLAM} methods using NeRFs and Gaussian Splatting for indoor robot navigation
    \item Integrated object segmentation, grasp-pose estimation, and 3D mapping on the Fetch robot via ROS, demonstrating
          a novel \textbf{tabletop object rearrangement} algorithm that reduced cost by \textbf{20\%} compared to the state-of-the-art
  \end{itemize}}

\newcommand{\Tomography}{%
  \entry{3D Tomography with Primal-Dual Neural Networks}{University College London}{May 2021 -- Jul 2023}
  \begin{itemize}
    \item Developed a stochastic neural-network architecture of primal-dual algorithm for online reconstruction of 3D
          volumes from tomographic projections, with \textbf{99.6\,\% structural similarity} in low-dosage conditions
    \item Built a Python library with custom gradient operators for reconstructing volumes in a single pass, reducing
          compute requirements by up to \textbf{5x} over SOTA learning-based approaches for cone vector tomography
  \end{itemize}}

% ---------------------------------------------------------------------------------------------------
% Projects
\newcommand{\OOXML}{%
  \entry{ooxml}{Open-source contributor \textbar\ \href{https://github.com/ChristopherVR/ooxml}{\color{accent}GitHub}}{Sep 2026 -- Present}
  \begin{itemize}
    \item Contributed to an open-source \textbf{TypeScript} engine that renders and edits PowerPoint and Word files in the browser,
          improving editing fidelity, \textbf{real-time collaboration} and load performance across React, Vue and Svelte viewers
  \end{itemize}}

\newcommand{\MirrorAI}{%
  \entry{Mirror AI: Deployable Personas}{Honorable mention in Supabase YC Hackathon}{Oct 2024 -- Dec 2024}
  \begin{itemize}
    \item Designed an \textbf{agentic LLM architecture} with LangGraph to mirror user personalities, creating interactive digital
          personas, and deployed it as a full-stack platform on Supabase and Vercel with one-click publishing
  \end{itemize}}

\newcommand{\LLMReasoning}{%
  \entry{Improving LLM Reasoning for Numerical Problems}{\report{https://sudhansh6.github.io/assets/pdf/MathPrompter++.pdf}}{Sep 2024 -- Dec 2024}
  \begin{itemize}
    \item Enhanced MathPrompter (ACL 2023) with CoT, achieving \textbf{10\% higher accuracy} on Llama 3.1 1B where prior methods failed,
          and reduced hallucination rates significantly by integrating multi-step validation
  \end{itemize}}

\newcommand{\InverseRendering}{%
  \entry{Inverse Rendering with 2D Gaussian Splatting}{\report{https://sudhansh6.github.io/assets/pdf/CSE252D_Report.pdf}}{Mar 2024 -- May 2024}
  \begin{itemize}
    \item Developed a novel \textbf{inverse rendering} framework in CUDA to recover PBR properties of a scene using 2D Gaussian
          Splatting, improving normal map MAE by \textbf{15\%} over the current SOTA methods
  \end{itemize}}

% ---------------------------------------------------------------------------------------------------
% Publications: title on the left, venue on the right, one line each.
\newcommand{\publine}[3]{\par\noindent\href{#1}{#2}\hfill\hspace{1em}\side{#3}\par\addvspace{1.5pt}}
\newcommand{\PubSaturation}{\publine{https://doi.org/10.1016/j.sigpro.2023.109349}{A likelihood based method for compressive signal recovery under Gaussian and saturation noise}{Signal Processing, 2024}}
\newcommand{\PubPermutation}{\publine{https://ieeexplore.ieee.org/document/10888147}{Identification and correction of permutation errors in compressed sensing-based group testing}{ICASSP 2025}}
\newcommand{\PubRearrangement}{\publine{https://arxiv.org/abs/2411.10899}{Planning for tabletop object rearrangement}{ICRA 2025}}

% ---------------------------------------------------------------------------------------------------
% Education, skills, service
% Each profile sets \UCSDCourses, \IITBCourses and \SkillRows before \input{resume-common}.
\providecommand{\UCSDCourses}{}
\providecommand{\IITBCourses}{}
\providecommand{\SkillRows}{}
\newcommand{\coursework}[1]{\ifblank{#1}{}{\par\noindent{\fontsize{9.5}{11.5}\selectfont{\color{muted}Coursework:} #1}\par}}
\newcommand{\Education}{%
  \section{Education}
  \entry{University of California San Diego}{M.S. in Computer Science and Engineering (AI)}{GPA 4.00/4.00 \quad Sep 2023 -- Mar 2025}
  \coursework{\UCSDCourses}
  \entry{Indian Institute of Technology Bombay}{B.Tech.\ with Honors in Computer Science}{CPI 9.66/10 \quad Jul 2019 -- Jul 2023}
  \coursework{\IITBCourses}}

% \skill{Label}{Items} is one row of the skills table.
\newcommand{\skill}[2]{{\headfont\bfseries #1} & #2 \\[1.5pt]}
\newcommand{\cpp}{C\kern-0.04em\raisebox{0.06ex}{\small++}}
\newcommand{\Skills}{%
  \section{Skills}
  \begin{tabularx}{\linewidth}{@{}l@{\hspace{12pt}}>{\raggedright\arraybackslash}X@{}}
    \SkillRows
  \end{tabularx}}
 to try another pairing.
% Default: Charter body with Source Sans headings. All-sans alternative:
%   \def\resumefont{\usepackage[default,semibold]{sourcesanspro}}\def\headfont{}
\providecommand{\resumefont}{\usepackage[scaled=0.96]{XCharter}\usepackage[semibold,scale=0.96]{sourcesanspro}}
\providecommand{\headfont}{\sffamily}
\providecommand{\namefont}{\rmfamily\scshape}
\resumefont
\usepackage[final]{microtype}
\usepackage[dvipsnames]{xcolor}
\usepackage{enumitem}
\usepackage{titlesec}
\usepackage{tabularx}
\usepackage{fontawesome5}
\usepackage{etoolbox}
\usepackage[hidelinks]{hyperref}

% The site's green, darkened for print.
\definecolor{accent}{HTML}{3A7A45}
\definecolor{muted}{HTML}{5C675D}

\pagestyle{empty}
\setlength{\parindent}{0pt}
\setlength{\tabcolsep}{0pt}
\providecommand{\leading}{1.1}
\linespread{\leading}

\hypersetup{pdfauthor={Sudhansh Peddabomma}}

% Type scale: name 26pt small caps, section titles 12.5pt caps, entry titles 10.8pt, body 10pt,
% side notes 9pt.
\newcommand{\side}[1]{{\fontsize{9}{11}\selectfont\color{muted}#1}}
\newcommand{\sectitle}[1]{\textls[60]{\MakeUppercase{#1}}}
\titleformat{\section}{\headfont\color{accent}\bfseries\fontsize{12.5}{14}\selectfont}{}{0pt}{\sectitle}[\vspace{1.5pt}{\color{accent}\titlerule[0.8pt]}\global\firstentrytrue]
% The first entry after a section title skips \entrygap, so only \rulegap sits under the rule.
\newif\iffirstentry

% Vertical rhythm, the same idea as the earlier hand-tuned résumés: bullets inside an entry sit
% tight together, entries are separated by a clear gap, and sections by a larger one. A profile
% may override these before % Shared preamble and content for Sudhansh Peddabomma's résumés.
% Each profile (SudhanshPeddabomma-AI.tex, SudhanshPeddabomma-Vision.tex) inputs this file and picks
% which blocks to show and in what order, so every bullet is written once.
% Layout rule: every bullet fills one line, or runs 1.5 to 2 lines. Check with: python3 check_lines.py <pdf>
\documentclass[10pt,letterpaper]{article}

\usepackage[margin=0.55in,top=0.5in,bottom=0.45in]{geometry}
\usepackage[T1]{fontenc}
% Fonts: \resumefont loads the body font; \headfont styles the name and section titles.
% Override either before \input{resume-common} to try another pairing.
% Default: Charter body with Source Sans headings. All-sans alternative:
%   \def\resumefont{\usepackage[default,semibold]{sourcesanspro}}\def\headfont{}
\providecommand{\resumefont}{\usepackage[scaled=0.96]{XCharter}\usepackage[semibold,scale=0.96]{sourcesanspro}}
\providecommand{\headfont}{\sffamily}
\providecommand{\namefont}{\rmfamily\scshape}
\resumefont
\usepackage[final]{microtype}
\usepackage[dvipsnames]{xcolor}
\usepackage{enumitem}
\usepackage{titlesec}
\usepackage{tabularx}
\usepackage{fontawesome5}
\usepackage{etoolbox}
\usepackage[hidelinks]{hyperref}

% The site's green, darkened for print.
\definecolor{accent}{HTML}{3A7A45}
\definecolor{muted}{HTML}{5C675D}

\pagestyle{empty}
\setlength{\parindent}{0pt}
\setlength{\tabcolsep}{0pt}
\providecommand{\leading}{1.1}
\linespread{\leading}

\hypersetup{pdfauthor={Sudhansh Peddabomma}}

% Type scale: name 26pt small caps, section titles 12.5pt caps, entry titles 10.8pt, body 10pt,
% side notes 9pt.
\newcommand{\side}[1]{{\fontsize{9}{11}\selectfont\color{muted}#1}}
\newcommand{\sectitle}[1]{\textls[60]{\MakeUppercase{#1}}}
\titleformat{\section}{\headfont\color{accent}\bfseries\fontsize{12.5}{14}\selectfont}{}{0pt}{\sectitle}[\vspace{1.5pt}{\color{accent}\titlerule[0.8pt]}\global\firstentrytrue]
% The first entry after a section title skips \entrygap, so only \rulegap sits under the rule.
\newif\iffirstentry

% Vertical rhythm, the same idea as the earlier hand-tuned résumés: bullets inside an entry sit
% tight together, entries are separated by a clear gap, and sections by a larger one. A profile
% may override these before \input{resume-common} to fit one page.
\providecommand{\secgap}{12pt}   % above a section title
\providecommand{\rulegap}{5.5pt}   % section rule to first entry
\providecommand{\entrygap}{7pt}  % between entries
\providecommand{\headgap}{1.5pt} % entry heading to its bullets
\providecommand{\itemgap}{0.6pt} % between bullets of one entry
\titlespacing*{\section}{0pt}{\secgap}{\rulegap}

\newcommand{\squarebullet}{\raisebox{0.32ex}{\color{accent}\rule{0.5ex}{0.5ex}}}
\setlist[itemize]{leftmargin=12pt, labelsep=6pt, topsep=\headgap, itemsep=\itemgap, parsep=0pt, partopsep=0pt,
  label=\squarebullet}

% \entry{Organisation or title}{Role or detail (may be empty)}{Dates}; \sub{Role or detail}{Right-hand note}
\newcommand{\entry}[3]{%
  \par\iffirstentry\global\firstentryfalse\else\addvspace{\entrygap}\fi\noindent{\headfont\fontsize{10.8}{13}\selectfont\textbf{#1}}\ifblank{#2}{}{\ {\color{muted}|}\ #2}\hfill\hspace{1em}\side{#3}\par}
\newcommand{\sub}[2]{%
  \noindent\textit{#1}\hfill{\color{muted}\textit{#2}}\par}
\newcommand{\report}[1]{\href{#1}{\color{accent}Report}}
\newcommand{\me}{\textbf{S. Peddabomma}}

% ---------------------------------------------------------------------------------------------------
% Header
\newcommand{\Header}{%
  \begin{center}
    {\namefont\fontsize{26}{28}\selectfont Sudhansh Peddabomma}\\[6pt]
    {\headfont\fontsize{9.5}{11}\selectfont\color{muted}
      \href{mailto:sudhansh6@gmail.com}{\faEnvelope[regular]\ sudhansh6@gmail.com} \quad
      \href{https://sudhansh6.github.io}{\faGlobe\ sudhansh6.github.io} \quad
      \href{https://www.linkedin.com/in/sudhansh-peddabomma-73a208192/}{\faLinkedin\ sudhansh-peddabomma} \quad
      \href{https://github.com/Sudhansh6}{\faGithub\ Sudhansh6}}
  \end{center}\vspace{-4pt}}

% ---------------------------------------------------------------------------------------------------
% Glean bullets, shared by both profiles
\newcommand{\GleanSlides}{\item Led \textbf{AI slide generation} in Glean Assistant, now creating \textbf{180K decks a week} for \textbf{50K weekly active users} at enterprises}
\newcommand{\GleanSlidesPerf}{\item Improved the slides harness, raising reliability to \textbf{97\%}, cutting median latency \textbf{26\%} and cost per deck by \textbf{half}
  in four months}
\newcommand{\GleanSlidesOrch}{\item Architected a \textbf{slide orchestration system} that reuses relevant slides from a company's existing decks and generates
  new ones in its brand and template, so teams can refine decks with each other and Glean in full context, without leaving the product}
\newcommand{\GleanAgents}{\item Moved Glean's \textbf{workflow agents} from one-shot planning to an \textbf{agentic tool-use loop}, greatly improving agentic
  capabilities and instruction following, with \textbf{18\%} higher answer completeness and \textbf{17\%} fewer tokens}
\newcommand{\GleanTeamChat}{\item Engineered the core of \textbf{Team chat}, letting teammates work with Glean in one shared conversation and co-edit its artifacts}
\newcommand{\GleanVideo}{\item Shipped \textbf{video artifacts} end to end, generating narrated videos in chat with streamed playback and admin rollout controls}
\newcommand{\GleanEvals}{\item Created the \textbf{evaluations and metrics} that guide model choices and rollouts across agents, slides and interactive artifacts}

% ---------------------------------------------------------------------------------------------------
% Experience
% \groupbreak separates related groups of bullets inside one long entry (slides, then platform work).
\providecommand{\groupgap}{3.5pt}
\newcommand{\groupbreak}{\vspace{\groupgap}}
\newcommand{\GleanAI}{%
  \entry{Glean}{Software Engineer}{Apr 2025 -- Present}
  \begin{itemize}
    \GleanSlides
    \GleanSlidesOrch
    \GleanSlidesPerf \groupbreak
    \GleanTeamChat
    \GleanVideo
    \GleanAgents
    \GleanEvals
  \end{itemize}}

\newcommand{\GleanVision}{%
  \entry{Glean}{Software Engineer}{Apr 2025 -- Present}
  \begin{itemize}
    \GleanSlides
    \GleanSlidesOrch
    \GleanSlidesPerf \groupbreak
    \GleanVideo
    \GleanEvals
  \end{itemize}}

\newcommand{\Duality}{%
  \entry{Duality AI}{Computer Vision Intern}{Jun 2024 -- Sep 2024}
  \begin{itemize}
    \item Designed automated \textbf{3D reconstruction} techniques for featureless objects, reducing digital-twin generation time
          by \textbf{40\%}
    \item Collaborated with \textbf{Autodesk} to validate Unreal Engine simulations for robotics tasks; leveraging structured
          domain randomization to reduce Sim2Real gap and increase mAP-50 by \textbf{15\%} for object detection and segmentation
  \end{itemize}}

\newcommand{\Microsoft}{%
  \entry{Microsoft}{Data and Applied Scientist Intern}{May 2022 -- Jul 2022}
  \begin{itemize}
    \item Developed a \textbf{decision-tree ranker} to recommend emails without user queries, improving \textbf{Outlook} search capabilities
    \item Proposed hierarchical feature-sets for the ranker, reducing latency for recommendations and improving recall
  \end{itemize}}

% ---------------------------------------------------------------------------------------------------
% Research
\newcommand{\HomeRobots}{%
  \entry{Real-time 3D Perception for Home Robots}{UC San Diego}{Sep 2023 -- Sep 2024}
  \begin{itemize}
    \item Investigated real-time dense \textbf{visual SLAM} methods using NeRFs and Gaussian Splatting for indoor robot navigation
    \item Integrated object segmentation, grasp-pose estimation, and 3D mapping on the Fetch robot via ROS, demonstrating
          a novel \textbf{tabletop object rearrangement} algorithm that reduced cost by \textbf{20\%} compared to the state-of-the-art
  \end{itemize}}

\newcommand{\Tomography}{%
  \entry{3D Tomography with Primal-Dual Neural Networks}{University College London}{May 2021 -- Jul 2023}
  \begin{itemize}
    \item Developed a stochastic neural-network architecture of primal-dual algorithm for online reconstruction of 3D
          volumes from tomographic projections, with \textbf{99.6\,\% structural similarity} in low-dosage conditions
    \item Built a Python library with custom gradient operators for reconstructing volumes in a single pass, reducing
          compute requirements by up to \textbf{5x} over SOTA learning-based approaches for cone vector tomography
  \end{itemize}}

% ---------------------------------------------------------------------------------------------------
% Projects
\newcommand{\OOXML}{%
  \entry{ooxml}{Open-source contributor \textbar\ \href{https://github.com/ChristopherVR/ooxml}{\color{accent}GitHub}}{Sep 2026 -- Present}
  \begin{itemize}
    \item Contributed to an open-source \textbf{TypeScript} engine that renders and edits PowerPoint and Word files in the browser,
          improving editing fidelity, \textbf{real-time collaboration} and load performance across React, Vue and Svelte viewers
  \end{itemize}}

\newcommand{\MirrorAI}{%
  \entry{Mirror AI: Deployable Personas}{Honorable mention in Supabase YC Hackathon}{Oct 2024 -- Dec 2024}
  \begin{itemize}
    \item Designed an \textbf{agentic LLM architecture} with LangGraph to mirror user personalities, creating interactive digital
          personas, and deployed it as a full-stack platform on Supabase and Vercel with one-click publishing
  \end{itemize}}

\newcommand{\LLMReasoning}{%
  \entry{Improving LLM Reasoning for Numerical Problems}{\report{https://sudhansh6.github.io/assets/pdf/MathPrompter++.pdf}}{Sep 2024 -- Dec 2024}
  \begin{itemize}
    \item Enhanced MathPrompter (ACL 2023) with CoT, achieving \textbf{10\% higher accuracy} on Llama 3.1 1B where prior methods failed,
          and reduced hallucination rates significantly by integrating multi-step validation
  \end{itemize}}

\newcommand{\InverseRendering}{%
  \entry{Inverse Rendering with 2D Gaussian Splatting}{\report{https://sudhansh6.github.io/assets/pdf/CSE252D_Report.pdf}}{Mar 2024 -- May 2024}
  \begin{itemize}
    \item Developed a novel \textbf{inverse rendering} framework in CUDA to recover PBR properties of a scene using 2D Gaussian
          Splatting, improving normal map MAE by \textbf{15\%} over the current SOTA methods
  \end{itemize}}

% ---------------------------------------------------------------------------------------------------
% Publications: title on the left, venue on the right, one line each.
\newcommand{\publine}[3]{\par\noindent\href{#1}{#2}\hfill\hspace{1em}\side{#3}\par\addvspace{1.5pt}}
\newcommand{\PubSaturation}{\publine{https://doi.org/10.1016/j.sigpro.2023.109349}{A likelihood based method for compressive signal recovery under Gaussian and saturation noise}{Signal Processing, 2024}}
\newcommand{\PubPermutation}{\publine{https://ieeexplore.ieee.org/document/10888147}{Identification and correction of permutation errors in compressed sensing-based group testing}{ICASSP 2025}}
\newcommand{\PubRearrangement}{\publine{https://arxiv.org/abs/2411.10899}{Planning for tabletop object rearrangement}{ICRA 2025}}

% ---------------------------------------------------------------------------------------------------
% Education, skills, service
% Each profile sets \UCSDCourses, \IITBCourses and \SkillRows before \input{resume-common}.
\providecommand{\UCSDCourses}{}
\providecommand{\IITBCourses}{}
\providecommand{\SkillRows}{}
\newcommand{\coursework}[1]{\ifblank{#1}{}{\par\noindent{\fontsize{9.5}{11.5}\selectfont{\color{muted}Coursework:} #1}\par}}
\newcommand{\Education}{%
  \section{Education}
  \entry{University of California San Diego}{M.S. in Computer Science and Engineering (AI)}{GPA 4.00/4.00 \quad Sep 2023 -- Mar 2025}
  \coursework{\UCSDCourses}
  \entry{Indian Institute of Technology Bombay}{B.Tech.\ with Honors in Computer Science}{CPI 9.66/10 \quad Jul 2019 -- Jul 2023}
  \coursework{\IITBCourses}}

% \skill{Label}{Items} is one row of the skills table.
\newcommand{\skill}[2]{{\headfont\bfseries #1} & #2 \\[1.5pt]}
\newcommand{\cpp}{C\kern-0.04em\raisebox{0.06ex}{\small++}}
\newcommand{\Skills}{%
  \section{Skills}
  \begin{tabularx}{\linewidth}{@{}l@{\hspace{12pt}}>{\raggedright\arraybackslash}X@{}}
    \SkillRows
  \end{tabularx}}
 to fit one page.
\providecommand{\secgap}{12pt}   % above a section title
\providecommand{\rulegap}{5.5pt}   % section rule to first entry
\providecommand{\entrygap}{7pt}  % between entries
\providecommand{\headgap}{1.5pt} % entry heading to its bullets
\providecommand{\itemgap}{0.6pt} % between bullets of one entry
\titlespacing*{\section}{0pt}{\secgap}{\rulegap}

\newcommand{\squarebullet}{\raisebox{0.32ex}{\color{accent}\rule{0.5ex}{0.5ex}}}
\setlist[itemize]{leftmargin=12pt, labelsep=6pt, topsep=\headgap, itemsep=\itemgap, parsep=0pt, partopsep=0pt,
  label=\squarebullet}

% \entry{Organisation or title}{Role or detail (may be empty)}{Dates}; \sub{Role or detail}{Right-hand note}
\newcommand{\entry}[3]{%
  \par\iffirstentry\global\firstentryfalse\else\addvspace{\entrygap}\fi\noindent{\headfont\fontsize{10.8}{13}\selectfont\textbf{#1}}\ifblank{#2}{}{\ {\color{muted}|}\ #2}\hfill\hspace{1em}\side{#3}\par}
\newcommand{\sub}[2]{%
  \noindent\textit{#1}\hfill{\color{muted}\textit{#2}}\par}
\newcommand{\report}[1]{\href{#1}{\color{accent}Report}}
\newcommand{\me}{\textbf{S. Peddabomma}}

% ---------------------------------------------------------------------------------------------------
% Header
\newcommand{\Header}{%
  \begin{center}
    {\namefont\fontsize{26}{28}\selectfont Sudhansh Peddabomma}\\[6pt]
    {\headfont\fontsize{9.5}{11}\selectfont\color{muted}
      \href{mailto:sudhansh6@gmail.com}{\faEnvelope[regular]\ sudhansh6@gmail.com} \quad
      \href{https://sudhansh6.github.io}{\faGlobe\ sudhansh6.github.io} \quad
      \href{https://www.linkedin.com/in/sudhansh-peddabomma-73a208192/}{\faLinkedin\ sudhansh-peddabomma} \quad
      \href{https://github.com/Sudhansh6}{\faGithub\ Sudhansh6}}
  \end{center}\vspace{-4pt}}

% ---------------------------------------------------------------------------------------------------
% Glean bullets, shared by both profiles
\newcommand{\GleanSlides}{\item Led \textbf{AI slide generation} in Glean Assistant, now creating \textbf{180K decks a week} for \textbf{50K weekly active users} at enterprises}
\newcommand{\GleanSlidesPerf}{\item Improved the slides harness, raising reliability to \textbf{97\%}, cutting median latency \textbf{26\%} and cost per deck by \textbf{half}
  in four months}
\newcommand{\GleanSlidesOrch}{\item Architected a \textbf{slide orchestration system} that reuses relevant slides from a company's existing decks and generates
  new ones in its brand and template, so teams can refine decks with each other and Glean in full context, without leaving the product}
\newcommand{\GleanAgents}{\item Moved Glean's \textbf{workflow agents} from one-shot planning to an \textbf{agentic tool-use loop}, greatly improving agentic
  capabilities and instruction following, with \textbf{18\%} higher answer completeness and \textbf{17\%} fewer tokens}
\newcommand{\GleanTeamChat}{\item Engineered the core of \textbf{Team chat}, letting teammates work with Glean in one shared conversation and co-edit its artifacts}
\newcommand{\GleanVideo}{\item Shipped \textbf{video artifacts} end to end, generating narrated videos in chat with streamed playback and admin rollout controls}
\newcommand{\GleanEvals}{\item Created the \textbf{evaluations and metrics} that guide model choices and rollouts across agents, slides and interactive artifacts}

% ---------------------------------------------------------------------------------------------------
% Experience
% \groupbreak separates related groups of bullets inside one long entry (slides, then platform work).
\providecommand{\groupgap}{3.5pt}
\newcommand{\groupbreak}{\vspace{\groupgap}}
\newcommand{\GleanAI}{%
  \entry{Glean}{Software Engineer}{Apr 2025 -- Present}
  \begin{itemize}
    \GleanSlides
    \GleanSlidesOrch
    \GleanSlidesPerf \groupbreak
    \GleanTeamChat
    \GleanVideo
    \GleanAgents
    \GleanEvals
  \end{itemize}}

\newcommand{\GleanVision}{%
  \entry{Glean}{Software Engineer}{Apr 2025 -- Present}
  \begin{itemize}
    \GleanSlides
    \GleanSlidesOrch
    \GleanSlidesPerf \groupbreak
    \GleanVideo
    \GleanEvals
  \end{itemize}}

\newcommand{\Duality}{%
  \entry{Duality AI}{Computer Vision Intern}{Jun 2024 -- Sep 2024}
  \begin{itemize}
    \item Designed automated \textbf{3D reconstruction} techniques for featureless objects, reducing digital-twin generation time
          by \textbf{40\%}
    \item Collaborated with \textbf{Autodesk} to validate Unreal Engine simulations for robotics tasks; leveraging structured
          domain randomization to reduce Sim2Real gap and increase mAP-50 by \textbf{15\%} for object detection and segmentation
  \end{itemize}}

\newcommand{\Microsoft}{%
  \entry{Microsoft}{Data and Applied Scientist Intern}{May 2022 -- Jul 2022}
  \begin{itemize}
    \item Developed a \textbf{decision-tree ranker} to recommend emails without user queries, improving \textbf{Outlook} search capabilities
    \item Proposed hierarchical feature-sets for the ranker, reducing latency for recommendations and improving recall
  \end{itemize}}

% ---------------------------------------------------------------------------------------------------
% Research
\newcommand{\HomeRobots}{%
  \entry{Real-time 3D Perception for Home Robots}{UC San Diego}{Sep 2023 -- Sep 2024}
  \begin{itemize}
    \item Investigated real-time dense \textbf{visual SLAM} methods using NeRFs and Gaussian Splatting for indoor robot navigation
    \item Integrated object segmentation, grasp-pose estimation, and 3D mapping on the Fetch robot via ROS, demonstrating
          a novel \textbf{tabletop object rearrangement} algorithm that reduced cost by \textbf{20\%} compared to the state-of-the-art
  \end{itemize}}

\newcommand{\Tomography}{%
  \entry{3D Tomography with Primal-Dual Neural Networks}{University College London}{May 2021 -- Jul 2023}
  \begin{itemize}
    \item Developed a stochastic neural-network architecture of primal-dual algorithm for online reconstruction of 3D
          volumes from tomographic projections, with \textbf{99.6\,\% structural similarity} in low-dosage conditions
    \item Built a Python library with custom gradient operators for reconstructing volumes in a single pass, reducing
          compute requirements by up to \textbf{5x} over SOTA learning-based approaches for cone vector tomography
  \end{itemize}}

% ---------------------------------------------------------------------------------------------------
% Projects
\newcommand{\OOXML}{%
  \entry{ooxml}{Open-source contributor \textbar\ \href{https://github.com/ChristopherVR/ooxml}{\color{accent}GitHub}}{Sep 2026 -- Present}
  \begin{itemize}
    \item Contributed to an open-source \textbf{TypeScript} engine that renders and edits PowerPoint and Word files in the browser,
          improving editing fidelity, \textbf{real-time collaboration} and load performance across React, Vue and Svelte viewers
  \end{itemize}}

\newcommand{\MirrorAI}{%
  \entry{Mirror AI: Deployable Personas}{Honorable mention in Supabase YC Hackathon}{Oct 2024 -- Dec 2024}
  \begin{itemize}
    \item Designed an \textbf{agentic LLM architecture} with LangGraph to mirror user personalities, creating interactive digital
          personas, and deployed it as a full-stack platform on Supabase and Vercel with one-click publishing
  \end{itemize}}

\newcommand{\LLMReasoning}{%
  \entry{Improving LLM Reasoning for Numerical Problems}{\report{https://sudhansh6.github.io/assets/pdf/MathPrompter++.pdf}}{Sep 2024 -- Dec 2024}
  \begin{itemize}
    \item Enhanced MathPrompter (ACL 2023) with CoT, achieving \textbf{10\% higher accuracy} on Llama 3.1 1B where prior methods failed,
          and reduced hallucination rates significantly by integrating multi-step validation
  \end{itemize}}

\newcommand{\InverseRendering}{%
  \entry{Inverse Rendering with 2D Gaussian Splatting}{\report{https://sudhansh6.github.io/assets/pdf/CSE252D_Report.pdf}}{Mar 2024 -- May 2024}
  \begin{itemize}
    \item Developed a novel \textbf{inverse rendering} framework in CUDA to recover PBR properties of a scene using 2D Gaussian
          Splatting, improving normal map MAE by \textbf{15\%} over the current SOTA methods
  \end{itemize}}

% ---------------------------------------------------------------------------------------------------
% Publications: title on the left, venue on the right, one line each.
\newcommand{\publine}[3]{\par\noindent\href{#1}{#2}\hfill\hspace{1em}\side{#3}\par\addvspace{1.5pt}}
\newcommand{\PubSaturation}{\publine{https://doi.org/10.1016/j.sigpro.2023.109349}{A likelihood based method for compressive signal recovery under Gaussian and saturation noise}{Signal Processing, 2024}}
\newcommand{\PubPermutation}{\publine{https://ieeexplore.ieee.org/document/10888147}{Identification and correction of permutation errors in compressed sensing-based group testing}{ICASSP 2025}}
\newcommand{\PubRearrangement}{\publine{https://arxiv.org/abs/2411.10899}{Planning for tabletop object rearrangement}{ICRA 2025}}

% ---------------------------------------------------------------------------------------------------
% Education, skills, service
% Each profile sets \UCSDCourses, \IITBCourses and \SkillRows before % Shared preamble and content for Sudhansh Peddabomma's résumés.
% Each profile (SudhanshPeddabomma-AI.tex, SudhanshPeddabomma-Vision.tex) inputs this file and picks
% which blocks to show and in what order, so every bullet is written once.
% Layout rule: every bullet fills one line, or runs 1.5 to 2 lines. Check with: python3 check_lines.py <pdf>
\documentclass[10pt,letterpaper]{article}

\usepackage[margin=0.55in,top=0.5in,bottom=0.45in]{geometry}
\usepackage[T1]{fontenc}
% Fonts: \resumefont loads the body font; \headfont styles the name and section titles.
% Override either before \input{resume-common} to try another pairing.
% Default: Charter body with Source Sans headings. All-sans alternative:
%   \def\resumefont{\usepackage[default,semibold]{sourcesanspro}}\def\headfont{}
\providecommand{\resumefont}{\usepackage[scaled=0.96]{XCharter}\usepackage[semibold,scale=0.96]{sourcesanspro}}
\providecommand{\headfont}{\sffamily}
\providecommand{\namefont}{\rmfamily\scshape}
\resumefont
\usepackage[final]{microtype}
\usepackage[dvipsnames]{xcolor}
\usepackage{enumitem}
\usepackage{titlesec}
\usepackage{tabularx}
\usepackage{fontawesome5}
\usepackage{etoolbox}
\usepackage[hidelinks]{hyperref}

% The site's green, darkened for print.
\definecolor{accent}{HTML}{3A7A45}
\definecolor{muted}{HTML}{5C675D}

\pagestyle{empty}
\setlength{\parindent}{0pt}
\setlength{\tabcolsep}{0pt}
\providecommand{\leading}{1.1}
\linespread{\leading}

\hypersetup{pdfauthor={Sudhansh Peddabomma}}

% Type scale: name 26pt small caps, section titles 12.5pt caps, entry titles 10.8pt, body 10pt,
% side notes 9pt.
\newcommand{\side}[1]{{\fontsize{9}{11}\selectfont\color{muted}#1}}
\newcommand{\sectitle}[1]{\textls[60]{\MakeUppercase{#1}}}
\titleformat{\section}{\headfont\color{accent}\bfseries\fontsize{12.5}{14}\selectfont}{}{0pt}{\sectitle}[\vspace{1.5pt}{\color{accent}\titlerule[0.8pt]}\global\firstentrytrue]
% The first entry after a section title skips \entrygap, so only \rulegap sits under the rule.
\newif\iffirstentry

% Vertical rhythm, the same idea as the earlier hand-tuned résumés: bullets inside an entry sit
% tight together, entries are separated by a clear gap, and sections by a larger one. A profile
% may override these before \input{resume-common} to fit one page.
\providecommand{\secgap}{12pt}   % above a section title
\providecommand{\rulegap}{5.5pt}   % section rule to first entry
\providecommand{\entrygap}{7pt}  % between entries
\providecommand{\headgap}{1.5pt} % entry heading to its bullets
\providecommand{\itemgap}{0.6pt} % between bullets of one entry
\titlespacing*{\section}{0pt}{\secgap}{\rulegap}

\newcommand{\squarebullet}{\raisebox{0.32ex}{\color{accent}\rule{0.5ex}{0.5ex}}}
\setlist[itemize]{leftmargin=12pt, labelsep=6pt, topsep=\headgap, itemsep=\itemgap, parsep=0pt, partopsep=0pt,
  label=\squarebullet}

% \entry{Organisation or title}{Role or detail (may be empty)}{Dates}; \sub{Role or detail}{Right-hand note}
\newcommand{\entry}[3]{%
  \par\iffirstentry\global\firstentryfalse\else\addvspace{\entrygap}\fi\noindent{\headfont\fontsize{10.8}{13}\selectfont\textbf{#1}}\ifblank{#2}{}{\ {\color{muted}|}\ #2}\hfill\hspace{1em}\side{#3}\par}
\newcommand{\sub}[2]{%
  \noindent\textit{#1}\hfill{\color{muted}\textit{#2}}\par}
\newcommand{\report}[1]{\href{#1}{\color{accent}Report}}
\newcommand{\me}{\textbf{S. Peddabomma}}

% ---------------------------------------------------------------------------------------------------
% Header
\newcommand{\Header}{%
  \begin{center}
    {\namefont\fontsize{26}{28}\selectfont Sudhansh Peddabomma}\\[6pt]
    {\headfont\fontsize{9.5}{11}\selectfont\color{muted}
      \href{mailto:sudhansh6@gmail.com}{\faEnvelope[regular]\ sudhansh6@gmail.com} \quad
      \href{https://sudhansh6.github.io}{\faGlobe\ sudhansh6.github.io} \quad
      \href{https://www.linkedin.com/in/sudhansh-peddabomma-73a208192/}{\faLinkedin\ sudhansh-peddabomma} \quad
      \href{https://github.com/Sudhansh6}{\faGithub\ Sudhansh6}}
  \end{center}\vspace{-4pt}}

% ---------------------------------------------------------------------------------------------------
% Glean bullets, shared by both profiles
\newcommand{\GleanSlides}{\item Led \textbf{AI slide generation} in Glean Assistant, now creating \textbf{180K decks a week} for \textbf{50K weekly active users} at enterprises}
\newcommand{\GleanSlidesPerf}{\item Improved the slides harness, raising reliability to \textbf{97\%}, cutting median latency \textbf{26\%} and cost per deck by \textbf{half}
  in four months}
\newcommand{\GleanSlidesOrch}{\item Architected a \textbf{slide orchestration system} that reuses relevant slides from a company's existing decks and generates
  new ones in its brand and template, so teams can refine decks with each other and Glean in full context, without leaving the product}
\newcommand{\GleanAgents}{\item Moved Glean's \textbf{workflow agents} from one-shot planning to an \textbf{agentic tool-use loop}, greatly improving agentic
  capabilities and instruction following, with \textbf{18\%} higher answer completeness and \textbf{17\%} fewer tokens}
\newcommand{\GleanTeamChat}{\item Engineered the core of \textbf{Team chat}, letting teammates work with Glean in one shared conversation and co-edit its artifacts}
\newcommand{\GleanVideo}{\item Shipped \textbf{video artifacts} end to end, generating narrated videos in chat with streamed playback and admin rollout controls}
\newcommand{\GleanEvals}{\item Created the \textbf{evaluations and metrics} that guide model choices and rollouts across agents, slides and interactive artifacts}

% ---------------------------------------------------------------------------------------------------
% Experience
% \groupbreak separates related groups of bullets inside one long entry (slides, then platform work).
\providecommand{\groupgap}{3.5pt}
\newcommand{\groupbreak}{\vspace{\groupgap}}
\newcommand{\GleanAI}{%
  \entry{Glean}{Software Engineer}{Apr 2025 -- Present}
  \begin{itemize}
    \GleanSlides
    \GleanSlidesOrch
    \GleanSlidesPerf \groupbreak
    \GleanTeamChat
    \GleanVideo
    \GleanAgents
    \GleanEvals
  \end{itemize}}

\newcommand{\GleanVision}{%
  \entry{Glean}{Software Engineer}{Apr 2025 -- Present}
  \begin{itemize}
    \GleanSlides
    \GleanSlidesOrch
    \GleanSlidesPerf \groupbreak
    \GleanVideo
    \GleanEvals
  \end{itemize}}

\newcommand{\Duality}{%
  \entry{Duality AI}{Computer Vision Intern}{Jun 2024 -- Sep 2024}
  \begin{itemize}
    \item Designed automated \textbf{3D reconstruction} techniques for featureless objects, reducing digital-twin generation time
          by \textbf{40\%}
    \item Collaborated with \textbf{Autodesk} to validate Unreal Engine simulations for robotics tasks; leveraging structured
          domain randomization to reduce Sim2Real gap and increase mAP-50 by \textbf{15\%} for object detection and segmentation
  \end{itemize}}

\newcommand{\Microsoft}{%
  \entry{Microsoft}{Data and Applied Scientist Intern}{May 2022 -- Jul 2022}
  \begin{itemize}
    \item Developed a \textbf{decision-tree ranker} to recommend emails without user queries, improving \textbf{Outlook} search capabilities
    \item Proposed hierarchical feature-sets for the ranker, reducing latency for recommendations and improving recall
  \end{itemize}}

% ---------------------------------------------------------------------------------------------------
% Research
\newcommand{\HomeRobots}{%
  \entry{Real-time 3D Perception for Home Robots}{UC San Diego}{Sep 2023 -- Sep 2024}
  \begin{itemize}
    \item Investigated real-time dense \textbf{visual SLAM} methods using NeRFs and Gaussian Splatting for indoor robot navigation
    \item Integrated object segmentation, grasp-pose estimation, and 3D mapping on the Fetch robot via ROS, demonstrating
          a novel \textbf{tabletop object rearrangement} algorithm that reduced cost by \textbf{20\%} compared to the state-of-the-art
  \end{itemize}}

\newcommand{\Tomography}{%
  \entry{3D Tomography with Primal-Dual Neural Networks}{University College London}{May 2021 -- Jul 2023}
  \begin{itemize}
    \item Developed a stochastic neural-network architecture of primal-dual algorithm for online reconstruction of 3D
          volumes from tomographic projections, with \textbf{99.6\,\% structural similarity} in low-dosage conditions
    \item Built a Python library with custom gradient operators for reconstructing volumes in a single pass, reducing
          compute requirements by up to \textbf{5x} over SOTA learning-based approaches for cone vector tomography
  \end{itemize}}

% ---------------------------------------------------------------------------------------------------
% Projects
\newcommand{\OOXML}{%
  \entry{ooxml}{Open-source contributor \textbar\ \href{https://github.com/ChristopherVR/ooxml}{\color{accent}GitHub}}{Sep 2026 -- Present}
  \begin{itemize}
    \item Contributed to an open-source \textbf{TypeScript} engine that renders and edits PowerPoint and Word files in the browser,
          improving editing fidelity, \textbf{real-time collaboration} and load performance across React, Vue and Svelte viewers
  \end{itemize}}

\newcommand{\MirrorAI}{%
  \entry{Mirror AI: Deployable Personas}{Honorable mention in Supabase YC Hackathon}{Oct 2024 -- Dec 2024}
  \begin{itemize}
    \item Designed an \textbf{agentic LLM architecture} with LangGraph to mirror user personalities, creating interactive digital
          personas, and deployed it as a full-stack platform on Supabase and Vercel with one-click publishing
  \end{itemize}}

\newcommand{\LLMReasoning}{%
  \entry{Improving LLM Reasoning for Numerical Problems}{\report{https://sudhansh6.github.io/assets/pdf/MathPrompter++.pdf}}{Sep 2024 -- Dec 2024}
  \begin{itemize}
    \item Enhanced MathPrompter (ACL 2023) with CoT, achieving \textbf{10\% higher accuracy} on Llama 3.1 1B where prior methods failed,
          and reduced hallucination rates significantly by integrating multi-step validation
  \end{itemize}}

\newcommand{\InverseRendering}{%
  \entry{Inverse Rendering with 2D Gaussian Splatting}{\report{https://sudhansh6.github.io/assets/pdf/CSE252D_Report.pdf}}{Mar 2024 -- May 2024}
  \begin{itemize}
    \item Developed a novel \textbf{inverse rendering} framework in CUDA to recover PBR properties of a scene using 2D Gaussian
          Splatting, improving normal map MAE by \textbf{15\%} over the current SOTA methods
  \end{itemize}}

% ---------------------------------------------------------------------------------------------------
% Publications: title on the left, venue on the right, one line each.
\newcommand{\publine}[3]{\par\noindent\href{#1}{#2}\hfill\hspace{1em}\side{#3}\par\addvspace{1.5pt}}
\newcommand{\PubSaturation}{\publine{https://doi.org/10.1016/j.sigpro.2023.109349}{A likelihood based method for compressive signal recovery under Gaussian and saturation noise}{Signal Processing, 2024}}
\newcommand{\PubPermutation}{\publine{https://ieeexplore.ieee.org/document/10888147}{Identification and correction of permutation errors in compressed sensing-based group testing}{ICASSP 2025}}
\newcommand{\PubRearrangement}{\publine{https://arxiv.org/abs/2411.10899}{Planning for tabletop object rearrangement}{ICRA 2025}}

% ---------------------------------------------------------------------------------------------------
% Education, skills, service
% Each profile sets \UCSDCourses, \IITBCourses and \SkillRows before \input{resume-common}.
\providecommand{\UCSDCourses}{}
\providecommand{\IITBCourses}{}
\providecommand{\SkillRows}{}
\newcommand{\coursework}[1]{\ifblank{#1}{}{\par\noindent{\fontsize{9.5}{11.5}\selectfont{\color{muted}Coursework:} #1}\par}}
\newcommand{\Education}{%
  \section{Education}
  \entry{University of California San Diego}{M.S. in Computer Science and Engineering (AI)}{GPA 4.00/4.00 \quad Sep 2023 -- Mar 2025}
  \coursework{\UCSDCourses}
  \entry{Indian Institute of Technology Bombay}{B.Tech.\ with Honors in Computer Science}{CPI 9.66/10 \quad Jul 2019 -- Jul 2023}
  \coursework{\IITBCourses}}

% \skill{Label}{Items} is one row of the skills table.
\newcommand{\skill}[2]{{\headfont\bfseries #1} & #2 \\[1.5pt]}
\newcommand{\cpp}{C\kern-0.04em\raisebox{0.06ex}{\small++}}
\newcommand{\Skills}{%
  \section{Skills}
  \begin{tabularx}{\linewidth}{@{}l@{\hspace{12pt}}>{\raggedright\arraybackslash}X@{}}
    \SkillRows
  \end{tabularx}}
.
\providecommand{\UCSDCourses}{}
\providecommand{\IITBCourses}{}
\providecommand{\SkillRows}{}
\newcommand{\coursework}[1]{\ifblank{#1}{}{\par\noindent{\fontsize{9.5}{11.5}\selectfont{\color{muted}Coursework:} #1}\par}}
\newcommand{\Education}{%
  \section{Education}
  \entry{University of California San Diego}{M.S. in Computer Science and Engineering (AI)}{GPA 4.00/4.00 \quad Sep 2023 -- Mar 2025}
  \coursework{\UCSDCourses}
  \entry{Indian Institute of Technology Bombay}{B.Tech.\ with Honors in Computer Science}{CPI 9.66/10 \quad Jul 2019 -- Jul 2023}
  \coursework{\IITBCourses}}

% \skill{Label}{Items} is one row of the skills table.
\newcommand{\skill}[2]{{\headfont\bfseries #1} & #2 \\[1.5pt]}
\newcommand{\cpp}{C\kern-0.04em\raisebox{0.06ex}{\small++}}
\newcommand{\Skills}{%
  \section{Skills}
  \begin{tabularx}{\linewidth}{@{}l@{\hspace{12pt}}>{\raggedright\arraybackslash}X@{}}
    \SkillRows
  \end{tabularx}}
 to try another pairing.
% Default: Charter body with Source Sans headings. All-sans alternative:
%   \def\resumefont{\usepackage[default,semibold]{sourcesanspro}}\def\headfont{}
\providecommand{\resumefont}{\usepackage[scaled=0.96]{XCharter}\usepackage[semibold,scale=0.96]{sourcesanspro}}
\providecommand{\headfont}{\sffamily}
\providecommand{\namefont}{\rmfamily\scshape}
\resumefont
\usepackage[final]{microtype}
\usepackage[dvipsnames]{xcolor}
\usepackage{enumitem}
\usepackage{titlesec}
\usepackage{tabularx}
\usepackage{fontawesome5}
\usepackage{etoolbox}
\usepackage[hidelinks]{hyperref}

% The site's green, darkened for print.
\definecolor{accent}{HTML}{3A7A45}
\definecolor{muted}{HTML}{5C675D}

\pagestyle{empty}
\setlength{\parindent}{0pt}
\setlength{\tabcolsep}{0pt}
\providecommand{\leading}{1.1}
\linespread{\leading}

\hypersetup{pdfauthor={Sudhansh Peddabomma}}

% Type scale: name 26pt small caps, section titles 12.5pt caps, entry titles 10.8pt, body 10pt,
% side notes 9pt.
\newcommand{\side}[1]{{\fontsize{9}{11}\selectfont\color{muted}#1}}
\newcommand{\sectitle}[1]{\textls[60]{\MakeUppercase{#1}}}
\titleformat{\section}{\headfont\color{accent}\bfseries\fontsize{12.5}{14}\selectfont}{}{0pt}{\sectitle}[\vspace{1.5pt}{\color{accent}\titlerule[0.8pt]}\global\firstentrytrue]
% The first entry after a section title skips \entrygap, so only \rulegap sits under the rule.
\newif\iffirstentry

% Vertical rhythm, the same idea as the earlier hand-tuned résumés: bullets inside an entry sit
% tight together, entries are separated by a clear gap, and sections by a larger one. A profile
% may override these before % Shared preamble and content for Sudhansh Peddabomma's résumés.
% Each profile (SudhanshPeddabomma-AI.tex, SudhanshPeddabomma-Vision.tex) inputs this file and picks
% which blocks to show and in what order, so every bullet is written once.
% Layout rule: every bullet fills one line, or runs 1.5 to 2 lines. Check with: python3 check_lines.py <pdf>
\documentclass[10pt,letterpaper]{article}

\usepackage[margin=0.55in,top=0.5in,bottom=0.45in]{geometry}
\usepackage[T1]{fontenc}
% Fonts: \resumefont loads the body font; \headfont styles the name and section titles.
% Override either before % Shared preamble and content for Sudhansh Peddabomma's résumés.
% Each profile (SudhanshPeddabomma-AI.tex, SudhanshPeddabomma-Vision.tex) inputs this file and picks
% which blocks to show and in what order, so every bullet is written once.
% Layout rule: every bullet fills one line, or runs 1.5 to 2 lines. Check with: python3 check_lines.py <pdf>
\documentclass[10pt,letterpaper]{article}

\usepackage[margin=0.55in,top=0.5in,bottom=0.45in]{geometry}
\usepackage[T1]{fontenc}
% Fonts: \resumefont loads the body font; \headfont styles the name and section titles.
% Override either before \input{resume-common} to try another pairing.
% Default: Charter body with Source Sans headings. All-sans alternative:
%   \def\resumefont{\usepackage[default,semibold]{sourcesanspro}}\def\headfont{}
\providecommand{\resumefont}{\usepackage[scaled=0.96]{XCharter}\usepackage[semibold,scale=0.96]{sourcesanspro}}
\providecommand{\headfont}{\sffamily}
\providecommand{\namefont}{\rmfamily\scshape}
\resumefont
\usepackage[final]{microtype}
\usepackage[dvipsnames]{xcolor}
\usepackage{enumitem}
\usepackage{titlesec}
\usepackage{tabularx}
\usepackage{fontawesome5}
\usepackage{etoolbox}
\usepackage[hidelinks]{hyperref}

% The site's green, darkened for print.
\definecolor{accent}{HTML}{3A7A45}
\definecolor{muted}{HTML}{5C675D}

\pagestyle{empty}
\setlength{\parindent}{0pt}
\setlength{\tabcolsep}{0pt}
\providecommand{\leading}{1.1}
\linespread{\leading}

\hypersetup{pdfauthor={Sudhansh Peddabomma}}

% Type scale: name 26pt small caps, section titles 12.5pt caps, entry titles 10.8pt, body 10pt,
% side notes 9pt.
\newcommand{\side}[1]{{\fontsize{9}{11}\selectfont\color{muted}#1}}
\newcommand{\sectitle}[1]{\textls[60]{\MakeUppercase{#1}}}
\titleformat{\section}{\headfont\color{accent}\bfseries\fontsize{12.5}{14}\selectfont}{}{0pt}{\sectitle}[\vspace{1.5pt}{\color{accent}\titlerule[0.8pt]}\global\firstentrytrue]
% The first entry after a section title skips \entrygap, so only \rulegap sits under the rule.
\newif\iffirstentry

% Vertical rhythm, the same idea as the earlier hand-tuned résumés: bullets inside an entry sit
% tight together, entries are separated by a clear gap, and sections by a larger one. A profile
% may override these before \input{resume-common} to fit one page.
\providecommand{\secgap}{12pt}   % above a section title
\providecommand{\rulegap}{5.5pt}   % section rule to first entry
\providecommand{\entrygap}{7pt}  % between entries
\providecommand{\headgap}{1.5pt} % entry heading to its bullets
\providecommand{\itemgap}{0.6pt} % between bullets of one entry
\titlespacing*{\section}{0pt}{\secgap}{\rulegap}

\newcommand{\squarebullet}{\raisebox{0.32ex}{\color{accent}\rule{0.5ex}{0.5ex}}}
\setlist[itemize]{leftmargin=12pt, labelsep=6pt, topsep=\headgap, itemsep=\itemgap, parsep=0pt, partopsep=0pt,
  label=\squarebullet}

% \entry{Organisation or title}{Role or detail (may be empty)}{Dates}; \sub{Role or detail}{Right-hand note}
\newcommand{\entry}[3]{%
  \par\iffirstentry\global\firstentryfalse\else\addvspace{\entrygap}\fi\noindent{\headfont\fontsize{10.8}{13}\selectfont\textbf{#1}}\ifblank{#2}{}{\ {\color{muted}|}\ #2}\hfill\hspace{1em}\side{#3}\par}
\newcommand{\sub}[2]{%
  \noindent\textit{#1}\hfill{\color{muted}\textit{#2}}\par}
\newcommand{\report}[1]{\href{#1}{\color{accent}Report}}
\newcommand{\me}{\textbf{S. Peddabomma}}

% ---------------------------------------------------------------------------------------------------
% Header
\newcommand{\Header}{%
  \begin{center}
    {\namefont\fontsize{26}{28}\selectfont Sudhansh Peddabomma}\\[6pt]
    {\headfont\fontsize{9.5}{11}\selectfont\color{muted}
      \href{mailto:sudhansh6@gmail.com}{\faEnvelope[regular]\ sudhansh6@gmail.com} \quad
      \href{https://sudhansh6.github.io}{\faGlobe\ sudhansh6.github.io} \quad
      \href{https://www.linkedin.com/in/sudhansh-peddabomma-73a208192/}{\faLinkedin\ sudhansh-peddabomma} \quad
      \href{https://github.com/Sudhansh6}{\faGithub\ Sudhansh6}}
  \end{center}\vspace{-4pt}}

% ---------------------------------------------------------------------------------------------------
% Glean bullets, shared by both profiles
\newcommand{\GleanSlides}{\item Led \textbf{AI slide generation} in Glean Assistant, now creating \textbf{180K decks a week} for \textbf{50K weekly active users} at enterprises}
\newcommand{\GleanSlidesPerf}{\item Improved the slides harness, raising reliability to \textbf{97\%}, cutting median latency \textbf{26\%} and cost per deck by \textbf{half}
  in four months}
\newcommand{\GleanSlidesOrch}{\item Architected a \textbf{slide orchestration system} that reuses relevant slides from a company's existing decks and generates
  new ones in its brand and template, so teams can refine decks with each other and Glean in full context, without leaving the product}
\newcommand{\GleanAgents}{\item Moved Glean's \textbf{workflow agents} from one-shot planning to an \textbf{agentic tool-use loop}, greatly improving agentic
  capabilities and instruction following, with \textbf{18\%} higher answer completeness and \textbf{17\%} fewer tokens}
\newcommand{\GleanTeamChat}{\item Engineered the core of \textbf{Team chat}, letting teammates work with Glean in one shared conversation and co-edit its artifacts}
\newcommand{\GleanVideo}{\item Shipped \textbf{video artifacts} end to end, generating narrated videos in chat with streamed playback and admin rollout controls}
\newcommand{\GleanEvals}{\item Created the \textbf{evaluations and metrics} that guide model choices and rollouts across agents, slides and interactive artifacts}

% ---------------------------------------------------------------------------------------------------
% Experience
% \groupbreak separates related groups of bullets inside one long entry (slides, then platform work).
\providecommand{\groupgap}{3.5pt}
\newcommand{\groupbreak}{\vspace{\groupgap}}
\newcommand{\GleanAI}{%
  \entry{Glean}{Software Engineer}{Apr 2025 -- Present}
  \begin{itemize}
    \GleanSlides
    \GleanSlidesOrch
    \GleanSlidesPerf \groupbreak
    \GleanTeamChat
    \GleanVideo
    \GleanAgents
    \GleanEvals
  \end{itemize}}

\newcommand{\GleanVision}{%
  \entry{Glean}{Software Engineer}{Apr 2025 -- Present}
  \begin{itemize}
    \GleanSlides
    \GleanSlidesOrch
    \GleanSlidesPerf \groupbreak
    \GleanVideo
    \GleanEvals
  \end{itemize}}

\newcommand{\Duality}{%
  \entry{Duality AI}{Computer Vision Intern}{Jun 2024 -- Sep 2024}
  \begin{itemize}
    \item Designed automated \textbf{3D reconstruction} techniques for featureless objects, reducing digital-twin generation time
          by \textbf{40\%}
    \item Collaborated with \textbf{Autodesk} to validate Unreal Engine simulations for robotics tasks; leveraging structured
          domain randomization to reduce Sim2Real gap and increase mAP-50 by \textbf{15\%} for object detection and segmentation
  \end{itemize}}

\newcommand{\Microsoft}{%
  \entry{Microsoft}{Data and Applied Scientist Intern}{May 2022 -- Jul 2022}
  \begin{itemize}
    \item Developed a \textbf{decision-tree ranker} to recommend emails without user queries, improving \textbf{Outlook} search capabilities
    \item Proposed hierarchical feature-sets for the ranker, reducing latency for recommendations and improving recall
  \end{itemize}}

% ---------------------------------------------------------------------------------------------------
% Research
\newcommand{\HomeRobots}{%
  \entry{Real-time 3D Perception for Home Robots}{UC San Diego}{Sep 2023 -- Sep 2024}
  \begin{itemize}
    \item Investigated real-time dense \textbf{visual SLAM} methods using NeRFs and Gaussian Splatting for indoor robot navigation
    \item Integrated object segmentation, grasp-pose estimation, and 3D mapping on the Fetch robot via ROS, demonstrating
          a novel \textbf{tabletop object rearrangement} algorithm that reduced cost by \textbf{20\%} compared to the state-of-the-art
  \end{itemize}}

\newcommand{\Tomography}{%
  \entry{3D Tomography with Primal-Dual Neural Networks}{University College London}{May 2021 -- Jul 2023}
  \begin{itemize}
    \item Developed a stochastic neural-network architecture of primal-dual algorithm for online reconstruction of 3D
          volumes from tomographic projections, with \textbf{99.6\,\% structural similarity} in low-dosage conditions
    \item Built a Python library with custom gradient operators for reconstructing volumes in a single pass, reducing
          compute requirements by up to \textbf{5x} over SOTA learning-based approaches for cone vector tomography
  \end{itemize}}

% ---------------------------------------------------------------------------------------------------
% Projects
\newcommand{\OOXML}{%
  \entry{ooxml}{Open-source contributor \textbar\ \href{https://github.com/ChristopherVR/ooxml}{\color{accent}GitHub}}{Sep 2026 -- Present}
  \begin{itemize}
    \item Contributed to an open-source \textbf{TypeScript} engine that renders and edits PowerPoint and Word files in the browser,
          improving editing fidelity, \textbf{real-time collaboration} and load performance across React, Vue and Svelte viewers
  \end{itemize}}

\newcommand{\MirrorAI}{%
  \entry{Mirror AI: Deployable Personas}{Honorable mention in Supabase YC Hackathon}{Oct 2024 -- Dec 2024}
  \begin{itemize}
    \item Designed an \textbf{agentic LLM architecture} with LangGraph to mirror user personalities, creating interactive digital
          personas, and deployed it as a full-stack platform on Supabase and Vercel with one-click publishing
  \end{itemize}}

\newcommand{\LLMReasoning}{%
  \entry{Improving LLM Reasoning for Numerical Problems}{\report{https://sudhansh6.github.io/assets/pdf/MathPrompter++.pdf}}{Sep 2024 -- Dec 2024}
  \begin{itemize}
    \item Enhanced MathPrompter (ACL 2023) with CoT, achieving \textbf{10\% higher accuracy} on Llama 3.1 1B where prior methods failed,
          and reduced hallucination rates significantly by integrating multi-step validation
  \end{itemize}}

\newcommand{\InverseRendering}{%
  \entry{Inverse Rendering with 2D Gaussian Splatting}{\report{https://sudhansh6.github.io/assets/pdf/CSE252D_Report.pdf}}{Mar 2024 -- May 2024}
  \begin{itemize}
    \item Developed a novel \textbf{inverse rendering} framework in CUDA to recover PBR properties of a scene using 2D Gaussian
          Splatting, improving normal map MAE by \textbf{15\%} over the current SOTA methods
  \end{itemize}}

% ---------------------------------------------------------------------------------------------------
% Publications: title on the left, venue on the right, one line each.
\newcommand{\publine}[3]{\par\noindent\href{#1}{#2}\hfill\hspace{1em}\side{#3}\par\addvspace{1.5pt}}
\newcommand{\PubSaturation}{\publine{https://doi.org/10.1016/j.sigpro.2023.109349}{A likelihood based method for compressive signal recovery under Gaussian and saturation noise}{Signal Processing, 2024}}
\newcommand{\PubPermutation}{\publine{https://ieeexplore.ieee.org/document/10888147}{Identification and correction of permutation errors in compressed sensing-based group testing}{ICASSP 2025}}
\newcommand{\PubRearrangement}{\publine{https://arxiv.org/abs/2411.10899}{Planning for tabletop object rearrangement}{ICRA 2025}}

% ---------------------------------------------------------------------------------------------------
% Education, skills, service
% Each profile sets \UCSDCourses, \IITBCourses and \SkillRows before \input{resume-common}.
\providecommand{\UCSDCourses}{}
\providecommand{\IITBCourses}{}
\providecommand{\SkillRows}{}
\newcommand{\coursework}[1]{\ifblank{#1}{}{\par\noindent{\fontsize{9.5}{11.5}\selectfont{\color{muted}Coursework:} #1}\par}}
\newcommand{\Education}{%
  \section{Education}
  \entry{University of California San Diego}{M.S. in Computer Science and Engineering (AI)}{GPA 4.00/4.00 \quad Sep 2023 -- Mar 2025}
  \coursework{\UCSDCourses}
  \entry{Indian Institute of Technology Bombay}{B.Tech.\ with Honors in Computer Science}{CPI 9.66/10 \quad Jul 2019 -- Jul 2023}
  \coursework{\IITBCourses}}

% \skill{Label}{Items} is one row of the skills table.
\newcommand{\skill}[2]{{\headfont\bfseries #1} & #2 \\[1.5pt]}
\newcommand{\cpp}{C\kern-0.04em\raisebox{0.06ex}{\small++}}
\newcommand{\Skills}{%
  \section{Skills}
  \begin{tabularx}{\linewidth}{@{}l@{\hspace{12pt}}>{\raggedright\arraybackslash}X@{}}
    \SkillRows
  \end{tabularx}}
 to try another pairing.
% Default: Charter body with Source Sans headings. All-sans alternative:
%   \def\resumefont{\usepackage[default,semibold]{sourcesanspro}}\def\headfont{}
\providecommand{\resumefont}{\usepackage[scaled=0.96]{XCharter}\usepackage[semibold,scale=0.96]{sourcesanspro}}
\providecommand{\headfont}{\sffamily}
\providecommand{\namefont}{\rmfamily\scshape}
\resumefont
\usepackage[final]{microtype}
\usepackage[dvipsnames]{xcolor}
\usepackage{enumitem}
\usepackage{titlesec}
\usepackage{tabularx}
\usepackage{fontawesome5}
\usepackage{etoolbox}
\usepackage[hidelinks]{hyperref}

% The site's green, darkened for print.
\definecolor{accent}{HTML}{3A7A45}
\definecolor{muted}{HTML}{5C675D}

\pagestyle{empty}
\setlength{\parindent}{0pt}
\setlength{\tabcolsep}{0pt}
\providecommand{\leading}{1.1}
\linespread{\leading}

\hypersetup{pdfauthor={Sudhansh Peddabomma}}

% Type scale: name 26pt small caps, section titles 12.5pt caps, entry titles 10.8pt, body 10pt,
% side notes 9pt.
\newcommand{\side}[1]{{\fontsize{9}{11}\selectfont\color{muted}#1}}
\newcommand{\sectitle}[1]{\textls[60]{\MakeUppercase{#1}}}
\titleformat{\section}{\headfont\color{accent}\bfseries\fontsize{12.5}{14}\selectfont}{}{0pt}{\sectitle}[\vspace{1.5pt}{\color{accent}\titlerule[0.8pt]}\global\firstentrytrue]
% The first entry after a section title skips \entrygap, so only \rulegap sits under the rule.
\newif\iffirstentry

% Vertical rhythm, the same idea as the earlier hand-tuned résumés: bullets inside an entry sit
% tight together, entries are separated by a clear gap, and sections by a larger one. A profile
% may override these before % Shared preamble and content for Sudhansh Peddabomma's résumés.
% Each profile (SudhanshPeddabomma-AI.tex, SudhanshPeddabomma-Vision.tex) inputs this file and picks
% which blocks to show and in what order, so every bullet is written once.
% Layout rule: every bullet fills one line, or runs 1.5 to 2 lines. Check with: python3 check_lines.py <pdf>
\documentclass[10pt,letterpaper]{article}

\usepackage[margin=0.55in,top=0.5in,bottom=0.45in]{geometry}
\usepackage[T1]{fontenc}
% Fonts: \resumefont loads the body font; \headfont styles the name and section titles.
% Override either before \input{resume-common} to try another pairing.
% Default: Charter body with Source Sans headings. All-sans alternative:
%   \def\resumefont{\usepackage[default,semibold]{sourcesanspro}}\def\headfont{}
\providecommand{\resumefont}{\usepackage[scaled=0.96]{XCharter}\usepackage[semibold,scale=0.96]{sourcesanspro}}
\providecommand{\headfont}{\sffamily}
\providecommand{\namefont}{\rmfamily\scshape}
\resumefont
\usepackage[final]{microtype}
\usepackage[dvipsnames]{xcolor}
\usepackage{enumitem}
\usepackage{titlesec}
\usepackage{tabularx}
\usepackage{fontawesome5}
\usepackage{etoolbox}
\usepackage[hidelinks]{hyperref}

% The site's green, darkened for print.
\definecolor{accent}{HTML}{3A7A45}
\definecolor{muted}{HTML}{5C675D}

\pagestyle{empty}
\setlength{\parindent}{0pt}
\setlength{\tabcolsep}{0pt}
\providecommand{\leading}{1.1}
\linespread{\leading}

\hypersetup{pdfauthor={Sudhansh Peddabomma}}

% Type scale: name 26pt small caps, section titles 12.5pt caps, entry titles 10.8pt, body 10pt,
% side notes 9pt.
\newcommand{\side}[1]{{\fontsize{9}{11}\selectfont\color{muted}#1}}
\newcommand{\sectitle}[1]{\textls[60]{\MakeUppercase{#1}}}
\titleformat{\section}{\headfont\color{accent}\bfseries\fontsize{12.5}{14}\selectfont}{}{0pt}{\sectitle}[\vspace{1.5pt}{\color{accent}\titlerule[0.8pt]}\global\firstentrytrue]
% The first entry after a section title skips \entrygap, so only \rulegap sits under the rule.
\newif\iffirstentry

% Vertical rhythm, the same idea as the earlier hand-tuned résumés: bullets inside an entry sit
% tight together, entries are separated by a clear gap, and sections by a larger one. A profile
% may override these before \input{resume-common} to fit one page.
\providecommand{\secgap}{12pt}   % above a section title
\providecommand{\rulegap}{5.5pt}   % section rule to first entry
\providecommand{\entrygap}{7pt}  % between entries
\providecommand{\headgap}{1.5pt} % entry heading to its bullets
\providecommand{\itemgap}{0.6pt} % between bullets of one entry
\titlespacing*{\section}{0pt}{\secgap}{\rulegap}

\newcommand{\squarebullet}{\raisebox{0.32ex}{\color{accent}\rule{0.5ex}{0.5ex}}}
\setlist[itemize]{leftmargin=12pt, labelsep=6pt, topsep=\headgap, itemsep=\itemgap, parsep=0pt, partopsep=0pt,
  label=\squarebullet}

% \entry{Organisation or title}{Role or detail (may be empty)}{Dates}; \sub{Role or detail}{Right-hand note}
\newcommand{\entry}[3]{%
  \par\iffirstentry\global\firstentryfalse\else\addvspace{\entrygap}\fi\noindent{\headfont\fontsize{10.8}{13}\selectfont\textbf{#1}}\ifblank{#2}{}{\ {\color{muted}|}\ #2}\hfill\hspace{1em}\side{#3}\par}
\newcommand{\sub}[2]{%
  \noindent\textit{#1}\hfill{\color{muted}\textit{#2}}\par}
\newcommand{\report}[1]{\href{#1}{\color{accent}Report}}
\newcommand{\me}{\textbf{S. Peddabomma}}

% ---------------------------------------------------------------------------------------------------
% Header
\newcommand{\Header}{%
  \begin{center}
    {\namefont\fontsize{26}{28}\selectfont Sudhansh Peddabomma}\\[6pt]
    {\headfont\fontsize{9.5}{11}\selectfont\color{muted}
      \href{mailto:sudhansh6@gmail.com}{\faEnvelope[regular]\ sudhansh6@gmail.com} \quad
      \href{https://sudhansh6.github.io}{\faGlobe\ sudhansh6.github.io} \quad
      \href{https://www.linkedin.com/in/sudhansh-peddabomma-73a208192/}{\faLinkedin\ sudhansh-peddabomma} \quad
      \href{https://github.com/Sudhansh6}{\faGithub\ Sudhansh6}}
  \end{center}\vspace{-4pt}}

% ---------------------------------------------------------------------------------------------------
% Glean bullets, shared by both profiles
\newcommand{\GleanSlides}{\item Led \textbf{AI slide generation} in Glean Assistant, now creating \textbf{180K decks a week} for \textbf{50K weekly active users} at enterprises}
\newcommand{\GleanSlidesPerf}{\item Improved the slides harness, raising reliability to \textbf{97\%}, cutting median latency \textbf{26\%} and cost per deck by \textbf{half}
  in four months}
\newcommand{\GleanSlidesOrch}{\item Architected a \textbf{slide orchestration system} that reuses relevant slides from a company's existing decks and generates
  new ones in its brand and template, so teams can refine decks with each other and Glean in full context, without leaving the product}
\newcommand{\GleanAgents}{\item Moved Glean's \textbf{workflow agents} from one-shot planning to an \textbf{agentic tool-use loop}, greatly improving agentic
  capabilities and instruction following, with \textbf{18\%} higher answer completeness and \textbf{17\%} fewer tokens}
\newcommand{\GleanTeamChat}{\item Engineered the core of \textbf{Team chat}, letting teammates work with Glean in one shared conversation and co-edit its artifacts}
\newcommand{\GleanVideo}{\item Shipped \textbf{video artifacts} end to end, generating narrated videos in chat with streamed playback and admin rollout controls}
\newcommand{\GleanEvals}{\item Created the \textbf{evaluations and metrics} that guide model choices and rollouts across agents, slides and interactive artifacts}

% ---------------------------------------------------------------------------------------------------
% Experience
% \groupbreak separates related groups of bullets inside one long entry (slides, then platform work).
\providecommand{\groupgap}{3.5pt}
\newcommand{\groupbreak}{\vspace{\groupgap}}
\newcommand{\GleanAI}{%
  \entry{Glean}{Software Engineer}{Apr 2025 -- Present}
  \begin{itemize}
    \GleanSlides
    \GleanSlidesOrch
    \GleanSlidesPerf \groupbreak
    \GleanTeamChat
    \GleanVideo
    \GleanAgents
    \GleanEvals
  \end{itemize}}

\newcommand{\GleanVision}{%
  \entry{Glean}{Software Engineer}{Apr 2025 -- Present}
  \begin{itemize}
    \GleanSlides
    \GleanSlidesOrch
    \GleanSlidesPerf \groupbreak
    \GleanVideo
    \GleanEvals
  \end{itemize}}

\newcommand{\Duality}{%
  \entry{Duality AI}{Computer Vision Intern}{Jun 2024 -- Sep 2024}
  \begin{itemize}
    \item Designed automated \textbf{3D reconstruction} techniques for featureless objects, reducing digital-twin generation time
          by \textbf{40\%}
    \item Collaborated with \textbf{Autodesk} to validate Unreal Engine simulations for robotics tasks; leveraging structured
          domain randomization to reduce Sim2Real gap and increase mAP-50 by \textbf{15\%} for object detection and segmentation
  \end{itemize}}

\newcommand{\Microsoft}{%
  \entry{Microsoft}{Data and Applied Scientist Intern}{May 2022 -- Jul 2022}
  \begin{itemize}
    \item Developed a \textbf{decision-tree ranker} to recommend emails without user queries, improving \textbf{Outlook} search capabilities
    \item Proposed hierarchical feature-sets for the ranker, reducing latency for recommendations and improving recall
  \end{itemize}}

% ---------------------------------------------------------------------------------------------------
% Research
\newcommand{\HomeRobots}{%
  \entry{Real-time 3D Perception for Home Robots}{UC San Diego}{Sep 2023 -- Sep 2024}
  \begin{itemize}
    \item Investigated real-time dense \textbf{visual SLAM} methods using NeRFs and Gaussian Splatting for indoor robot navigation
    \item Integrated object segmentation, grasp-pose estimation, and 3D mapping on the Fetch robot via ROS, demonstrating
          a novel \textbf{tabletop object rearrangement} algorithm that reduced cost by \textbf{20\%} compared to the state-of-the-art
  \end{itemize}}

\newcommand{\Tomography}{%
  \entry{3D Tomography with Primal-Dual Neural Networks}{University College London}{May 2021 -- Jul 2023}
  \begin{itemize}
    \item Developed a stochastic neural-network architecture of primal-dual algorithm for online reconstruction of 3D
          volumes from tomographic projections, with \textbf{99.6\,\% structural similarity} in low-dosage conditions
    \item Built a Python library with custom gradient operators for reconstructing volumes in a single pass, reducing
          compute requirements by up to \textbf{5x} over SOTA learning-based approaches for cone vector tomography
  \end{itemize}}

% ---------------------------------------------------------------------------------------------------
% Projects
\newcommand{\OOXML}{%
  \entry{ooxml}{Open-source contributor \textbar\ \href{https://github.com/ChristopherVR/ooxml}{\color{accent}GitHub}}{Sep 2026 -- Present}
  \begin{itemize}
    \item Contributed to an open-source \textbf{TypeScript} engine that renders and edits PowerPoint and Word files in the browser,
          improving editing fidelity, \textbf{real-time collaboration} and load performance across React, Vue and Svelte viewers
  \end{itemize}}

\newcommand{\MirrorAI}{%
  \entry{Mirror AI: Deployable Personas}{Honorable mention in Supabase YC Hackathon}{Oct 2024 -- Dec 2024}
  \begin{itemize}
    \item Designed an \textbf{agentic LLM architecture} with LangGraph to mirror user personalities, creating interactive digital
          personas, and deployed it as a full-stack platform on Supabase and Vercel with one-click publishing
  \end{itemize}}

\newcommand{\LLMReasoning}{%
  \entry{Improving LLM Reasoning for Numerical Problems}{\report{https://sudhansh6.github.io/assets/pdf/MathPrompter++.pdf}}{Sep 2024 -- Dec 2024}
  \begin{itemize}
    \item Enhanced MathPrompter (ACL 2023) with CoT, achieving \textbf{10\% higher accuracy} on Llama 3.1 1B where prior methods failed,
          and reduced hallucination rates significantly by integrating multi-step validation
  \end{itemize}}

\newcommand{\InverseRendering}{%
  \entry{Inverse Rendering with 2D Gaussian Splatting}{\report{https://sudhansh6.github.io/assets/pdf/CSE252D_Report.pdf}}{Mar 2024 -- May 2024}
  \begin{itemize}
    \item Developed a novel \textbf{inverse rendering} framework in CUDA to recover PBR properties of a scene using 2D Gaussian
          Splatting, improving normal map MAE by \textbf{15\%} over the current SOTA methods
  \end{itemize}}

% ---------------------------------------------------------------------------------------------------
% Publications: title on the left, venue on the right, one line each.
\newcommand{\publine}[3]{\par\noindent\href{#1}{#2}\hfill\hspace{1em}\side{#3}\par\addvspace{1.5pt}}
\newcommand{\PubSaturation}{\publine{https://doi.org/10.1016/j.sigpro.2023.109349}{A likelihood based method for compressive signal recovery under Gaussian and saturation noise}{Signal Processing, 2024}}
\newcommand{\PubPermutation}{\publine{https://ieeexplore.ieee.org/document/10888147}{Identification and correction of permutation errors in compressed sensing-based group testing}{ICASSP 2025}}
\newcommand{\PubRearrangement}{\publine{https://arxiv.org/abs/2411.10899}{Planning for tabletop object rearrangement}{ICRA 2025}}

% ---------------------------------------------------------------------------------------------------
% Education, skills, service
% Each profile sets \UCSDCourses, \IITBCourses and \SkillRows before \input{resume-common}.
\providecommand{\UCSDCourses}{}
\providecommand{\IITBCourses}{}
\providecommand{\SkillRows}{}
\newcommand{\coursework}[1]{\ifblank{#1}{}{\par\noindent{\fontsize{9.5}{11.5}\selectfont{\color{muted}Coursework:} #1}\par}}
\newcommand{\Education}{%
  \section{Education}
  \entry{University of California San Diego}{M.S. in Computer Science and Engineering (AI)}{GPA 4.00/4.00 \quad Sep 2023 -- Mar 2025}
  \coursework{\UCSDCourses}
  \entry{Indian Institute of Technology Bombay}{B.Tech.\ with Honors in Computer Science}{CPI 9.66/10 \quad Jul 2019 -- Jul 2023}
  \coursework{\IITBCourses}}

% \skill{Label}{Items} is one row of the skills table.
\newcommand{\skill}[2]{{\headfont\bfseries #1} & #2 \\[1.5pt]}
\newcommand{\cpp}{C\kern-0.04em\raisebox{0.06ex}{\small++}}
\newcommand{\Skills}{%
  \section{Skills}
  \begin{tabularx}{\linewidth}{@{}l@{\hspace{12pt}}>{\raggedright\arraybackslash}X@{}}
    \SkillRows
  \end{tabularx}}
 to fit one page.
\providecommand{\secgap}{12pt}   % above a section title
\providecommand{\rulegap}{5.5pt}   % section rule to first entry
\providecommand{\entrygap}{7pt}  % between entries
\providecommand{\headgap}{1.5pt} % entry heading to its bullets
\providecommand{\itemgap}{0.6pt} % between bullets of one entry
\titlespacing*{\section}{0pt}{\secgap}{\rulegap}

\newcommand{\squarebullet}{\raisebox{0.32ex}{\color{accent}\rule{0.5ex}{0.5ex}}}
\setlist[itemize]{leftmargin=12pt, labelsep=6pt, topsep=\headgap, itemsep=\itemgap, parsep=0pt, partopsep=0pt,
  label=\squarebullet}

% \entry{Organisation or title}{Role or detail (may be empty)}{Dates}; \sub{Role or detail}{Right-hand note}
\newcommand{\entry}[3]{%
  \par\iffirstentry\global\firstentryfalse\else\addvspace{\entrygap}\fi\noindent{\headfont\fontsize{10.8}{13}\selectfont\textbf{#1}}\ifblank{#2}{}{\ {\color{muted}|}\ #2}\hfill\hspace{1em}\side{#3}\par}
\newcommand{\sub}[2]{%
  \noindent\textit{#1}\hfill{\color{muted}\textit{#2}}\par}
\newcommand{\report}[1]{\href{#1}{\color{accent}Report}}
\newcommand{\me}{\textbf{S. Peddabomma}}

% ---------------------------------------------------------------------------------------------------
% Header
\newcommand{\Header}{%
  \begin{center}
    {\namefont\fontsize{26}{28}\selectfont Sudhansh Peddabomma}\\[6pt]
    {\headfont\fontsize{9.5}{11}\selectfont\color{muted}
      \href{mailto:sudhansh6@gmail.com}{\faEnvelope[regular]\ sudhansh6@gmail.com} \quad
      \href{https://sudhansh6.github.io}{\faGlobe\ sudhansh6.github.io} \quad
      \href{https://www.linkedin.com/in/sudhansh-peddabomma-73a208192/}{\faLinkedin\ sudhansh-peddabomma} \quad
      \href{https://github.com/Sudhansh6}{\faGithub\ Sudhansh6}}
  \end{center}\vspace{-4pt}}

% ---------------------------------------------------------------------------------------------------
% Glean bullets, shared by both profiles
\newcommand{\GleanSlides}{\item Led \textbf{AI slide generation} in Glean Assistant, now creating \textbf{180K decks a week} for \textbf{50K weekly active users} at enterprises}
\newcommand{\GleanSlidesPerf}{\item Improved the slides harness, raising reliability to \textbf{97\%}, cutting median latency \textbf{26\%} and cost per deck by \textbf{half}
  in four months}
\newcommand{\GleanSlidesOrch}{\item Architected a \textbf{slide orchestration system} that reuses relevant slides from a company's existing decks and generates
  new ones in its brand and template, so teams can refine decks with each other and Glean in full context, without leaving the product}
\newcommand{\GleanAgents}{\item Moved Glean's \textbf{workflow agents} from one-shot planning to an \textbf{agentic tool-use loop}, greatly improving agentic
  capabilities and instruction following, with \textbf{18\%} higher answer completeness and \textbf{17\%} fewer tokens}
\newcommand{\GleanTeamChat}{\item Engineered the core of \textbf{Team chat}, letting teammates work with Glean in one shared conversation and co-edit its artifacts}
\newcommand{\GleanVideo}{\item Shipped \textbf{video artifacts} end to end, generating narrated videos in chat with streamed playback and admin rollout controls}
\newcommand{\GleanEvals}{\item Created the \textbf{evaluations and metrics} that guide model choices and rollouts across agents, slides and interactive artifacts}

% ---------------------------------------------------------------------------------------------------
% Experience
% \groupbreak separates related groups of bullets inside one long entry (slides, then platform work).
\providecommand{\groupgap}{3.5pt}
\newcommand{\groupbreak}{\vspace{\groupgap}}
\newcommand{\GleanAI}{%
  \entry{Glean}{Software Engineer}{Apr 2025 -- Present}
  \begin{itemize}
    \GleanSlides
    \GleanSlidesOrch
    \GleanSlidesPerf \groupbreak
    \GleanTeamChat
    \GleanVideo
    \GleanAgents
    \GleanEvals
  \end{itemize}}

\newcommand{\GleanVision}{%
  \entry{Glean}{Software Engineer}{Apr 2025 -- Present}
  \begin{itemize}
    \GleanSlides
    \GleanSlidesOrch
    \GleanSlidesPerf \groupbreak
    \GleanVideo
    \GleanEvals
  \end{itemize}}

\newcommand{\Duality}{%
  \entry{Duality AI}{Computer Vision Intern}{Jun 2024 -- Sep 2024}
  \begin{itemize}
    \item Designed automated \textbf{3D reconstruction} techniques for featureless objects, reducing digital-twin generation time
          by \textbf{40\%}
    \item Collaborated with \textbf{Autodesk} to validate Unreal Engine simulations for robotics tasks; leveraging structured
          domain randomization to reduce Sim2Real gap and increase mAP-50 by \textbf{15\%} for object detection and segmentation
  \end{itemize}}

\newcommand{\Microsoft}{%
  \entry{Microsoft}{Data and Applied Scientist Intern}{May 2022 -- Jul 2022}
  \begin{itemize}
    \item Developed a \textbf{decision-tree ranker} to recommend emails without user queries, improving \textbf{Outlook} search capabilities
    \item Proposed hierarchical feature-sets for the ranker, reducing latency for recommendations and improving recall
  \end{itemize}}

% ---------------------------------------------------------------------------------------------------
% Research
\newcommand{\HomeRobots}{%
  \entry{Real-time 3D Perception for Home Robots}{UC San Diego}{Sep 2023 -- Sep 2024}
  \begin{itemize}
    \item Investigated real-time dense \textbf{visual SLAM} methods using NeRFs and Gaussian Splatting for indoor robot navigation
    \item Integrated object segmentation, grasp-pose estimation, and 3D mapping on the Fetch robot via ROS, demonstrating
          a novel \textbf{tabletop object rearrangement} algorithm that reduced cost by \textbf{20\%} compared to the state-of-the-art
  \end{itemize}}

\newcommand{\Tomography}{%
  \entry{3D Tomography with Primal-Dual Neural Networks}{University College London}{May 2021 -- Jul 2023}
  \begin{itemize}
    \item Developed a stochastic neural-network architecture of primal-dual algorithm for online reconstruction of 3D
          volumes from tomographic projections, with \textbf{99.6\,\% structural similarity} in low-dosage conditions
    \item Built a Python library with custom gradient operators for reconstructing volumes in a single pass, reducing
          compute requirements by up to \textbf{5x} over SOTA learning-based approaches for cone vector tomography
  \end{itemize}}

% ---------------------------------------------------------------------------------------------------
% Projects
\newcommand{\OOXML}{%
  \entry{ooxml}{Open-source contributor \textbar\ \href{https://github.com/ChristopherVR/ooxml}{\color{accent}GitHub}}{Sep 2026 -- Present}
  \begin{itemize}
    \item Contributed to an open-source \textbf{TypeScript} engine that renders and edits PowerPoint and Word files in the browser,
          improving editing fidelity, \textbf{real-time collaboration} and load performance across React, Vue and Svelte viewers
  \end{itemize}}

\newcommand{\MirrorAI}{%
  \entry{Mirror AI: Deployable Personas}{Honorable mention in Supabase YC Hackathon}{Oct 2024 -- Dec 2024}
  \begin{itemize}
    \item Designed an \textbf{agentic LLM architecture} with LangGraph to mirror user personalities, creating interactive digital
          personas, and deployed it as a full-stack platform on Supabase and Vercel with one-click publishing
  \end{itemize}}

\newcommand{\LLMReasoning}{%
  \entry{Improving LLM Reasoning for Numerical Problems}{\report{https://sudhansh6.github.io/assets/pdf/MathPrompter++.pdf}}{Sep 2024 -- Dec 2024}
  \begin{itemize}
    \item Enhanced MathPrompter (ACL 2023) with CoT, achieving \textbf{10\% higher accuracy} on Llama 3.1 1B where prior methods failed,
          and reduced hallucination rates significantly by integrating multi-step validation
  \end{itemize}}

\newcommand{\InverseRendering}{%
  \entry{Inverse Rendering with 2D Gaussian Splatting}{\report{https://sudhansh6.github.io/assets/pdf/CSE252D_Report.pdf}}{Mar 2024 -- May 2024}
  \begin{itemize}
    \item Developed a novel \textbf{inverse rendering} framework in CUDA to recover PBR properties of a scene using 2D Gaussian
          Splatting, improving normal map MAE by \textbf{15\%} over the current SOTA methods
  \end{itemize}}

% ---------------------------------------------------------------------------------------------------
% Publications: title on the left, venue on the right, one line each.
\newcommand{\publine}[3]{\par\noindent\href{#1}{#2}\hfill\hspace{1em}\side{#3}\par\addvspace{1.5pt}}
\newcommand{\PubSaturation}{\publine{https://doi.org/10.1016/j.sigpro.2023.109349}{A likelihood based method for compressive signal recovery under Gaussian and saturation noise}{Signal Processing, 2024}}
\newcommand{\PubPermutation}{\publine{https://ieeexplore.ieee.org/document/10888147}{Identification and correction of permutation errors in compressed sensing-based group testing}{ICASSP 2025}}
\newcommand{\PubRearrangement}{\publine{https://arxiv.org/abs/2411.10899}{Planning for tabletop object rearrangement}{ICRA 2025}}

% ---------------------------------------------------------------------------------------------------
% Education, skills, service
% Each profile sets \UCSDCourses, \IITBCourses and \SkillRows before % Shared preamble and content for Sudhansh Peddabomma's résumés.
% Each profile (SudhanshPeddabomma-AI.tex, SudhanshPeddabomma-Vision.tex) inputs this file and picks
% which blocks to show and in what order, so every bullet is written once.
% Layout rule: every bullet fills one line, or runs 1.5 to 2 lines. Check with: python3 check_lines.py <pdf>
\documentclass[10pt,letterpaper]{article}

\usepackage[margin=0.55in,top=0.5in,bottom=0.45in]{geometry}
\usepackage[T1]{fontenc}
% Fonts: \resumefont loads the body font; \headfont styles the name and section titles.
% Override either before \input{resume-common} to try another pairing.
% Default: Charter body with Source Sans headings. All-sans alternative:
%   \def\resumefont{\usepackage[default,semibold]{sourcesanspro}}\def\headfont{}
\providecommand{\resumefont}{\usepackage[scaled=0.96]{XCharter}\usepackage[semibold,scale=0.96]{sourcesanspro}}
\providecommand{\headfont}{\sffamily}
\providecommand{\namefont}{\rmfamily\scshape}
\resumefont
\usepackage[final]{microtype}
\usepackage[dvipsnames]{xcolor}
\usepackage{enumitem}
\usepackage{titlesec}
\usepackage{tabularx}
\usepackage{fontawesome5}
\usepackage{etoolbox}
\usepackage[hidelinks]{hyperref}

% The site's green, darkened for print.
\definecolor{accent}{HTML}{3A7A45}
\definecolor{muted}{HTML}{5C675D}

\pagestyle{empty}
\setlength{\parindent}{0pt}
\setlength{\tabcolsep}{0pt}
\providecommand{\leading}{1.1}
\linespread{\leading}

\hypersetup{pdfauthor={Sudhansh Peddabomma}}

% Type scale: name 26pt small caps, section titles 12.5pt caps, entry titles 10.8pt, body 10pt,
% side notes 9pt.
\newcommand{\side}[1]{{\fontsize{9}{11}\selectfont\color{muted}#1}}
\newcommand{\sectitle}[1]{\textls[60]{\MakeUppercase{#1}}}
\titleformat{\section}{\headfont\color{accent}\bfseries\fontsize{12.5}{14}\selectfont}{}{0pt}{\sectitle}[\vspace{1.5pt}{\color{accent}\titlerule[0.8pt]}\global\firstentrytrue]
% The first entry after a section title skips \entrygap, so only \rulegap sits under the rule.
\newif\iffirstentry

% Vertical rhythm, the same idea as the earlier hand-tuned résumés: bullets inside an entry sit
% tight together, entries are separated by a clear gap, and sections by a larger one. A profile
% may override these before \input{resume-common} to fit one page.
\providecommand{\secgap}{12pt}   % above a section title
\providecommand{\rulegap}{5.5pt}   % section rule to first entry
\providecommand{\entrygap}{7pt}  % between entries
\providecommand{\headgap}{1.5pt} % entry heading to its bullets
\providecommand{\itemgap}{0.6pt} % between bullets of one entry
\titlespacing*{\section}{0pt}{\secgap}{\rulegap}

\newcommand{\squarebullet}{\raisebox{0.32ex}{\color{accent}\rule{0.5ex}{0.5ex}}}
\setlist[itemize]{leftmargin=12pt, labelsep=6pt, topsep=\headgap, itemsep=\itemgap, parsep=0pt, partopsep=0pt,
  label=\squarebullet}

% \entry{Organisation or title}{Role or detail (may be empty)}{Dates}; \sub{Role or detail}{Right-hand note}
\newcommand{\entry}[3]{%
  \par\iffirstentry\global\firstentryfalse\else\addvspace{\entrygap}\fi\noindent{\headfont\fontsize{10.8}{13}\selectfont\textbf{#1}}\ifblank{#2}{}{\ {\color{muted}|}\ #2}\hfill\hspace{1em}\side{#3}\par}
\newcommand{\sub}[2]{%
  \noindent\textit{#1}\hfill{\color{muted}\textit{#2}}\par}
\newcommand{\report}[1]{\href{#1}{\color{accent}Report}}
\newcommand{\me}{\textbf{S. Peddabomma}}

% ---------------------------------------------------------------------------------------------------
% Header
\newcommand{\Header}{%
  \begin{center}
    {\namefont\fontsize{26}{28}\selectfont Sudhansh Peddabomma}\\[6pt]
    {\headfont\fontsize{9.5}{11}\selectfont\color{muted}
      \href{mailto:sudhansh6@gmail.com}{\faEnvelope[regular]\ sudhansh6@gmail.com} \quad
      \href{https://sudhansh6.github.io}{\faGlobe\ sudhansh6.github.io} \quad
      \href{https://www.linkedin.com/in/sudhansh-peddabomma-73a208192/}{\faLinkedin\ sudhansh-peddabomma} \quad
      \href{https://github.com/Sudhansh6}{\faGithub\ Sudhansh6}}
  \end{center}\vspace{-4pt}}

% ---------------------------------------------------------------------------------------------------
% Glean bullets, shared by both profiles
\newcommand{\GleanSlides}{\item Led \textbf{AI slide generation} in Glean Assistant, now creating \textbf{180K decks a week} for \textbf{50K weekly active users} at enterprises}
\newcommand{\GleanSlidesPerf}{\item Improved the slides harness, raising reliability to \textbf{97\%}, cutting median latency \textbf{26\%} and cost per deck by \textbf{half}
  in four months}
\newcommand{\GleanSlidesOrch}{\item Architected a \textbf{slide orchestration system} that reuses relevant slides from a company's existing decks and generates
  new ones in its brand and template, so teams can refine decks with each other and Glean in full context, without leaving the product}
\newcommand{\GleanAgents}{\item Moved Glean's \textbf{workflow agents} from one-shot planning to an \textbf{agentic tool-use loop}, greatly improving agentic
  capabilities and instruction following, with \textbf{18\%} higher answer completeness and \textbf{17\%} fewer tokens}
\newcommand{\GleanTeamChat}{\item Engineered the core of \textbf{Team chat}, letting teammates work with Glean in one shared conversation and co-edit its artifacts}
\newcommand{\GleanVideo}{\item Shipped \textbf{video artifacts} end to end, generating narrated videos in chat with streamed playback and admin rollout controls}
\newcommand{\GleanEvals}{\item Created the \textbf{evaluations and metrics} that guide model choices and rollouts across agents, slides and interactive artifacts}

% ---------------------------------------------------------------------------------------------------
% Experience
% \groupbreak separates related groups of bullets inside one long entry (slides, then platform work).
\providecommand{\groupgap}{3.5pt}
\newcommand{\groupbreak}{\vspace{\groupgap}}
\newcommand{\GleanAI}{%
  \entry{Glean}{Software Engineer}{Apr 2025 -- Present}
  \begin{itemize}
    \GleanSlides
    \GleanSlidesOrch
    \GleanSlidesPerf \groupbreak
    \GleanTeamChat
    \GleanVideo
    \GleanAgents
    \GleanEvals
  \end{itemize}}

\newcommand{\GleanVision}{%
  \entry{Glean}{Software Engineer}{Apr 2025 -- Present}
  \begin{itemize}
    \GleanSlides
    \GleanSlidesOrch
    \GleanSlidesPerf \groupbreak
    \GleanVideo
    \GleanEvals
  \end{itemize}}

\newcommand{\Duality}{%
  \entry{Duality AI}{Computer Vision Intern}{Jun 2024 -- Sep 2024}
  \begin{itemize}
    \item Designed automated \textbf{3D reconstruction} techniques for featureless objects, reducing digital-twin generation time
          by \textbf{40\%}
    \item Collaborated with \textbf{Autodesk} to validate Unreal Engine simulations for robotics tasks; leveraging structured
          domain randomization to reduce Sim2Real gap and increase mAP-50 by \textbf{15\%} for object detection and segmentation
  \end{itemize}}

\newcommand{\Microsoft}{%
  \entry{Microsoft}{Data and Applied Scientist Intern}{May 2022 -- Jul 2022}
  \begin{itemize}
    \item Developed a \textbf{decision-tree ranker} to recommend emails without user queries, improving \textbf{Outlook} search capabilities
    \item Proposed hierarchical feature-sets for the ranker, reducing latency for recommendations and improving recall
  \end{itemize}}

% ---------------------------------------------------------------------------------------------------
% Research
\newcommand{\HomeRobots}{%
  \entry{Real-time 3D Perception for Home Robots}{UC San Diego}{Sep 2023 -- Sep 2024}
  \begin{itemize}
    \item Investigated real-time dense \textbf{visual SLAM} methods using NeRFs and Gaussian Splatting for indoor robot navigation
    \item Integrated object segmentation, grasp-pose estimation, and 3D mapping on the Fetch robot via ROS, demonstrating
          a novel \textbf{tabletop object rearrangement} algorithm that reduced cost by \textbf{20\%} compared to the state-of-the-art
  \end{itemize}}

\newcommand{\Tomography}{%
  \entry{3D Tomography with Primal-Dual Neural Networks}{University College London}{May 2021 -- Jul 2023}
  \begin{itemize}
    \item Developed a stochastic neural-network architecture of primal-dual algorithm for online reconstruction of 3D
          volumes from tomographic projections, with \textbf{99.6\,\% structural similarity} in low-dosage conditions
    \item Built a Python library with custom gradient operators for reconstructing volumes in a single pass, reducing
          compute requirements by up to \textbf{5x} over SOTA learning-based approaches for cone vector tomography
  \end{itemize}}

% ---------------------------------------------------------------------------------------------------
% Projects
\newcommand{\OOXML}{%
  \entry{ooxml}{Open-source contributor \textbar\ \href{https://github.com/ChristopherVR/ooxml}{\color{accent}GitHub}}{Sep 2026 -- Present}
  \begin{itemize}
    \item Contributed to an open-source \textbf{TypeScript} engine that renders and edits PowerPoint and Word files in the browser,
          improving editing fidelity, \textbf{real-time collaboration} and load performance across React, Vue and Svelte viewers
  \end{itemize}}

\newcommand{\MirrorAI}{%
  \entry{Mirror AI: Deployable Personas}{Honorable mention in Supabase YC Hackathon}{Oct 2024 -- Dec 2024}
  \begin{itemize}
    \item Designed an \textbf{agentic LLM architecture} with LangGraph to mirror user personalities, creating interactive digital
          personas, and deployed it as a full-stack platform on Supabase and Vercel with one-click publishing
  \end{itemize}}

\newcommand{\LLMReasoning}{%
  \entry{Improving LLM Reasoning for Numerical Problems}{\report{https://sudhansh6.github.io/assets/pdf/MathPrompter++.pdf}}{Sep 2024 -- Dec 2024}
  \begin{itemize}
    \item Enhanced MathPrompter (ACL 2023) with CoT, achieving \textbf{10\% higher accuracy} on Llama 3.1 1B where prior methods failed,
          and reduced hallucination rates significantly by integrating multi-step validation
  \end{itemize}}

\newcommand{\InverseRendering}{%
  \entry{Inverse Rendering with 2D Gaussian Splatting}{\report{https://sudhansh6.github.io/assets/pdf/CSE252D_Report.pdf}}{Mar 2024 -- May 2024}
  \begin{itemize}
    \item Developed a novel \textbf{inverse rendering} framework in CUDA to recover PBR properties of a scene using 2D Gaussian
          Splatting, improving normal map MAE by \textbf{15\%} over the current SOTA methods
  \end{itemize}}

% ---------------------------------------------------------------------------------------------------
% Publications: title on the left, venue on the right, one line each.
\newcommand{\publine}[3]{\par\noindent\href{#1}{#2}\hfill\hspace{1em}\side{#3}\par\addvspace{1.5pt}}
\newcommand{\PubSaturation}{\publine{https://doi.org/10.1016/j.sigpro.2023.109349}{A likelihood based method for compressive signal recovery under Gaussian and saturation noise}{Signal Processing, 2024}}
\newcommand{\PubPermutation}{\publine{https://ieeexplore.ieee.org/document/10888147}{Identification and correction of permutation errors in compressed sensing-based group testing}{ICASSP 2025}}
\newcommand{\PubRearrangement}{\publine{https://arxiv.org/abs/2411.10899}{Planning for tabletop object rearrangement}{ICRA 2025}}

% ---------------------------------------------------------------------------------------------------
% Education, skills, service
% Each profile sets \UCSDCourses, \IITBCourses and \SkillRows before \input{resume-common}.
\providecommand{\UCSDCourses}{}
\providecommand{\IITBCourses}{}
\providecommand{\SkillRows}{}
\newcommand{\coursework}[1]{\ifblank{#1}{}{\par\noindent{\fontsize{9.5}{11.5}\selectfont{\color{muted}Coursework:} #1}\par}}
\newcommand{\Education}{%
  \section{Education}
  \entry{University of California San Diego}{M.S. in Computer Science and Engineering (AI)}{GPA 4.00/4.00 \quad Sep 2023 -- Mar 2025}
  \coursework{\UCSDCourses}
  \entry{Indian Institute of Technology Bombay}{B.Tech.\ with Honors in Computer Science}{CPI 9.66/10 \quad Jul 2019 -- Jul 2023}
  \coursework{\IITBCourses}}

% \skill{Label}{Items} is one row of the skills table.
\newcommand{\skill}[2]{{\headfont\bfseries #1} & #2 \\[1.5pt]}
\newcommand{\cpp}{C\kern-0.04em\raisebox{0.06ex}{\small++}}
\newcommand{\Skills}{%
  \section{Skills}
  \begin{tabularx}{\linewidth}{@{}l@{\hspace{12pt}}>{\raggedright\arraybackslash}X@{}}
    \SkillRows
  \end{tabularx}}
.
\providecommand{\UCSDCourses}{}
\providecommand{\IITBCourses}{}
\providecommand{\SkillRows}{}
\newcommand{\coursework}[1]{\ifblank{#1}{}{\par\noindent{\fontsize{9.5}{11.5}\selectfont{\color{muted}Coursework:} #1}\par}}
\newcommand{\Education}{%
  \section{Education}
  \entry{University of California San Diego}{M.S. in Computer Science and Engineering (AI)}{GPA 4.00/4.00 \quad Sep 2023 -- Mar 2025}
  \coursework{\UCSDCourses}
  \entry{Indian Institute of Technology Bombay}{B.Tech.\ with Honors in Computer Science}{CPI 9.66/10 \quad Jul 2019 -- Jul 2023}
  \coursework{\IITBCourses}}

% \skill{Label}{Items} is one row of the skills table.
\newcommand{\skill}[2]{{\headfont\bfseries #1} & #2 \\[1.5pt]}
\newcommand{\cpp}{C\kern-0.04em\raisebox{0.06ex}{\small++}}
\newcommand{\Skills}{%
  \section{Skills}
  \begin{tabularx}{\linewidth}{@{}l@{\hspace{12pt}}>{\raggedright\arraybackslash}X@{}}
    \SkillRows
  \end{tabularx}}
 to fit one page.
\providecommand{\secgap}{12pt}   % above a section title
\providecommand{\rulegap}{5.5pt}   % section rule to first entry
\providecommand{\entrygap}{7pt}  % between entries
\providecommand{\headgap}{1.5pt} % entry heading to its bullets
\providecommand{\itemgap}{0.6pt} % between bullets of one entry
\titlespacing*{\section}{0pt}{\secgap}{\rulegap}

\newcommand{\squarebullet}{\raisebox{0.32ex}{\color{accent}\rule{0.5ex}{0.5ex}}}
\setlist[itemize]{leftmargin=12pt, labelsep=6pt, topsep=\headgap, itemsep=\itemgap, parsep=0pt, partopsep=0pt,
  label=\squarebullet}

% \entry{Organisation or title}{Role or detail (may be empty)}{Dates}; \sub{Role or detail}{Right-hand note}
\newcommand{\entry}[3]{%
  \par\iffirstentry\global\firstentryfalse\else\addvspace{\entrygap}\fi\noindent{\headfont\fontsize{10.8}{13}\selectfont\textbf{#1}}\ifblank{#2}{}{\ {\color{muted}|}\ #2}\hfill\hspace{1em}\side{#3}\par}
\newcommand{\sub}[2]{%
  \noindent\textit{#1}\hfill{\color{muted}\textit{#2}}\par}
\newcommand{\report}[1]{\href{#1}{\color{accent}Report}}
\newcommand{\me}{\textbf{S. Peddabomma}}

% ---------------------------------------------------------------------------------------------------
% Header
\newcommand{\Header}{%
  \begin{center}
    {\namefont\fontsize{26}{28}\selectfont Sudhansh Peddabomma}\\[6pt]
    {\headfont\fontsize{9.5}{11}\selectfont\color{muted}
      \href{mailto:sudhansh6@gmail.com}{\faEnvelope[regular]\ sudhansh6@gmail.com} \quad
      \href{https://sudhansh6.github.io}{\faGlobe\ sudhansh6.github.io} \quad
      \href{https://www.linkedin.com/in/sudhansh-peddabomma-73a208192/}{\faLinkedin\ sudhansh-peddabomma} \quad
      \href{https://github.com/Sudhansh6}{\faGithub\ Sudhansh6}}
  \end{center}\vspace{-4pt}}

% ---------------------------------------------------------------------------------------------------
% Glean bullets, shared by both profiles
\newcommand{\GleanSlides}{\item Led \textbf{AI slide generation} in Glean Assistant, now creating \textbf{180K decks a week} for \textbf{50K weekly active users} at enterprises}
\newcommand{\GleanSlidesPerf}{\item Improved the slides harness, raising reliability to \textbf{97\%}, cutting median latency \textbf{26\%} and cost per deck by \textbf{half}
  in four months}
\newcommand{\GleanSlidesOrch}{\item Architected a \textbf{slide orchestration system} that reuses relevant slides from a company's existing decks and generates
  new ones in its brand and template, so teams can refine decks with each other and Glean in full context, without leaving the product}
\newcommand{\GleanAgents}{\item Moved Glean's \textbf{workflow agents} from one-shot planning to an \textbf{agentic tool-use loop}, greatly improving agentic
  capabilities and instruction following, with \textbf{18\%} higher answer completeness and \textbf{17\%} fewer tokens}
\newcommand{\GleanTeamChat}{\item Engineered the core of \textbf{Team chat}, letting teammates work with Glean in one shared conversation and co-edit its artifacts}
\newcommand{\GleanVideo}{\item Shipped \textbf{video artifacts} end to end, generating narrated videos in chat with streamed playback and admin rollout controls}
\newcommand{\GleanEvals}{\item Created the \textbf{evaluations and metrics} that guide model choices and rollouts across agents, slides and interactive artifacts}

% ---------------------------------------------------------------------------------------------------
% Experience
% \groupbreak separates related groups of bullets inside one long entry (slides, then platform work).
\providecommand{\groupgap}{3.5pt}
\newcommand{\groupbreak}{\vspace{\groupgap}}
\newcommand{\GleanAI}{%
  \entry{Glean}{Software Engineer}{Apr 2025 -- Present}
  \begin{itemize}
    \GleanSlides
    \GleanSlidesOrch
    \GleanSlidesPerf \groupbreak
    \GleanTeamChat
    \GleanVideo
    \GleanAgents
    \GleanEvals
  \end{itemize}}

\newcommand{\GleanVision}{%
  \entry{Glean}{Software Engineer}{Apr 2025 -- Present}
  \begin{itemize}
    \GleanSlides
    \GleanSlidesOrch
    \GleanSlidesPerf \groupbreak
    \GleanVideo
    \GleanEvals
  \end{itemize}}

\newcommand{\Duality}{%
  \entry{Duality AI}{Computer Vision Intern}{Jun 2024 -- Sep 2024}
  \begin{itemize}
    \item Designed automated \textbf{3D reconstruction} techniques for featureless objects, reducing digital-twin generation time
          by \textbf{40\%}
    \item Collaborated with \textbf{Autodesk} to validate Unreal Engine simulations for robotics tasks; leveraging structured
          domain randomization to reduce Sim2Real gap and increase mAP-50 by \textbf{15\%} for object detection and segmentation
  \end{itemize}}

\newcommand{\Microsoft}{%
  \entry{Microsoft}{Data and Applied Scientist Intern}{May 2022 -- Jul 2022}
  \begin{itemize}
    \item Developed a \textbf{decision-tree ranker} to recommend emails without user queries, improving \textbf{Outlook} search capabilities
    \item Proposed hierarchical feature-sets for the ranker, reducing latency for recommendations and improving recall
  \end{itemize}}

% ---------------------------------------------------------------------------------------------------
% Research
\newcommand{\HomeRobots}{%
  \entry{Real-time 3D Perception for Home Robots}{UC San Diego}{Sep 2023 -- Sep 2024}
  \begin{itemize}
    \item Investigated real-time dense \textbf{visual SLAM} methods using NeRFs and Gaussian Splatting for indoor robot navigation
    \item Integrated object segmentation, grasp-pose estimation, and 3D mapping on the Fetch robot via ROS, demonstrating
          a novel \textbf{tabletop object rearrangement} algorithm that reduced cost by \textbf{20\%} compared to the state-of-the-art
  \end{itemize}}

\newcommand{\Tomography}{%
  \entry{3D Tomography with Primal-Dual Neural Networks}{University College London}{May 2021 -- Jul 2023}
  \begin{itemize}
    \item Developed a stochastic neural-network architecture of primal-dual algorithm for online reconstruction of 3D
          volumes from tomographic projections, with \textbf{99.6\,\% structural similarity} in low-dosage conditions
    \item Built a Python library with custom gradient operators for reconstructing volumes in a single pass, reducing
          compute requirements by up to \textbf{5x} over SOTA learning-based approaches for cone vector tomography
  \end{itemize}}

% ---------------------------------------------------------------------------------------------------
% Projects
\newcommand{\OOXML}{%
  \entry{ooxml}{Open-source contributor \textbar\ \href{https://github.com/ChristopherVR/ooxml}{\color{accent}GitHub}}{Sep 2026 -- Present}
  \begin{itemize}
    \item Contributed to an open-source \textbf{TypeScript} engine that renders and edits PowerPoint and Word files in the browser,
          improving editing fidelity, \textbf{real-time collaboration} and load performance across React, Vue and Svelte viewers
  \end{itemize}}

\newcommand{\MirrorAI}{%
  \entry{Mirror AI: Deployable Personas}{Honorable mention in Supabase YC Hackathon}{Oct 2024 -- Dec 2024}
  \begin{itemize}
    \item Designed an \textbf{agentic LLM architecture} with LangGraph to mirror user personalities, creating interactive digital
          personas, and deployed it as a full-stack platform on Supabase and Vercel with one-click publishing
  \end{itemize}}

\newcommand{\LLMReasoning}{%
  \entry{Improving LLM Reasoning for Numerical Problems}{\report{https://sudhansh6.github.io/assets/pdf/MathPrompter++.pdf}}{Sep 2024 -- Dec 2024}
  \begin{itemize}
    \item Enhanced MathPrompter (ACL 2023) with CoT, achieving \textbf{10\% higher accuracy} on Llama 3.1 1B where prior methods failed,
          and reduced hallucination rates significantly by integrating multi-step validation
  \end{itemize}}

\newcommand{\InverseRendering}{%
  \entry{Inverse Rendering with 2D Gaussian Splatting}{\report{https://sudhansh6.github.io/assets/pdf/CSE252D_Report.pdf}}{Mar 2024 -- May 2024}
  \begin{itemize}
    \item Developed a novel \textbf{inverse rendering} framework in CUDA to recover PBR properties of a scene using 2D Gaussian
          Splatting, improving normal map MAE by \textbf{15\%} over the current SOTA methods
  \end{itemize}}

% ---------------------------------------------------------------------------------------------------
% Publications: title on the left, venue on the right, one line each.
\newcommand{\publine}[3]{\par\noindent\href{#1}{#2}\hfill\hspace{1em}\side{#3}\par\addvspace{1.5pt}}
\newcommand{\PubSaturation}{\publine{https://doi.org/10.1016/j.sigpro.2023.109349}{A likelihood based method for compressive signal recovery under Gaussian and saturation noise}{Signal Processing, 2024}}
\newcommand{\PubPermutation}{\publine{https://ieeexplore.ieee.org/document/10888147}{Identification and correction of permutation errors in compressed sensing-based group testing}{ICASSP 2025}}
\newcommand{\PubRearrangement}{\publine{https://arxiv.org/abs/2411.10899}{Planning for tabletop object rearrangement}{ICRA 2025}}

% ---------------------------------------------------------------------------------------------------
% Education, skills, service
% Each profile sets \UCSDCourses, \IITBCourses and \SkillRows before % Shared preamble and content for Sudhansh Peddabomma's résumés.
% Each profile (SudhanshPeddabomma-AI.tex, SudhanshPeddabomma-Vision.tex) inputs this file and picks
% which blocks to show and in what order, so every bullet is written once.
% Layout rule: every bullet fills one line, or runs 1.5 to 2 lines. Check with: python3 check_lines.py <pdf>
\documentclass[10pt,letterpaper]{article}

\usepackage[margin=0.55in,top=0.5in,bottom=0.45in]{geometry}
\usepackage[T1]{fontenc}
% Fonts: \resumefont loads the body font; \headfont styles the name and section titles.
% Override either before % Shared preamble and content for Sudhansh Peddabomma's résumés.
% Each profile (SudhanshPeddabomma-AI.tex, SudhanshPeddabomma-Vision.tex) inputs this file and picks
% which blocks to show and in what order, so every bullet is written once.
% Layout rule: every bullet fills one line, or runs 1.5 to 2 lines. Check with: python3 check_lines.py <pdf>
\documentclass[10pt,letterpaper]{article}

\usepackage[margin=0.55in,top=0.5in,bottom=0.45in]{geometry}
\usepackage[T1]{fontenc}
% Fonts: \resumefont loads the body font; \headfont styles the name and section titles.
% Override either before \input{resume-common} to try another pairing.
% Default: Charter body with Source Sans headings. All-sans alternative:
%   \def\resumefont{\usepackage[default,semibold]{sourcesanspro}}\def\headfont{}
\providecommand{\resumefont}{\usepackage[scaled=0.96]{XCharter}\usepackage[semibold,scale=0.96]{sourcesanspro}}
\providecommand{\headfont}{\sffamily}
\providecommand{\namefont}{\rmfamily\scshape}
\resumefont
\usepackage[final]{microtype}
\usepackage[dvipsnames]{xcolor}
\usepackage{enumitem}
\usepackage{titlesec}
\usepackage{tabularx}
\usepackage{fontawesome5}
\usepackage{etoolbox}
\usepackage[hidelinks]{hyperref}

% The site's green, darkened for print.
\definecolor{accent}{HTML}{3A7A45}
\definecolor{muted}{HTML}{5C675D}

\pagestyle{empty}
\setlength{\parindent}{0pt}
\setlength{\tabcolsep}{0pt}
\providecommand{\leading}{1.1}
\linespread{\leading}

\hypersetup{pdfauthor={Sudhansh Peddabomma}}

% Type scale: name 26pt small caps, section titles 12.5pt caps, entry titles 10.8pt, body 10pt,
% side notes 9pt.
\newcommand{\side}[1]{{\fontsize{9}{11}\selectfont\color{muted}#1}}
\newcommand{\sectitle}[1]{\textls[60]{\MakeUppercase{#1}}}
\titleformat{\section}{\headfont\color{accent}\bfseries\fontsize{12.5}{14}\selectfont}{}{0pt}{\sectitle}[\vspace{1.5pt}{\color{accent}\titlerule[0.8pt]}\global\firstentrytrue]
% The first entry after a section title skips \entrygap, so only \rulegap sits under the rule.
\newif\iffirstentry

% Vertical rhythm, the same idea as the earlier hand-tuned résumés: bullets inside an entry sit
% tight together, entries are separated by a clear gap, and sections by a larger one. A profile
% may override these before \input{resume-common} to fit one page.
\providecommand{\secgap}{12pt}   % above a section title
\providecommand{\rulegap}{5.5pt}   % section rule to first entry
\providecommand{\entrygap}{7pt}  % between entries
\providecommand{\headgap}{1.5pt} % entry heading to its bullets
\providecommand{\itemgap}{0.6pt} % between bullets of one entry
\titlespacing*{\section}{0pt}{\secgap}{\rulegap}

\newcommand{\squarebullet}{\raisebox{0.32ex}{\color{accent}\rule{0.5ex}{0.5ex}}}
\setlist[itemize]{leftmargin=12pt, labelsep=6pt, topsep=\headgap, itemsep=\itemgap, parsep=0pt, partopsep=0pt,
  label=\squarebullet}

% \entry{Organisation or title}{Role or detail (may be empty)}{Dates}; \sub{Role or detail}{Right-hand note}
\newcommand{\entry}[3]{%
  \par\iffirstentry\global\firstentryfalse\else\addvspace{\entrygap}\fi\noindent{\headfont\fontsize{10.8}{13}\selectfont\textbf{#1}}\ifblank{#2}{}{\ {\color{muted}|}\ #2}\hfill\hspace{1em}\side{#3}\par}
\newcommand{\sub}[2]{%
  \noindent\textit{#1}\hfill{\color{muted}\textit{#2}}\par}
\newcommand{\report}[1]{\href{#1}{\color{accent}Report}}
\newcommand{\me}{\textbf{S. Peddabomma}}

% ---------------------------------------------------------------------------------------------------
% Header
\newcommand{\Header}{%
  \begin{center}
    {\namefont\fontsize{26}{28}\selectfont Sudhansh Peddabomma}\\[6pt]
    {\headfont\fontsize{9.5}{11}\selectfont\color{muted}
      \href{mailto:sudhansh6@gmail.com}{\faEnvelope[regular]\ sudhansh6@gmail.com} \quad
      \href{https://sudhansh6.github.io}{\faGlobe\ sudhansh6.github.io} \quad
      \href{https://www.linkedin.com/in/sudhansh-peddabomma-73a208192/}{\faLinkedin\ sudhansh-peddabomma} \quad
      \href{https://github.com/Sudhansh6}{\faGithub\ Sudhansh6}}
  \end{center}\vspace{-4pt}}

% ---------------------------------------------------------------------------------------------------
% Glean bullets, shared by both profiles
\newcommand{\GleanSlides}{\item Led \textbf{AI slide generation} in Glean Assistant, now creating \textbf{180K decks a week} for \textbf{50K weekly active users} at enterprises}
\newcommand{\GleanSlidesPerf}{\item Improved the slides harness, raising reliability to \textbf{97\%}, cutting median latency \textbf{26\%} and cost per deck by \textbf{half}
  in four months}
\newcommand{\GleanSlidesOrch}{\item Architected a \textbf{slide orchestration system} that reuses relevant slides from a company's existing decks and generates
  new ones in its brand and template, so teams can refine decks with each other and Glean in full context, without leaving the product}
\newcommand{\GleanAgents}{\item Moved Glean's \textbf{workflow agents} from one-shot planning to an \textbf{agentic tool-use loop}, greatly improving agentic
  capabilities and instruction following, with \textbf{18\%} higher answer completeness and \textbf{17\%} fewer tokens}
\newcommand{\GleanTeamChat}{\item Engineered the core of \textbf{Team chat}, letting teammates work with Glean in one shared conversation and co-edit its artifacts}
\newcommand{\GleanVideo}{\item Shipped \textbf{video artifacts} end to end, generating narrated videos in chat with streamed playback and admin rollout controls}
\newcommand{\GleanEvals}{\item Created the \textbf{evaluations and metrics} that guide model choices and rollouts across agents, slides and interactive artifacts}

% ---------------------------------------------------------------------------------------------------
% Experience
% \groupbreak separates related groups of bullets inside one long entry (slides, then platform work).
\providecommand{\groupgap}{3.5pt}
\newcommand{\groupbreak}{\vspace{\groupgap}}
\newcommand{\GleanAI}{%
  \entry{Glean}{Software Engineer}{Apr 2025 -- Present}
  \begin{itemize}
    \GleanSlides
    \GleanSlidesOrch
    \GleanSlidesPerf \groupbreak
    \GleanTeamChat
    \GleanVideo
    \GleanAgents
    \GleanEvals
  \end{itemize}}

\newcommand{\GleanVision}{%
  \entry{Glean}{Software Engineer}{Apr 2025 -- Present}
  \begin{itemize}
    \GleanSlides
    \GleanSlidesOrch
    \GleanSlidesPerf \groupbreak
    \GleanVideo
    \GleanEvals
  \end{itemize}}

\newcommand{\Duality}{%
  \entry{Duality AI}{Computer Vision Intern}{Jun 2024 -- Sep 2024}
  \begin{itemize}
    \item Designed automated \textbf{3D reconstruction} techniques for featureless objects, reducing digital-twin generation time
          by \textbf{40\%}
    \item Collaborated with \textbf{Autodesk} to validate Unreal Engine simulations for robotics tasks; leveraging structured
          domain randomization to reduce Sim2Real gap and increase mAP-50 by \textbf{15\%} for object detection and segmentation
  \end{itemize}}

\newcommand{\Microsoft}{%
  \entry{Microsoft}{Data and Applied Scientist Intern}{May 2022 -- Jul 2022}
  \begin{itemize}
    \item Developed a \textbf{decision-tree ranker} to recommend emails without user queries, improving \textbf{Outlook} search capabilities
    \item Proposed hierarchical feature-sets for the ranker, reducing latency for recommendations and improving recall
  \end{itemize}}

% ---------------------------------------------------------------------------------------------------
% Research
\newcommand{\HomeRobots}{%
  \entry{Real-time 3D Perception for Home Robots}{UC San Diego}{Sep 2023 -- Sep 2024}
  \begin{itemize}
    \item Investigated real-time dense \textbf{visual SLAM} methods using NeRFs and Gaussian Splatting for indoor robot navigation
    \item Integrated object segmentation, grasp-pose estimation, and 3D mapping on the Fetch robot via ROS, demonstrating
          a novel \textbf{tabletop object rearrangement} algorithm that reduced cost by \textbf{20\%} compared to the state-of-the-art
  \end{itemize}}

\newcommand{\Tomography}{%
  \entry{3D Tomography with Primal-Dual Neural Networks}{University College London}{May 2021 -- Jul 2023}
  \begin{itemize}
    \item Developed a stochastic neural-network architecture of primal-dual algorithm for online reconstruction of 3D
          volumes from tomographic projections, with \textbf{99.6\,\% structural similarity} in low-dosage conditions
    \item Built a Python library with custom gradient operators for reconstructing volumes in a single pass, reducing
          compute requirements by up to \textbf{5x} over SOTA learning-based approaches for cone vector tomography
  \end{itemize}}

% ---------------------------------------------------------------------------------------------------
% Projects
\newcommand{\OOXML}{%
  \entry{ooxml}{Open-source contributor \textbar\ \href{https://github.com/ChristopherVR/ooxml}{\color{accent}GitHub}}{Sep 2026 -- Present}
  \begin{itemize}
    \item Contributed to an open-source \textbf{TypeScript} engine that renders and edits PowerPoint and Word files in the browser,
          improving editing fidelity, \textbf{real-time collaboration} and load performance across React, Vue and Svelte viewers
  \end{itemize}}

\newcommand{\MirrorAI}{%
  \entry{Mirror AI: Deployable Personas}{Honorable mention in Supabase YC Hackathon}{Oct 2024 -- Dec 2024}
  \begin{itemize}
    \item Designed an \textbf{agentic LLM architecture} with LangGraph to mirror user personalities, creating interactive digital
          personas, and deployed it as a full-stack platform on Supabase and Vercel with one-click publishing
  \end{itemize}}

\newcommand{\LLMReasoning}{%
  \entry{Improving LLM Reasoning for Numerical Problems}{\report{https://sudhansh6.github.io/assets/pdf/MathPrompter++.pdf}}{Sep 2024 -- Dec 2024}
  \begin{itemize}
    \item Enhanced MathPrompter (ACL 2023) with CoT, achieving \textbf{10\% higher accuracy} on Llama 3.1 1B where prior methods failed,
          and reduced hallucination rates significantly by integrating multi-step validation
  \end{itemize}}

\newcommand{\InverseRendering}{%
  \entry{Inverse Rendering with 2D Gaussian Splatting}{\report{https://sudhansh6.github.io/assets/pdf/CSE252D_Report.pdf}}{Mar 2024 -- May 2024}
  \begin{itemize}
    \item Developed a novel \textbf{inverse rendering} framework in CUDA to recover PBR properties of a scene using 2D Gaussian
          Splatting, improving normal map MAE by \textbf{15\%} over the current SOTA methods
  \end{itemize}}

% ---------------------------------------------------------------------------------------------------
% Publications: title on the left, venue on the right, one line each.
\newcommand{\publine}[3]{\par\noindent\href{#1}{#2}\hfill\hspace{1em}\side{#3}\par\addvspace{1.5pt}}
\newcommand{\PubSaturation}{\publine{https://doi.org/10.1016/j.sigpro.2023.109349}{A likelihood based method for compressive signal recovery under Gaussian and saturation noise}{Signal Processing, 2024}}
\newcommand{\PubPermutation}{\publine{https://ieeexplore.ieee.org/document/10888147}{Identification and correction of permutation errors in compressed sensing-based group testing}{ICASSP 2025}}
\newcommand{\PubRearrangement}{\publine{https://arxiv.org/abs/2411.10899}{Planning for tabletop object rearrangement}{ICRA 2025}}

% ---------------------------------------------------------------------------------------------------
% Education, skills, service
% Each profile sets \UCSDCourses, \IITBCourses and \SkillRows before \input{resume-common}.
\providecommand{\UCSDCourses}{}
\providecommand{\IITBCourses}{}
\providecommand{\SkillRows}{}
\newcommand{\coursework}[1]{\ifblank{#1}{}{\par\noindent{\fontsize{9.5}{11.5}\selectfont{\color{muted}Coursework:} #1}\par}}
\newcommand{\Education}{%
  \section{Education}
  \entry{University of California San Diego}{M.S. in Computer Science and Engineering (AI)}{GPA 4.00/4.00 \quad Sep 2023 -- Mar 2025}
  \coursework{\UCSDCourses}
  \entry{Indian Institute of Technology Bombay}{B.Tech.\ with Honors in Computer Science}{CPI 9.66/10 \quad Jul 2019 -- Jul 2023}
  \coursework{\IITBCourses}}

% \skill{Label}{Items} is one row of the skills table.
\newcommand{\skill}[2]{{\headfont\bfseries #1} & #2 \\[1.5pt]}
\newcommand{\cpp}{C\kern-0.04em\raisebox{0.06ex}{\small++}}
\newcommand{\Skills}{%
  \section{Skills}
  \begin{tabularx}{\linewidth}{@{}l@{\hspace{12pt}}>{\raggedright\arraybackslash}X@{}}
    \SkillRows
  \end{tabularx}}
 to try another pairing.
% Default: Charter body with Source Sans headings. All-sans alternative:
%   \def\resumefont{\usepackage[default,semibold]{sourcesanspro}}\def\headfont{}
\providecommand{\resumefont}{\usepackage[scaled=0.96]{XCharter}\usepackage[semibold,scale=0.96]{sourcesanspro}}
\providecommand{\headfont}{\sffamily}
\providecommand{\namefont}{\rmfamily\scshape}
\resumefont
\usepackage[final]{microtype}
\usepackage[dvipsnames]{xcolor}
\usepackage{enumitem}
\usepackage{titlesec}
\usepackage{tabularx}
\usepackage{fontawesome5}
\usepackage{etoolbox}
\usepackage[hidelinks]{hyperref}

% The site's green, darkened for print.
\definecolor{accent}{HTML}{3A7A45}
\definecolor{muted}{HTML}{5C675D}

\pagestyle{empty}
\setlength{\parindent}{0pt}
\setlength{\tabcolsep}{0pt}
\providecommand{\leading}{1.1}
\linespread{\leading}

\hypersetup{pdfauthor={Sudhansh Peddabomma}}

% Type scale: name 26pt small caps, section titles 12.5pt caps, entry titles 10.8pt, body 10pt,
% side notes 9pt.
\newcommand{\side}[1]{{\fontsize{9}{11}\selectfont\color{muted}#1}}
\newcommand{\sectitle}[1]{\textls[60]{\MakeUppercase{#1}}}
\titleformat{\section}{\headfont\color{accent}\bfseries\fontsize{12.5}{14}\selectfont}{}{0pt}{\sectitle}[\vspace{1.5pt}{\color{accent}\titlerule[0.8pt]}\global\firstentrytrue]
% The first entry after a section title skips \entrygap, so only \rulegap sits under the rule.
\newif\iffirstentry

% Vertical rhythm, the same idea as the earlier hand-tuned résumés: bullets inside an entry sit
% tight together, entries are separated by a clear gap, and sections by a larger one. A profile
% may override these before % Shared preamble and content for Sudhansh Peddabomma's résumés.
% Each profile (SudhanshPeddabomma-AI.tex, SudhanshPeddabomma-Vision.tex) inputs this file and picks
% which blocks to show and in what order, so every bullet is written once.
% Layout rule: every bullet fills one line, or runs 1.5 to 2 lines. Check with: python3 check_lines.py <pdf>
\documentclass[10pt,letterpaper]{article}

\usepackage[margin=0.55in,top=0.5in,bottom=0.45in]{geometry}
\usepackage[T1]{fontenc}
% Fonts: \resumefont loads the body font; \headfont styles the name and section titles.
% Override either before \input{resume-common} to try another pairing.
% Default: Charter body with Source Sans headings. All-sans alternative:
%   \def\resumefont{\usepackage[default,semibold]{sourcesanspro}}\def\headfont{}
\providecommand{\resumefont}{\usepackage[scaled=0.96]{XCharter}\usepackage[semibold,scale=0.96]{sourcesanspro}}
\providecommand{\headfont}{\sffamily}
\providecommand{\namefont}{\rmfamily\scshape}
\resumefont
\usepackage[final]{microtype}
\usepackage[dvipsnames]{xcolor}
\usepackage{enumitem}
\usepackage{titlesec}
\usepackage{tabularx}
\usepackage{fontawesome5}
\usepackage{etoolbox}
\usepackage[hidelinks]{hyperref}

% The site's green, darkened for print.
\definecolor{accent}{HTML}{3A7A45}
\definecolor{muted}{HTML}{5C675D}

\pagestyle{empty}
\setlength{\parindent}{0pt}
\setlength{\tabcolsep}{0pt}
\providecommand{\leading}{1.1}
\linespread{\leading}

\hypersetup{pdfauthor={Sudhansh Peddabomma}}

% Type scale: name 26pt small caps, section titles 12.5pt caps, entry titles 10.8pt, body 10pt,
% side notes 9pt.
\newcommand{\side}[1]{{\fontsize{9}{11}\selectfont\color{muted}#1}}
\newcommand{\sectitle}[1]{\textls[60]{\MakeUppercase{#1}}}
\titleformat{\section}{\headfont\color{accent}\bfseries\fontsize{12.5}{14}\selectfont}{}{0pt}{\sectitle}[\vspace{1.5pt}{\color{accent}\titlerule[0.8pt]}\global\firstentrytrue]
% The first entry after a section title skips \entrygap, so only \rulegap sits under the rule.
\newif\iffirstentry

% Vertical rhythm, the same idea as the earlier hand-tuned résumés: bullets inside an entry sit
% tight together, entries are separated by a clear gap, and sections by a larger one. A profile
% may override these before \input{resume-common} to fit one page.
\providecommand{\secgap}{12pt}   % above a section title
\providecommand{\rulegap}{5.5pt}   % section rule to first entry
\providecommand{\entrygap}{7pt}  % between entries
\providecommand{\headgap}{1.5pt} % entry heading to its bullets
\providecommand{\itemgap}{0.6pt} % between bullets of one entry
\titlespacing*{\section}{0pt}{\secgap}{\rulegap}

\newcommand{\squarebullet}{\raisebox{0.32ex}{\color{accent}\rule{0.5ex}{0.5ex}}}
\setlist[itemize]{leftmargin=12pt, labelsep=6pt, topsep=\headgap, itemsep=\itemgap, parsep=0pt, partopsep=0pt,
  label=\squarebullet}

% \entry{Organisation or title}{Role or detail (may be empty)}{Dates}; \sub{Role or detail}{Right-hand note}
\newcommand{\entry}[3]{%
  \par\iffirstentry\global\firstentryfalse\else\addvspace{\entrygap}\fi\noindent{\headfont\fontsize{10.8}{13}\selectfont\textbf{#1}}\ifblank{#2}{}{\ {\color{muted}|}\ #2}\hfill\hspace{1em}\side{#3}\par}
\newcommand{\sub}[2]{%
  \noindent\textit{#1}\hfill{\color{muted}\textit{#2}}\par}
\newcommand{\report}[1]{\href{#1}{\color{accent}Report}}
\newcommand{\me}{\textbf{S. Peddabomma}}

% ---------------------------------------------------------------------------------------------------
% Header
\newcommand{\Header}{%
  \begin{center}
    {\namefont\fontsize{26}{28}\selectfont Sudhansh Peddabomma}\\[6pt]
    {\headfont\fontsize{9.5}{11}\selectfont\color{muted}
      \href{mailto:sudhansh6@gmail.com}{\faEnvelope[regular]\ sudhansh6@gmail.com} \quad
      \href{https://sudhansh6.github.io}{\faGlobe\ sudhansh6.github.io} \quad
      \href{https://www.linkedin.com/in/sudhansh-peddabomma-73a208192/}{\faLinkedin\ sudhansh-peddabomma} \quad
      \href{https://github.com/Sudhansh6}{\faGithub\ Sudhansh6}}
  \end{center}\vspace{-4pt}}

% ---------------------------------------------------------------------------------------------------
% Glean bullets, shared by both profiles
\newcommand{\GleanSlides}{\item Led \textbf{AI slide generation} in Glean Assistant, now creating \textbf{180K decks a week} for \textbf{50K weekly active users} at enterprises}
\newcommand{\GleanSlidesPerf}{\item Improved the slides harness, raising reliability to \textbf{97\%}, cutting median latency \textbf{26\%} and cost per deck by \textbf{half}
  in four months}
\newcommand{\GleanSlidesOrch}{\item Architected a \textbf{slide orchestration system} that reuses relevant slides from a company's existing decks and generates
  new ones in its brand and template, so teams can refine decks with each other and Glean in full context, without leaving the product}
\newcommand{\GleanAgents}{\item Moved Glean's \textbf{workflow agents} from one-shot planning to an \textbf{agentic tool-use loop}, greatly improving agentic
  capabilities and instruction following, with \textbf{18\%} higher answer completeness and \textbf{17\%} fewer tokens}
\newcommand{\GleanTeamChat}{\item Engineered the core of \textbf{Team chat}, letting teammates work with Glean in one shared conversation and co-edit its artifacts}
\newcommand{\GleanVideo}{\item Shipped \textbf{video artifacts} end to end, generating narrated videos in chat with streamed playback and admin rollout controls}
\newcommand{\GleanEvals}{\item Created the \textbf{evaluations and metrics} that guide model choices and rollouts across agents, slides and interactive artifacts}

% ---------------------------------------------------------------------------------------------------
% Experience
% \groupbreak separates related groups of bullets inside one long entry (slides, then platform work).
\providecommand{\groupgap}{3.5pt}
\newcommand{\groupbreak}{\vspace{\groupgap}}
\newcommand{\GleanAI}{%
  \entry{Glean}{Software Engineer}{Apr 2025 -- Present}
  \begin{itemize}
    \GleanSlides
    \GleanSlidesOrch
    \GleanSlidesPerf \groupbreak
    \GleanTeamChat
    \GleanVideo
    \GleanAgents
    \GleanEvals
  \end{itemize}}

\newcommand{\GleanVision}{%
  \entry{Glean}{Software Engineer}{Apr 2025 -- Present}
  \begin{itemize}
    \GleanSlides
    \GleanSlidesOrch
    \GleanSlidesPerf \groupbreak
    \GleanVideo
    \GleanEvals
  \end{itemize}}

\newcommand{\Duality}{%
  \entry{Duality AI}{Computer Vision Intern}{Jun 2024 -- Sep 2024}
  \begin{itemize}
    \item Designed automated \textbf{3D reconstruction} techniques for featureless objects, reducing digital-twin generation time
          by \textbf{40\%}
    \item Collaborated with \textbf{Autodesk} to validate Unreal Engine simulations for robotics tasks; leveraging structured
          domain randomization to reduce Sim2Real gap and increase mAP-50 by \textbf{15\%} for object detection and segmentation
  \end{itemize}}

\newcommand{\Microsoft}{%
  \entry{Microsoft}{Data and Applied Scientist Intern}{May 2022 -- Jul 2022}
  \begin{itemize}
    \item Developed a \textbf{decision-tree ranker} to recommend emails without user queries, improving \textbf{Outlook} search capabilities
    \item Proposed hierarchical feature-sets for the ranker, reducing latency for recommendations and improving recall
  \end{itemize}}

% ---------------------------------------------------------------------------------------------------
% Research
\newcommand{\HomeRobots}{%
  \entry{Real-time 3D Perception for Home Robots}{UC San Diego}{Sep 2023 -- Sep 2024}
  \begin{itemize}
    \item Investigated real-time dense \textbf{visual SLAM} methods using NeRFs and Gaussian Splatting for indoor robot navigation
    \item Integrated object segmentation, grasp-pose estimation, and 3D mapping on the Fetch robot via ROS, demonstrating
          a novel \textbf{tabletop object rearrangement} algorithm that reduced cost by \textbf{20\%} compared to the state-of-the-art
  \end{itemize}}

\newcommand{\Tomography}{%
  \entry{3D Tomography with Primal-Dual Neural Networks}{University College London}{May 2021 -- Jul 2023}
  \begin{itemize}
    \item Developed a stochastic neural-network architecture of primal-dual algorithm for online reconstruction of 3D
          volumes from tomographic projections, with \textbf{99.6\,\% structural similarity} in low-dosage conditions
    \item Built a Python library with custom gradient operators for reconstructing volumes in a single pass, reducing
          compute requirements by up to \textbf{5x} over SOTA learning-based approaches for cone vector tomography
  \end{itemize}}

% ---------------------------------------------------------------------------------------------------
% Projects
\newcommand{\OOXML}{%
  \entry{ooxml}{Open-source contributor \textbar\ \href{https://github.com/ChristopherVR/ooxml}{\color{accent}GitHub}}{Sep 2026 -- Present}
  \begin{itemize}
    \item Contributed to an open-source \textbf{TypeScript} engine that renders and edits PowerPoint and Word files in the browser,
          improving editing fidelity, \textbf{real-time collaboration} and load performance across React, Vue and Svelte viewers
  \end{itemize}}

\newcommand{\MirrorAI}{%
  \entry{Mirror AI: Deployable Personas}{Honorable mention in Supabase YC Hackathon}{Oct 2024 -- Dec 2024}
  \begin{itemize}
    \item Designed an \textbf{agentic LLM architecture} with LangGraph to mirror user personalities, creating interactive digital
          personas, and deployed it as a full-stack platform on Supabase and Vercel with one-click publishing
  \end{itemize}}

\newcommand{\LLMReasoning}{%
  \entry{Improving LLM Reasoning for Numerical Problems}{\report{https://sudhansh6.github.io/assets/pdf/MathPrompter++.pdf}}{Sep 2024 -- Dec 2024}
  \begin{itemize}
    \item Enhanced MathPrompter (ACL 2023) with CoT, achieving \textbf{10\% higher accuracy} on Llama 3.1 1B where prior methods failed,
          and reduced hallucination rates significantly by integrating multi-step validation
  \end{itemize}}

\newcommand{\InverseRendering}{%
  \entry{Inverse Rendering with 2D Gaussian Splatting}{\report{https://sudhansh6.github.io/assets/pdf/CSE252D_Report.pdf}}{Mar 2024 -- May 2024}
  \begin{itemize}
    \item Developed a novel \textbf{inverse rendering} framework in CUDA to recover PBR properties of a scene using 2D Gaussian
          Splatting, improving normal map MAE by \textbf{15\%} over the current SOTA methods
  \end{itemize}}

% ---------------------------------------------------------------------------------------------------
% Publications: title on the left, venue on the right, one line each.
\newcommand{\publine}[3]{\par\noindent\href{#1}{#2}\hfill\hspace{1em}\side{#3}\par\addvspace{1.5pt}}
\newcommand{\PubSaturation}{\publine{https://doi.org/10.1016/j.sigpro.2023.109349}{A likelihood based method for compressive signal recovery under Gaussian and saturation noise}{Signal Processing, 2024}}
\newcommand{\PubPermutation}{\publine{https://ieeexplore.ieee.org/document/10888147}{Identification and correction of permutation errors in compressed sensing-based group testing}{ICASSP 2025}}
\newcommand{\PubRearrangement}{\publine{https://arxiv.org/abs/2411.10899}{Planning for tabletop object rearrangement}{ICRA 2025}}

% ---------------------------------------------------------------------------------------------------
% Education, skills, service
% Each profile sets \UCSDCourses, \IITBCourses and \SkillRows before \input{resume-common}.
\providecommand{\UCSDCourses}{}
\providecommand{\IITBCourses}{}
\providecommand{\SkillRows}{}
\newcommand{\coursework}[1]{\ifblank{#1}{}{\par\noindent{\fontsize{9.5}{11.5}\selectfont{\color{muted}Coursework:} #1}\par}}
\newcommand{\Education}{%
  \section{Education}
  \entry{University of California San Diego}{M.S. in Computer Science and Engineering (AI)}{GPA 4.00/4.00 \quad Sep 2023 -- Mar 2025}
  \coursework{\UCSDCourses}
  \entry{Indian Institute of Technology Bombay}{B.Tech.\ with Honors in Computer Science}{CPI 9.66/10 \quad Jul 2019 -- Jul 2023}
  \coursework{\IITBCourses}}

% \skill{Label}{Items} is one row of the skills table.
\newcommand{\skill}[2]{{\headfont\bfseries #1} & #2 \\[1.5pt]}
\newcommand{\cpp}{C\kern-0.04em\raisebox{0.06ex}{\small++}}
\newcommand{\Skills}{%
  \section{Skills}
  \begin{tabularx}{\linewidth}{@{}l@{\hspace{12pt}}>{\raggedright\arraybackslash}X@{}}
    \SkillRows
  \end{tabularx}}
 to fit one page.
\providecommand{\secgap}{12pt}   % above a section title
\providecommand{\rulegap}{5.5pt}   % section rule to first entry
\providecommand{\entrygap}{7pt}  % between entries
\providecommand{\headgap}{1.5pt} % entry heading to its bullets
\providecommand{\itemgap}{0.6pt} % between bullets of one entry
\titlespacing*{\section}{0pt}{\secgap}{\rulegap}

\newcommand{\squarebullet}{\raisebox{0.32ex}{\color{accent}\rule{0.5ex}{0.5ex}}}
\setlist[itemize]{leftmargin=12pt, labelsep=6pt, topsep=\headgap, itemsep=\itemgap, parsep=0pt, partopsep=0pt,
  label=\squarebullet}

% \entry{Organisation or title}{Role or detail (may be empty)}{Dates}; \sub{Role or detail}{Right-hand note}
\newcommand{\entry}[3]{%
  \par\iffirstentry\global\firstentryfalse\else\addvspace{\entrygap}\fi\noindent{\headfont\fontsize{10.8}{13}\selectfont\textbf{#1}}\ifblank{#2}{}{\ {\color{muted}|}\ #2}\hfill\hspace{1em}\side{#3}\par}
\newcommand{\sub}[2]{%
  \noindent\textit{#1}\hfill{\color{muted}\textit{#2}}\par}
\newcommand{\report}[1]{\href{#1}{\color{accent}Report}}
\newcommand{\me}{\textbf{S. Peddabomma}}

% ---------------------------------------------------------------------------------------------------
% Header
\newcommand{\Header}{%
  \begin{center}
    {\namefont\fontsize{26}{28}\selectfont Sudhansh Peddabomma}\\[6pt]
    {\headfont\fontsize{9.5}{11}\selectfont\color{muted}
      \href{mailto:sudhansh6@gmail.com}{\faEnvelope[regular]\ sudhansh6@gmail.com} \quad
      \href{https://sudhansh6.github.io}{\faGlobe\ sudhansh6.github.io} \quad
      \href{https://www.linkedin.com/in/sudhansh-peddabomma-73a208192/}{\faLinkedin\ sudhansh-peddabomma} \quad
      \href{https://github.com/Sudhansh6}{\faGithub\ Sudhansh6}}
  \end{center}\vspace{-4pt}}

% ---------------------------------------------------------------------------------------------------
% Glean bullets, shared by both profiles
\newcommand{\GleanSlides}{\item Led \textbf{AI slide generation} in Glean Assistant, now creating \textbf{180K decks a week} for \textbf{50K weekly active users} at enterprises}
\newcommand{\GleanSlidesPerf}{\item Improved the slides harness, raising reliability to \textbf{97\%}, cutting median latency \textbf{26\%} and cost per deck by \textbf{half}
  in four months}
\newcommand{\GleanSlidesOrch}{\item Architected a \textbf{slide orchestration system} that reuses relevant slides from a company's existing decks and generates
  new ones in its brand and template, so teams can refine decks with each other and Glean in full context, without leaving the product}
\newcommand{\GleanAgents}{\item Moved Glean's \textbf{workflow agents} from one-shot planning to an \textbf{agentic tool-use loop}, greatly improving agentic
  capabilities and instruction following, with \textbf{18\%} higher answer completeness and \textbf{17\%} fewer tokens}
\newcommand{\GleanTeamChat}{\item Engineered the core of \textbf{Team chat}, letting teammates work with Glean in one shared conversation and co-edit its artifacts}
\newcommand{\GleanVideo}{\item Shipped \textbf{video artifacts} end to end, generating narrated videos in chat with streamed playback and admin rollout controls}
\newcommand{\GleanEvals}{\item Created the \textbf{evaluations and metrics} that guide model choices and rollouts across agents, slides and interactive artifacts}

% ---------------------------------------------------------------------------------------------------
% Experience
% \groupbreak separates related groups of bullets inside one long entry (slides, then platform work).
\providecommand{\groupgap}{3.5pt}
\newcommand{\groupbreak}{\vspace{\groupgap}}
\newcommand{\GleanAI}{%
  \entry{Glean}{Software Engineer}{Apr 2025 -- Present}
  \begin{itemize}
    \GleanSlides
    \GleanSlidesOrch
    \GleanSlidesPerf \groupbreak
    \GleanTeamChat
    \GleanVideo
    \GleanAgents
    \GleanEvals
  \end{itemize}}

\newcommand{\GleanVision}{%
  \entry{Glean}{Software Engineer}{Apr 2025 -- Present}
  \begin{itemize}
    \GleanSlides
    \GleanSlidesOrch
    \GleanSlidesPerf \groupbreak
    \GleanVideo
    \GleanEvals
  \end{itemize}}

\newcommand{\Duality}{%
  \entry{Duality AI}{Computer Vision Intern}{Jun 2024 -- Sep 2024}
  \begin{itemize}
    \item Designed automated \textbf{3D reconstruction} techniques for featureless objects, reducing digital-twin generation time
          by \textbf{40\%}
    \item Collaborated with \textbf{Autodesk} to validate Unreal Engine simulations for robotics tasks; leveraging structured
          domain randomization to reduce Sim2Real gap and increase mAP-50 by \textbf{15\%} for object detection and segmentation
  \end{itemize}}

\newcommand{\Microsoft}{%
  \entry{Microsoft}{Data and Applied Scientist Intern}{May 2022 -- Jul 2022}
  \begin{itemize}
    \item Developed a \textbf{decision-tree ranker} to recommend emails without user queries, improving \textbf{Outlook} search capabilities
    \item Proposed hierarchical feature-sets for the ranker, reducing latency for recommendations and improving recall
  \end{itemize}}

% ---------------------------------------------------------------------------------------------------
% Research
\newcommand{\HomeRobots}{%
  \entry{Real-time 3D Perception for Home Robots}{UC San Diego}{Sep 2023 -- Sep 2024}
  \begin{itemize}
    \item Investigated real-time dense \textbf{visual SLAM} methods using NeRFs and Gaussian Splatting for indoor robot navigation
    \item Integrated object segmentation, grasp-pose estimation, and 3D mapping on the Fetch robot via ROS, demonstrating
          a novel \textbf{tabletop object rearrangement} algorithm that reduced cost by \textbf{20\%} compared to the state-of-the-art
  \end{itemize}}

\newcommand{\Tomography}{%
  \entry{3D Tomography with Primal-Dual Neural Networks}{University College London}{May 2021 -- Jul 2023}
  \begin{itemize}
    \item Developed a stochastic neural-network architecture of primal-dual algorithm for online reconstruction of 3D
          volumes from tomographic projections, with \textbf{99.6\,\% structural similarity} in low-dosage conditions
    \item Built a Python library with custom gradient operators for reconstructing volumes in a single pass, reducing
          compute requirements by up to \textbf{5x} over SOTA learning-based approaches for cone vector tomography
  \end{itemize}}

% ---------------------------------------------------------------------------------------------------
% Projects
\newcommand{\OOXML}{%
  \entry{ooxml}{Open-source contributor \textbar\ \href{https://github.com/ChristopherVR/ooxml}{\color{accent}GitHub}}{Sep 2026 -- Present}
  \begin{itemize}
    \item Contributed to an open-source \textbf{TypeScript} engine that renders and edits PowerPoint and Word files in the browser,
          improving editing fidelity, \textbf{real-time collaboration} and load performance across React, Vue and Svelte viewers
  \end{itemize}}

\newcommand{\MirrorAI}{%
  \entry{Mirror AI: Deployable Personas}{Honorable mention in Supabase YC Hackathon}{Oct 2024 -- Dec 2024}
  \begin{itemize}
    \item Designed an \textbf{agentic LLM architecture} with LangGraph to mirror user personalities, creating interactive digital
          personas, and deployed it as a full-stack platform on Supabase and Vercel with one-click publishing
  \end{itemize}}

\newcommand{\LLMReasoning}{%
  \entry{Improving LLM Reasoning for Numerical Problems}{\report{https://sudhansh6.github.io/assets/pdf/MathPrompter++.pdf}}{Sep 2024 -- Dec 2024}
  \begin{itemize}
    \item Enhanced MathPrompter (ACL 2023) with CoT, achieving \textbf{10\% higher accuracy} on Llama 3.1 1B where prior methods failed,
          and reduced hallucination rates significantly by integrating multi-step validation
  \end{itemize}}

\newcommand{\InverseRendering}{%
  \entry{Inverse Rendering with 2D Gaussian Splatting}{\report{https://sudhansh6.github.io/assets/pdf/CSE252D_Report.pdf}}{Mar 2024 -- May 2024}
  \begin{itemize}
    \item Developed a novel \textbf{inverse rendering} framework in CUDA to recover PBR properties of a scene using 2D Gaussian
          Splatting, improving normal map MAE by \textbf{15\%} over the current SOTA methods
  \end{itemize}}

% ---------------------------------------------------------------------------------------------------
% Publications: title on the left, venue on the right, one line each.
\newcommand{\publine}[3]{\par\noindent\href{#1}{#2}\hfill\hspace{1em}\side{#3}\par\addvspace{1.5pt}}
\newcommand{\PubSaturation}{\publine{https://doi.org/10.1016/j.sigpro.2023.109349}{A likelihood based method for compressive signal recovery under Gaussian and saturation noise}{Signal Processing, 2024}}
\newcommand{\PubPermutation}{\publine{https://ieeexplore.ieee.org/document/10888147}{Identification and correction of permutation errors in compressed sensing-based group testing}{ICASSP 2025}}
\newcommand{\PubRearrangement}{\publine{https://arxiv.org/abs/2411.10899}{Planning for tabletop object rearrangement}{ICRA 2025}}

% ---------------------------------------------------------------------------------------------------
% Education, skills, service
% Each profile sets \UCSDCourses, \IITBCourses and \SkillRows before % Shared preamble and content for Sudhansh Peddabomma's résumés.
% Each profile (SudhanshPeddabomma-AI.tex, SudhanshPeddabomma-Vision.tex) inputs this file and picks
% which blocks to show and in what order, so every bullet is written once.
% Layout rule: every bullet fills one line, or runs 1.5 to 2 lines. Check with: python3 check_lines.py <pdf>
\documentclass[10pt,letterpaper]{article}

\usepackage[margin=0.55in,top=0.5in,bottom=0.45in]{geometry}
\usepackage[T1]{fontenc}
% Fonts: \resumefont loads the body font; \headfont styles the name and section titles.
% Override either before \input{resume-common} to try another pairing.
% Default: Charter body with Source Sans headings. All-sans alternative:
%   \def\resumefont{\usepackage[default,semibold]{sourcesanspro}}\def\headfont{}
\providecommand{\resumefont}{\usepackage[scaled=0.96]{XCharter}\usepackage[semibold,scale=0.96]{sourcesanspro}}
\providecommand{\headfont}{\sffamily}
\providecommand{\namefont}{\rmfamily\scshape}
\resumefont
\usepackage[final]{microtype}
\usepackage[dvipsnames]{xcolor}
\usepackage{enumitem}
\usepackage{titlesec}
\usepackage{tabularx}
\usepackage{fontawesome5}
\usepackage{etoolbox}
\usepackage[hidelinks]{hyperref}

% The site's green, darkened for print.
\definecolor{accent}{HTML}{3A7A45}
\definecolor{muted}{HTML}{5C675D}

\pagestyle{empty}
\setlength{\parindent}{0pt}
\setlength{\tabcolsep}{0pt}
\providecommand{\leading}{1.1}
\linespread{\leading}

\hypersetup{pdfauthor={Sudhansh Peddabomma}}

% Type scale: name 26pt small caps, section titles 12.5pt caps, entry titles 10.8pt, body 10pt,
% side notes 9pt.
\newcommand{\side}[1]{{\fontsize{9}{11}\selectfont\color{muted}#1}}
\newcommand{\sectitle}[1]{\textls[60]{\MakeUppercase{#1}}}
\titleformat{\section}{\headfont\color{accent}\bfseries\fontsize{12.5}{14}\selectfont}{}{0pt}{\sectitle}[\vspace{1.5pt}{\color{accent}\titlerule[0.8pt]}\global\firstentrytrue]
% The first entry after a section title skips \entrygap, so only \rulegap sits under the rule.
\newif\iffirstentry

% Vertical rhythm, the same idea as the earlier hand-tuned résumés: bullets inside an entry sit
% tight together, entries are separated by a clear gap, and sections by a larger one. A profile
% may override these before \input{resume-common} to fit one page.
\providecommand{\secgap}{12pt}   % above a section title
\providecommand{\rulegap}{5.5pt}   % section rule to first entry
\providecommand{\entrygap}{7pt}  % between entries
\providecommand{\headgap}{1.5pt} % entry heading to its bullets
\providecommand{\itemgap}{0.6pt} % between bullets of one entry
\titlespacing*{\section}{0pt}{\secgap}{\rulegap}

\newcommand{\squarebullet}{\raisebox{0.32ex}{\color{accent}\rule{0.5ex}{0.5ex}}}
\setlist[itemize]{leftmargin=12pt, labelsep=6pt, topsep=\headgap, itemsep=\itemgap, parsep=0pt, partopsep=0pt,
  label=\squarebullet}

% \entry{Organisation or title}{Role or detail (may be empty)}{Dates}; \sub{Role or detail}{Right-hand note}
\newcommand{\entry}[3]{%
  \par\iffirstentry\global\firstentryfalse\else\addvspace{\entrygap}\fi\noindent{\headfont\fontsize{10.8}{13}\selectfont\textbf{#1}}\ifblank{#2}{}{\ {\color{muted}|}\ #2}\hfill\hspace{1em}\side{#3}\par}
\newcommand{\sub}[2]{%
  \noindent\textit{#1}\hfill{\color{muted}\textit{#2}}\par}
\newcommand{\report}[1]{\href{#1}{\color{accent}Report}}
\newcommand{\me}{\textbf{S. Peddabomma}}

% ---------------------------------------------------------------------------------------------------
% Header
\newcommand{\Header}{%
  \begin{center}
    {\namefont\fontsize{26}{28}\selectfont Sudhansh Peddabomma}\\[6pt]
    {\headfont\fontsize{9.5}{11}\selectfont\color{muted}
      \href{mailto:sudhansh6@gmail.com}{\faEnvelope[regular]\ sudhansh6@gmail.com} \quad
      \href{https://sudhansh6.github.io}{\faGlobe\ sudhansh6.github.io} \quad
      \href{https://www.linkedin.com/in/sudhansh-peddabomma-73a208192/}{\faLinkedin\ sudhansh-peddabomma} \quad
      \href{https://github.com/Sudhansh6}{\faGithub\ Sudhansh6}}
  \end{center}\vspace{-4pt}}

% ---------------------------------------------------------------------------------------------------
% Glean bullets, shared by both profiles
\newcommand{\GleanSlides}{\item Led \textbf{AI slide generation} in Glean Assistant, now creating \textbf{180K decks a week} for \textbf{50K weekly active users} at enterprises}
\newcommand{\GleanSlidesPerf}{\item Improved the slides harness, raising reliability to \textbf{97\%}, cutting median latency \textbf{26\%} and cost per deck by \textbf{half}
  in four months}
\newcommand{\GleanSlidesOrch}{\item Architected a \textbf{slide orchestration system} that reuses relevant slides from a company's existing decks and generates
  new ones in its brand and template, so teams can refine decks with each other and Glean in full context, without leaving the product}
\newcommand{\GleanAgents}{\item Moved Glean's \textbf{workflow agents} from one-shot planning to an \textbf{agentic tool-use loop}, greatly improving agentic
  capabilities and instruction following, with \textbf{18\%} higher answer completeness and \textbf{17\%} fewer tokens}
\newcommand{\GleanTeamChat}{\item Engineered the core of \textbf{Team chat}, letting teammates work with Glean in one shared conversation and co-edit its artifacts}
\newcommand{\GleanVideo}{\item Shipped \textbf{video artifacts} end to end, generating narrated videos in chat with streamed playback and admin rollout controls}
\newcommand{\GleanEvals}{\item Created the \textbf{evaluations and metrics} that guide model choices and rollouts across agents, slides and interactive artifacts}

% ---------------------------------------------------------------------------------------------------
% Experience
% \groupbreak separates related groups of bullets inside one long entry (slides, then platform work).
\providecommand{\groupgap}{3.5pt}
\newcommand{\groupbreak}{\vspace{\groupgap}}
\newcommand{\GleanAI}{%
  \entry{Glean}{Software Engineer}{Apr 2025 -- Present}
  \begin{itemize}
    \GleanSlides
    \GleanSlidesOrch
    \GleanSlidesPerf \groupbreak
    \GleanTeamChat
    \GleanVideo
    \GleanAgents
    \GleanEvals
  \end{itemize}}

\newcommand{\GleanVision}{%
  \entry{Glean}{Software Engineer}{Apr 2025 -- Present}
  \begin{itemize}
    \GleanSlides
    \GleanSlidesOrch
    \GleanSlidesPerf \groupbreak
    \GleanVideo
    \GleanEvals
  \end{itemize}}

\newcommand{\Duality}{%
  \entry{Duality AI}{Computer Vision Intern}{Jun 2024 -- Sep 2024}
  \begin{itemize}
    \item Designed automated \textbf{3D reconstruction} techniques for featureless objects, reducing digital-twin generation time
          by \textbf{40\%}
    \item Collaborated with \textbf{Autodesk} to validate Unreal Engine simulations for robotics tasks; leveraging structured
          domain randomization to reduce Sim2Real gap and increase mAP-50 by \textbf{15\%} for object detection and segmentation
  \end{itemize}}

\newcommand{\Microsoft}{%
  \entry{Microsoft}{Data and Applied Scientist Intern}{May 2022 -- Jul 2022}
  \begin{itemize}
    \item Developed a \textbf{decision-tree ranker} to recommend emails without user queries, improving \textbf{Outlook} search capabilities
    \item Proposed hierarchical feature-sets for the ranker, reducing latency for recommendations and improving recall
  \end{itemize}}

% ---------------------------------------------------------------------------------------------------
% Research
\newcommand{\HomeRobots}{%
  \entry{Real-time 3D Perception for Home Robots}{UC San Diego}{Sep 2023 -- Sep 2024}
  \begin{itemize}
    \item Investigated real-time dense \textbf{visual SLAM} methods using NeRFs and Gaussian Splatting for indoor robot navigation
    \item Integrated object segmentation, grasp-pose estimation, and 3D mapping on the Fetch robot via ROS, demonstrating
          a novel \textbf{tabletop object rearrangement} algorithm that reduced cost by \textbf{20\%} compared to the state-of-the-art
  \end{itemize}}

\newcommand{\Tomography}{%
  \entry{3D Tomography with Primal-Dual Neural Networks}{University College London}{May 2021 -- Jul 2023}
  \begin{itemize}
    \item Developed a stochastic neural-network architecture of primal-dual algorithm for online reconstruction of 3D
          volumes from tomographic projections, with \textbf{99.6\,\% structural similarity} in low-dosage conditions
    \item Built a Python library with custom gradient operators for reconstructing volumes in a single pass, reducing
          compute requirements by up to \textbf{5x} over SOTA learning-based approaches for cone vector tomography
  \end{itemize}}

% ---------------------------------------------------------------------------------------------------
% Projects
\newcommand{\OOXML}{%
  \entry{ooxml}{Open-source contributor \textbar\ \href{https://github.com/ChristopherVR/ooxml}{\color{accent}GitHub}}{Sep 2026 -- Present}
  \begin{itemize}
    \item Contributed to an open-source \textbf{TypeScript} engine that renders and edits PowerPoint and Word files in the browser,
          improving editing fidelity, \textbf{real-time collaboration} and load performance across React, Vue and Svelte viewers
  \end{itemize}}

\newcommand{\MirrorAI}{%
  \entry{Mirror AI: Deployable Personas}{Honorable mention in Supabase YC Hackathon}{Oct 2024 -- Dec 2024}
  \begin{itemize}
    \item Designed an \textbf{agentic LLM architecture} with LangGraph to mirror user personalities, creating interactive digital
          personas, and deployed it as a full-stack platform on Supabase and Vercel with one-click publishing
  \end{itemize}}

\newcommand{\LLMReasoning}{%
  \entry{Improving LLM Reasoning for Numerical Problems}{\report{https://sudhansh6.github.io/assets/pdf/MathPrompter++.pdf}}{Sep 2024 -- Dec 2024}
  \begin{itemize}
    \item Enhanced MathPrompter (ACL 2023) with CoT, achieving \textbf{10\% higher accuracy} on Llama 3.1 1B where prior methods failed,
          and reduced hallucination rates significantly by integrating multi-step validation
  \end{itemize}}

\newcommand{\InverseRendering}{%
  \entry{Inverse Rendering with 2D Gaussian Splatting}{\report{https://sudhansh6.github.io/assets/pdf/CSE252D_Report.pdf}}{Mar 2024 -- May 2024}
  \begin{itemize}
    \item Developed a novel \textbf{inverse rendering} framework in CUDA to recover PBR properties of a scene using 2D Gaussian
          Splatting, improving normal map MAE by \textbf{15\%} over the current SOTA methods
  \end{itemize}}

% ---------------------------------------------------------------------------------------------------
% Publications: title on the left, venue on the right, one line each.
\newcommand{\publine}[3]{\par\noindent\href{#1}{#2}\hfill\hspace{1em}\side{#3}\par\addvspace{1.5pt}}
\newcommand{\PubSaturation}{\publine{https://doi.org/10.1016/j.sigpro.2023.109349}{A likelihood based method for compressive signal recovery under Gaussian and saturation noise}{Signal Processing, 2024}}
\newcommand{\PubPermutation}{\publine{https://ieeexplore.ieee.org/document/10888147}{Identification and correction of permutation errors in compressed sensing-based group testing}{ICASSP 2025}}
\newcommand{\PubRearrangement}{\publine{https://arxiv.org/abs/2411.10899}{Planning for tabletop object rearrangement}{ICRA 2025}}

% ---------------------------------------------------------------------------------------------------
% Education, skills, service
% Each profile sets \UCSDCourses, \IITBCourses and \SkillRows before \input{resume-common}.
\providecommand{\UCSDCourses}{}
\providecommand{\IITBCourses}{}
\providecommand{\SkillRows}{}
\newcommand{\coursework}[1]{\ifblank{#1}{}{\par\noindent{\fontsize{9.5}{11.5}\selectfont{\color{muted}Coursework:} #1}\par}}
\newcommand{\Education}{%
  \section{Education}
  \entry{University of California San Diego}{M.S. in Computer Science and Engineering (AI)}{GPA 4.00/4.00 \quad Sep 2023 -- Mar 2025}
  \coursework{\UCSDCourses}
  \entry{Indian Institute of Technology Bombay}{B.Tech.\ with Honors in Computer Science}{CPI 9.66/10 \quad Jul 2019 -- Jul 2023}
  \coursework{\IITBCourses}}

% \skill{Label}{Items} is one row of the skills table.
\newcommand{\skill}[2]{{\headfont\bfseries #1} & #2 \\[1.5pt]}
\newcommand{\cpp}{C\kern-0.04em\raisebox{0.06ex}{\small++}}
\newcommand{\Skills}{%
  \section{Skills}
  \begin{tabularx}{\linewidth}{@{}l@{\hspace{12pt}}>{\raggedright\arraybackslash}X@{}}
    \SkillRows
  \end{tabularx}}
.
\providecommand{\UCSDCourses}{}
\providecommand{\IITBCourses}{}
\providecommand{\SkillRows}{}
\newcommand{\coursework}[1]{\ifblank{#1}{}{\par\noindent{\fontsize{9.5}{11.5}\selectfont{\color{muted}Coursework:} #1}\par}}
\newcommand{\Education}{%
  \section{Education}
  \entry{University of California San Diego}{M.S. in Computer Science and Engineering (AI)}{GPA 4.00/4.00 \quad Sep 2023 -- Mar 2025}
  \coursework{\UCSDCourses}
  \entry{Indian Institute of Technology Bombay}{B.Tech.\ with Honors in Computer Science}{CPI 9.66/10 \quad Jul 2019 -- Jul 2023}
  \coursework{\IITBCourses}}

% \skill{Label}{Items} is one row of the skills table.
\newcommand{\skill}[2]{{\headfont\bfseries #1} & #2 \\[1.5pt]}
\newcommand{\cpp}{C\kern-0.04em\raisebox{0.06ex}{\small++}}
\newcommand{\Skills}{%
  \section{Skills}
  \begin{tabularx}{\linewidth}{@{}l@{\hspace{12pt}}>{\raggedright\arraybackslash}X@{}}
    \SkillRows
  \end{tabularx}}
.
\providecommand{\UCSDCourses}{}
\providecommand{\IITBCourses}{}
\providecommand{\SkillRows}{}
\newcommand{\coursework}[1]{\ifblank{#1}{}{\par\noindent{\fontsize{9.5}{11.5}\selectfont{\color{muted}Coursework:} #1}\par}}
\newcommand{\Education}{%
  \section{Education}
  \entry{University of California San Diego}{M.S. in Computer Science and Engineering (AI)}{GPA 4.00/4.00 \quad Sep 2023 -- Mar 2025}
  \coursework{\UCSDCourses}
  \entry{Indian Institute of Technology Bombay}{B.Tech.\ with Honors in Computer Science}{CPI 9.66/10 \quad Jul 2019 -- Jul 2023}
  \coursework{\IITBCourses}}

% \skill{Label}{Items} is one row of the skills table.
\newcommand{\skill}[2]{{\headfont\bfseries #1} & #2 \\[1.5pt]}
\newcommand{\cpp}{C\kern-0.04em\raisebox{0.06ex}{\small++}}
\newcommand{\Skills}{%
  \section{Skills}
  \begin{tabularx}{\linewidth}{@{}l@{\hspace{12pt}}>{\raggedright\arraybackslash}X@{}}
    \SkillRows
  \end{tabularx}}
 to fit one page.
\providecommand{\secgap}{12pt}   % above a section title
\providecommand{\rulegap}{5.5pt}   % section rule to first entry
\providecommand{\entrygap}{7pt}  % between entries
\providecommand{\headgap}{1.5pt} % entry heading to its bullets
\providecommand{\itemgap}{0.6pt} % between bullets of one entry
\titlespacing*{\section}{0pt}{\secgap}{\rulegap}

\newcommand{\squarebullet}{\raisebox{0.32ex}{\color{accent}\rule{0.5ex}{0.5ex}}}
\setlist[itemize]{leftmargin=12pt, labelsep=6pt, topsep=\headgap, itemsep=\itemgap, parsep=0pt, partopsep=0pt,
  label=\squarebullet}

% \entry{Organisation or title}{Role or detail (may be empty)}{Dates}; \sub{Role or detail}{Right-hand note}
\newcommand{\entry}[3]{%
  \par\iffirstentry\global\firstentryfalse\else\addvspace{\entrygap}\fi\noindent{\headfont\fontsize{10.8}{13}\selectfont\textbf{#1}}\ifblank{#2}{}{\ {\color{muted}|}\ #2}\hfill\hspace{1em}\side{#3}\par}
\newcommand{\sub}[2]{%
  \noindent\textit{#1}\hfill{\color{muted}\textit{#2}}\par}
\newcommand{\report}[1]{\href{#1}{\color{accent}Report}}
\newcommand{\me}{\textbf{S. Peddabomma}}

% ---------------------------------------------------------------------------------------------------
% Header
\newcommand{\Header}{%
  \begin{center}
    {\namefont\fontsize{26}{28}\selectfont Sudhansh Peddabomma}\\[6pt]
    {\headfont\fontsize{9.5}{11}\selectfont\color{muted}
      \href{mailto:sudhansh6@gmail.com}{\faEnvelope[regular]\ sudhansh6@gmail.com} \quad
      \href{https://sudhansh6.github.io}{\faGlobe\ sudhansh6.github.io} \quad
      \href{https://www.linkedin.com/in/sudhansh-peddabomma-73a208192/}{\faLinkedin\ sudhansh-peddabomma} \quad
      \href{https://github.com/Sudhansh6}{\faGithub\ Sudhansh6}}
  \end{center}\vspace{-4pt}}

% ---------------------------------------------------------------------------------------------------
% Glean bullets, shared by both profiles
\newcommand{\GleanSlides}{\item Led \textbf{AI slide generation} in Glean Assistant, now creating \textbf{180K decks a week} for \textbf{50K weekly active users} at enterprises}
\newcommand{\GleanSlidesPerf}{\item Improved the slides harness, raising reliability to \textbf{97\%}, cutting median latency \textbf{26\%} and cost per deck by \textbf{half}
  in four months}
\newcommand{\GleanSlidesOrch}{\item Architected a \textbf{slide orchestration system} that reuses relevant slides from a company's existing decks and generates
  new ones in its brand and template, so teams can refine decks with each other and Glean in full context, without leaving the product}
\newcommand{\GleanAgents}{\item Moved Glean's \textbf{workflow agents} from one-shot planning to an \textbf{agentic tool-use loop}, greatly improving agentic
  capabilities and instruction following, with \textbf{18\%} higher answer completeness and \textbf{17\%} fewer tokens}
\newcommand{\GleanTeamChat}{\item Engineered the core of \textbf{Team chat}, letting teammates work with Glean in one shared conversation and co-edit its artifacts}
\newcommand{\GleanVideo}{\item Shipped \textbf{video artifacts} end to end, generating narrated videos in chat with streamed playback and admin rollout controls}
\newcommand{\GleanEvals}{\item Created the \textbf{evaluations and metrics} that guide model choices and rollouts across agents, slides and interactive artifacts}

% ---------------------------------------------------------------------------------------------------
% Experience
% \groupbreak separates related groups of bullets inside one long entry (slides, then platform work).
\providecommand{\groupgap}{3.5pt}
\newcommand{\groupbreak}{\vspace{\groupgap}}
\newcommand{\GleanAI}{%
  \entry{Glean}{Software Engineer}{Apr 2025 -- Present}
  \begin{itemize}
    \GleanSlides
    \GleanSlidesOrch
    \GleanSlidesPerf \groupbreak
    \GleanTeamChat
    \GleanVideo
    \GleanAgents
    \GleanEvals
  \end{itemize}}

\newcommand{\GleanVision}{%
  \entry{Glean}{Software Engineer}{Apr 2025 -- Present}
  \begin{itemize}
    \GleanSlides
    \GleanSlidesOrch
    \GleanSlidesPerf \groupbreak
    \GleanVideo
    \GleanEvals
  \end{itemize}}

\newcommand{\Duality}{%
  \entry{Duality AI}{Computer Vision Intern}{Jun 2024 -- Sep 2024}
  \begin{itemize}
    \item Designed automated \textbf{3D reconstruction} techniques for featureless objects, reducing digital-twin generation time
          by \textbf{40\%}
    \item Collaborated with \textbf{Autodesk} to validate Unreal Engine simulations for robotics tasks; leveraging structured
          domain randomization to reduce Sim2Real gap and increase mAP-50 by \textbf{15\%} for object detection and segmentation
  \end{itemize}}

\newcommand{\Microsoft}{%
  \entry{Microsoft}{Data and Applied Scientist Intern}{May 2022 -- Jul 2022}
  \begin{itemize}
    \item Developed a \textbf{decision-tree ranker} to recommend emails without user queries, improving \textbf{Outlook} search capabilities
    \item Proposed hierarchical feature-sets for the ranker, reducing latency for recommendations and improving recall
  \end{itemize}}

% ---------------------------------------------------------------------------------------------------
% Research
\newcommand{\HomeRobots}{%
  \entry{Real-time 3D Perception for Home Robots}{UC San Diego}{Sep 2023 -- Sep 2024}
  \begin{itemize}
    \item Investigated real-time dense \textbf{visual SLAM} methods using NeRFs and Gaussian Splatting for indoor robot navigation
    \item Integrated object segmentation, grasp-pose estimation, and 3D mapping on the Fetch robot via ROS, demonstrating
          a novel \textbf{tabletop object rearrangement} algorithm that reduced cost by \textbf{20\%} compared to the state-of-the-art
  \end{itemize}}

\newcommand{\Tomography}{%
  \entry{3D Tomography with Primal-Dual Neural Networks}{University College London}{May 2021 -- Jul 2023}
  \begin{itemize}
    \item Developed a stochastic neural-network architecture of primal-dual algorithm for online reconstruction of 3D
          volumes from tomographic projections, with \textbf{99.6\,\% structural similarity} in low-dosage conditions
    \item Built a Python library with custom gradient operators for reconstructing volumes in a single pass, reducing
          compute requirements by up to \textbf{5x} over SOTA learning-based approaches for cone vector tomography
  \end{itemize}}

% ---------------------------------------------------------------------------------------------------
% Projects
\newcommand{\OOXML}{%
  \entry{ooxml}{Open-source contributor \textbar\ \href{https://github.com/ChristopherVR/ooxml}{\color{accent}GitHub}}{Sep 2026 -- Present}
  \begin{itemize}
    \item Contributed to an open-source \textbf{TypeScript} engine that renders and edits PowerPoint and Word files in the browser,
          improving editing fidelity, \textbf{real-time collaboration} and load performance across React, Vue and Svelte viewers
  \end{itemize}}

\newcommand{\MirrorAI}{%
  \entry{Mirror AI: Deployable Personas}{Honorable mention in Supabase YC Hackathon}{Oct 2024 -- Dec 2024}
  \begin{itemize}
    \item Designed an \textbf{agentic LLM architecture} with LangGraph to mirror user personalities, creating interactive digital
          personas, and deployed it as a full-stack platform on Supabase and Vercel with one-click publishing
  \end{itemize}}

\newcommand{\LLMReasoning}{%
  \entry{Improving LLM Reasoning for Numerical Problems}{\report{https://sudhansh6.github.io/assets/pdf/MathPrompter++.pdf}}{Sep 2024 -- Dec 2024}
  \begin{itemize}
    \item Enhanced MathPrompter (ACL 2023) with CoT, achieving \textbf{10\% higher accuracy} on Llama 3.1 1B where prior methods failed,
          and reduced hallucination rates significantly by integrating multi-step validation
  \end{itemize}}

\newcommand{\InverseRendering}{%
  \entry{Inverse Rendering with 2D Gaussian Splatting}{\report{https://sudhansh6.github.io/assets/pdf/CSE252D_Report.pdf}}{Mar 2024 -- May 2024}
  \begin{itemize}
    \item Developed a novel \textbf{inverse rendering} framework in CUDA to recover PBR properties of a scene using 2D Gaussian
          Splatting, improving normal map MAE by \textbf{15\%} over the current SOTA methods
  \end{itemize}}

% ---------------------------------------------------------------------------------------------------
% Publications: title on the left, venue on the right, one line each.
\newcommand{\publine}[3]{\par\noindent\href{#1}{#2}\hfill\hspace{1em}\side{#3}\par\addvspace{1.5pt}}
\newcommand{\PubSaturation}{\publine{https://doi.org/10.1016/j.sigpro.2023.109349}{A likelihood based method for compressive signal recovery under Gaussian and saturation noise}{Signal Processing, 2024}}
\newcommand{\PubPermutation}{\publine{https://ieeexplore.ieee.org/document/10888147}{Identification and correction of permutation errors in compressed sensing-based group testing}{ICASSP 2025}}
\newcommand{\PubRearrangement}{\publine{https://arxiv.org/abs/2411.10899}{Planning for tabletop object rearrangement}{ICRA 2025}}

% ---------------------------------------------------------------------------------------------------
% Education, skills, service
% Each profile sets \UCSDCourses, \IITBCourses and \SkillRows before % Shared preamble and content for Sudhansh Peddabomma's résumés.
% Each profile (SudhanshPeddabomma-AI.tex, SudhanshPeddabomma-Vision.tex) inputs this file and picks
% which blocks to show and in what order, so every bullet is written once.
% Layout rule: every bullet fills one line, or runs 1.5 to 2 lines. Check with: python3 check_lines.py <pdf>
\documentclass[10pt,letterpaper]{article}

\usepackage[margin=0.55in,top=0.5in,bottom=0.45in]{geometry}
\usepackage[T1]{fontenc}
% Fonts: \resumefont loads the body font; \headfont styles the name and section titles.
% Override either before % Shared preamble and content for Sudhansh Peddabomma's résumés.
% Each profile (SudhanshPeddabomma-AI.tex, SudhanshPeddabomma-Vision.tex) inputs this file and picks
% which blocks to show and in what order, so every bullet is written once.
% Layout rule: every bullet fills one line, or runs 1.5 to 2 lines. Check with: python3 check_lines.py <pdf>
\documentclass[10pt,letterpaper]{article}

\usepackage[margin=0.55in,top=0.5in,bottom=0.45in]{geometry}
\usepackage[T1]{fontenc}
% Fonts: \resumefont loads the body font; \headfont styles the name and section titles.
% Override either before % Shared preamble and content for Sudhansh Peddabomma's résumés.
% Each profile (SudhanshPeddabomma-AI.tex, SudhanshPeddabomma-Vision.tex) inputs this file and picks
% which blocks to show and in what order, so every bullet is written once.
% Layout rule: every bullet fills one line, or runs 1.5 to 2 lines. Check with: python3 check_lines.py <pdf>
\documentclass[10pt,letterpaper]{article}

\usepackage[margin=0.55in,top=0.5in,bottom=0.45in]{geometry}
\usepackage[T1]{fontenc}
% Fonts: \resumefont loads the body font; \headfont styles the name and section titles.
% Override either before \input{resume-common} to try another pairing.
% Default: Charter body with Source Sans headings. All-sans alternative:
%   \def\resumefont{\usepackage[default,semibold]{sourcesanspro}}\def\headfont{}
\providecommand{\resumefont}{\usepackage[scaled=0.96]{XCharter}\usepackage[semibold,scale=0.96]{sourcesanspro}}
\providecommand{\headfont}{\sffamily}
\providecommand{\namefont}{\rmfamily\scshape}
\resumefont
\usepackage[final]{microtype}
\usepackage[dvipsnames]{xcolor}
\usepackage{enumitem}
\usepackage{titlesec}
\usepackage{tabularx}
\usepackage{fontawesome5}
\usepackage{etoolbox}
\usepackage[hidelinks]{hyperref}

% The site's green, darkened for print.
\definecolor{accent}{HTML}{3A7A45}
\definecolor{muted}{HTML}{5C675D}

\pagestyle{empty}
\setlength{\parindent}{0pt}
\setlength{\tabcolsep}{0pt}
\providecommand{\leading}{1.1}
\linespread{\leading}

\hypersetup{pdfauthor={Sudhansh Peddabomma}}

% Type scale: name 26pt small caps, section titles 12.5pt caps, entry titles 10.8pt, body 10pt,
% side notes 9pt.
\newcommand{\side}[1]{{\fontsize{9}{11}\selectfont\color{muted}#1}}
\newcommand{\sectitle}[1]{\textls[60]{\MakeUppercase{#1}}}
\titleformat{\section}{\headfont\color{accent}\bfseries\fontsize{12.5}{14}\selectfont}{}{0pt}{\sectitle}[\vspace{1.5pt}{\color{accent}\titlerule[0.8pt]}\global\firstentrytrue]
% The first entry after a section title skips \entrygap, so only \rulegap sits under the rule.
\newif\iffirstentry

% Vertical rhythm, the same idea as the earlier hand-tuned résumés: bullets inside an entry sit
% tight together, entries are separated by a clear gap, and sections by a larger one. A profile
% may override these before \input{resume-common} to fit one page.
\providecommand{\secgap}{12pt}   % above a section title
\providecommand{\rulegap}{5.5pt}   % section rule to first entry
\providecommand{\entrygap}{7pt}  % between entries
\providecommand{\headgap}{1.5pt} % entry heading to its bullets
\providecommand{\itemgap}{0.6pt} % between bullets of one entry
\titlespacing*{\section}{0pt}{\secgap}{\rulegap}

\newcommand{\squarebullet}{\raisebox{0.32ex}{\color{accent}\rule{0.5ex}{0.5ex}}}
\setlist[itemize]{leftmargin=12pt, labelsep=6pt, topsep=\headgap, itemsep=\itemgap, parsep=0pt, partopsep=0pt,
  label=\squarebullet}

% \entry{Organisation or title}{Role or detail (may be empty)}{Dates}; \sub{Role or detail}{Right-hand note}
\newcommand{\entry}[3]{%
  \par\iffirstentry\global\firstentryfalse\else\addvspace{\entrygap}\fi\noindent{\headfont\fontsize{10.8}{13}\selectfont\textbf{#1}}\ifblank{#2}{}{\ {\color{muted}|}\ #2}\hfill\hspace{1em}\side{#3}\par}
\newcommand{\sub}[2]{%
  \noindent\textit{#1}\hfill{\color{muted}\textit{#2}}\par}
\newcommand{\report}[1]{\href{#1}{\color{accent}Report}}
\newcommand{\me}{\textbf{S. Peddabomma}}

% ---------------------------------------------------------------------------------------------------
% Header
\newcommand{\Header}{%
  \begin{center}
    {\namefont\fontsize{26}{28}\selectfont Sudhansh Peddabomma}\\[6pt]
    {\headfont\fontsize{9.5}{11}\selectfont\color{muted}
      \href{mailto:sudhansh6@gmail.com}{\faEnvelope[regular]\ sudhansh6@gmail.com} \quad
      \href{https://sudhansh6.github.io}{\faGlobe\ sudhansh6.github.io} \quad
      \href{https://www.linkedin.com/in/sudhansh-peddabomma-73a208192/}{\faLinkedin\ sudhansh-peddabomma} \quad
      \href{https://github.com/Sudhansh6}{\faGithub\ Sudhansh6}}
  \end{center}\vspace{-4pt}}

% ---------------------------------------------------------------------------------------------------
% Glean bullets, shared by both profiles
\newcommand{\GleanSlides}{\item Led \textbf{AI slide generation} in Glean Assistant, now creating \textbf{180K decks a week} for \textbf{50K weekly active users} at enterprises}
\newcommand{\GleanSlidesPerf}{\item Improved the slides harness, raising reliability to \textbf{97\%}, cutting median latency \textbf{26\%} and cost per deck by \textbf{half}
  in four months}
\newcommand{\GleanSlidesOrch}{\item Architected a \textbf{slide orchestration system} that reuses relevant slides from a company's existing decks and generates
  new ones in its brand and template, so teams can refine decks with each other and Glean in full context, without leaving the product}
\newcommand{\GleanAgents}{\item Moved Glean's \textbf{workflow agents} from one-shot planning to an \textbf{agentic tool-use loop}, greatly improving agentic
  capabilities and instruction following, with \textbf{18\%} higher answer completeness and \textbf{17\%} fewer tokens}
\newcommand{\GleanTeamChat}{\item Engineered the core of \textbf{Team chat}, letting teammates work with Glean in one shared conversation and co-edit its artifacts}
\newcommand{\GleanVideo}{\item Shipped \textbf{video artifacts} end to end, generating narrated videos in chat with streamed playback and admin rollout controls}
\newcommand{\GleanEvals}{\item Created the \textbf{evaluations and metrics} that guide model choices and rollouts across agents, slides and interactive artifacts}

% ---------------------------------------------------------------------------------------------------
% Experience
% \groupbreak separates related groups of bullets inside one long entry (slides, then platform work).
\providecommand{\groupgap}{3.5pt}
\newcommand{\groupbreak}{\vspace{\groupgap}}
\newcommand{\GleanAI}{%
  \entry{Glean}{Software Engineer}{Apr 2025 -- Present}
  \begin{itemize}
    \GleanSlides
    \GleanSlidesOrch
    \GleanSlidesPerf \groupbreak
    \GleanTeamChat
    \GleanVideo
    \GleanAgents
    \GleanEvals
  \end{itemize}}

\newcommand{\GleanVision}{%
  \entry{Glean}{Software Engineer}{Apr 2025 -- Present}
  \begin{itemize}
    \GleanSlides
    \GleanSlidesOrch
    \GleanSlidesPerf \groupbreak
    \GleanVideo
    \GleanEvals
  \end{itemize}}

\newcommand{\Duality}{%
  \entry{Duality AI}{Computer Vision Intern}{Jun 2024 -- Sep 2024}
  \begin{itemize}
    \item Designed automated \textbf{3D reconstruction} techniques for featureless objects, reducing digital-twin generation time
          by \textbf{40\%}
    \item Collaborated with \textbf{Autodesk} to validate Unreal Engine simulations for robotics tasks; leveraging structured
          domain randomization to reduce Sim2Real gap and increase mAP-50 by \textbf{15\%} for object detection and segmentation
  \end{itemize}}

\newcommand{\Microsoft}{%
  \entry{Microsoft}{Data and Applied Scientist Intern}{May 2022 -- Jul 2022}
  \begin{itemize}
    \item Developed a \textbf{decision-tree ranker} to recommend emails without user queries, improving \textbf{Outlook} search capabilities
    \item Proposed hierarchical feature-sets for the ranker, reducing latency for recommendations and improving recall
  \end{itemize}}

% ---------------------------------------------------------------------------------------------------
% Research
\newcommand{\HomeRobots}{%
  \entry{Real-time 3D Perception for Home Robots}{UC San Diego}{Sep 2023 -- Sep 2024}
  \begin{itemize}
    \item Investigated real-time dense \textbf{visual SLAM} methods using NeRFs and Gaussian Splatting for indoor robot navigation
    \item Integrated object segmentation, grasp-pose estimation, and 3D mapping on the Fetch robot via ROS, demonstrating
          a novel \textbf{tabletop object rearrangement} algorithm that reduced cost by \textbf{20\%} compared to the state-of-the-art
  \end{itemize}}

\newcommand{\Tomography}{%
  \entry{3D Tomography with Primal-Dual Neural Networks}{University College London}{May 2021 -- Jul 2023}
  \begin{itemize}
    \item Developed a stochastic neural-network architecture of primal-dual algorithm for online reconstruction of 3D
          volumes from tomographic projections, with \textbf{99.6\,\% structural similarity} in low-dosage conditions
    \item Built a Python library with custom gradient operators for reconstructing volumes in a single pass, reducing
          compute requirements by up to \textbf{5x} over SOTA learning-based approaches for cone vector tomography
  \end{itemize}}

% ---------------------------------------------------------------------------------------------------
% Projects
\newcommand{\OOXML}{%
  \entry{ooxml}{Open-source contributor \textbar\ \href{https://github.com/ChristopherVR/ooxml}{\color{accent}GitHub}}{Sep 2026 -- Present}
  \begin{itemize}
    \item Contributed to an open-source \textbf{TypeScript} engine that renders and edits PowerPoint and Word files in the browser,
          improving editing fidelity, \textbf{real-time collaboration} and load performance across React, Vue and Svelte viewers
  \end{itemize}}

\newcommand{\MirrorAI}{%
  \entry{Mirror AI: Deployable Personas}{Honorable mention in Supabase YC Hackathon}{Oct 2024 -- Dec 2024}
  \begin{itemize}
    \item Designed an \textbf{agentic LLM architecture} with LangGraph to mirror user personalities, creating interactive digital
          personas, and deployed it as a full-stack platform on Supabase and Vercel with one-click publishing
  \end{itemize}}

\newcommand{\LLMReasoning}{%
  \entry{Improving LLM Reasoning for Numerical Problems}{\report{https://sudhansh6.github.io/assets/pdf/MathPrompter++.pdf}}{Sep 2024 -- Dec 2024}
  \begin{itemize}
    \item Enhanced MathPrompter (ACL 2023) with CoT, achieving \textbf{10\% higher accuracy} on Llama 3.1 1B where prior methods failed,
          and reduced hallucination rates significantly by integrating multi-step validation
  \end{itemize}}

\newcommand{\InverseRendering}{%
  \entry{Inverse Rendering with 2D Gaussian Splatting}{\report{https://sudhansh6.github.io/assets/pdf/CSE252D_Report.pdf}}{Mar 2024 -- May 2024}
  \begin{itemize}
    \item Developed a novel \textbf{inverse rendering} framework in CUDA to recover PBR properties of a scene using 2D Gaussian
          Splatting, improving normal map MAE by \textbf{15\%} over the current SOTA methods
  \end{itemize}}

% ---------------------------------------------------------------------------------------------------
% Publications: title on the left, venue on the right, one line each.
\newcommand{\publine}[3]{\par\noindent\href{#1}{#2}\hfill\hspace{1em}\side{#3}\par\addvspace{1.5pt}}
\newcommand{\PubSaturation}{\publine{https://doi.org/10.1016/j.sigpro.2023.109349}{A likelihood based method for compressive signal recovery under Gaussian and saturation noise}{Signal Processing, 2024}}
\newcommand{\PubPermutation}{\publine{https://ieeexplore.ieee.org/document/10888147}{Identification and correction of permutation errors in compressed sensing-based group testing}{ICASSP 2025}}
\newcommand{\PubRearrangement}{\publine{https://arxiv.org/abs/2411.10899}{Planning for tabletop object rearrangement}{ICRA 2025}}

% ---------------------------------------------------------------------------------------------------
% Education, skills, service
% Each profile sets \UCSDCourses, \IITBCourses and \SkillRows before \input{resume-common}.
\providecommand{\UCSDCourses}{}
\providecommand{\IITBCourses}{}
\providecommand{\SkillRows}{}
\newcommand{\coursework}[1]{\ifblank{#1}{}{\par\noindent{\fontsize{9.5}{11.5}\selectfont{\color{muted}Coursework:} #1}\par}}
\newcommand{\Education}{%
  \section{Education}
  \entry{University of California San Diego}{M.S. in Computer Science and Engineering (AI)}{GPA 4.00/4.00 \quad Sep 2023 -- Mar 2025}
  \coursework{\UCSDCourses}
  \entry{Indian Institute of Technology Bombay}{B.Tech.\ with Honors in Computer Science}{CPI 9.66/10 \quad Jul 2019 -- Jul 2023}
  \coursework{\IITBCourses}}

% \skill{Label}{Items} is one row of the skills table.
\newcommand{\skill}[2]{{\headfont\bfseries #1} & #2 \\[1.5pt]}
\newcommand{\cpp}{C\kern-0.04em\raisebox{0.06ex}{\small++}}
\newcommand{\Skills}{%
  \section{Skills}
  \begin{tabularx}{\linewidth}{@{}l@{\hspace{12pt}}>{\raggedright\arraybackslash}X@{}}
    \SkillRows
  \end{tabularx}}
 to try another pairing.
% Default: Charter body with Source Sans headings. All-sans alternative:
%   \def\resumefont{\usepackage[default,semibold]{sourcesanspro}}\def\headfont{}
\providecommand{\resumefont}{\usepackage[scaled=0.96]{XCharter}\usepackage[semibold,scale=0.96]{sourcesanspro}}
\providecommand{\headfont}{\sffamily}
\providecommand{\namefont}{\rmfamily\scshape}
\resumefont
\usepackage[final]{microtype}
\usepackage[dvipsnames]{xcolor}
\usepackage{enumitem}
\usepackage{titlesec}
\usepackage{tabularx}
\usepackage{fontawesome5}
\usepackage{etoolbox}
\usepackage[hidelinks]{hyperref}

% The site's green, darkened for print.
\definecolor{accent}{HTML}{3A7A45}
\definecolor{muted}{HTML}{5C675D}

\pagestyle{empty}
\setlength{\parindent}{0pt}
\setlength{\tabcolsep}{0pt}
\providecommand{\leading}{1.1}
\linespread{\leading}

\hypersetup{pdfauthor={Sudhansh Peddabomma}}

% Type scale: name 26pt small caps, section titles 12.5pt caps, entry titles 10.8pt, body 10pt,
% side notes 9pt.
\newcommand{\side}[1]{{\fontsize{9}{11}\selectfont\color{muted}#1}}
\newcommand{\sectitle}[1]{\textls[60]{\MakeUppercase{#1}}}
\titleformat{\section}{\headfont\color{accent}\bfseries\fontsize{12.5}{14}\selectfont}{}{0pt}{\sectitle}[\vspace{1.5pt}{\color{accent}\titlerule[0.8pt]}\global\firstentrytrue]
% The first entry after a section title skips \entrygap, so only \rulegap sits under the rule.
\newif\iffirstentry

% Vertical rhythm, the same idea as the earlier hand-tuned résumés: bullets inside an entry sit
% tight together, entries are separated by a clear gap, and sections by a larger one. A profile
% may override these before % Shared preamble and content for Sudhansh Peddabomma's résumés.
% Each profile (SudhanshPeddabomma-AI.tex, SudhanshPeddabomma-Vision.tex) inputs this file and picks
% which blocks to show and in what order, so every bullet is written once.
% Layout rule: every bullet fills one line, or runs 1.5 to 2 lines. Check with: python3 check_lines.py <pdf>
\documentclass[10pt,letterpaper]{article}

\usepackage[margin=0.55in,top=0.5in,bottom=0.45in]{geometry}
\usepackage[T1]{fontenc}
% Fonts: \resumefont loads the body font; \headfont styles the name and section titles.
% Override either before \input{resume-common} to try another pairing.
% Default: Charter body with Source Sans headings. All-sans alternative:
%   \def\resumefont{\usepackage[default,semibold]{sourcesanspro}}\def\headfont{}
\providecommand{\resumefont}{\usepackage[scaled=0.96]{XCharter}\usepackage[semibold,scale=0.96]{sourcesanspro}}
\providecommand{\headfont}{\sffamily}
\providecommand{\namefont}{\rmfamily\scshape}
\resumefont
\usepackage[final]{microtype}
\usepackage[dvipsnames]{xcolor}
\usepackage{enumitem}
\usepackage{titlesec}
\usepackage{tabularx}
\usepackage{fontawesome5}
\usepackage{etoolbox}
\usepackage[hidelinks]{hyperref}

% The site's green, darkened for print.
\definecolor{accent}{HTML}{3A7A45}
\definecolor{muted}{HTML}{5C675D}

\pagestyle{empty}
\setlength{\parindent}{0pt}
\setlength{\tabcolsep}{0pt}
\providecommand{\leading}{1.1}
\linespread{\leading}

\hypersetup{pdfauthor={Sudhansh Peddabomma}}

% Type scale: name 26pt small caps, section titles 12.5pt caps, entry titles 10.8pt, body 10pt,
% side notes 9pt.
\newcommand{\side}[1]{{\fontsize{9}{11}\selectfont\color{muted}#1}}
\newcommand{\sectitle}[1]{\textls[60]{\MakeUppercase{#1}}}
\titleformat{\section}{\headfont\color{accent}\bfseries\fontsize{12.5}{14}\selectfont}{}{0pt}{\sectitle}[\vspace{1.5pt}{\color{accent}\titlerule[0.8pt]}\global\firstentrytrue]
% The first entry after a section title skips \entrygap, so only \rulegap sits under the rule.
\newif\iffirstentry

% Vertical rhythm, the same idea as the earlier hand-tuned résumés: bullets inside an entry sit
% tight together, entries are separated by a clear gap, and sections by a larger one. A profile
% may override these before \input{resume-common} to fit one page.
\providecommand{\secgap}{12pt}   % above a section title
\providecommand{\rulegap}{5.5pt}   % section rule to first entry
\providecommand{\entrygap}{7pt}  % between entries
\providecommand{\headgap}{1.5pt} % entry heading to its bullets
\providecommand{\itemgap}{0.6pt} % between bullets of one entry
\titlespacing*{\section}{0pt}{\secgap}{\rulegap}

\newcommand{\squarebullet}{\raisebox{0.32ex}{\color{accent}\rule{0.5ex}{0.5ex}}}
\setlist[itemize]{leftmargin=12pt, labelsep=6pt, topsep=\headgap, itemsep=\itemgap, parsep=0pt, partopsep=0pt,
  label=\squarebullet}

% \entry{Organisation or title}{Role or detail (may be empty)}{Dates}; \sub{Role or detail}{Right-hand note}
\newcommand{\entry}[3]{%
  \par\iffirstentry\global\firstentryfalse\else\addvspace{\entrygap}\fi\noindent{\headfont\fontsize{10.8}{13}\selectfont\textbf{#1}}\ifblank{#2}{}{\ {\color{muted}|}\ #2}\hfill\hspace{1em}\side{#3}\par}
\newcommand{\sub}[2]{%
  \noindent\textit{#1}\hfill{\color{muted}\textit{#2}}\par}
\newcommand{\report}[1]{\href{#1}{\color{accent}Report}}
\newcommand{\me}{\textbf{S. Peddabomma}}

% ---------------------------------------------------------------------------------------------------
% Header
\newcommand{\Header}{%
  \begin{center}
    {\namefont\fontsize{26}{28}\selectfont Sudhansh Peddabomma}\\[6pt]
    {\headfont\fontsize{9.5}{11}\selectfont\color{muted}
      \href{mailto:sudhansh6@gmail.com}{\faEnvelope[regular]\ sudhansh6@gmail.com} \quad
      \href{https://sudhansh6.github.io}{\faGlobe\ sudhansh6.github.io} \quad
      \href{https://www.linkedin.com/in/sudhansh-peddabomma-73a208192/}{\faLinkedin\ sudhansh-peddabomma} \quad
      \href{https://github.com/Sudhansh6}{\faGithub\ Sudhansh6}}
  \end{center}\vspace{-4pt}}

% ---------------------------------------------------------------------------------------------------
% Glean bullets, shared by both profiles
\newcommand{\GleanSlides}{\item Led \textbf{AI slide generation} in Glean Assistant, now creating \textbf{180K decks a week} for \textbf{50K weekly active users} at enterprises}
\newcommand{\GleanSlidesPerf}{\item Improved the slides harness, raising reliability to \textbf{97\%}, cutting median latency \textbf{26\%} and cost per deck by \textbf{half}
  in four months}
\newcommand{\GleanSlidesOrch}{\item Architected a \textbf{slide orchestration system} that reuses relevant slides from a company's existing decks and generates
  new ones in its brand and template, so teams can refine decks with each other and Glean in full context, without leaving the product}
\newcommand{\GleanAgents}{\item Moved Glean's \textbf{workflow agents} from one-shot planning to an \textbf{agentic tool-use loop}, greatly improving agentic
  capabilities and instruction following, with \textbf{18\%} higher answer completeness and \textbf{17\%} fewer tokens}
\newcommand{\GleanTeamChat}{\item Engineered the core of \textbf{Team chat}, letting teammates work with Glean in one shared conversation and co-edit its artifacts}
\newcommand{\GleanVideo}{\item Shipped \textbf{video artifacts} end to end, generating narrated videos in chat with streamed playback and admin rollout controls}
\newcommand{\GleanEvals}{\item Created the \textbf{evaluations and metrics} that guide model choices and rollouts across agents, slides and interactive artifacts}

% ---------------------------------------------------------------------------------------------------
% Experience
% \groupbreak separates related groups of bullets inside one long entry (slides, then platform work).
\providecommand{\groupgap}{3.5pt}
\newcommand{\groupbreak}{\vspace{\groupgap}}
\newcommand{\GleanAI}{%
  \entry{Glean}{Software Engineer}{Apr 2025 -- Present}
  \begin{itemize}
    \GleanSlides
    \GleanSlidesOrch
    \GleanSlidesPerf \groupbreak
    \GleanTeamChat
    \GleanVideo
    \GleanAgents
    \GleanEvals
  \end{itemize}}

\newcommand{\GleanVision}{%
  \entry{Glean}{Software Engineer}{Apr 2025 -- Present}
  \begin{itemize}
    \GleanSlides
    \GleanSlidesOrch
    \GleanSlidesPerf \groupbreak
    \GleanVideo
    \GleanEvals
  \end{itemize}}

\newcommand{\Duality}{%
  \entry{Duality AI}{Computer Vision Intern}{Jun 2024 -- Sep 2024}
  \begin{itemize}
    \item Designed automated \textbf{3D reconstruction} techniques for featureless objects, reducing digital-twin generation time
          by \textbf{40\%}
    \item Collaborated with \textbf{Autodesk} to validate Unreal Engine simulations for robotics tasks; leveraging structured
          domain randomization to reduce Sim2Real gap and increase mAP-50 by \textbf{15\%} for object detection and segmentation
  \end{itemize}}

\newcommand{\Microsoft}{%
  \entry{Microsoft}{Data and Applied Scientist Intern}{May 2022 -- Jul 2022}
  \begin{itemize}
    \item Developed a \textbf{decision-tree ranker} to recommend emails without user queries, improving \textbf{Outlook} search capabilities
    \item Proposed hierarchical feature-sets for the ranker, reducing latency for recommendations and improving recall
  \end{itemize}}

% ---------------------------------------------------------------------------------------------------
% Research
\newcommand{\HomeRobots}{%
  \entry{Real-time 3D Perception for Home Robots}{UC San Diego}{Sep 2023 -- Sep 2024}
  \begin{itemize}
    \item Investigated real-time dense \textbf{visual SLAM} methods using NeRFs and Gaussian Splatting for indoor robot navigation
    \item Integrated object segmentation, grasp-pose estimation, and 3D mapping on the Fetch robot via ROS, demonstrating
          a novel \textbf{tabletop object rearrangement} algorithm that reduced cost by \textbf{20\%} compared to the state-of-the-art
  \end{itemize}}

\newcommand{\Tomography}{%
  \entry{3D Tomography with Primal-Dual Neural Networks}{University College London}{May 2021 -- Jul 2023}
  \begin{itemize}
    \item Developed a stochastic neural-network architecture of primal-dual algorithm for online reconstruction of 3D
          volumes from tomographic projections, with \textbf{99.6\,\% structural similarity} in low-dosage conditions
    \item Built a Python library with custom gradient operators for reconstructing volumes in a single pass, reducing
          compute requirements by up to \textbf{5x} over SOTA learning-based approaches for cone vector tomography
  \end{itemize}}

% ---------------------------------------------------------------------------------------------------
% Projects
\newcommand{\OOXML}{%
  \entry{ooxml}{Open-source contributor \textbar\ \href{https://github.com/ChristopherVR/ooxml}{\color{accent}GitHub}}{Sep 2026 -- Present}
  \begin{itemize}
    \item Contributed to an open-source \textbf{TypeScript} engine that renders and edits PowerPoint and Word files in the browser,
          improving editing fidelity, \textbf{real-time collaboration} and load performance across React, Vue and Svelte viewers
  \end{itemize}}

\newcommand{\MirrorAI}{%
  \entry{Mirror AI: Deployable Personas}{Honorable mention in Supabase YC Hackathon}{Oct 2024 -- Dec 2024}
  \begin{itemize}
    \item Designed an \textbf{agentic LLM architecture} with LangGraph to mirror user personalities, creating interactive digital
          personas, and deployed it as a full-stack platform on Supabase and Vercel with one-click publishing
  \end{itemize}}

\newcommand{\LLMReasoning}{%
  \entry{Improving LLM Reasoning for Numerical Problems}{\report{https://sudhansh6.github.io/assets/pdf/MathPrompter++.pdf}}{Sep 2024 -- Dec 2024}
  \begin{itemize}
    \item Enhanced MathPrompter (ACL 2023) with CoT, achieving \textbf{10\% higher accuracy} on Llama 3.1 1B where prior methods failed,
          and reduced hallucination rates significantly by integrating multi-step validation
  \end{itemize}}

\newcommand{\InverseRendering}{%
  \entry{Inverse Rendering with 2D Gaussian Splatting}{\report{https://sudhansh6.github.io/assets/pdf/CSE252D_Report.pdf}}{Mar 2024 -- May 2024}
  \begin{itemize}
    \item Developed a novel \textbf{inverse rendering} framework in CUDA to recover PBR properties of a scene using 2D Gaussian
          Splatting, improving normal map MAE by \textbf{15\%} over the current SOTA methods
  \end{itemize}}

% ---------------------------------------------------------------------------------------------------
% Publications: title on the left, venue on the right, one line each.
\newcommand{\publine}[3]{\par\noindent\href{#1}{#2}\hfill\hspace{1em}\side{#3}\par\addvspace{1.5pt}}
\newcommand{\PubSaturation}{\publine{https://doi.org/10.1016/j.sigpro.2023.109349}{A likelihood based method for compressive signal recovery under Gaussian and saturation noise}{Signal Processing, 2024}}
\newcommand{\PubPermutation}{\publine{https://ieeexplore.ieee.org/document/10888147}{Identification and correction of permutation errors in compressed sensing-based group testing}{ICASSP 2025}}
\newcommand{\PubRearrangement}{\publine{https://arxiv.org/abs/2411.10899}{Planning for tabletop object rearrangement}{ICRA 2025}}

% ---------------------------------------------------------------------------------------------------
% Education, skills, service
% Each profile sets \UCSDCourses, \IITBCourses and \SkillRows before \input{resume-common}.
\providecommand{\UCSDCourses}{}
\providecommand{\IITBCourses}{}
\providecommand{\SkillRows}{}
\newcommand{\coursework}[1]{\ifblank{#1}{}{\par\noindent{\fontsize{9.5}{11.5}\selectfont{\color{muted}Coursework:} #1}\par}}
\newcommand{\Education}{%
  \section{Education}
  \entry{University of California San Diego}{M.S. in Computer Science and Engineering (AI)}{GPA 4.00/4.00 \quad Sep 2023 -- Mar 2025}
  \coursework{\UCSDCourses}
  \entry{Indian Institute of Technology Bombay}{B.Tech.\ with Honors in Computer Science}{CPI 9.66/10 \quad Jul 2019 -- Jul 2023}
  \coursework{\IITBCourses}}

% \skill{Label}{Items} is one row of the skills table.
\newcommand{\skill}[2]{{\headfont\bfseries #1} & #2 \\[1.5pt]}
\newcommand{\cpp}{C\kern-0.04em\raisebox{0.06ex}{\small++}}
\newcommand{\Skills}{%
  \section{Skills}
  \begin{tabularx}{\linewidth}{@{}l@{\hspace{12pt}}>{\raggedright\arraybackslash}X@{}}
    \SkillRows
  \end{tabularx}}
 to fit one page.
\providecommand{\secgap}{12pt}   % above a section title
\providecommand{\rulegap}{5.5pt}   % section rule to first entry
\providecommand{\entrygap}{7pt}  % between entries
\providecommand{\headgap}{1.5pt} % entry heading to its bullets
\providecommand{\itemgap}{0.6pt} % between bullets of one entry
\titlespacing*{\section}{0pt}{\secgap}{\rulegap}

\newcommand{\squarebullet}{\raisebox{0.32ex}{\color{accent}\rule{0.5ex}{0.5ex}}}
\setlist[itemize]{leftmargin=12pt, labelsep=6pt, topsep=\headgap, itemsep=\itemgap, parsep=0pt, partopsep=0pt,
  label=\squarebullet}

% \entry{Organisation or title}{Role or detail (may be empty)}{Dates}; \sub{Role or detail}{Right-hand note}
\newcommand{\entry}[3]{%
  \par\iffirstentry\global\firstentryfalse\else\addvspace{\entrygap}\fi\noindent{\headfont\fontsize{10.8}{13}\selectfont\textbf{#1}}\ifblank{#2}{}{\ {\color{muted}|}\ #2}\hfill\hspace{1em}\side{#3}\par}
\newcommand{\sub}[2]{%
  \noindent\textit{#1}\hfill{\color{muted}\textit{#2}}\par}
\newcommand{\report}[1]{\href{#1}{\color{accent}Report}}
\newcommand{\me}{\textbf{S. Peddabomma}}

% ---------------------------------------------------------------------------------------------------
% Header
\newcommand{\Header}{%
  \begin{center}
    {\namefont\fontsize{26}{28}\selectfont Sudhansh Peddabomma}\\[6pt]
    {\headfont\fontsize{9.5}{11}\selectfont\color{muted}
      \href{mailto:sudhansh6@gmail.com}{\faEnvelope[regular]\ sudhansh6@gmail.com} \quad
      \href{https://sudhansh6.github.io}{\faGlobe\ sudhansh6.github.io} \quad
      \href{https://www.linkedin.com/in/sudhansh-peddabomma-73a208192/}{\faLinkedin\ sudhansh-peddabomma} \quad
      \href{https://github.com/Sudhansh6}{\faGithub\ Sudhansh6}}
  \end{center}\vspace{-4pt}}

% ---------------------------------------------------------------------------------------------------
% Glean bullets, shared by both profiles
\newcommand{\GleanSlides}{\item Led \textbf{AI slide generation} in Glean Assistant, now creating \textbf{180K decks a week} for \textbf{50K weekly active users} at enterprises}
\newcommand{\GleanSlidesPerf}{\item Improved the slides harness, raising reliability to \textbf{97\%}, cutting median latency \textbf{26\%} and cost per deck by \textbf{half}
  in four months}
\newcommand{\GleanSlidesOrch}{\item Architected a \textbf{slide orchestration system} that reuses relevant slides from a company's existing decks and generates
  new ones in its brand and template, so teams can refine decks with each other and Glean in full context, without leaving the product}
\newcommand{\GleanAgents}{\item Moved Glean's \textbf{workflow agents} from one-shot planning to an \textbf{agentic tool-use loop}, greatly improving agentic
  capabilities and instruction following, with \textbf{18\%} higher answer completeness and \textbf{17\%} fewer tokens}
\newcommand{\GleanTeamChat}{\item Engineered the core of \textbf{Team chat}, letting teammates work with Glean in one shared conversation and co-edit its artifacts}
\newcommand{\GleanVideo}{\item Shipped \textbf{video artifacts} end to end, generating narrated videos in chat with streamed playback and admin rollout controls}
\newcommand{\GleanEvals}{\item Created the \textbf{evaluations and metrics} that guide model choices and rollouts across agents, slides and interactive artifacts}

% ---------------------------------------------------------------------------------------------------
% Experience
% \groupbreak separates related groups of bullets inside one long entry (slides, then platform work).
\providecommand{\groupgap}{3.5pt}
\newcommand{\groupbreak}{\vspace{\groupgap}}
\newcommand{\GleanAI}{%
  \entry{Glean}{Software Engineer}{Apr 2025 -- Present}
  \begin{itemize}
    \GleanSlides
    \GleanSlidesOrch
    \GleanSlidesPerf \groupbreak
    \GleanTeamChat
    \GleanVideo
    \GleanAgents
    \GleanEvals
  \end{itemize}}

\newcommand{\GleanVision}{%
  \entry{Glean}{Software Engineer}{Apr 2025 -- Present}
  \begin{itemize}
    \GleanSlides
    \GleanSlidesOrch
    \GleanSlidesPerf \groupbreak
    \GleanVideo
    \GleanEvals
  \end{itemize}}

\newcommand{\Duality}{%
  \entry{Duality AI}{Computer Vision Intern}{Jun 2024 -- Sep 2024}
  \begin{itemize}
    \item Designed automated \textbf{3D reconstruction} techniques for featureless objects, reducing digital-twin generation time
          by \textbf{40\%}
    \item Collaborated with \textbf{Autodesk} to validate Unreal Engine simulations for robotics tasks; leveraging structured
          domain randomization to reduce Sim2Real gap and increase mAP-50 by \textbf{15\%} for object detection and segmentation
  \end{itemize}}

\newcommand{\Microsoft}{%
  \entry{Microsoft}{Data and Applied Scientist Intern}{May 2022 -- Jul 2022}
  \begin{itemize}
    \item Developed a \textbf{decision-tree ranker} to recommend emails without user queries, improving \textbf{Outlook} search capabilities
    \item Proposed hierarchical feature-sets for the ranker, reducing latency for recommendations and improving recall
  \end{itemize}}

% ---------------------------------------------------------------------------------------------------
% Research
\newcommand{\HomeRobots}{%
  \entry{Real-time 3D Perception for Home Robots}{UC San Diego}{Sep 2023 -- Sep 2024}
  \begin{itemize}
    \item Investigated real-time dense \textbf{visual SLAM} methods using NeRFs and Gaussian Splatting for indoor robot navigation
    \item Integrated object segmentation, grasp-pose estimation, and 3D mapping on the Fetch robot via ROS, demonstrating
          a novel \textbf{tabletop object rearrangement} algorithm that reduced cost by \textbf{20\%} compared to the state-of-the-art
  \end{itemize}}

\newcommand{\Tomography}{%
  \entry{3D Tomography with Primal-Dual Neural Networks}{University College London}{May 2021 -- Jul 2023}
  \begin{itemize}
    \item Developed a stochastic neural-network architecture of primal-dual algorithm for online reconstruction of 3D
          volumes from tomographic projections, with \textbf{99.6\,\% structural similarity} in low-dosage conditions
    \item Built a Python library with custom gradient operators for reconstructing volumes in a single pass, reducing
          compute requirements by up to \textbf{5x} over SOTA learning-based approaches for cone vector tomography
  \end{itemize}}

% ---------------------------------------------------------------------------------------------------
% Projects
\newcommand{\OOXML}{%
  \entry{ooxml}{Open-source contributor \textbar\ \href{https://github.com/ChristopherVR/ooxml}{\color{accent}GitHub}}{Sep 2026 -- Present}
  \begin{itemize}
    \item Contributed to an open-source \textbf{TypeScript} engine that renders and edits PowerPoint and Word files in the browser,
          improving editing fidelity, \textbf{real-time collaboration} and load performance across React, Vue and Svelte viewers
  \end{itemize}}

\newcommand{\MirrorAI}{%
  \entry{Mirror AI: Deployable Personas}{Honorable mention in Supabase YC Hackathon}{Oct 2024 -- Dec 2024}
  \begin{itemize}
    \item Designed an \textbf{agentic LLM architecture} with LangGraph to mirror user personalities, creating interactive digital
          personas, and deployed it as a full-stack platform on Supabase and Vercel with one-click publishing
  \end{itemize}}

\newcommand{\LLMReasoning}{%
  \entry{Improving LLM Reasoning for Numerical Problems}{\report{https://sudhansh6.github.io/assets/pdf/MathPrompter++.pdf}}{Sep 2024 -- Dec 2024}
  \begin{itemize}
    \item Enhanced MathPrompter (ACL 2023) with CoT, achieving \textbf{10\% higher accuracy} on Llama 3.1 1B where prior methods failed,
          and reduced hallucination rates significantly by integrating multi-step validation
  \end{itemize}}

\newcommand{\InverseRendering}{%
  \entry{Inverse Rendering with 2D Gaussian Splatting}{\report{https://sudhansh6.github.io/assets/pdf/CSE252D_Report.pdf}}{Mar 2024 -- May 2024}
  \begin{itemize}
    \item Developed a novel \textbf{inverse rendering} framework in CUDA to recover PBR properties of a scene using 2D Gaussian
          Splatting, improving normal map MAE by \textbf{15\%} over the current SOTA methods
  \end{itemize}}

% ---------------------------------------------------------------------------------------------------
% Publications: title on the left, venue on the right, one line each.
\newcommand{\publine}[3]{\par\noindent\href{#1}{#2}\hfill\hspace{1em}\side{#3}\par\addvspace{1.5pt}}
\newcommand{\PubSaturation}{\publine{https://doi.org/10.1016/j.sigpro.2023.109349}{A likelihood based method for compressive signal recovery under Gaussian and saturation noise}{Signal Processing, 2024}}
\newcommand{\PubPermutation}{\publine{https://ieeexplore.ieee.org/document/10888147}{Identification and correction of permutation errors in compressed sensing-based group testing}{ICASSP 2025}}
\newcommand{\PubRearrangement}{\publine{https://arxiv.org/abs/2411.10899}{Planning for tabletop object rearrangement}{ICRA 2025}}

% ---------------------------------------------------------------------------------------------------
% Education, skills, service
% Each profile sets \UCSDCourses, \IITBCourses and \SkillRows before % Shared preamble and content for Sudhansh Peddabomma's résumés.
% Each profile (SudhanshPeddabomma-AI.tex, SudhanshPeddabomma-Vision.tex) inputs this file and picks
% which blocks to show and in what order, so every bullet is written once.
% Layout rule: every bullet fills one line, or runs 1.5 to 2 lines. Check with: python3 check_lines.py <pdf>
\documentclass[10pt,letterpaper]{article}

\usepackage[margin=0.55in,top=0.5in,bottom=0.45in]{geometry}
\usepackage[T1]{fontenc}
% Fonts: \resumefont loads the body font; \headfont styles the name and section titles.
% Override either before \input{resume-common} to try another pairing.
% Default: Charter body with Source Sans headings. All-sans alternative:
%   \def\resumefont{\usepackage[default,semibold]{sourcesanspro}}\def\headfont{}
\providecommand{\resumefont}{\usepackage[scaled=0.96]{XCharter}\usepackage[semibold,scale=0.96]{sourcesanspro}}
\providecommand{\headfont}{\sffamily}
\providecommand{\namefont}{\rmfamily\scshape}
\resumefont
\usepackage[final]{microtype}
\usepackage[dvipsnames]{xcolor}
\usepackage{enumitem}
\usepackage{titlesec}
\usepackage{tabularx}
\usepackage{fontawesome5}
\usepackage{etoolbox}
\usepackage[hidelinks]{hyperref}

% The site's green, darkened for print.
\definecolor{accent}{HTML}{3A7A45}
\definecolor{muted}{HTML}{5C675D}

\pagestyle{empty}
\setlength{\parindent}{0pt}
\setlength{\tabcolsep}{0pt}
\providecommand{\leading}{1.1}
\linespread{\leading}

\hypersetup{pdfauthor={Sudhansh Peddabomma}}

% Type scale: name 26pt small caps, section titles 12.5pt caps, entry titles 10.8pt, body 10pt,
% side notes 9pt.
\newcommand{\side}[1]{{\fontsize{9}{11}\selectfont\color{muted}#1}}
\newcommand{\sectitle}[1]{\textls[60]{\MakeUppercase{#1}}}
\titleformat{\section}{\headfont\color{accent}\bfseries\fontsize{12.5}{14}\selectfont}{}{0pt}{\sectitle}[\vspace{1.5pt}{\color{accent}\titlerule[0.8pt]}\global\firstentrytrue]
% The first entry after a section title skips \entrygap, so only \rulegap sits under the rule.
\newif\iffirstentry

% Vertical rhythm, the same idea as the earlier hand-tuned résumés: bullets inside an entry sit
% tight together, entries are separated by a clear gap, and sections by a larger one. A profile
% may override these before \input{resume-common} to fit one page.
\providecommand{\secgap}{12pt}   % above a section title
\providecommand{\rulegap}{5.5pt}   % section rule to first entry
\providecommand{\entrygap}{7pt}  % between entries
\providecommand{\headgap}{1.5pt} % entry heading to its bullets
\providecommand{\itemgap}{0.6pt} % between bullets of one entry
\titlespacing*{\section}{0pt}{\secgap}{\rulegap}

\newcommand{\squarebullet}{\raisebox{0.32ex}{\color{accent}\rule{0.5ex}{0.5ex}}}
\setlist[itemize]{leftmargin=12pt, labelsep=6pt, topsep=\headgap, itemsep=\itemgap, parsep=0pt, partopsep=0pt,
  label=\squarebullet}

% \entry{Organisation or title}{Role or detail (may be empty)}{Dates}; \sub{Role or detail}{Right-hand note}
\newcommand{\entry}[3]{%
  \par\iffirstentry\global\firstentryfalse\else\addvspace{\entrygap}\fi\noindent{\headfont\fontsize{10.8}{13}\selectfont\textbf{#1}}\ifblank{#2}{}{\ {\color{muted}|}\ #2}\hfill\hspace{1em}\side{#3}\par}
\newcommand{\sub}[2]{%
  \noindent\textit{#1}\hfill{\color{muted}\textit{#2}}\par}
\newcommand{\report}[1]{\href{#1}{\color{accent}Report}}
\newcommand{\me}{\textbf{S. Peddabomma}}

% ---------------------------------------------------------------------------------------------------
% Header
\newcommand{\Header}{%
  \begin{center}
    {\namefont\fontsize{26}{28}\selectfont Sudhansh Peddabomma}\\[6pt]
    {\headfont\fontsize{9.5}{11}\selectfont\color{muted}
      \href{mailto:sudhansh6@gmail.com}{\faEnvelope[regular]\ sudhansh6@gmail.com} \quad
      \href{https://sudhansh6.github.io}{\faGlobe\ sudhansh6.github.io} \quad
      \href{https://www.linkedin.com/in/sudhansh-peddabomma-73a208192/}{\faLinkedin\ sudhansh-peddabomma} \quad
      \href{https://github.com/Sudhansh6}{\faGithub\ Sudhansh6}}
  \end{center}\vspace{-4pt}}

% ---------------------------------------------------------------------------------------------------
% Glean bullets, shared by both profiles
\newcommand{\GleanSlides}{\item Led \textbf{AI slide generation} in Glean Assistant, now creating \textbf{180K decks a week} for \textbf{50K weekly active users} at enterprises}
\newcommand{\GleanSlidesPerf}{\item Improved the slides harness, raising reliability to \textbf{97\%}, cutting median latency \textbf{26\%} and cost per deck by \textbf{half}
  in four months}
\newcommand{\GleanSlidesOrch}{\item Architected a \textbf{slide orchestration system} that reuses relevant slides from a company's existing decks and generates
  new ones in its brand and template, so teams can refine decks with each other and Glean in full context, without leaving the product}
\newcommand{\GleanAgents}{\item Moved Glean's \textbf{workflow agents} from one-shot planning to an \textbf{agentic tool-use loop}, greatly improving agentic
  capabilities and instruction following, with \textbf{18\%} higher answer completeness and \textbf{17\%} fewer tokens}
\newcommand{\GleanTeamChat}{\item Engineered the core of \textbf{Team chat}, letting teammates work with Glean in one shared conversation and co-edit its artifacts}
\newcommand{\GleanVideo}{\item Shipped \textbf{video artifacts} end to end, generating narrated videos in chat with streamed playback and admin rollout controls}
\newcommand{\GleanEvals}{\item Created the \textbf{evaluations and metrics} that guide model choices and rollouts across agents, slides and interactive artifacts}

% ---------------------------------------------------------------------------------------------------
% Experience
% \groupbreak separates related groups of bullets inside one long entry (slides, then platform work).
\providecommand{\groupgap}{3.5pt}
\newcommand{\groupbreak}{\vspace{\groupgap}}
\newcommand{\GleanAI}{%
  \entry{Glean}{Software Engineer}{Apr 2025 -- Present}
  \begin{itemize}
    \GleanSlides
    \GleanSlidesOrch
    \GleanSlidesPerf \groupbreak
    \GleanTeamChat
    \GleanVideo
    \GleanAgents
    \GleanEvals
  \end{itemize}}

\newcommand{\GleanVision}{%
  \entry{Glean}{Software Engineer}{Apr 2025 -- Present}
  \begin{itemize}
    \GleanSlides
    \GleanSlidesOrch
    \GleanSlidesPerf \groupbreak
    \GleanVideo
    \GleanEvals
  \end{itemize}}

\newcommand{\Duality}{%
  \entry{Duality AI}{Computer Vision Intern}{Jun 2024 -- Sep 2024}
  \begin{itemize}
    \item Designed automated \textbf{3D reconstruction} techniques for featureless objects, reducing digital-twin generation time
          by \textbf{40\%}
    \item Collaborated with \textbf{Autodesk} to validate Unreal Engine simulations for robotics tasks; leveraging structured
          domain randomization to reduce Sim2Real gap and increase mAP-50 by \textbf{15\%} for object detection and segmentation
  \end{itemize}}

\newcommand{\Microsoft}{%
  \entry{Microsoft}{Data and Applied Scientist Intern}{May 2022 -- Jul 2022}
  \begin{itemize}
    \item Developed a \textbf{decision-tree ranker} to recommend emails without user queries, improving \textbf{Outlook} search capabilities
    \item Proposed hierarchical feature-sets for the ranker, reducing latency for recommendations and improving recall
  \end{itemize}}

% ---------------------------------------------------------------------------------------------------
% Research
\newcommand{\HomeRobots}{%
  \entry{Real-time 3D Perception for Home Robots}{UC San Diego}{Sep 2023 -- Sep 2024}
  \begin{itemize}
    \item Investigated real-time dense \textbf{visual SLAM} methods using NeRFs and Gaussian Splatting for indoor robot navigation
    \item Integrated object segmentation, grasp-pose estimation, and 3D mapping on the Fetch robot via ROS, demonstrating
          a novel \textbf{tabletop object rearrangement} algorithm that reduced cost by \textbf{20\%} compared to the state-of-the-art
  \end{itemize}}

\newcommand{\Tomography}{%
  \entry{3D Tomography with Primal-Dual Neural Networks}{University College London}{May 2021 -- Jul 2023}
  \begin{itemize}
    \item Developed a stochastic neural-network architecture of primal-dual algorithm for online reconstruction of 3D
          volumes from tomographic projections, with \textbf{99.6\,\% structural similarity} in low-dosage conditions
    \item Built a Python library with custom gradient operators for reconstructing volumes in a single pass, reducing
          compute requirements by up to \textbf{5x} over SOTA learning-based approaches for cone vector tomography
  \end{itemize}}

% ---------------------------------------------------------------------------------------------------
% Projects
\newcommand{\OOXML}{%
  \entry{ooxml}{Open-source contributor \textbar\ \href{https://github.com/ChristopherVR/ooxml}{\color{accent}GitHub}}{Sep 2026 -- Present}
  \begin{itemize}
    \item Contributed to an open-source \textbf{TypeScript} engine that renders and edits PowerPoint and Word files in the browser,
          improving editing fidelity, \textbf{real-time collaboration} and load performance across React, Vue and Svelte viewers
  \end{itemize}}

\newcommand{\MirrorAI}{%
  \entry{Mirror AI: Deployable Personas}{Honorable mention in Supabase YC Hackathon}{Oct 2024 -- Dec 2024}
  \begin{itemize}
    \item Designed an \textbf{agentic LLM architecture} with LangGraph to mirror user personalities, creating interactive digital
          personas, and deployed it as a full-stack platform on Supabase and Vercel with one-click publishing
  \end{itemize}}

\newcommand{\LLMReasoning}{%
  \entry{Improving LLM Reasoning for Numerical Problems}{\report{https://sudhansh6.github.io/assets/pdf/MathPrompter++.pdf}}{Sep 2024 -- Dec 2024}
  \begin{itemize}
    \item Enhanced MathPrompter (ACL 2023) with CoT, achieving \textbf{10\% higher accuracy} on Llama 3.1 1B where prior methods failed,
          and reduced hallucination rates significantly by integrating multi-step validation
  \end{itemize}}

\newcommand{\InverseRendering}{%
  \entry{Inverse Rendering with 2D Gaussian Splatting}{\report{https://sudhansh6.github.io/assets/pdf/CSE252D_Report.pdf}}{Mar 2024 -- May 2024}
  \begin{itemize}
    \item Developed a novel \textbf{inverse rendering} framework in CUDA to recover PBR properties of a scene using 2D Gaussian
          Splatting, improving normal map MAE by \textbf{15\%} over the current SOTA methods
  \end{itemize}}

% ---------------------------------------------------------------------------------------------------
% Publications: title on the left, venue on the right, one line each.
\newcommand{\publine}[3]{\par\noindent\href{#1}{#2}\hfill\hspace{1em}\side{#3}\par\addvspace{1.5pt}}
\newcommand{\PubSaturation}{\publine{https://doi.org/10.1016/j.sigpro.2023.109349}{A likelihood based method for compressive signal recovery under Gaussian and saturation noise}{Signal Processing, 2024}}
\newcommand{\PubPermutation}{\publine{https://ieeexplore.ieee.org/document/10888147}{Identification and correction of permutation errors in compressed sensing-based group testing}{ICASSP 2025}}
\newcommand{\PubRearrangement}{\publine{https://arxiv.org/abs/2411.10899}{Planning for tabletop object rearrangement}{ICRA 2025}}

% ---------------------------------------------------------------------------------------------------
% Education, skills, service
% Each profile sets \UCSDCourses, \IITBCourses and \SkillRows before \input{resume-common}.
\providecommand{\UCSDCourses}{}
\providecommand{\IITBCourses}{}
\providecommand{\SkillRows}{}
\newcommand{\coursework}[1]{\ifblank{#1}{}{\par\noindent{\fontsize{9.5}{11.5}\selectfont{\color{muted}Coursework:} #1}\par}}
\newcommand{\Education}{%
  \section{Education}
  \entry{University of California San Diego}{M.S. in Computer Science and Engineering (AI)}{GPA 4.00/4.00 \quad Sep 2023 -- Mar 2025}
  \coursework{\UCSDCourses}
  \entry{Indian Institute of Technology Bombay}{B.Tech.\ with Honors in Computer Science}{CPI 9.66/10 \quad Jul 2019 -- Jul 2023}
  \coursework{\IITBCourses}}

% \skill{Label}{Items} is one row of the skills table.
\newcommand{\skill}[2]{{\headfont\bfseries #1} & #2 \\[1.5pt]}
\newcommand{\cpp}{C\kern-0.04em\raisebox{0.06ex}{\small++}}
\newcommand{\Skills}{%
  \section{Skills}
  \begin{tabularx}{\linewidth}{@{}l@{\hspace{12pt}}>{\raggedright\arraybackslash}X@{}}
    \SkillRows
  \end{tabularx}}
.
\providecommand{\UCSDCourses}{}
\providecommand{\IITBCourses}{}
\providecommand{\SkillRows}{}
\newcommand{\coursework}[1]{\ifblank{#1}{}{\par\noindent{\fontsize{9.5}{11.5}\selectfont{\color{muted}Coursework:} #1}\par}}
\newcommand{\Education}{%
  \section{Education}
  \entry{University of California San Diego}{M.S. in Computer Science and Engineering (AI)}{GPA 4.00/4.00 \quad Sep 2023 -- Mar 2025}
  \coursework{\UCSDCourses}
  \entry{Indian Institute of Technology Bombay}{B.Tech.\ with Honors in Computer Science}{CPI 9.66/10 \quad Jul 2019 -- Jul 2023}
  \coursework{\IITBCourses}}

% \skill{Label}{Items} is one row of the skills table.
\newcommand{\skill}[2]{{\headfont\bfseries #1} & #2 \\[1.5pt]}
\newcommand{\cpp}{C\kern-0.04em\raisebox{0.06ex}{\small++}}
\newcommand{\Skills}{%
  \section{Skills}
  \begin{tabularx}{\linewidth}{@{}l@{\hspace{12pt}}>{\raggedright\arraybackslash}X@{}}
    \SkillRows
  \end{tabularx}}
 to try another pairing.
% Default: Charter body with Source Sans headings. All-sans alternative:
%   \def\resumefont{\usepackage[default,semibold]{sourcesanspro}}\def\headfont{}
\providecommand{\resumefont}{\usepackage[scaled=0.96]{XCharter}\usepackage[semibold,scale=0.96]{sourcesanspro}}
\providecommand{\headfont}{\sffamily}
\providecommand{\namefont}{\rmfamily\scshape}
\resumefont
\usepackage[final]{microtype}
\usepackage[dvipsnames]{xcolor}
\usepackage{enumitem}
\usepackage{titlesec}
\usepackage{tabularx}
\usepackage{fontawesome5}
\usepackage{etoolbox}
\usepackage[hidelinks]{hyperref}

% The site's green, darkened for print.
\definecolor{accent}{HTML}{3A7A45}
\definecolor{muted}{HTML}{5C675D}

\pagestyle{empty}
\setlength{\parindent}{0pt}
\setlength{\tabcolsep}{0pt}
\providecommand{\leading}{1.1}
\linespread{\leading}

\hypersetup{pdfauthor={Sudhansh Peddabomma}}

% Type scale: name 26pt small caps, section titles 12.5pt caps, entry titles 10.8pt, body 10pt,
% side notes 9pt.
\newcommand{\side}[1]{{\fontsize{9}{11}\selectfont\color{muted}#1}}
\newcommand{\sectitle}[1]{\textls[60]{\MakeUppercase{#1}}}
\titleformat{\section}{\headfont\color{accent}\bfseries\fontsize{12.5}{14}\selectfont}{}{0pt}{\sectitle}[\vspace{1.5pt}{\color{accent}\titlerule[0.8pt]}\global\firstentrytrue]
% The first entry after a section title skips \entrygap, so only \rulegap sits under the rule.
\newif\iffirstentry

% Vertical rhythm, the same idea as the earlier hand-tuned résumés: bullets inside an entry sit
% tight together, entries are separated by a clear gap, and sections by a larger one. A profile
% may override these before % Shared preamble and content for Sudhansh Peddabomma's résumés.
% Each profile (SudhanshPeddabomma-AI.tex, SudhanshPeddabomma-Vision.tex) inputs this file and picks
% which blocks to show and in what order, so every bullet is written once.
% Layout rule: every bullet fills one line, or runs 1.5 to 2 lines. Check with: python3 check_lines.py <pdf>
\documentclass[10pt,letterpaper]{article}

\usepackage[margin=0.55in,top=0.5in,bottom=0.45in]{geometry}
\usepackage[T1]{fontenc}
% Fonts: \resumefont loads the body font; \headfont styles the name and section titles.
% Override either before % Shared preamble and content for Sudhansh Peddabomma's résumés.
% Each profile (SudhanshPeddabomma-AI.tex, SudhanshPeddabomma-Vision.tex) inputs this file and picks
% which blocks to show and in what order, so every bullet is written once.
% Layout rule: every bullet fills one line, or runs 1.5 to 2 lines. Check with: python3 check_lines.py <pdf>
\documentclass[10pt,letterpaper]{article}

\usepackage[margin=0.55in,top=0.5in,bottom=0.45in]{geometry}
\usepackage[T1]{fontenc}
% Fonts: \resumefont loads the body font; \headfont styles the name and section titles.
% Override either before \input{resume-common} to try another pairing.
% Default: Charter body with Source Sans headings. All-sans alternative:
%   \def\resumefont{\usepackage[default,semibold]{sourcesanspro}}\def\headfont{}
\providecommand{\resumefont}{\usepackage[scaled=0.96]{XCharter}\usepackage[semibold,scale=0.96]{sourcesanspro}}
\providecommand{\headfont}{\sffamily}
\providecommand{\namefont}{\rmfamily\scshape}
\resumefont
\usepackage[final]{microtype}
\usepackage[dvipsnames]{xcolor}
\usepackage{enumitem}
\usepackage{titlesec}
\usepackage{tabularx}
\usepackage{fontawesome5}
\usepackage{etoolbox}
\usepackage[hidelinks]{hyperref}

% The site's green, darkened for print.
\definecolor{accent}{HTML}{3A7A45}
\definecolor{muted}{HTML}{5C675D}

\pagestyle{empty}
\setlength{\parindent}{0pt}
\setlength{\tabcolsep}{0pt}
\providecommand{\leading}{1.1}
\linespread{\leading}

\hypersetup{pdfauthor={Sudhansh Peddabomma}}

% Type scale: name 26pt small caps, section titles 12.5pt caps, entry titles 10.8pt, body 10pt,
% side notes 9pt.
\newcommand{\side}[1]{{\fontsize{9}{11}\selectfont\color{muted}#1}}
\newcommand{\sectitle}[1]{\textls[60]{\MakeUppercase{#1}}}
\titleformat{\section}{\headfont\color{accent}\bfseries\fontsize{12.5}{14}\selectfont}{}{0pt}{\sectitle}[\vspace{1.5pt}{\color{accent}\titlerule[0.8pt]}\global\firstentrytrue]
% The first entry after a section title skips \entrygap, so only \rulegap sits under the rule.
\newif\iffirstentry

% Vertical rhythm, the same idea as the earlier hand-tuned résumés: bullets inside an entry sit
% tight together, entries are separated by a clear gap, and sections by a larger one. A profile
% may override these before \input{resume-common} to fit one page.
\providecommand{\secgap}{12pt}   % above a section title
\providecommand{\rulegap}{5.5pt}   % section rule to first entry
\providecommand{\entrygap}{7pt}  % between entries
\providecommand{\headgap}{1.5pt} % entry heading to its bullets
\providecommand{\itemgap}{0.6pt} % between bullets of one entry
\titlespacing*{\section}{0pt}{\secgap}{\rulegap}

\newcommand{\squarebullet}{\raisebox{0.32ex}{\color{accent}\rule{0.5ex}{0.5ex}}}
\setlist[itemize]{leftmargin=12pt, labelsep=6pt, topsep=\headgap, itemsep=\itemgap, parsep=0pt, partopsep=0pt,
  label=\squarebullet}

% \entry{Organisation or title}{Role or detail (may be empty)}{Dates}; \sub{Role or detail}{Right-hand note}
\newcommand{\entry}[3]{%
  \par\iffirstentry\global\firstentryfalse\else\addvspace{\entrygap}\fi\noindent{\headfont\fontsize{10.8}{13}\selectfont\textbf{#1}}\ifblank{#2}{}{\ {\color{muted}|}\ #2}\hfill\hspace{1em}\side{#3}\par}
\newcommand{\sub}[2]{%
  \noindent\textit{#1}\hfill{\color{muted}\textit{#2}}\par}
\newcommand{\report}[1]{\href{#1}{\color{accent}Report}}
\newcommand{\me}{\textbf{S. Peddabomma}}

% ---------------------------------------------------------------------------------------------------
% Header
\newcommand{\Header}{%
  \begin{center}
    {\namefont\fontsize{26}{28}\selectfont Sudhansh Peddabomma}\\[6pt]
    {\headfont\fontsize{9.5}{11}\selectfont\color{muted}
      \href{mailto:sudhansh6@gmail.com}{\faEnvelope[regular]\ sudhansh6@gmail.com} \quad
      \href{https://sudhansh6.github.io}{\faGlobe\ sudhansh6.github.io} \quad
      \href{https://www.linkedin.com/in/sudhansh-peddabomma-73a208192/}{\faLinkedin\ sudhansh-peddabomma} \quad
      \href{https://github.com/Sudhansh6}{\faGithub\ Sudhansh6}}
  \end{center}\vspace{-4pt}}

% ---------------------------------------------------------------------------------------------------
% Glean bullets, shared by both profiles
\newcommand{\GleanSlides}{\item Led \textbf{AI slide generation} in Glean Assistant, now creating \textbf{180K decks a week} for \textbf{50K weekly active users} at enterprises}
\newcommand{\GleanSlidesPerf}{\item Improved the slides harness, raising reliability to \textbf{97\%}, cutting median latency \textbf{26\%} and cost per deck by \textbf{half}
  in four months}
\newcommand{\GleanSlidesOrch}{\item Architected a \textbf{slide orchestration system} that reuses relevant slides from a company's existing decks and generates
  new ones in its brand and template, so teams can refine decks with each other and Glean in full context, without leaving the product}
\newcommand{\GleanAgents}{\item Moved Glean's \textbf{workflow agents} from one-shot planning to an \textbf{agentic tool-use loop}, greatly improving agentic
  capabilities and instruction following, with \textbf{18\%} higher answer completeness and \textbf{17\%} fewer tokens}
\newcommand{\GleanTeamChat}{\item Engineered the core of \textbf{Team chat}, letting teammates work with Glean in one shared conversation and co-edit its artifacts}
\newcommand{\GleanVideo}{\item Shipped \textbf{video artifacts} end to end, generating narrated videos in chat with streamed playback and admin rollout controls}
\newcommand{\GleanEvals}{\item Created the \textbf{evaluations and metrics} that guide model choices and rollouts across agents, slides and interactive artifacts}

% ---------------------------------------------------------------------------------------------------
% Experience
% \groupbreak separates related groups of bullets inside one long entry (slides, then platform work).
\providecommand{\groupgap}{3.5pt}
\newcommand{\groupbreak}{\vspace{\groupgap}}
\newcommand{\GleanAI}{%
  \entry{Glean}{Software Engineer}{Apr 2025 -- Present}
  \begin{itemize}
    \GleanSlides
    \GleanSlidesOrch
    \GleanSlidesPerf \groupbreak
    \GleanTeamChat
    \GleanVideo
    \GleanAgents
    \GleanEvals
  \end{itemize}}

\newcommand{\GleanVision}{%
  \entry{Glean}{Software Engineer}{Apr 2025 -- Present}
  \begin{itemize}
    \GleanSlides
    \GleanSlidesOrch
    \GleanSlidesPerf \groupbreak
    \GleanVideo
    \GleanEvals
  \end{itemize}}

\newcommand{\Duality}{%
  \entry{Duality AI}{Computer Vision Intern}{Jun 2024 -- Sep 2024}
  \begin{itemize}
    \item Designed automated \textbf{3D reconstruction} techniques for featureless objects, reducing digital-twin generation time
          by \textbf{40\%}
    \item Collaborated with \textbf{Autodesk} to validate Unreal Engine simulations for robotics tasks; leveraging structured
          domain randomization to reduce Sim2Real gap and increase mAP-50 by \textbf{15\%} for object detection and segmentation
  \end{itemize}}

\newcommand{\Microsoft}{%
  \entry{Microsoft}{Data and Applied Scientist Intern}{May 2022 -- Jul 2022}
  \begin{itemize}
    \item Developed a \textbf{decision-tree ranker} to recommend emails without user queries, improving \textbf{Outlook} search capabilities
    \item Proposed hierarchical feature-sets for the ranker, reducing latency for recommendations and improving recall
  \end{itemize}}

% ---------------------------------------------------------------------------------------------------
% Research
\newcommand{\HomeRobots}{%
  \entry{Real-time 3D Perception for Home Robots}{UC San Diego}{Sep 2023 -- Sep 2024}
  \begin{itemize}
    \item Investigated real-time dense \textbf{visual SLAM} methods using NeRFs and Gaussian Splatting for indoor robot navigation
    \item Integrated object segmentation, grasp-pose estimation, and 3D mapping on the Fetch robot via ROS, demonstrating
          a novel \textbf{tabletop object rearrangement} algorithm that reduced cost by \textbf{20\%} compared to the state-of-the-art
  \end{itemize}}

\newcommand{\Tomography}{%
  \entry{3D Tomography with Primal-Dual Neural Networks}{University College London}{May 2021 -- Jul 2023}
  \begin{itemize}
    \item Developed a stochastic neural-network architecture of primal-dual algorithm for online reconstruction of 3D
          volumes from tomographic projections, with \textbf{99.6\,\% structural similarity} in low-dosage conditions
    \item Built a Python library with custom gradient operators for reconstructing volumes in a single pass, reducing
          compute requirements by up to \textbf{5x} over SOTA learning-based approaches for cone vector tomography
  \end{itemize}}

% ---------------------------------------------------------------------------------------------------
% Projects
\newcommand{\OOXML}{%
  \entry{ooxml}{Open-source contributor \textbar\ \href{https://github.com/ChristopherVR/ooxml}{\color{accent}GitHub}}{Sep 2026 -- Present}
  \begin{itemize}
    \item Contributed to an open-source \textbf{TypeScript} engine that renders and edits PowerPoint and Word files in the browser,
          improving editing fidelity, \textbf{real-time collaboration} and load performance across React, Vue and Svelte viewers
  \end{itemize}}

\newcommand{\MirrorAI}{%
  \entry{Mirror AI: Deployable Personas}{Honorable mention in Supabase YC Hackathon}{Oct 2024 -- Dec 2024}
  \begin{itemize}
    \item Designed an \textbf{agentic LLM architecture} with LangGraph to mirror user personalities, creating interactive digital
          personas, and deployed it as a full-stack platform on Supabase and Vercel with one-click publishing
  \end{itemize}}

\newcommand{\LLMReasoning}{%
  \entry{Improving LLM Reasoning for Numerical Problems}{\report{https://sudhansh6.github.io/assets/pdf/MathPrompter++.pdf}}{Sep 2024 -- Dec 2024}
  \begin{itemize}
    \item Enhanced MathPrompter (ACL 2023) with CoT, achieving \textbf{10\% higher accuracy} on Llama 3.1 1B where prior methods failed,
          and reduced hallucination rates significantly by integrating multi-step validation
  \end{itemize}}

\newcommand{\InverseRendering}{%
  \entry{Inverse Rendering with 2D Gaussian Splatting}{\report{https://sudhansh6.github.io/assets/pdf/CSE252D_Report.pdf}}{Mar 2024 -- May 2024}
  \begin{itemize}
    \item Developed a novel \textbf{inverse rendering} framework in CUDA to recover PBR properties of a scene using 2D Gaussian
          Splatting, improving normal map MAE by \textbf{15\%} over the current SOTA methods
  \end{itemize}}

% ---------------------------------------------------------------------------------------------------
% Publications: title on the left, venue on the right, one line each.
\newcommand{\publine}[3]{\par\noindent\href{#1}{#2}\hfill\hspace{1em}\side{#3}\par\addvspace{1.5pt}}
\newcommand{\PubSaturation}{\publine{https://doi.org/10.1016/j.sigpro.2023.109349}{A likelihood based method for compressive signal recovery under Gaussian and saturation noise}{Signal Processing, 2024}}
\newcommand{\PubPermutation}{\publine{https://ieeexplore.ieee.org/document/10888147}{Identification and correction of permutation errors in compressed sensing-based group testing}{ICASSP 2025}}
\newcommand{\PubRearrangement}{\publine{https://arxiv.org/abs/2411.10899}{Planning for tabletop object rearrangement}{ICRA 2025}}

% ---------------------------------------------------------------------------------------------------
% Education, skills, service
% Each profile sets \UCSDCourses, \IITBCourses and \SkillRows before \input{resume-common}.
\providecommand{\UCSDCourses}{}
\providecommand{\IITBCourses}{}
\providecommand{\SkillRows}{}
\newcommand{\coursework}[1]{\ifblank{#1}{}{\par\noindent{\fontsize{9.5}{11.5}\selectfont{\color{muted}Coursework:} #1}\par}}
\newcommand{\Education}{%
  \section{Education}
  \entry{University of California San Diego}{M.S. in Computer Science and Engineering (AI)}{GPA 4.00/4.00 \quad Sep 2023 -- Mar 2025}
  \coursework{\UCSDCourses}
  \entry{Indian Institute of Technology Bombay}{B.Tech.\ with Honors in Computer Science}{CPI 9.66/10 \quad Jul 2019 -- Jul 2023}
  \coursework{\IITBCourses}}

% \skill{Label}{Items} is one row of the skills table.
\newcommand{\skill}[2]{{\headfont\bfseries #1} & #2 \\[1.5pt]}
\newcommand{\cpp}{C\kern-0.04em\raisebox{0.06ex}{\small++}}
\newcommand{\Skills}{%
  \section{Skills}
  \begin{tabularx}{\linewidth}{@{}l@{\hspace{12pt}}>{\raggedright\arraybackslash}X@{}}
    \SkillRows
  \end{tabularx}}
 to try another pairing.
% Default: Charter body with Source Sans headings. All-sans alternative:
%   \def\resumefont{\usepackage[default,semibold]{sourcesanspro}}\def\headfont{}
\providecommand{\resumefont}{\usepackage[scaled=0.96]{XCharter}\usepackage[semibold,scale=0.96]{sourcesanspro}}
\providecommand{\headfont}{\sffamily}
\providecommand{\namefont}{\rmfamily\scshape}
\resumefont
\usepackage[final]{microtype}
\usepackage[dvipsnames]{xcolor}
\usepackage{enumitem}
\usepackage{titlesec}
\usepackage{tabularx}
\usepackage{fontawesome5}
\usepackage{etoolbox}
\usepackage[hidelinks]{hyperref}

% The site's green, darkened for print.
\definecolor{accent}{HTML}{3A7A45}
\definecolor{muted}{HTML}{5C675D}

\pagestyle{empty}
\setlength{\parindent}{0pt}
\setlength{\tabcolsep}{0pt}
\providecommand{\leading}{1.1}
\linespread{\leading}

\hypersetup{pdfauthor={Sudhansh Peddabomma}}

% Type scale: name 26pt small caps, section titles 12.5pt caps, entry titles 10.8pt, body 10pt,
% side notes 9pt.
\newcommand{\side}[1]{{\fontsize{9}{11}\selectfont\color{muted}#1}}
\newcommand{\sectitle}[1]{\textls[60]{\MakeUppercase{#1}}}
\titleformat{\section}{\headfont\color{accent}\bfseries\fontsize{12.5}{14}\selectfont}{}{0pt}{\sectitle}[\vspace{1.5pt}{\color{accent}\titlerule[0.8pt]}\global\firstentrytrue]
% The first entry after a section title skips \entrygap, so only \rulegap sits under the rule.
\newif\iffirstentry

% Vertical rhythm, the same idea as the earlier hand-tuned résumés: bullets inside an entry sit
% tight together, entries are separated by a clear gap, and sections by a larger one. A profile
% may override these before % Shared preamble and content for Sudhansh Peddabomma's résumés.
% Each profile (SudhanshPeddabomma-AI.tex, SudhanshPeddabomma-Vision.tex) inputs this file and picks
% which blocks to show and in what order, so every bullet is written once.
% Layout rule: every bullet fills one line, or runs 1.5 to 2 lines. Check with: python3 check_lines.py <pdf>
\documentclass[10pt,letterpaper]{article}

\usepackage[margin=0.55in,top=0.5in,bottom=0.45in]{geometry}
\usepackage[T1]{fontenc}
% Fonts: \resumefont loads the body font; \headfont styles the name and section titles.
% Override either before \input{resume-common} to try another pairing.
% Default: Charter body with Source Sans headings. All-sans alternative:
%   \def\resumefont{\usepackage[default,semibold]{sourcesanspro}}\def\headfont{}
\providecommand{\resumefont}{\usepackage[scaled=0.96]{XCharter}\usepackage[semibold,scale=0.96]{sourcesanspro}}
\providecommand{\headfont}{\sffamily}
\providecommand{\namefont}{\rmfamily\scshape}
\resumefont
\usepackage[final]{microtype}
\usepackage[dvipsnames]{xcolor}
\usepackage{enumitem}
\usepackage{titlesec}
\usepackage{tabularx}
\usepackage{fontawesome5}
\usepackage{etoolbox}
\usepackage[hidelinks]{hyperref}

% The site's green, darkened for print.
\definecolor{accent}{HTML}{3A7A45}
\definecolor{muted}{HTML}{5C675D}

\pagestyle{empty}
\setlength{\parindent}{0pt}
\setlength{\tabcolsep}{0pt}
\providecommand{\leading}{1.1}
\linespread{\leading}

\hypersetup{pdfauthor={Sudhansh Peddabomma}}

% Type scale: name 26pt small caps, section titles 12.5pt caps, entry titles 10.8pt, body 10pt,
% side notes 9pt.
\newcommand{\side}[1]{{\fontsize{9}{11}\selectfont\color{muted}#1}}
\newcommand{\sectitle}[1]{\textls[60]{\MakeUppercase{#1}}}
\titleformat{\section}{\headfont\color{accent}\bfseries\fontsize{12.5}{14}\selectfont}{}{0pt}{\sectitle}[\vspace{1.5pt}{\color{accent}\titlerule[0.8pt]}\global\firstentrytrue]
% The first entry after a section title skips \entrygap, so only \rulegap sits under the rule.
\newif\iffirstentry

% Vertical rhythm, the same idea as the earlier hand-tuned résumés: bullets inside an entry sit
% tight together, entries are separated by a clear gap, and sections by a larger one. A profile
% may override these before \input{resume-common} to fit one page.
\providecommand{\secgap}{12pt}   % above a section title
\providecommand{\rulegap}{5.5pt}   % section rule to first entry
\providecommand{\entrygap}{7pt}  % between entries
\providecommand{\headgap}{1.5pt} % entry heading to its bullets
\providecommand{\itemgap}{0.6pt} % between bullets of one entry
\titlespacing*{\section}{0pt}{\secgap}{\rulegap}

\newcommand{\squarebullet}{\raisebox{0.32ex}{\color{accent}\rule{0.5ex}{0.5ex}}}
\setlist[itemize]{leftmargin=12pt, labelsep=6pt, topsep=\headgap, itemsep=\itemgap, parsep=0pt, partopsep=0pt,
  label=\squarebullet}

% \entry{Organisation or title}{Role or detail (may be empty)}{Dates}; \sub{Role or detail}{Right-hand note}
\newcommand{\entry}[3]{%
  \par\iffirstentry\global\firstentryfalse\else\addvspace{\entrygap}\fi\noindent{\headfont\fontsize{10.8}{13}\selectfont\textbf{#1}}\ifblank{#2}{}{\ {\color{muted}|}\ #2}\hfill\hspace{1em}\side{#3}\par}
\newcommand{\sub}[2]{%
  \noindent\textit{#1}\hfill{\color{muted}\textit{#2}}\par}
\newcommand{\report}[1]{\href{#1}{\color{accent}Report}}
\newcommand{\me}{\textbf{S. Peddabomma}}

% ---------------------------------------------------------------------------------------------------
% Header
\newcommand{\Header}{%
  \begin{center}
    {\namefont\fontsize{26}{28}\selectfont Sudhansh Peddabomma}\\[6pt]
    {\headfont\fontsize{9.5}{11}\selectfont\color{muted}
      \href{mailto:sudhansh6@gmail.com}{\faEnvelope[regular]\ sudhansh6@gmail.com} \quad
      \href{https://sudhansh6.github.io}{\faGlobe\ sudhansh6.github.io} \quad
      \href{https://www.linkedin.com/in/sudhansh-peddabomma-73a208192/}{\faLinkedin\ sudhansh-peddabomma} \quad
      \href{https://github.com/Sudhansh6}{\faGithub\ Sudhansh6}}
  \end{center}\vspace{-4pt}}

% ---------------------------------------------------------------------------------------------------
% Glean bullets, shared by both profiles
\newcommand{\GleanSlides}{\item Led \textbf{AI slide generation} in Glean Assistant, now creating \textbf{180K decks a week} for \textbf{50K weekly active users} at enterprises}
\newcommand{\GleanSlidesPerf}{\item Improved the slides harness, raising reliability to \textbf{97\%}, cutting median latency \textbf{26\%} and cost per deck by \textbf{half}
  in four months}
\newcommand{\GleanSlidesOrch}{\item Architected a \textbf{slide orchestration system} that reuses relevant slides from a company's existing decks and generates
  new ones in its brand and template, so teams can refine decks with each other and Glean in full context, without leaving the product}
\newcommand{\GleanAgents}{\item Moved Glean's \textbf{workflow agents} from one-shot planning to an \textbf{agentic tool-use loop}, greatly improving agentic
  capabilities and instruction following, with \textbf{18\%} higher answer completeness and \textbf{17\%} fewer tokens}
\newcommand{\GleanTeamChat}{\item Engineered the core of \textbf{Team chat}, letting teammates work with Glean in one shared conversation and co-edit its artifacts}
\newcommand{\GleanVideo}{\item Shipped \textbf{video artifacts} end to end, generating narrated videos in chat with streamed playback and admin rollout controls}
\newcommand{\GleanEvals}{\item Created the \textbf{evaluations and metrics} that guide model choices and rollouts across agents, slides and interactive artifacts}

% ---------------------------------------------------------------------------------------------------
% Experience
% \groupbreak separates related groups of bullets inside one long entry (slides, then platform work).
\providecommand{\groupgap}{3.5pt}
\newcommand{\groupbreak}{\vspace{\groupgap}}
\newcommand{\GleanAI}{%
  \entry{Glean}{Software Engineer}{Apr 2025 -- Present}
  \begin{itemize}
    \GleanSlides
    \GleanSlidesOrch
    \GleanSlidesPerf \groupbreak
    \GleanTeamChat
    \GleanVideo
    \GleanAgents
    \GleanEvals
  \end{itemize}}

\newcommand{\GleanVision}{%
  \entry{Glean}{Software Engineer}{Apr 2025 -- Present}
  \begin{itemize}
    \GleanSlides
    \GleanSlidesOrch
    \GleanSlidesPerf \groupbreak
    \GleanVideo
    \GleanEvals
  \end{itemize}}

\newcommand{\Duality}{%
  \entry{Duality AI}{Computer Vision Intern}{Jun 2024 -- Sep 2024}
  \begin{itemize}
    \item Designed automated \textbf{3D reconstruction} techniques for featureless objects, reducing digital-twin generation time
          by \textbf{40\%}
    \item Collaborated with \textbf{Autodesk} to validate Unreal Engine simulations for robotics tasks; leveraging structured
          domain randomization to reduce Sim2Real gap and increase mAP-50 by \textbf{15\%} for object detection and segmentation
  \end{itemize}}

\newcommand{\Microsoft}{%
  \entry{Microsoft}{Data and Applied Scientist Intern}{May 2022 -- Jul 2022}
  \begin{itemize}
    \item Developed a \textbf{decision-tree ranker} to recommend emails without user queries, improving \textbf{Outlook} search capabilities
    \item Proposed hierarchical feature-sets for the ranker, reducing latency for recommendations and improving recall
  \end{itemize}}

% ---------------------------------------------------------------------------------------------------
% Research
\newcommand{\HomeRobots}{%
  \entry{Real-time 3D Perception for Home Robots}{UC San Diego}{Sep 2023 -- Sep 2024}
  \begin{itemize}
    \item Investigated real-time dense \textbf{visual SLAM} methods using NeRFs and Gaussian Splatting for indoor robot navigation
    \item Integrated object segmentation, grasp-pose estimation, and 3D mapping on the Fetch robot via ROS, demonstrating
          a novel \textbf{tabletop object rearrangement} algorithm that reduced cost by \textbf{20\%} compared to the state-of-the-art
  \end{itemize}}

\newcommand{\Tomography}{%
  \entry{3D Tomography with Primal-Dual Neural Networks}{University College London}{May 2021 -- Jul 2023}
  \begin{itemize}
    \item Developed a stochastic neural-network architecture of primal-dual algorithm for online reconstruction of 3D
          volumes from tomographic projections, with \textbf{99.6\,\% structural similarity} in low-dosage conditions
    \item Built a Python library with custom gradient operators for reconstructing volumes in a single pass, reducing
          compute requirements by up to \textbf{5x} over SOTA learning-based approaches for cone vector tomography
  \end{itemize}}

% ---------------------------------------------------------------------------------------------------
% Projects
\newcommand{\OOXML}{%
  \entry{ooxml}{Open-source contributor \textbar\ \href{https://github.com/ChristopherVR/ooxml}{\color{accent}GitHub}}{Sep 2026 -- Present}
  \begin{itemize}
    \item Contributed to an open-source \textbf{TypeScript} engine that renders and edits PowerPoint and Word files in the browser,
          improving editing fidelity, \textbf{real-time collaboration} and load performance across React, Vue and Svelte viewers
  \end{itemize}}

\newcommand{\MirrorAI}{%
  \entry{Mirror AI: Deployable Personas}{Honorable mention in Supabase YC Hackathon}{Oct 2024 -- Dec 2024}
  \begin{itemize}
    \item Designed an \textbf{agentic LLM architecture} with LangGraph to mirror user personalities, creating interactive digital
          personas, and deployed it as a full-stack platform on Supabase and Vercel with one-click publishing
  \end{itemize}}

\newcommand{\LLMReasoning}{%
  \entry{Improving LLM Reasoning for Numerical Problems}{\report{https://sudhansh6.github.io/assets/pdf/MathPrompter++.pdf}}{Sep 2024 -- Dec 2024}
  \begin{itemize}
    \item Enhanced MathPrompter (ACL 2023) with CoT, achieving \textbf{10\% higher accuracy} on Llama 3.1 1B where prior methods failed,
          and reduced hallucination rates significantly by integrating multi-step validation
  \end{itemize}}

\newcommand{\InverseRendering}{%
  \entry{Inverse Rendering with 2D Gaussian Splatting}{\report{https://sudhansh6.github.io/assets/pdf/CSE252D_Report.pdf}}{Mar 2024 -- May 2024}
  \begin{itemize}
    \item Developed a novel \textbf{inverse rendering} framework in CUDA to recover PBR properties of a scene using 2D Gaussian
          Splatting, improving normal map MAE by \textbf{15\%} over the current SOTA methods
  \end{itemize}}

% ---------------------------------------------------------------------------------------------------
% Publications: title on the left, venue on the right, one line each.
\newcommand{\publine}[3]{\par\noindent\href{#1}{#2}\hfill\hspace{1em}\side{#3}\par\addvspace{1.5pt}}
\newcommand{\PubSaturation}{\publine{https://doi.org/10.1016/j.sigpro.2023.109349}{A likelihood based method for compressive signal recovery under Gaussian and saturation noise}{Signal Processing, 2024}}
\newcommand{\PubPermutation}{\publine{https://ieeexplore.ieee.org/document/10888147}{Identification and correction of permutation errors in compressed sensing-based group testing}{ICASSP 2025}}
\newcommand{\PubRearrangement}{\publine{https://arxiv.org/abs/2411.10899}{Planning for tabletop object rearrangement}{ICRA 2025}}

% ---------------------------------------------------------------------------------------------------
% Education, skills, service
% Each profile sets \UCSDCourses, \IITBCourses and \SkillRows before \input{resume-common}.
\providecommand{\UCSDCourses}{}
\providecommand{\IITBCourses}{}
\providecommand{\SkillRows}{}
\newcommand{\coursework}[1]{\ifblank{#1}{}{\par\noindent{\fontsize{9.5}{11.5}\selectfont{\color{muted}Coursework:} #1}\par}}
\newcommand{\Education}{%
  \section{Education}
  \entry{University of California San Diego}{M.S. in Computer Science and Engineering (AI)}{GPA 4.00/4.00 \quad Sep 2023 -- Mar 2025}
  \coursework{\UCSDCourses}
  \entry{Indian Institute of Technology Bombay}{B.Tech.\ with Honors in Computer Science}{CPI 9.66/10 \quad Jul 2019 -- Jul 2023}
  \coursework{\IITBCourses}}

% \skill{Label}{Items} is one row of the skills table.
\newcommand{\skill}[2]{{\headfont\bfseries #1} & #2 \\[1.5pt]}
\newcommand{\cpp}{C\kern-0.04em\raisebox{0.06ex}{\small++}}
\newcommand{\Skills}{%
  \section{Skills}
  \begin{tabularx}{\linewidth}{@{}l@{\hspace{12pt}}>{\raggedright\arraybackslash}X@{}}
    \SkillRows
  \end{tabularx}}
 to fit one page.
\providecommand{\secgap}{12pt}   % above a section title
\providecommand{\rulegap}{5.5pt}   % section rule to first entry
\providecommand{\entrygap}{7pt}  % between entries
\providecommand{\headgap}{1.5pt} % entry heading to its bullets
\providecommand{\itemgap}{0.6pt} % between bullets of one entry
\titlespacing*{\section}{0pt}{\secgap}{\rulegap}

\newcommand{\squarebullet}{\raisebox{0.32ex}{\color{accent}\rule{0.5ex}{0.5ex}}}
\setlist[itemize]{leftmargin=12pt, labelsep=6pt, topsep=\headgap, itemsep=\itemgap, parsep=0pt, partopsep=0pt,
  label=\squarebullet}

% \entry{Organisation or title}{Role or detail (may be empty)}{Dates}; \sub{Role or detail}{Right-hand note}
\newcommand{\entry}[3]{%
  \par\iffirstentry\global\firstentryfalse\else\addvspace{\entrygap}\fi\noindent{\headfont\fontsize{10.8}{13}\selectfont\textbf{#1}}\ifblank{#2}{}{\ {\color{muted}|}\ #2}\hfill\hspace{1em}\side{#3}\par}
\newcommand{\sub}[2]{%
  \noindent\textit{#1}\hfill{\color{muted}\textit{#2}}\par}
\newcommand{\report}[1]{\href{#1}{\color{accent}Report}}
\newcommand{\me}{\textbf{S. Peddabomma}}

% ---------------------------------------------------------------------------------------------------
% Header
\newcommand{\Header}{%
  \begin{center}
    {\namefont\fontsize{26}{28}\selectfont Sudhansh Peddabomma}\\[6pt]
    {\headfont\fontsize{9.5}{11}\selectfont\color{muted}
      \href{mailto:sudhansh6@gmail.com}{\faEnvelope[regular]\ sudhansh6@gmail.com} \quad
      \href{https://sudhansh6.github.io}{\faGlobe\ sudhansh6.github.io} \quad
      \href{https://www.linkedin.com/in/sudhansh-peddabomma-73a208192/}{\faLinkedin\ sudhansh-peddabomma} \quad
      \href{https://github.com/Sudhansh6}{\faGithub\ Sudhansh6}}
  \end{center}\vspace{-4pt}}

% ---------------------------------------------------------------------------------------------------
% Glean bullets, shared by both profiles
\newcommand{\GleanSlides}{\item Led \textbf{AI slide generation} in Glean Assistant, now creating \textbf{180K decks a week} for \textbf{50K weekly active users} at enterprises}
\newcommand{\GleanSlidesPerf}{\item Improved the slides harness, raising reliability to \textbf{97\%}, cutting median latency \textbf{26\%} and cost per deck by \textbf{half}
  in four months}
\newcommand{\GleanSlidesOrch}{\item Architected a \textbf{slide orchestration system} that reuses relevant slides from a company's existing decks and generates
  new ones in its brand and template, so teams can refine decks with each other and Glean in full context, without leaving the product}
\newcommand{\GleanAgents}{\item Moved Glean's \textbf{workflow agents} from one-shot planning to an \textbf{agentic tool-use loop}, greatly improving agentic
  capabilities and instruction following, with \textbf{18\%} higher answer completeness and \textbf{17\%} fewer tokens}
\newcommand{\GleanTeamChat}{\item Engineered the core of \textbf{Team chat}, letting teammates work with Glean in one shared conversation and co-edit its artifacts}
\newcommand{\GleanVideo}{\item Shipped \textbf{video artifacts} end to end, generating narrated videos in chat with streamed playback and admin rollout controls}
\newcommand{\GleanEvals}{\item Created the \textbf{evaluations and metrics} that guide model choices and rollouts across agents, slides and interactive artifacts}

% ---------------------------------------------------------------------------------------------------
% Experience
% \groupbreak separates related groups of bullets inside one long entry (slides, then platform work).
\providecommand{\groupgap}{3.5pt}
\newcommand{\groupbreak}{\vspace{\groupgap}}
\newcommand{\GleanAI}{%
  \entry{Glean}{Software Engineer}{Apr 2025 -- Present}
  \begin{itemize}
    \GleanSlides
    \GleanSlidesOrch
    \GleanSlidesPerf \groupbreak
    \GleanTeamChat
    \GleanVideo
    \GleanAgents
    \GleanEvals
  \end{itemize}}

\newcommand{\GleanVision}{%
  \entry{Glean}{Software Engineer}{Apr 2025 -- Present}
  \begin{itemize}
    \GleanSlides
    \GleanSlidesOrch
    \GleanSlidesPerf \groupbreak
    \GleanVideo
    \GleanEvals
  \end{itemize}}

\newcommand{\Duality}{%
  \entry{Duality AI}{Computer Vision Intern}{Jun 2024 -- Sep 2024}
  \begin{itemize}
    \item Designed automated \textbf{3D reconstruction} techniques for featureless objects, reducing digital-twin generation time
          by \textbf{40\%}
    \item Collaborated with \textbf{Autodesk} to validate Unreal Engine simulations for robotics tasks; leveraging structured
          domain randomization to reduce Sim2Real gap and increase mAP-50 by \textbf{15\%} for object detection and segmentation
  \end{itemize}}

\newcommand{\Microsoft}{%
  \entry{Microsoft}{Data and Applied Scientist Intern}{May 2022 -- Jul 2022}
  \begin{itemize}
    \item Developed a \textbf{decision-tree ranker} to recommend emails without user queries, improving \textbf{Outlook} search capabilities
    \item Proposed hierarchical feature-sets for the ranker, reducing latency for recommendations and improving recall
  \end{itemize}}

% ---------------------------------------------------------------------------------------------------
% Research
\newcommand{\HomeRobots}{%
  \entry{Real-time 3D Perception for Home Robots}{UC San Diego}{Sep 2023 -- Sep 2024}
  \begin{itemize}
    \item Investigated real-time dense \textbf{visual SLAM} methods using NeRFs and Gaussian Splatting for indoor robot navigation
    \item Integrated object segmentation, grasp-pose estimation, and 3D mapping on the Fetch robot via ROS, demonstrating
          a novel \textbf{tabletop object rearrangement} algorithm that reduced cost by \textbf{20\%} compared to the state-of-the-art
  \end{itemize}}

\newcommand{\Tomography}{%
  \entry{3D Tomography with Primal-Dual Neural Networks}{University College London}{May 2021 -- Jul 2023}
  \begin{itemize}
    \item Developed a stochastic neural-network architecture of primal-dual algorithm for online reconstruction of 3D
          volumes from tomographic projections, with \textbf{99.6\,\% structural similarity} in low-dosage conditions
    \item Built a Python library with custom gradient operators for reconstructing volumes in a single pass, reducing
          compute requirements by up to \textbf{5x} over SOTA learning-based approaches for cone vector tomography
  \end{itemize}}

% ---------------------------------------------------------------------------------------------------
% Projects
\newcommand{\OOXML}{%
  \entry{ooxml}{Open-source contributor \textbar\ \href{https://github.com/ChristopherVR/ooxml}{\color{accent}GitHub}}{Sep 2026 -- Present}
  \begin{itemize}
    \item Contributed to an open-source \textbf{TypeScript} engine that renders and edits PowerPoint and Word files in the browser,
          improving editing fidelity, \textbf{real-time collaboration} and load performance across React, Vue and Svelte viewers
  \end{itemize}}

\newcommand{\MirrorAI}{%
  \entry{Mirror AI: Deployable Personas}{Honorable mention in Supabase YC Hackathon}{Oct 2024 -- Dec 2024}
  \begin{itemize}
    \item Designed an \textbf{agentic LLM architecture} with LangGraph to mirror user personalities, creating interactive digital
          personas, and deployed it as a full-stack platform on Supabase and Vercel with one-click publishing
  \end{itemize}}

\newcommand{\LLMReasoning}{%
  \entry{Improving LLM Reasoning for Numerical Problems}{\report{https://sudhansh6.github.io/assets/pdf/MathPrompter++.pdf}}{Sep 2024 -- Dec 2024}
  \begin{itemize}
    \item Enhanced MathPrompter (ACL 2023) with CoT, achieving \textbf{10\% higher accuracy} on Llama 3.1 1B where prior methods failed,
          and reduced hallucination rates significantly by integrating multi-step validation
  \end{itemize}}

\newcommand{\InverseRendering}{%
  \entry{Inverse Rendering with 2D Gaussian Splatting}{\report{https://sudhansh6.github.io/assets/pdf/CSE252D_Report.pdf}}{Mar 2024 -- May 2024}
  \begin{itemize}
    \item Developed a novel \textbf{inverse rendering} framework in CUDA to recover PBR properties of a scene using 2D Gaussian
          Splatting, improving normal map MAE by \textbf{15\%} over the current SOTA methods
  \end{itemize}}

% ---------------------------------------------------------------------------------------------------
% Publications: title on the left, venue on the right, one line each.
\newcommand{\publine}[3]{\par\noindent\href{#1}{#2}\hfill\hspace{1em}\side{#3}\par\addvspace{1.5pt}}
\newcommand{\PubSaturation}{\publine{https://doi.org/10.1016/j.sigpro.2023.109349}{A likelihood based method for compressive signal recovery under Gaussian and saturation noise}{Signal Processing, 2024}}
\newcommand{\PubPermutation}{\publine{https://ieeexplore.ieee.org/document/10888147}{Identification and correction of permutation errors in compressed sensing-based group testing}{ICASSP 2025}}
\newcommand{\PubRearrangement}{\publine{https://arxiv.org/abs/2411.10899}{Planning for tabletop object rearrangement}{ICRA 2025}}

% ---------------------------------------------------------------------------------------------------
% Education, skills, service
% Each profile sets \UCSDCourses, \IITBCourses and \SkillRows before % Shared preamble and content for Sudhansh Peddabomma's résumés.
% Each profile (SudhanshPeddabomma-AI.tex, SudhanshPeddabomma-Vision.tex) inputs this file and picks
% which blocks to show and in what order, so every bullet is written once.
% Layout rule: every bullet fills one line, or runs 1.5 to 2 lines. Check with: python3 check_lines.py <pdf>
\documentclass[10pt,letterpaper]{article}

\usepackage[margin=0.55in,top=0.5in,bottom=0.45in]{geometry}
\usepackage[T1]{fontenc}
% Fonts: \resumefont loads the body font; \headfont styles the name and section titles.
% Override either before \input{resume-common} to try another pairing.
% Default: Charter body with Source Sans headings. All-sans alternative:
%   \def\resumefont{\usepackage[default,semibold]{sourcesanspro}}\def\headfont{}
\providecommand{\resumefont}{\usepackage[scaled=0.96]{XCharter}\usepackage[semibold,scale=0.96]{sourcesanspro}}
\providecommand{\headfont}{\sffamily}
\providecommand{\namefont}{\rmfamily\scshape}
\resumefont
\usepackage[final]{microtype}
\usepackage[dvipsnames]{xcolor}
\usepackage{enumitem}
\usepackage{titlesec}
\usepackage{tabularx}
\usepackage{fontawesome5}
\usepackage{etoolbox}
\usepackage[hidelinks]{hyperref}

% The site's green, darkened for print.
\definecolor{accent}{HTML}{3A7A45}
\definecolor{muted}{HTML}{5C675D}

\pagestyle{empty}
\setlength{\parindent}{0pt}
\setlength{\tabcolsep}{0pt}
\providecommand{\leading}{1.1}
\linespread{\leading}

\hypersetup{pdfauthor={Sudhansh Peddabomma}}

% Type scale: name 26pt small caps, section titles 12.5pt caps, entry titles 10.8pt, body 10pt,
% side notes 9pt.
\newcommand{\side}[1]{{\fontsize{9}{11}\selectfont\color{muted}#1}}
\newcommand{\sectitle}[1]{\textls[60]{\MakeUppercase{#1}}}
\titleformat{\section}{\headfont\color{accent}\bfseries\fontsize{12.5}{14}\selectfont}{}{0pt}{\sectitle}[\vspace{1.5pt}{\color{accent}\titlerule[0.8pt]}\global\firstentrytrue]
% The first entry after a section title skips \entrygap, so only \rulegap sits under the rule.
\newif\iffirstentry

% Vertical rhythm, the same idea as the earlier hand-tuned résumés: bullets inside an entry sit
% tight together, entries are separated by a clear gap, and sections by a larger one. A profile
% may override these before \input{resume-common} to fit one page.
\providecommand{\secgap}{12pt}   % above a section title
\providecommand{\rulegap}{5.5pt}   % section rule to first entry
\providecommand{\entrygap}{7pt}  % between entries
\providecommand{\headgap}{1.5pt} % entry heading to its bullets
\providecommand{\itemgap}{0.6pt} % between bullets of one entry
\titlespacing*{\section}{0pt}{\secgap}{\rulegap}

\newcommand{\squarebullet}{\raisebox{0.32ex}{\color{accent}\rule{0.5ex}{0.5ex}}}
\setlist[itemize]{leftmargin=12pt, labelsep=6pt, topsep=\headgap, itemsep=\itemgap, parsep=0pt, partopsep=0pt,
  label=\squarebullet}

% \entry{Organisation or title}{Role or detail (may be empty)}{Dates}; \sub{Role or detail}{Right-hand note}
\newcommand{\entry}[3]{%
  \par\iffirstentry\global\firstentryfalse\else\addvspace{\entrygap}\fi\noindent{\headfont\fontsize{10.8}{13}\selectfont\textbf{#1}}\ifblank{#2}{}{\ {\color{muted}|}\ #2}\hfill\hspace{1em}\side{#3}\par}
\newcommand{\sub}[2]{%
  \noindent\textit{#1}\hfill{\color{muted}\textit{#2}}\par}
\newcommand{\report}[1]{\href{#1}{\color{accent}Report}}
\newcommand{\me}{\textbf{S. Peddabomma}}

% ---------------------------------------------------------------------------------------------------
% Header
\newcommand{\Header}{%
  \begin{center}
    {\namefont\fontsize{26}{28}\selectfont Sudhansh Peddabomma}\\[6pt]
    {\headfont\fontsize{9.5}{11}\selectfont\color{muted}
      \href{mailto:sudhansh6@gmail.com}{\faEnvelope[regular]\ sudhansh6@gmail.com} \quad
      \href{https://sudhansh6.github.io}{\faGlobe\ sudhansh6.github.io} \quad
      \href{https://www.linkedin.com/in/sudhansh-peddabomma-73a208192/}{\faLinkedin\ sudhansh-peddabomma} \quad
      \href{https://github.com/Sudhansh6}{\faGithub\ Sudhansh6}}
  \end{center}\vspace{-4pt}}

% ---------------------------------------------------------------------------------------------------
% Glean bullets, shared by both profiles
\newcommand{\GleanSlides}{\item Led \textbf{AI slide generation} in Glean Assistant, now creating \textbf{180K decks a week} for \textbf{50K weekly active users} at enterprises}
\newcommand{\GleanSlidesPerf}{\item Improved the slides harness, raising reliability to \textbf{97\%}, cutting median latency \textbf{26\%} and cost per deck by \textbf{half}
  in four months}
\newcommand{\GleanSlidesOrch}{\item Architected a \textbf{slide orchestration system} that reuses relevant slides from a company's existing decks and generates
  new ones in its brand and template, so teams can refine decks with each other and Glean in full context, without leaving the product}
\newcommand{\GleanAgents}{\item Moved Glean's \textbf{workflow agents} from one-shot planning to an \textbf{agentic tool-use loop}, greatly improving agentic
  capabilities and instruction following, with \textbf{18\%} higher answer completeness and \textbf{17\%} fewer tokens}
\newcommand{\GleanTeamChat}{\item Engineered the core of \textbf{Team chat}, letting teammates work with Glean in one shared conversation and co-edit its artifacts}
\newcommand{\GleanVideo}{\item Shipped \textbf{video artifacts} end to end, generating narrated videos in chat with streamed playback and admin rollout controls}
\newcommand{\GleanEvals}{\item Created the \textbf{evaluations and metrics} that guide model choices and rollouts across agents, slides and interactive artifacts}

% ---------------------------------------------------------------------------------------------------
% Experience
% \groupbreak separates related groups of bullets inside one long entry (slides, then platform work).
\providecommand{\groupgap}{3.5pt}
\newcommand{\groupbreak}{\vspace{\groupgap}}
\newcommand{\GleanAI}{%
  \entry{Glean}{Software Engineer}{Apr 2025 -- Present}
  \begin{itemize}
    \GleanSlides
    \GleanSlidesOrch
    \GleanSlidesPerf \groupbreak
    \GleanTeamChat
    \GleanVideo
    \GleanAgents
    \GleanEvals
  \end{itemize}}

\newcommand{\GleanVision}{%
  \entry{Glean}{Software Engineer}{Apr 2025 -- Present}
  \begin{itemize}
    \GleanSlides
    \GleanSlidesOrch
    \GleanSlidesPerf \groupbreak
    \GleanVideo
    \GleanEvals
  \end{itemize}}

\newcommand{\Duality}{%
  \entry{Duality AI}{Computer Vision Intern}{Jun 2024 -- Sep 2024}
  \begin{itemize}
    \item Designed automated \textbf{3D reconstruction} techniques for featureless objects, reducing digital-twin generation time
          by \textbf{40\%}
    \item Collaborated with \textbf{Autodesk} to validate Unreal Engine simulations for robotics tasks; leveraging structured
          domain randomization to reduce Sim2Real gap and increase mAP-50 by \textbf{15\%} for object detection and segmentation
  \end{itemize}}

\newcommand{\Microsoft}{%
  \entry{Microsoft}{Data and Applied Scientist Intern}{May 2022 -- Jul 2022}
  \begin{itemize}
    \item Developed a \textbf{decision-tree ranker} to recommend emails without user queries, improving \textbf{Outlook} search capabilities
    \item Proposed hierarchical feature-sets for the ranker, reducing latency for recommendations and improving recall
  \end{itemize}}

% ---------------------------------------------------------------------------------------------------
% Research
\newcommand{\HomeRobots}{%
  \entry{Real-time 3D Perception for Home Robots}{UC San Diego}{Sep 2023 -- Sep 2024}
  \begin{itemize}
    \item Investigated real-time dense \textbf{visual SLAM} methods using NeRFs and Gaussian Splatting for indoor robot navigation
    \item Integrated object segmentation, grasp-pose estimation, and 3D mapping on the Fetch robot via ROS, demonstrating
          a novel \textbf{tabletop object rearrangement} algorithm that reduced cost by \textbf{20\%} compared to the state-of-the-art
  \end{itemize}}

\newcommand{\Tomography}{%
  \entry{3D Tomography with Primal-Dual Neural Networks}{University College London}{May 2021 -- Jul 2023}
  \begin{itemize}
    \item Developed a stochastic neural-network architecture of primal-dual algorithm for online reconstruction of 3D
          volumes from tomographic projections, with \textbf{99.6\,\% structural similarity} in low-dosage conditions
    \item Built a Python library with custom gradient operators for reconstructing volumes in a single pass, reducing
          compute requirements by up to \textbf{5x} over SOTA learning-based approaches for cone vector tomography
  \end{itemize}}

% ---------------------------------------------------------------------------------------------------
% Projects
\newcommand{\OOXML}{%
  \entry{ooxml}{Open-source contributor \textbar\ \href{https://github.com/ChristopherVR/ooxml}{\color{accent}GitHub}}{Sep 2026 -- Present}
  \begin{itemize}
    \item Contributed to an open-source \textbf{TypeScript} engine that renders and edits PowerPoint and Word files in the browser,
          improving editing fidelity, \textbf{real-time collaboration} and load performance across React, Vue and Svelte viewers
  \end{itemize}}

\newcommand{\MirrorAI}{%
  \entry{Mirror AI: Deployable Personas}{Honorable mention in Supabase YC Hackathon}{Oct 2024 -- Dec 2024}
  \begin{itemize}
    \item Designed an \textbf{agentic LLM architecture} with LangGraph to mirror user personalities, creating interactive digital
          personas, and deployed it as a full-stack platform on Supabase and Vercel with one-click publishing
  \end{itemize}}

\newcommand{\LLMReasoning}{%
  \entry{Improving LLM Reasoning for Numerical Problems}{\report{https://sudhansh6.github.io/assets/pdf/MathPrompter++.pdf}}{Sep 2024 -- Dec 2024}
  \begin{itemize}
    \item Enhanced MathPrompter (ACL 2023) with CoT, achieving \textbf{10\% higher accuracy} on Llama 3.1 1B where prior methods failed,
          and reduced hallucination rates significantly by integrating multi-step validation
  \end{itemize}}

\newcommand{\InverseRendering}{%
  \entry{Inverse Rendering with 2D Gaussian Splatting}{\report{https://sudhansh6.github.io/assets/pdf/CSE252D_Report.pdf}}{Mar 2024 -- May 2024}
  \begin{itemize}
    \item Developed a novel \textbf{inverse rendering} framework in CUDA to recover PBR properties of a scene using 2D Gaussian
          Splatting, improving normal map MAE by \textbf{15\%} over the current SOTA methods
  \end{itemize}}

% ---------------------------------------------------------------------------------------------------
% Publications: title on the left, venue on the right, one line each.
\newcommand{\publine}[3]{\par\noindent\href{#1}{#2}\hfill\hspace{1em}\side{#3}\par\addvspace{1.5pt}}
\newcommand{\PubSaturation}{\publine{https://doi.org/10.1016/j.sigpro.2023.109349}{A likelihood based method for compressive signal recovery under Gaussian and saturation noise}{Signal Processing, 2024}}
\newcommand{\PubPermutation}{\publine{https://ieeexplore.ieee.org/document/10888147}{Identification and correction of permutation errors in compressed sensing-based group testing}{ICASSP 2025}}
\newcommand{\PubRearrangement}{\publine{https://arxiv.org/abs/2411.10899}{Planning for tabletop object rearrangement}{ICRA 2025}}

% ---------------------------------------------------------------------------------------------------
% Education, skills, service
% Each profile sets \UCSDCourses, \IITBCourses and \SkillRows before \input{resume-common}.
\providecommand{\UCSDCourses}{}
\providecommand{\IITBCourses}{}
\providecommand{\SkillRows}{}
\newcommand{\coursework}[1]{\ifblank{#1}{}{\par\noindent{\fontsize{9.5}{11.5}\selectfont{\color{muted}Coursework:} #1}\par}}
\newcommand{\Education}{%
  \section{Education}
  \entry{University of California San Diego}{M.S. in Computer Science and Engineering (AI)}{GPA 4.00/4.00 \quad Sep 2023 -- Mar 2025}
  \coursework{\UCSDCourses}
  \entry{Indian Institute of Technology Bombay}{B.Tech.\ with Honors in Computer Science}{CPI 9.66/10 \quad Jul 2019 -- Jul 2023}
  \coursework{\IITBCourses}}

% \skill{Label}{Items} is one row of the skills table.
\newcommand{\skill}[2]{{\headfont\bfseries #1} & #2 \\[1.5pt]}
\newcommand{\cpp}{C\kern-0.04em\raisebox{0.06ex}{\small++}}
\newcommand{\Skills}{%
  \section{Skills}
  \begin{tabularx}{\linewidth}{@{}l@{\hspace{12pt}}>{\raggedright\arraybackslash}X@{}}
    \SkillRows
  \end{tabularx}}
.
\providecommand{\UCSDCourses}{}
\providecommand{\IITBCourses}{}
\providecommand{\SkillRows}{}
\newcommand{\coursework}[1]{\ifblank{#1}{}{\par\noindent{\fontsize{9.5}{11.5}\selectfont{\color{muted}Coursework:} #1}\par}}
\newcommand{\Education}{%
  \section{Education}
  \entry{University of California San Diego}{M.S. in Computer Science and Engineering (AI)}{GPA 4.00/4.00 \quad Sep 2023 -- Mar 2025}
  \coursework{\UCSDCourses}
  \entry{Indian Institute of Technology Bombay}{B.Tech.\ with Honors in Computer Science}{CPI 9.66/10 \quad Jul 2019 -- Jul 2023}
  \coursework{\IITBCourses}}

% \skill{Label}{Items} is one row of the skills table.
\newcommand{\skill}[2]{{\headfont\bfseries #1} & #2 \\[1.5pt]}
\newcommand{\cpp}{C\kern-0.04em\raisebox{0.06ex}{\small++}}
\newcommand{\Skills}{%
  \section{Skills}
  \begin{tabularx}{\linewidth}{@{}l@{\hspace{12pt}}>{\raggedright\arraybackslash}X@{}}
    \SkillRows
  \end{tabularx}}
 to fit one page.
\providecommand{\secgap}{12pt}   % above a section title
\providecommand{\rulegap}{5.5pt}   % section rule to first entry
\providecommand{\entrygap}{7pt}  % between entries
\providecommand{\headgap}{1.5pt} % entry heading to its bullets
\providecommand{\itemgap}{0.6pt} % between bullets of one entry
\titlespacing*{\section}{0pt}{\secgap}{\rulegap}

\newcommand{\squarebullet}{\raisebox{0.32ex}{\color{accent}\rule{0.5ex}{0.5ex}}}
\setlist[itemize]{leftmargin=12pt, labelsep=6pt, topsep=\headgap, itemsep=\itemgap, parsep=0pt, partopsep=0pt,
  label=\squarebullet}

% \entry{Organisation or title}{Role or detail (may be empty)}{Dates}; \sub{Role or detail}{Right-hand note}
\newcommand{\entry}[3]{%
  \par\iffirstentry\global\firstentryfalse\else\addvspace{\entrygap}\fi\noindent{\headfont\fontsize{10.8}{13}\selectfont\textbf{#1}}\ifblank{#2}{}{\ {\color{muted}|}\ #2}\hfill\hspace{1em}\side{#3}\par}
\newcommand{\sub}[2]{%
  \noindent\textit{#1}\hfill{\color{muted}\textit{#2}}\par}
\newcommand{\report}[1]{\href{#1}{\color{accent}Report}}
\newcommand{\me}{\textbf{S. Peddabomma}}

% ---------------------------------------------------------------------------------------------------
% Header
\newcommand{\Header}{%
  \begin{center}
    {\namefont\fontsize{26}{28}\selectfont Sudhansh Peddabomma}\\[6pt]
    {\headfont\fontsize{9.5}{11}\selectfont\color{muted}
      \href{mailto:sudhansh6@gmail.com}{\faEnvelope[regular]\ sudhansh6@gmail.com} \quad
      \href{https://sudhansh6.github.io}{\faGlobe\ sudhansh6.github.io} \quad
      \href{https://www.linkedin.com/in/sudhansh-peddabomma-73a208192/}{\faLinkedin\ sudhansh-peddabomma} \quad
      \href{https://github.com/Sudhansh6}{\faGithub\ Sudhansh6}}
  \end{center}\vspace{-4pt}}

% ---------------------------------------------------------------------------------------------------
% Glean bullets, shared by both profiles
\newcommand{\GleanSlides}{\item Led \textbf{AI slide generation} in Glean Assistant, now creating \textbf{180K decks a week} for \textbf{50K weekly active users} at enterprises}
\newcommand{\GleanSlidesPerf}{\item Improved the slides harness, raising reliability to \textbf{97\%}, cutting median latency \textbf{26\%} and cost per deck by \textbf{half}
  in four months}
\newcommand{\GleanSlidesOrch}{\item Architected a \textbf{slide orchestration system} that reuses relevant slides from a company's existing decks and generates
  new ones in its brand and template, so teams can refine decks with each other and Glean in full context, without leaving the product}
\newcommand{\GleanAgents}{\item Moved Glean's \textbf{workflow agents} from one-shot planning to an \textbf{agentic tool-use loop}, greatly improving agentic
  capabilities and instruction following, with \textbf{18\%} higher answer completeness and \textbf{17\%} fewer tokens}
\newcommand{\GleanTeamChat}{\item Engineered the core of \textbf{Team chat}, letting teammates work with Glean in one shared conversation and co-edit its artifacts}
\newcommand{\GleanVideo}{\item Shipped \textbf{video artifacts} end to end, generating narrated videos in chat with streamed playback and admin rollout controls}
\newcommand{\GleanEvals}{\item Created the \textbf{evaluations and metrics} that guide model choices and rollouts across agents, slides and interactive artifacts}

% ---------------------------------------------------------------------------------------------------
% Experience
% \groupbreak separates related groups of bullets inside one long entry (slides, then platform work).
\providecommand{\groupgap}{3.5pt}
\newcommand{\groupbreak}{\vspace{\groupgap}}
\newcommand{\GleanAI}{%
  \entry{Glean}{Software Engineer}{Apr 2025 -- Present}
  \begin{itemize}
    \GleanSlides
    \GleanSlidesOrch
    \GleanSlidesPerf \groupbreak
    \GleanTeamChat
    \GleanVideo
    \GleanAgents
    \GleanEvals
  \end{itemize}}

\newcommand{\GleanVision}{%
  \entry{Glean}{Software Engineer}{Apr 2025 -- Present}
  \begin{itemize}
    \GleanSlides
    \GleanSlidesOrch
    \GleanSlidesPerf \groupbreak
    \GleanVideo
    \GleanEvals
  \end{itemize}}

\newcommand{\Duality}{%
  \entry{Duality AI}{Computer Vision Intern}{Jun 2024 -- Sep 2024}
  \begin{itemize}
    \item Designed automated \textbf{3D reconstruction} techniques for featureless objects, reducing digital-twin generation time
          by \textbf{40\%}
    \item Collaborated with \textbf{Autodesk} to validate Unreal Engine simulations for robotics tasks; leveraging structured
          domain randomization to reduce Sim2Real gap and increase mAP-50 by \textbf{15\%} for object detection and segmentation
  \end{itemize}}

\newcommand{\Microsoft}{%
  \entry{Microsoft}{Data and Applied Scientist Intern}{May 2022 -- Jul 2022}
  \begin{itemize}
    \item Developed a \textbf{decision-tree ranker} to recommend emails without user queries, improving \textbf{Outlook} search capabilities
    \item Proposed hierarchical feature-sets for the ranker, reducing latency for recommendations and improving recall
  \end{itemize}}

% ---------------------------------------------------------------------------------------------------
% Research
\newcommand{\HomeRobots}{%
  \entry{Real-time 3D Perception for Home Robots}{UC San Diego}{Sep 2023 -- Sep 2024}
  \begin{itemize}
    \item Investigated real-time dense \textbf{visual SLAM} methods using NeRFs and Gaussian Splatting for indoor robot navigation
    \item Integrated object segmentation, grasp-pose estimation, and 3D mapping on the Fetch robot via ROS, demonstrating
          a novel \textbf{tabletop object rearrangement} algorithm that reduced cost by \textbf{20\%} compared to the state-of-the-art
  \end{itemize}}

\newcommand{\Tomography}{%
  \entry{3D Tomography with Primal-Dual Neural Networks}{University College London}{May 2021 -- Jul 2023}
  \begin{itemize}
    \item Developed a stochastic neural-network architecture of primal-dual algorithm for online reconstruction of 3D
          volumes from tomographic projections, with \textbf{99.6\,\% structural similarity} in low-dosage conditions
    \item Built a Python library with custom gradient operators for reconstructing volumes in a single pass, reducing
          compute requirements by up to \textbf{5x} over SOTA learning-based approaches for cone vector tomography
  \end{itemize}}

% ---------------------------------------------------------------------------------------------------
% Projects
\newcommand{\OOXML}{%
  \entry{ooxml}{Open-source contributor \textbar\ \href{https://github.com/ChristopherVR/ooxml}{\color{accent}GitHub}}{Sep 2026 -- Present}
  \begin{itemize}
    \item Contributed to an open-source \textbf{TypeScript} engine that renders and edits PowerPoint and Word files in the browser,
          improving editing fidelity, \textbf{real-time collaboration} and load performance across React, Vue and Svelte viewers
  \end{itemize}}

\newcommand{\MirrorAI}{%
  \entry{Mirror AI: Deployable Personas}{Honorable mention in Supabase YC Hackathon}{Oct 2024 -- Dec 2024}
  \begin{itemize}
    \item Designed an \textbf{agentic LLM architecture} with LangGraph to mirror user personalities, creating interactive digital
          personas, and deployed it as a full-stack platform on Supabase and Vercel with one-click publishing
  \end{itemize}}

\newcommand{\LLMReasoning}{%
  \entry{Improving LLM Reasoning for Numerical Problems}{\report{https://sudhansh6.github.io/assets/pdf/MathPrompter++.pdf}}{Sep 2024 -- Dec 2024}
  \begin{itemize}
    \item Enhanced MathPrompter (ACL 2023) with CoT, achieving \textbf{10\% higher accuracy} on Llama 3.1 1B where prior methods failed,
          and reduced hallucination rates significantly by integrating multi-step validation
  \end{itemize}}

\newcommand{\InverseRendering}{%
  \entry{Inverse Rendering with 2D Gaussian Splatting}{\report{https://sudhansh6.github.io/assets/pdf/CSE252D_Report.pdf}}{Mar 2024 -- May 2024}
  \begin{itemize}
    \item Developed a novel \textbf{inverse rendering} framework in CUDA to recover PBR properties of a scene using 2D Gaussian
          Splatting, improving normal map MAE by \textbf{15\%} over the current SOTA methods
  \end{itemize}}

% ---------------------------------------------------------------------------------------------------
% Publications: title on the left, venue on the right, one line each.
\newcommand{\publine}[3]{\par\noindent\href{#1}{#2}\hfill\hspace{1em}\side{#3}\par\addvspace{1.5pt}}
\newcommand{\PubSaturation}{\publine{https://doi.org/10.1016/j.sigpro.2023.109349}{A likelihood based method for compressive signal recovery under Gaussian and saturation noise}{Signal Processing, 2024}}
\newcommand{\PubPermutation}{\publine{https://ieeexplore.ieee.org/document/10888147}{Identification and correction of permutation errors in compressed sensing-based group testing}{ICASSP 2025}}
\newcommand{\PubRearrangement}{\publine{https://arxiv.org/abs/2411.10899}{Planning for tabletop object rearrangement}{ICRA 2025}}

% ---------------------------------------------------------------------------------------------------
% Education, skills, service
% Each profile sets \UCSDCourses, \IITBCourses and \SkillRows before % Shared preamble and content for Sudhansh Peddabomma's résumés.
% Each profile (SudhanshPeddabomma-AI.tex, SudhanshPeddabomma-Vision.tex) inputs this file and picks
% which blocks to show and in what order, so every bullet is written once.
% Layout rule: every bullet fills one line, or runs 1.5 to 2 lines. Check with: python3 check_lines.py <pdf>
\documentclass[10pt,letterpaper]{article}

\usepackage[margin=0.55in,top=0.5in,bottom=0.45in]{geometry}
\usepackage[T1]{fontenc}
% Fonts: \resumefont loads the body font; \headfont styles the name and section titles.
% Override either before % Shared preamble and content for Sudhansh Peddabomma's résumés.
% Each profile (SudhanshPeddabomma-AI.tex, SudhanshPeddabomma-Vision.tex) inputs this file and picks
% which blocks to show and in what order, so every bullet is written once.
% Layout rule: every bullet fills one line, or runs 1.5 to 2 lines. Check with: python3 check_lines.py <pdf>
\documentclass[10pt,letterpaper]{article}

\usepackage[margin=0.55in,top=0.5in,bottom=0.45in]{geometry}
\usepackage[T1]{fontenc}
% Fonts: \resumefont loads the body font; \headfont styles the name and section titles.
% Override either before \input{resume-common} to try another pairing.
% Default: Charter body with Source Sans headings. All-sans alternative:
%   \def\resumefont{\usepackage[default,semibold]{sourcesanspro}}\def\headfont{}
\providecommand{\resumefont}{\usepackage[scaled=0.96]{XCharter}\usepackage[semibold,scale=0.96]{sourcesanspro}}
\providecommand{\headfont}{\sffamily}
\providecommand{\namefont}{\rmfamily\scshape}
\resumefont
\usepackage[final]{microtype}
\usepackage[dvipsnames]{xcolor}
\usepackage{enumitem}
\usepackage{titlesec}
\usepackage{tabularx}
\usepackage{fontawesome5}
\usepackage{etoolbox}
\usepackage[hidelinks]{hyperref}

% The site's green, darkened for print.
\definecolor{accent}{HTML}{3A7A45}
\definecolor{muted}{HTML}{5C675D}

\pagestyle{empty}
\setlength{\parindent}{0pt}
\setlength{\tabcolsep}{0pt}
\providecommand{\leading}{1.1}
\linespread{\leading}

\hypersetup{pdfauthor={Sudhansh Peddabomma}}

% Type scale: name 26pt small caps, section titles 12.5pt caps, entry titles 10.8pt, body 10pt,
% side notes 9pt.
\newcommand{\side}[1]{{\fontsize{9}{11}\selectfont\color{muted}#1}}
\newcommand{\sectitle}[1]{\textls[60]{\MakeUppercase{#1}}}
\titleformat{\section}{\headfont\color{accent}\bfseries\fontsize{12.5}{14}\selectfont}{}{0pt}{\sectitle}[\vspace{1.5pt}{\color{accent}\titlerule[0.8pt]}\global\firstentrytrue]
% The first entry after a section title skips \entrygap, so only \rulegap sits under the rule.
\newif\iffirstentry

% Vertical rhythm, the same idea as the earlier hand-tuned résumés: bullets inside an entry sit
% tight together, entries are separated by a clear gap, and sections by a larger one. A profile
% may override these before \input{resume-common} to fit one page.
\providecommand{\secgap}{12pt}   % above a section title
\providecommand{\rulegap}{5.5pt}   % section rule to first entry
\providecommand{\entrygap}{7pt}  % between entries
\providecommand{\headgap}{1.5pt} % entry heading to its bullets
\providecommand{\itemgap}{0.6pt} % between bullets of one entry
\titlespacing*{\section}{0pt}{\secgap}{\rulegap}

\newcommand{\squarebullet}{\raisebox{0.32ex}{\color{accent}\rule{0.5ex}{0.5ex}}}
\setlist[itemize]{leftmargin=12pt, labelsep=6pt, topsep=\headgap, itemsep=\itemgap, parsep=0pt, partopsep=0pt,
  label=\squarebullet}

% \entry{Organisation or title}{Role or detail (may be empty)}{Dates}; \sub{Role or detail}{Right-hand note}
\newcommand{\entry}[3]{%
  \par\iffirstentry\global\firstentryfalse\else\addvspace{\entrygap}\fi\noindent{\headfont\fontsize{10.8}{13}\selectfont\textbf{#1}}\ifblank{#2}{}{\ {\color{muted}|}\ #2}\hfill\hspace{1em}\side{#3}\par}
\newcommand{\sub}[2]{%
  \noindent\textit{#1}\hfill{\color{muted}\textit{#2}}\par}
\newcommand{\report}[1]{\href{#1}{\color{accent}Report}}
\newcommand{\me}{\textbf{S. Peddabomma}}

% ---------------------------------------------------------------------------------------------------
% Header
\newcommand{\Header}{%
  \begin{center}
    {\namefont\fontsize{26}{28}\selectfont Sudhansh Peddabomma}\\[6pt]
    {\headfont\fontsize{9.5}{11}\selectfont\color{muted}
      \href{mailto:sudhansh6@gmail.com}{\faEnvelope[regular]\ sudhansh6@gmail.com} \quad
      \href{https://sudhansh6.github.io}{\faGlobe\ sudhansh6.github.io} \quad
      \href{https://www.linkedin.com/in/sudhansh-peddabomma-73a208192/}{\faLinkedin\ sudhansh-peddabomma} \quad
      \href{https://github.com/Sudhansh6}{\faGithub\ Sudhansh6}}
  \end{center}\vspace{-4pt}}

% ---------------------------------------------------------------------------------------------------
% Glean bullets, shared by both profiles
\newcommand{\GleanSlides}{\item Led \textbf{AI slide generation} in Glean Assistant, now creating \textbf{180K decks a week} for \textbf{50K weekly active users} at enterprises}
\newcommand{\GleanSlidesPerf}{\item Improved the slides harness, raising reliability to \textbf{97\%}, cutting median latency \textbf{26\%} and cost per deck by \textbf{half}
  in four months}
\newcommand{\GleanSlidesOrch}{\item Architected a \textbf{slide orchestration system} that reuses relevant slides from a company's existing decks and generates
  new ones in its brand and template, so teams can refine decks with each other and Glean in full context, without leaving the product}
\newcommand{\GleanAgents}{\item Moved Glean's \textbf{workflow agents} from one-shot planning to an \textbf{agentic tool-use loop}, greatly improving agentic
  capabilities and instruction following, with \textbf{18\%} higher answer completeness and \textbf{17\%} fewer tokens}
\newcommand{\GleanTeamChat}{\item Engineered the core of \textbf{Team chat}, letting teammates work with Glean in one shared conversation and co-edit its artifacts}
\newcommand{\GleanVideo}{\item Shipped \textbf{video artifacts} end to end, generating narrated videos in chat with streamed playback and admin rollout controls}
\newcommand{\GleanEvals}{\item Created the \textbf{evaluations and metrics} that guide model choices and rollouts across agents, slides and interactive artifacts}

% ---------------------------------------------------------------------------------------------------
% Experience
% \groupbreak separates related groups of bullets inside one long entry (slides, then platform work).
\providecommand{\groupgap}{3.5pt}
\newcommand{\groupbreak}{\vspace{\groupgap}}
\newcommand{\GleanAI}{%
  \entry{Glean}{Software Engineer}{Apr 2025 -- Present}
  \begin{itemize}
    \GleanSlides
    \GleanSlidesOrch
    \GleanSlidesPerf \groupbreak
    \GleanTeamChat
    \GleanVideo
    \GleanAgents
    \GleanEvals
  \end{itemize}}

\newcommand{\GleanVision}{%
  \entry{Glean}{Software Engineer}{Apr 2025 -- Present}
  \begin{itemize}
    \GleanSlides
    \GleanSlidesOrch
    \GleanSlidesPerf \groupbreak
    \GleanVideo
    \GleanEvals
  \end{itemize}}

\newcommand{\Duality}{%
  \entry{Duality AI}{Computer Vision Intern}{Jun 2024 -- Sep 2024}
  \begin{itemize}
    \item Designed automated \textbf{3D reconstruction} techniques for featureless objects, reducing digital-twin generation time
          by \textbf{40\%}
    \item Collaborated with \textbf{Autodesk} to validate Unreal Engine simulations for robotics tasks; leveraging structured
          domain randomization to reduce Sim2Real gap and increase mAP-50 by \textbf{15\%} for object detection and segmentation
  \end{itemize}}

\newcommand{\Microsoft}{%
  \entry{Microsoft}{Data and Applied Scientist Intern}{May 2022 -- Jul 2022}
  \begin{itemize}
    \item Developed a \textbf{decision-tree ranker} to recommend emails without user queries, improving \textbf{Outlook} search capabilities
    \item Proposed hierarchical feature-sets for the ranker, reducing latency for recommendations and improving recall
  \end{itemize}}

% ---------------------------------------------------------------------------------------------------
% Research
\newcommand{\HomeRobots}{%
  \entry{Real-time 3D Perception for Home Robots}{UC San Diego}{Sep 2023 -- Sep 2024}
  \begin{itemize}
    \item Investigated real-time dense \textbf{visual SLAM} methods using NeRFs and Gaussian Splatting for indoor robot navigation
    \item Integrated object segmentation, grasp-pose estimation, and 3D mapping on the Fetch robot via ROS, demonstrating
          a novel \textbf{tabletop object rearrangement} algorithm that reduced cost by \textbf{20\%} compared to the state-of-the-art
  \end{itemize}}

\newcommand{\Tomography}{%
  \entry{3D Tomography with Primal-Dual Neural Networks}{University College London}{May 2021 -- Jul 2023}
  \begin{itemize}
    \item Developed a stochastic neural-network architecture of primal-dual algorithm for online reconstruction of 3D
          volumes from tomographic projections, with \textbf{99.6\,\% structural similarity} in low-dosage conditions
    \item Built a Python library with custom gradient operators for reconstructing volumes in a single pass, reducing
          compute requirements by up to \textbf{5x} over SOTA learning-based approaches for cone vector tomography
  \end{itemize}}

% ---------------------------------------------------------------------------------------------------
% Projects
\newcommand{\OOXML}{%
  \entry{ooxml}{Open-source contributor \textbar\ \href{https://github.com/ChristopherVR/ooxml}{\color{accent}GitHub}}{Sep 2026 -- Present}
  \begin{itemize}
    \item Contributed to an open-source \textbf{TypeScript} engine that renders and edits PowerPoint and Word files in the browser,
          improving editing fidelity, \textbf{real-time collaboration} and load performance across React, Vue and Svelte viewers
  \end{itemize}}

\newcommand{\MirrorAI}{%
  \entry{Mirror AI: Deployable Personas}{Honorable mention in Supabase YC Hackathon}{Oct 2024 -- Dec 2024}
  \begin{itemize}
    \item Designed an \textbf{agentic LLM architecture} with LangGraph to mirror user personalities, creating interactive digital
          personas, and deployed it as a full-stack platform on Supabase and Vercel with one-click publishing
  \end{itemize}}

\newcommand{\LLMReasoning}{%
  \entry{Improving LLM Reasoning for Numerical Problems}{\report{https://sudhansh6.github.io/assets/pdf/MathPrompter++.pdf}}{Sep 2024 -- Dec 2024}
  \begin{itemize}
    \item Enhanced MathPrompter (ACL 2023) with CoT, achieving \textbf{10\% higher accuracy} on Llama 3.1 1B where prior methods failed,
          and reduced hallucination rates significantly by integrating multi-step validation
  \end{itemize}}

\newcommand{\InverseRendering}{%
  \entry{Inverse Rendering with 2D Gaussian Splatting}{\report{https://sudhansh6.github.io/assets/pdf/CSE252D_Report.pdf}}{Mar 2024 -- May 2024}
  \begin{itemize}
    \item Developed a novel \textbf{inverse rendering} framework in CUDA to recover PBR properties of a scene using 2D Gaussian
          Splatting, improving normal map MAE by \textbf{15\%} over the current SOTA methods
  \end{itemize}}

% ---------------------------------------------------------------------------------------------------
% Publications: title on the left, venue on the right, one line each.
\newcommand{\publine}[3]{\par\noindent\href{#1}{#2}\hfill\hspace{1em}\side{#3}\par\addvspace{1.5pt}}
\newcommand{\PubSaturation}{\publine{https://doi.org/10.1016/j.sigpro.2023.109349}{A likelihood based method for compressive signal recovery under Gaussian and saturation noise}{Signal Processing, 2024}}
\newcommand{\PubPermutation}{\publine{https://ieeexplore.ieee.org/document/10888147}{Identification and correction of permutation errors in compressed sensing-based group testing}{ICASSP 2025}}
\newcommand{\PubRearrangement}{\publine{https://arxiv.org/abs/2411.10899}{Planning for tabletop object rearrangement}{ICRA 2025}}

% ---------------------------------------------------------------------------------------------------
% Education, skills, service
% Each profile sets \UCSDCourses, \IITBCourses and \SkillRows before \input{resume-common}.
\providecommand{\UCSDCourses}{}
\providecommand{\IITBCourses}{}
\providecommand{\SkillRows}{}
\newcommand{\coursework}[1]{\ifblank{#1}{}{\par\noindent{\fontsize{9.5}{11.5}\selectfont{\color{muted}Coursework:} #1}\par}}
\newcommand{\Education}{%
  \section{Education}
  \entry{University of California San Diego}{M.S. in Computer Science and Engineering (AI)}{GPA 4.00/4.00 \quad Sep 2023 -- Mar 2025}
  \coursework{\UCSDCourses}
  \entry{Indian Institute of Technology Bombay}{B.Tech.\ with Honors in Computer Science}{CPI 9.66/10 \quad Jul 2019 -- Jul 2023}
  \coursework{\IITBCourses}}

% \skill{Label}{Items} is one row of the skills table.
\newcommand{\skill}[2]{{\headfont\bfseries #1} & #2 \\[1.5pt]}
\newcommand{\cpp}{C\kern-0.04em\raisebox{0.06ex}{\small++}}
\newcommand{\Skills}{%
  \section{Skills}
  \begin{tabularx}{\linewidth}{@{}l@{\hspace{12pt}}>{\raggedright\arraybackslash}X@{}}
    \SkillRows
  \end{tabularx}}
 to try another pairing.
% Default: Charter body with Source Sans headings. All-sans alternative:
%   \def\resumefont{\usepackage[default,semibold]{sourcesanspro}}\def\headfont{}
\providecommand{\resumefont}{\usepackage[scaled=0.96]{XCharter}\usepackage[semibold,scale=0.96]{sourcesanspro}}
\providecommand{\headfont}{\sffamily}
\providecommand{\namefont}{\rmfamily\scshape}
\resumefont
\usepackage[final]{microtype}
\usepackage[dvipsnames]{xcolor}
\usepackage{enumitem}
\usepackage{titlesec}
\usepackage{tabularx}
\usepackage{fontawesome5}
\usepackage{etoolbox}
\usepackage[hidelinks]{hyperref}

% The site's green, darkened for print.
\definecolor{accent}{HTML}{3A7A45}
\definecolor{muted}{HTML}{5C675D}

\pagestyle{empty}
\setlength{\parindent}{0pt}
\setlength{\tabcolsep}{0pt}
\providecommand{\leading}{1.1}
\linespread{\leading}

\hypersetup{pdfauthor={Sudhansh Peddabomma}}

% Type scale: name 26pt small caps, section titles 12.5pt caps, entry titles 10.8pt, body 10pt,
% side notes 9pt.
\newcommand{\side}[1]{{\fontsize{9}{11}\selectfont\color{muted}#1}}
\newcommand{\sectitle}[1]{\textls[60]{\MakeUppercase{#1}}}
\titleformat{\section}{\headfont\color{accent}\bfseries\fontsize{12.5}{14}\selectfont}{}{0pt}{\sectitle}[\vspace{1.5pt}{\color{accent}\titlerule[0.8pt]}\global\firstentrytrue]
% The first entry after a section title skips \entrygap, so only \rulegap sits under the rule.
\newif\iffirstentry

% Vertical rhythm, the same idea as the earlier hand-tuned résumés: bullets inside an entry sit
% tight together, entries are separated by a clear gap, and sections by a larger one. A profile
% may override these before % Shared preamble and content for Sudhansh Peddabomma's résumés.
% Each profile (SudhanshPeddabomma-AI.tex, SudhanshPeddabomma-Vision.tex) inputs this file and picks
% which blocks to show and in what order, so every bullet is written once.
% Layout rule: every bullet fills one line, or runs 1.5 to 2 lines. Check with: python3 check_lines.py <pdf>
\documentclass[10pt,letterpaper]{article}

\usepackage[margin=0.55in,top=0.5in,bottom=0.45in]{geometry}
\usepackage[T1]{fontenc}
% Fonts: \resumefont loads the body font; \headfont styles the name and section titles.
% Override either before \input{resume-common} to try another pairing.
% Default: Charter body with Source Sans headings. All-sans alternative:
%   \def\resumefont{\usepackage[default,semibold]{sourcesanspro}}\def\headfont{}
\providecommand{\resumefont}{\usepackage[scaled=0.96]{XCharter}\usepackage[semibold,scale=0.96]{sourcesanspro}}
\providecommand{\headfont}{\sffamily}
\providecommand{\namefont}{\rmfamily\scshape}
\resumefont
\usepackage[final]{microtype}
\usepackage[dvipsnames]{xcolor}
\usepackage{enumitem}
\usepackage{titlesec}
\usepackage{tabularx}
\usepackage{fontawesome5}
\usepackage{etoolbox}
\usepackage[hidelinks]{hyperref}

% The site's green, darkened for print.
\definecolor{accent}{HTML}{3A7A45}
\definecolor{muted}{HTML}{5C675D}

\pagestyle{empty}
\setlength{\parindent}{0pt}
\setlength{\tabcolsep}{0pt}
\providecommand{\leading}{1.1}
\linespread{\leading}

\hypersetup{pdfauthor={Sudhansh Peddabomma}}

% Type scale: name 26pt small caps, section titles 12.5pt caps, entry titles 10.8pt, body 10pt,
% side notes 9pt.
\newcommand{\side}[1]{{\fontsize{9}{11}\selectfont\color{muted}#1}}
\newcommand{\sectitle}[1]{\textls[60]{\MakeUppercase{#1}}}
\titleformat{\section}{\headfont\color{accent}\bfseries\fontsize{12.5}{14}\selectfont}{}{0pt}{\sectitle}[\vspace{1.5pt}{\color{accent}\titlerule[0.8pt]}\global\firstentrytrue]
% The first entry after a section title skips \entrygap, so only \rulegap sits under the rule.
\newif\iffirstentry

% Vertical rhythm, the same idea as the earlier hand-tuned résumés: bullets inside an entry sit
% tight together, entries are separated by a clear gap, and sections by a larger one. A profile
% may override these before \input{resume-common} to fit one page.
\providecommand{\secgap}{12pt}   % above a section title
\providecommand{\rulegap}{5.5pt}   % section rule to first entry
\providecommand{\entrygap}{7pt}  % between entries
\providecommand{\headgap}{1.5pt} % entry heading to its bullets
\providecommand{\itemgap}{0.6pt} % between bullets of one entry
\titlespacing*{\section}{0pt}{\secgap}{\rulegap}

\newcommand{\squarebullet}{\raisebox{0.32ex}{\color{accent}\rule{0.5ex}{0.5ex}}}
\setlist[itemize]{leftmargin=12pt, labelsep=6pt, topsep=\headgap, itemsep=\itemgap, parsep=0pt, partopsep=0pt,
  label=\squarebullet}

% \entry{Organisation or title}{Role or detail (may be empty)}{Dates}; \sub{Role or detail}{Right-hand note}
\newcommand{\entry}[3]{%
  \par\iffirstentry\global\firstentryfalse\else\addvspace{\entrygap}\fi\noindent{\headfont\fontsize{10.8}{13}\selectfont\textbf{#1}}\ifblank{#2}{}{\ {\color{muted}|}\ #2}\hfill\hspace{1em}\side{#3}\par}
\newcommand{\sub}[2]{%
  \noindent\textit{#1}\hfill{\color{muted}\textit{#2}}\par}
\newcommand{\report}[1]{\href{#1}{\color{accent}Report}}
\newcommand{\me}{\textbf{S. Peddabomma}}

% ---------------------------------------------------------------------------------------------------
% Header
\newcommand{\Header}{%
  \begin{center}
    {\namefont\fontsize{26}{28}\selectfont Sudhansh Peddabomma}\\[6pt]
    {\headfont\fontsize{9.5}{11}\selectfont\color{muted}
      \href{mailto:sudhansh6@gmail.com}{\faEnvelope[regular]\ sudhansh6@gmail.com} \quad
      \href{https://sudhansh6.github.io}{\faGlobe\ sudhansh6.github.io} \quad
      \href{https://www.linkedin.com/in/sudhansh-peddabomma-73a208192/}{\faLinkedin\ sudhansh-peddabomma} \quad
      \href{https://github.com/Sudhansh6}{\faGithub\ Sudhansh6}}
  \end{center}\vspace{-4pt}}

% ---------------------------------------------------------------------------------------------------
% Glean bullets, shared by both profiles
\newcommand{\GleanSlides}{\item Led \textbf{AI slide generation} in Glean Assistant, now creating \textbf{180K decks a week} for \textbf{50K weekly active users} at enterprises}
\newcommand{\GleanSlidesPerf}{\item Improved the slides harness, raising reliability to \textbf{97\%}, cutting median latency \textbf{26\%} and cost per deck by \textbf{half}
  in four months}
\newcommand{\GleanSlidesOrch}{\item Architected a \textbf{slide orchestration system} that reuses relevant slides from a company's existing decks and generates
  new ones in its brand and template, so teams can refine decks with each other and Glean in full context, without leaving the product}
\newcommand{\GleanAgents}{\item Moved Glean's \textbf{workflow agents} from one-shot planning to an \textbf{agentic tool-use loop}, greatly improving agentic
  capabilities and instruction following, with \textbf{18\%} higher answer completeness and \textbf{17\%} fewer tokens}
\newcommand{\GleanTeamChat}{\item Engineered the core of \textbf{Team chat}, letting teammates work with Glean in one shared conversation and co-edit its artifacts}
\newcommand{\GleanVideo}{\item Shipped \textbf{video artifacts} end to end, generating narrated videos in chat with streamed playback and admin rollout controls}
\newcommand{\GleanEvals}{\item Created the \textbf{evaluations and metrics} that guide model choices and rollouts across agents, slides and interactive artifacts}

% ---------------------------------------------------------------------------------------------------
% Experience
% \groupbreak separates related groups of bullets inside one long entry (slides, then platform work).
\providecommand{\groupgap}{3.5pt}
\newcommand{\groupbreak}{\vspace{\groupgap}}
\newcommand{\GleanAI}{%
  \entry{Glean}{Software Engineer}{Apr 2025 -- Present}
  \begin{itemize}
    \GleanSlides
    \GleanSlidesOrch
    \GleanSlidesPerf \groupbreak
    \GleanTeamChat
    \GleanVideo
    \GleanAgents
    \GleanEvals
  \end{itemize}}

\newcommand{\GleanVision}{%
  \entry{Glean}{Software Engineer}{Apr 2025 -- Present}
  \begin{itemize}
    \GleanSlides
    \GleanSlidesOrch
    \GleanSlidesPerf \groupbreak
    \GleanVideo
    \GleanEvals
  \end{itemize}}

\newcommand{\Duality}{%
  \entry{Duality AI}{Computer Vision Intern}{Jun 2024 -- Sep 2024}
  \begin{itemize}
    \item Designed automated \textbf{3D reconstruction} techniques for featureless objects, reducing digital-twin generation time
          by \textbf{40\%}
    \item Collaborated with \textbf{Autodesk} to validate Unreal Engine simulations for robotics tasks; leveraging structured
          domain randomization to reduce Sim2Real gap and increase mAP-50 by \textbf{15\%} for object detection and segmentation
  \end{itemize}}

\newcommand{\Microsoft}{%
  \entry{Microsoft}{Data and Applied Scientist Intern}{May 2022 -- Jul 2022}
  \begin{itemize}
    \item Developed a \textbf{decision-tree ranker} to recommend emails without user queries, improving \textbf{Outlook} search capabilities
    \item Proposed hierarchical feature-sets for the ranker, reducing latency for recommendations and improving recall
  \end{itemize}}

% ---------------------------------------------------------------------------------------------------
% Research
\newcommand{\HomeRobots}{%
  \entry{Real-time 3D Perception for Home Robots}{UC San Diego}{Sep 2023 -- Sep 2024}
  \begin{itemize}
    \item Investigated real-time dense \textbf{visual SLAM} methods using NeRFs and Gaussian Splatting for indoor robot navigation
    \item Integrated object segmentation, grasp-pose estimation, and 3D mapping on the Fetch robot via ROS, demonstrating
          a novel \textbf{tabletop object rearrangement} algorithm that reduced cost by \textbf{20\%} compared to the state-of-the-art
  \end{itemize}}

\newcommand{\Tomography}{%
  \entry{3D Tomography with Primal-Dual Neural Networks}{University College London}{May 2021 -- Jul 2023}
  \begin{itemize}
    \item Developed a stochastic neural-network architecture of primal-dual algorithm for online reconstruction of 3D
          volumes from tomographic projections, with \textbf{99.6\,\% structural similarity} in low-dosage conditions
    \item Built a Python library with custom gradient operators for reconstructing volumes in a single pass, reducing
          compute requirements by up to \textbf{5x} over SOTA learning-based approaches for cone vector tomography
  \end{itemize}}

% ---------------------------------------------------------------------------------------------------
% Projects
\newcommand{\OOXML}{%
  \entry{ooxml}{Open-source contributor \textbar\ \href{https://github.com/ChristopherVR/ooxml}{\color{accent}GitHub}}{Sep 2026 -- Present}
  \begin{itemize}
    \item Contributed to an open-source \textbf{TypeScript} engine that renders and edits PowerPoint and Word files in the browser,
          improving editing fidelity, \textbf{real-time collaboration} and load performance across React, Vue and Svelte viewers
  \end{itemize}}

\newcommand{\MirrorAI}{%
  \entry{Mirror AI: Deployable Personas}{Honorable mention in Supabase YC Hackathon}{Oct 2024 -- Dec 2024}
  \begin{itemize}
    \item Designed an \textbf{agentic LLM architecture} with LangGraph to mirror user personalities, creating interactive digital
          personas, and deployed it as a full-stack platform on Supabase and Vercel with one-click publishing
  \end{itemize}}

\newcommand{\LLMReasoning}{%
  \entry{Improving LLM Reasoning for Numerical Problems}{\report{https://sudhansh6.github.io/assets/pdf/MathPrompter++.pdf}}{Sep 2024 -- Dec 2024}
  \begin{itemize}
    \item Enhanced MathPrompter (ACL 2023) with CoT, achieving \textbf{10\% higher accuracy} on Llama 3.1 1B where prior methods failed,
          and reduced hallucination rates significantly by integrating multi-step validation
  \end{itemize}}

\newcommand{\InverseRendering}{%
  \entry{Inverse Rendering with 2D Gaussian Splatting}{\report{https://sudhansh6.github.io/assets/pdf/CSE252D_Report.pdf}}{Mar 2024 -- May 2024}
  \begin{itemize}
    \item Developed a novel \textbf{inverse rendering} framework in CUDA to recover PBR properties of a scene using 2D Gaussian
          Splatting, improving normal map MAE by \textbf{15\%} over the current SOTA methods
  \end{itemize}}

% ---------------------------------------------------------------------------------------------------
% Publications: title on the left, venue on the right, one line each.
\newcommand{\publine}[3]{\par\noindent\href{#1}{#2}\hfill\hspace{1em}\side{#3}\par\addvspace{1.5pt}}
\newcommand{\PubSaturation}{\publine{https://doi.org/10.1016/j.sigpro.2023.109349}{A likelihood based method for compressive signal recovery under Gaussian and saturation noise}{Signal Processing, 2024}}
\newcommand{\PubPermutation}{\publine{https://ieeexplore.ieee.org/document/10888147}{Identification and correction of permutation errors in compressed sensing-based group testing}{ICASSP 2025}}
\newcommand{\PubRearrangement}{\publine{https://arxiv.org/abs/2411.10899}{Planning for tabletop object rearrangement}{ICRA 2025}}

% ---------------------------------------------------------------------------------------------------
% Education, skills, service
% Each profile sets \UCSDCourses, \IITBCourses and \SkillRows before \input{resume-common}.
\providecommand{\UCSDCourses}{}
\providecommand{\IITBCourses}{}
\providecommand{\SkillRows}{}
\newcommand{\coursework}[1]{\ifblank{#1}{}{\par\noindent{\fontsize{9.5}{11.5}\selectfont{\color{muted}Coursework:} #1}\par}}
\newcommand{\Education}{%
  \section{Education}
  \entry{University of California San Diego}{M.S. in Computer Science and Engineering (AI)}{GPA 4.00/4.00 \quad Sep 2023 -- Mar 2025}
  \coursework{\UCSDCourses}
  \entry{Indian Institute of Technology Bombay}{B.Tech.\ with Honors in Computer Science}{CPI 9.66/10 \quad Jul 2019 -- Jul 2023}
  \coursework{\IITBCourses}}

% \skill{Label}{Items} is one row of the skills table.
\newcommand{\skill}[2]{{\headfont\bfseries #1} & #2 \\[1.5pt]}
\newcommand{\cpp}{C\kern-0.04em\raisebox{0.06ex}{\small++}}
\newcommand{\Skills}{%
  \section{Skills}
  \begin{tabularx}{\linewidth}{@{}l@{\hspace{12pt}}>{\raggedright\arraybackslash}X@{}}
    \SkillRows
  \end{tabularx}}
 to fit one page.
\providecommand{\secgap}{12pt}   % above a section title
\providecommand{\rulegap}{5.5pt}   % section rule to first entry
\providecommand{\entrygap}{7pt}  % between entries
\providecommand{\headgap}{1.5pt} % entry heading to its bullets
\providecommand{\itemgap}{0.6pt} % between bullets of one entry
\titlespacing*{\section}{0pt}{\secgap}{\rulegap}

\newcommand{\squarebullet}{\raisebox{0.32ex}{\color{accent}\rule{0.5ex}{0.5ex}}}
\setlist[itemize]{leftmargin=12pt, labelsep=6pt, topsep=\headgap, itemsep=\itemgap, parsep=0pt, partopsep=0pt,
  label=\squarebullet}

% \entry{Organisation or title}{Role or detail (may be empty)}{Dates}; \sub{Role or detail}{Right-hand note}
\newcommand{\entry}[3]{%
  \par\iffirstentry\global\firstentryfalse\else\addvspace{\entrygap}\fi\noindent{\headfont\fontsize{10.8}{13}\selectfont\textbf{#1}}\ifblank{#2}{}{\ {\color{muted}|}\ #2}\hfill\hspace{1em}\side{#3}\par}
\newcommand{\sub}[2]{%
  \noindent\textit{#1}\hfill{\color{muted}\textit{#2}}\par}
\newcommand{\report}[1]{\href{#1}{\color{accent}Report}}
\newcommand{\me}{\textbf{S. Peddabomma}}

% ---------------------------------------------------------------------------------------------------
% Header
\newcommand{\Header}{%
  \begin{center}
    {\namefont\fontsize{26}{28}\selectfont Sudhansh Peddabomma}\\[6pt]
    {\headfont\fontsize{9.5}{11}\selectfont\color{muted}
      \href{mailto:sudhansh6@gmail.com}{\faEnvelope[regular]\ sudhansh6@gmail.com} \quad
      \href{https://sudhansh6.github.io}{\faGlobe\ sudhansh6.github.io} \quad
      \href{https://www.linkedin.com/in/sudhansh-peddabomma-73a208192/}{\faLinkedin\ sudhansh-peddabomma} \quad
      \href{https://github.com/Sudhansh6}{\faGithub\ Sudhansh6}}
  \end{center}\vspace{-4pt}}

% ---------------------------------------------------------------------------------------------------
% Glean bullets, shared by both profiles
\newcommand{\GleanSlides}{\item Led \textbf{AI slide generation} in Glean Assistant, now creating \textbf{180K decks a week} for \textbf{50K weekly active users} at enterprises}
\newcommand{\GleanSlidesPerf}{\item Improved the slides harness, raising reliability to \textbf{97\%}, cutting median latency \textbf{26\%} and cost per deck by \textbf{half}
  in four months}
\newcommand{\GleanSlidesOrch}{\item Architected a \textbf{slide orchestration system} that reuses relevant slides from a company's existing decks and generates
  new ones in its brand and template, so teams can refine decks with each other and Glean in full context, without leaving the product}
\newcommand{\GleanAgents}{\item Moved Glean's \textbf{workflow agents} from one-shot planning to an \textbf{agentic tool-use loop}, greatly improving agentic
  capabilities and instruction following, with \textbf{18\%} higher answer completeness and \textbf{17\%} fewer tokens}
\newcommand{\GleanTeamChat}{\item Engineered the core of \textbf{Team chat}, letting teammates work with Glean in one shared conversation and co-edit its artifacts}
\newcommand{\GleanVideo}{\item Shipped \textbf{video artifacts} end to end, generating narrated videos in chat with streamed playback and admin rollout controls}
\newcommand{\GleanEvals}{\item Created the \textbf{evaluations and metrics} that guide model choices and rollouts across agents, slides and interactive artifacts}

% ---------------------------------------------------------------------------------------------------
% Experience
% \groupbreak separates related groups of bullets inside one long entry (slides, then platform work).
\providecommand{\groupgap}{3.5pt}
\newcommand{\groupbreak}{\vspace{\groupgap}}
\newcommand{\GleanAI}{%
  \entry{Glean}{Software Engineer}{Apr 2025 -- Present}
  \begin{itemize}
    \GleanSlides
    \GleanSlidesOrch
    \GleanSlidesPerf \groupbreak
    \GleanTeamChat
    \GleanVideo
    \GleanAgents
    \GleanEvals
  \end{itemize}}

\newcommand{\GleanVision}{%
  \entry{Glean}{Software Engineer}{Apr 2025 -- Present}
  \begin{itemize}
    \GleanSlides
    \GleanSlidesOrch
    \GleanSlidesPerf \groupbreak
    \GleanVideo
    \GleanEvals
  \end{itemize}}

\newcommand{\Duality}{%
  \entry{Duality AI}{Computer Vision Intern}{Jun 2024 -- Sep 2024}
  \begin{itemize}
    \item Designed automated \textbf{3D reconstruction} techniques for featureless objects, reducing digital-twin generation time
          by \textbf{40\%}
    \item Collaborated with \textbf{Autodesk} to validate Unreal Engine simulations for robotics tasks; leveraging structured
          domain randomization to reduce Sim2Real gap and increase mAP-50 by \textbf{15\%} for object detection and segmentation
  \end{itemize}}

\newcommand{\Microsoft}{%
  \entry{Microsoft}{Data and Applied Scientist Intern}{May 2022 -- Jul 2022}
  \begin{itemize}
    \item Developed a \textbf{decision-tree ranker} to recommend emails without user queries, improving \textbf{Outlook} search capabilities
    \item Proposed hierarchical feature-sets for the ranker, reducing latency for recommendations and improving recall
  \end{itemize}}

% ---------------------------------------------------------------------------------------------------
% Research
\newcommand{\HomeRobots}{%
  \entry{Real-time 3D Perception for Home Robots}{UC San Diego}{Sep 2023 -- Sep 2024}
  \begin{itemize}
    \item Investigated real-time dense \textbf{visual SLAM} methods using NeRFs and Gaussian Splatting for indoor robot navigation
    \item Integrated object segmentation, grasp-pose estimation, and 3D mapping on the Fetch robot via ROS, demonstrating
          a novel \textbf{tabletop object rearrangement} algorithm that reduced cost by \textbf{20\%} compared to the state-of-the-art
  \end{itemize}}

\newcommand{\Tomography}{%
  \entry{3D Tomography with Primal-Dual Neural Networks}{University College London}{May 2021 -- Jul 2023}
  \begin{itemize}
    \item Developed a stochastic neural-network architecture of primal-dual algorithm for online reconstruction of 3D
          volumes from tomographic projections, with \textbf{99.6\,\% structural similarity} in low-dosage conditions
    \item Built a Python library with custom gradient operators for reconstructing volumes in a single pass, reducing
          compute requirements by up to \textbf{5x} over SOTA learning-based approaches for cone vector tomography
  \end{itemize}}

% ---------------------------------------------------------------------------------------------------
% Projects
\newcommand{\OOXML}{%
  \entry{ooxml}{Open-source contributor \textbar\ \href{https://github.com/ChristopherVR/ooxml}{\color{accent}GitHub}}{Sep 2026 -- Present}
  \begin{itemize}
    \item Contributed to an open-source \textbf{TypeScript} engine that renders and edits PowerPoint and Word files in the browser,
          improving editing fidelity, \textbf{real-time collaboration} and load performance across React, Vue and Svelte viewers
  \end{itemize}}

\newcommand{\MirrorAI}{%
  \entry{Mirror AI: Deployable Personas}{Honorable mention in Supabase YC Hackathon}{Oct 2024 -- Dec 2024}
  \begin{itemize}
    \item Designed an \textbf{agentic LLM architecture} with LangGraph to mirror user personalities, creating interactive digital
          personas, and deployed it as a full-stack platform on Supabase and Vercel with one-click publishing
  \end{itemize}}

\newcommand{\LLMReasoning}{%
  \entry{Improving LLM Reasoning for Numerical Problems}{\report{https://sudhansh6.github.io/assets/pdf/MathPrompter++.pdf}}{Sep 2024 -- Dec 2024}
  \begin{itemize}
    \item Enhanced MathPrompter (ACL 2023) with CoT, achieving \textbf{10\% higher accuracy} on Llama 3.1 1B where prior methods failed,
          and reduced hallucination rates significantly by integrating multi-step validation
  \end{itemize}}

\newcommand{\InverseRendering}{%
  \entry{Inverse Rendering with 2D Gaussian Splatting}{\report{https://sudhansh6.github.io/assets/pdf/CSE252D_Report.pdf}}{Mar 2024 -- May 2024}
  \begin{itemize}
    \item Developed a novel \textbf{inverse rendering} framework in CUDA to recover PBR properties of a scene using 2D Gaussian
          Splatting, improving normal map MAE by \textbf{15\%} over the current SOTA methods
  \end{itemize}}

% ---------------------------------------------------------------------------------------------------
% Publications: title on the left, venue on the right, one line each.
\newcommand{\publine}[3]{\par\noindent\href{#1}{#2}\hfill\hspace{1em}\side{#3}\par\addvspace{1.5pt}}
\newcommand{\PubSaturation}{\publine{https://doi.org/10.1016/j.sigpro.2023.109349}{A likelihood based method for compressive signal recovery under Gaussian and saturation noise}{Signal Processing, 2024}}
\newcommand{\PubPermutation}{\publine{https://ieeexplore.ieee.org/document/10888147}{Identification and correction of permutation errors in compressed sensing-based group testing}{ICASSP 2025}}
\newcommand{\PubRearrangement}{\publine{https://arxiv.org/abs/2411.10899}{Planning for tabletop object rearrangement}{ICRA 2025}}

% ---------------------------------------------------------------------------------------------------
% Education, skills, service
% Each profile sets \UCSDCourses, \IITBCourses and \SkillRows before % Shared preamble and content for Sudhansh Peddabomma's résumés.
% Each profile (SudhanshPeddabomma-AI.tex, SudhanshPeddabomma-Vision.tex) inputs this file and picks
% which blocks to show and in what order, so every bullet is written once.
% Layout rule: every bullet fills one line, or runs 1.5 to 2 lines. Check with: python3 check_lines.py <pdf>
\documentclass[10pt,letterpaper]{article}

\usepackage[margin=0.55in,top=0.5in,bottom=0.45in]{geometry}
\usepackage[T1]{fontenc}
% Fonts: \resumefont loads the body font; \headfont styles the name and section titles.
% Override either before \input{resume-common} to try another pairing.
% Default: Charter body with Source Sans headings. All-sans alternative:
%   \def\resumefont{\usepackage[default,semibold]{sourcesanspro}}\def\headfont{}
\providecommand{\resumefont}{\usepackage[scaled=0.96]{XCharter}\usepackage[semibold,scale=0.96]{sourcesanspro}}
\providecommand{\headfont}{\sffamily}
\providecommand{\namefont}{\rmfamily\scshape}
\resumefont
\usepackage[final]{microtype}
\usepackage[dvipsnames]{xcolor}
\usepackage{enumitem}
\usepackage{titlesec}
\usepackage{tabularx}
\usepackage{fontawesome5}
\usepackage{etoolbox}
\usepackage[hidelinks]{hyperref}

% The site's green, darkened for print.
\definecolor{accent}{HTML}{3A7A45}
\definecolor{muted}{HTML}{5C675D}

\pagestyle{empty}
\setlength{\parindent}{0pt}
\setlength{\tabcolsep}{0pt}
\providecommand{\leading}{1.1}
\linespread{\leading}

\hypersetup{pdfauthor={Sudhansh Peddabomma}}

% Type scale: name 26pt small caps, section titles 12.5pt caps, entry titles 10.8pt, body 10pt,
% side notes 9pt.
\newcommand{\side}[1]{{\fontsize{9}{11}\selectfont\color{muted}#1}}
\newcommand{\sectitle}[1]{\textls[60]{\MakeUppercase{#1}}}
\titleformat{\section}{\headfont\color{accent}\bfseries\fontsize{12.5}{14}\selectfont}{}{0pt}{\sectitle}[\vspace{1.5pt}{\color{accent}\titlerule[0.8pt]}\global\firstentrytrue]
% The first entry after a section title skips \entrygap, so only \rulegap sits under the rule.
\newif\iffirstentry

% Vertical rhythm, the same idea as the earlier hand-tuned résumés: bullets inside an entry sit
% tight together, entries are separated by a clear gap, and sections by a larger one. A profile
% may override these before \input{resume-common} to fit one page.
\providecommand{\secgap}{12pt}   % above a section title
\providecommand{\rulegap}{5.5pt}   % section rule to first entry
\providecommand{\entrygap}{7pt}  % between entries
\providecommand{\headgap}{1.5pt} % entry heading to its bullets
\providecommand{\itemgap}{0.6pt} % between bullets of one entry
\titlespacing*{\section}{0pt}{\secgap}{\rulegap}

\newcommand{\squarebullet}{\raisebox{0.32ex}{\color{accent}\rule{0.5ex}{0.5ex}}}
\setlist[itemize]{leftmargin=12pt, labelsep=6pt, topsep=\headgap, itemsep=\itemgap, parsep=0pt, partopsep=0pt,
  label=\squarebullet}

% \entry{Organisation or title}{Role or detail (may be empty)}{Dates}; \sub{Role or detail}{Right-hand note}
\newcommand{\entry}[3]{%
  \par\iffirstentry\global\firstentryfalse\else\addvspace{\entrygap}\fi\noindent{\headfont\fontsize{10.8}{13}\selectfont\textbf{#1}}\ifblank{#2}{}{\ {\color{muted}|}\ #2}\hfill\hspace{1em}\side{#3}\par}
\newcommand{\sub}[2]{%
  \noindent\textit{#1}\hfill{\color{muted}\textit{#2}}\par}
\newcommand{\report}[1]{\href{#1}{\color{accent}Report}}
\newcommand{\me}{\textbf{S. Peddabomma}}

% ---------------------------------------------------------------------------------------------------
% Header
\newcommand{\Header}{%
  \begin{center}
    {\namefont\fontsize{26}{28}\selectfont Sudhansh Peddabomma}\\[6pt]
    {\headfont\fontsize{9.5}{11}\selectfont\color{muted}
      \href{mailto:sudhansh6@gmail.com}{\faEnvelope[regular]\ sudhansh6@gmail.com} \quad
      \href{https://sudhansh6.github.io}{\faGlobe\ sudhansh6.github.io} \quad
      \href{https://www.linkedin.com/in/sudhansh-peddabomma-73a208192/}{\faLinkedin\ sudhansh-peddabomma} \quad
      \href{https://github.com/Sudhansh6}{\faGithub\ Sudhansh6}}
  \end{center}\vspace{-4pt}}

% ---------------------------------------------------------------------------------------------------
% Glean bullets, shared by both profiles
\newcommand{\GleanSlides}{\item Led \textbf{AI slide generation} in Glean Assistant, now creating \textbf{180K decks a week} for \textbf{50K weekly active users} at enterprises}
\newcommand{\GleanSlidesPerf}{\item Improved the slides harness, raising reliability to \textbf{97\%}, cutting median latency \textbf{26\%} and cost per deck by \textbf{half}
  in four months}
\newcommand{\GleanSlidesOrch}{\item Architected a \textbf{slide orchestration system} that reuses relevant slides from a company's existing decks and generates
  new ones in its brand and template, so teams can refine decks with each other and Glean in full context, without leaving the product}
\newcommand{\GleanAgents}{\item Moved Glean's \textbf{workflow agents} from one-shot planning to an \textbf{agentic tool-use loop}, greatly improving agentic
  capabilities and instruction following, with \textbf{18\%} higher answer completeness and \textbf{17\%} fewer tokens}
\newcommand{\GleanTeamChat}{\item Engineered the core of \textbf{Team chat}, letting teammates work with Glean in one shared conversation and co-edit its artifacts}
\newcommand{\GleanVideo}{\item Shipped \textbf{video artifacts} end to end, generating narrated videos in chat with streamed playback and admin rollout controls}
\newcommand{\GleanEvals}{\item Created the \textbf{evaluations and metrics} that guide model choices and rollouts across agents, slides and interactive artifacts}

% ---------------------------------------------------------------------------------------------------
% Experience
% \groupbreak separates related groups of bullets inside one long entry (slides, then platform work).
\providecommand{\groupgap}{3.5pt}
\newcommand{\groupbreak}{\vspace{\groupgap}}
\newcommand{\GleanAI}{%
  \entry{Glean}{Software Engineer}{Apr 2025 -- Present}
  \begin{itemize}
    \GleanSlides
    \GleanSlidesOrch
    \GleanSlidesPerf \groupbreak
    \GleanTeamChat
    \GleanVideo
    \GleanAgents
    \GleanEvals
  \end{itemize}}

\newcommand{\GleanVision}{%
  \entry{Glean}{Software Engineer}{Apr 2025 -- Present}
  \begin{itemize}
    \GleanSlides
    \GleanSlidesOrch
    \GleanSlidesPerf \groupbreak
    \GleanVideo
    \GleanEvals
  \end{itemize}}

\newcommand{\Duality}{%
  \entry{Duality AI}{Computer Vision Intern}{Jun 2024 -- Sep 2024}
  \begin{itemize}
    \item Designed automated \textbf{3D reconstruction} techniques for featureless objects, reducing digital-twin generation time
          by \textbf{40\%}
    \item Collaborated with \textbf{Autodesk} to validate Unreal Engine simulations for robotics tasks; leveraging structured
          domain randomization to reduce Sim2Real gap and increase mAP-50 by \textbf{15\%} for object detection and segmentation
  \end{itemize}}

\newcommand{\Microsoft}{%
  \entry{Microsoft}{Data and Applied Scientist Intern}{May 2022 -- Jul 2022}
  \begin{itemize}
    \item Developed a \textbf{decision-tree ranker} to recommend emails without user queries, improving \textbf{Outlook} search capabilities
    \item Proposed hierarchical feature-sets for the ranker, reducing latency for recommendations and improving recall
  \end{itemize}}

% ---------------------------------------------------------------------------------------------------
% Research
\newcommand{\HomeRobots}{%
  \entry{Real-time 3D Perception for Home Robots}{UC San Diego}{Sep 2023 -- Sep 2024}
  \begin{itemize}
    \item Investigated real-time dense \textbf{visual SLAM} methods using NeRFs and Gaussian Splatting for indoor robot navigation
    \item Integrated object segmentation, grasp-pose estimation, and 3D mapping on the Fetch robot via ROS, demonstrating
          a novel \textbf{tabletop object rearrangement} algorithm that reduced cost by \textbf{20\%} compared to the state-of-the-art
  \end{itemize}}

\newcommand{\Tomography}{%
  \entry{3D Tomography with Primal-Dual Neural Networks}{University College London}{May 2021 -- Jul 2023}
  \begin{itemize}
    \item Developed a stochastic neural-network architecture of primal-dual algorithm for online reconstruction of 3D
          volumes from tomographic projections, with \textbf{99.6\,\% structural similarity} in low-dosage conditions
    \item Built a Python library with custom gradient operators for reconstructing volumes in a single pass, reducing
          compute requirements by up to \textbf{5x} over SOTA learning-based approaches for cone vector tomography
  \end{itemize}}

% ---------------------------------------------------------------------------------------------------
% Projects
\newcommand{\OOXML}{%
  \entry{ooxml}{Open-source contributor \textbar\ \href{https://github.com/ChristopherVR/ooxml}{\color{accent}GitHub}}{Sep 2026 -- Present}
  \begin{itemize}
    \item Contributed to an open-source \textbf{TypeScript} engine that renders and edits PowerPoint and Word files in the browser,
          improving editing fidelity, \textbf{real-time collaboration} and load performance across React, Vue and Svelte viewers
  \end{itemize}}

\newcommand{\MirrorAI}{%
  \entry{Mirror AI: Deployable Personas}{Honorable mention in Supabase YC Hackathon}{Oct 2024 -- Dec 2024}
  \begin{itemize}
    \item Designed an \textbf{agentic LLM architecture} with LangGraph to mirror user personalities, creating interactive digital
          personas, and deployed it as a full-stack platform on Supabase and Vercel with one-click publishing
  \end{itemize}}

\newcommand{\LLMReasoning}{%
  \entry{Improving LLM Reasoning for Numerical Problems}{\report{https://sudhansh6.github.io/assets/pdf/MathPrompter++.pdf}}{Sep 2024 -- Dec 2024}
  \begin{itemize}
    \item Enhanced MathPrompter (ACL 2023) with CoT, achieving \textbf{10\% higher accuracy} on Llama 3.1 1B where prior methods failed,
          and reduced hallucination rates significantly by integrating multi-step validation
  \end{itemize}}

\newcommand{\InverseRendering}{%
  \entry{Inverse Rendering with 2D Gaussian Splatting}{\report{https://sudhansh6.github.io/assets/pdf/CSE252D_Report.pdf}}{Mar 2024 -- May 2024}
  \begin{itemize}
    \item Developed a novel \textbf{inverse rendering} framework in CUDA to recover PBR properties of a scene using 2D Gaussian
          Splatting, improving normal map MAE by \textbf{15\%} over the current SOTA methods
  \end{itemize}}

% ---------------------------------------------------------------------------------------------------
% Publications: title on the left, venue on the right, one line each.
\newcommand{\publine}[3]{\par\noindent\href{#1}{#2}\hfill\hspace{1em}\side{#3}\par\addvspace{1.5pt}}
\newcommand{\PubSaturation}{\publine{https://doi.org/10.1016/j.sigpro.2023.109349}{A likelihood based method for compressive signal recovery under Gaussian and saturation noise}{Signal Processing, 2024}}
\newcommand{\PubPermutation}{\publine{https://ieeexplore.ieee.org/document/10888147}{Identification and correction of permutation errors in compressed sensing-based group testing}{ICASSP 2025}}
\newcommand{\PubRearrangement}{\publine{https://arxiv.org/abs/2411.10899}{Planning for tabletop object rearrangement}{ICRA 2025}}

% ---------------------------------------------------------------------------------------------------
% Education, skills, service
% Each profile sets \UCSDCourses, \IITBCourses and \SkillRows before \input{resume-common}.
\providecommand{\UCSDCourses}{}
\providecommand{\IITBCourses}{}
\providecommand{\SkillRows}{}
\newcommand{\coursework}[1]{\ifblank{#1}{}{\par\noindent{\fontsize{9.5}{11.5}\selectfont{\color{muted}Coursework:} #1}\par}}
\newcommand{\Education}{%
  \section{Education}
  \entry{University of California San Diego}{M.S. in Computer Science and Engineering (AI)}{GPA 4.00/4.00 \quad Sep 2023 -- Mar 2025}
  \coursework{\UCSDCourses}
  \entry{Indian Institute of Technology Bombay}{B.Tech.\ with Honors in Computer Science}{CPI 9.66/10 \quad Jul 2019 -- Jul 2023}
  \coursework{\IITBCourses}}

% \skill{Label}{Items} is one row of the skills table.
\newcommand{\skill}[2]{{\headfont\bfseries #1} & #2 \\[1.5pt]}
\newcommand{\cpp}{C\kern-0.04em\raisebox{0.06ex}{\small++}}
\newcommand{\Skills}{%
  \section{Skills}
  \begin{tabularx}{\linewidth}{@{}l@{\hspace{12pt}}>{\raggedright\arraybackslash}X@{}}
    \SkillRows
  \end{tabularx}}
.
\providecommand{\UCSDCourses}{}
\providecommand{\IITBCourses}{}
\providecommand{\SkillRows}{}
\newcommand{\coursework}[1]{\ifblank{#1}{}{\par\noindent{\fontsize{9.5}{11.5}\selectfont{\color{muted}Coursework:} #1}\par}}
\newcommand{\Education}{%
  \section{Education}
  \entry{University of California San Diego}{M.S. in Computer Science and Engineering (AI)}{GPA 4.00/4.00 \quad Sep 2023 -- Mar 2025}
  \coursework{\UCSDCourses}
  \entry{Indian Institute of Technology Bombay}{B.Tech.\ with Honors in Computer Science}{CPI 9.66/10 \quad Jul 2019 -- Jul 2023}
  \coursework{\IITBCourses}}

% \skill{Label}{Items} is one row of the skills table.
\newcommand{\skill}[2]{{\headfont\bfseries #1} & #2 \\[1.5pt]}
\newcommand{\cpp}{C\kern-0.04em\raisebox{0.06ex}{\small++}}
\newcommand{\Skills}{%
  \section{Skills}
  \begin{tabularx}{\linewidth}{@{}l@{\hspace{12pt}}>{\raggedright\arraybackslash}X@{}}
    \SkillRows
  \end{tabularx}}
.
\providecommand{\UCSDCourses}{}
\providecommand{\IITBCourses}{}
\providecommand{\SkillRows}{}
\newcommand{\coursework}[1]{\ifblank{#1}{}{\par\noindent{\fontsize{9.5}{11.5}\selectfont{\color{muted}Coursework:} #1}\par}}
\newcommand{\Education}{%
  \section{Education}
  \entry{University of California San Diego}{M.S. in Computer Science and Engineering (AI)}{GPA 4.00/4.00 \quad Sep 2023 -- Mar 2025}
  \coursework{\UCSDCourses}
  \entry{Indian Institute of Technology Bombay}{B.Tech.\ with Honors in Computer Science}{CPI 9.66/10 \quad Jul 2019 -- Jul 2023}
  \coursework{\IITBCourses}}

% \skill{Label}{Items} is one row of the skills table.
\newcommand{\skill}[2]{{\headfont\bfseries #1} & #2 \\[1.5pt]}
\newcommand{\cpp}{C\kern-0.04em\raisebox{0.06ex}{\small++}}
\newcommand{\Skills}{%
  \section{Skills}
  \begin{tabularx}{\linewidth}{@{}l@{\hspace{12pt}}>{\raggedright\arraybackslash}X@{}}
    \SkillRows
  \end{tabularx}}
.
\providecommand{\UCSDCourses}{}
\providecommand{\IITBCourses}{}
\providecommand{\SkillRows}{}
\newcommand{\coursework}[1]{\ifblank{#1}{}{\par\noindent{\fontsize{9.5}{11.5}\selectfont{\color{muted}Coursework:} #1}\par}}
\newcommand{\Education}{%
  \section{Education}
  \entry{University of California San Diego}{M.S. in Computer Science and Engineering (AI)}{GPA 4.00/4.00 \quad Sep 2023 -- Mar 2025}
  \coursework{\UCSDCourses}
  \entry{Indian Institute of Technology Bombay}{B.Tech.\ with Honors in Computer Science}{CPI 9.66/10 \quad Jul 2019 -- Jul 2023}
  \coursework{\IITBCourses}}

% \skill{Label}{Items} is one row of the skills table.
\newcommand{\skill}[2]{{\headfont\bfseries #1} & #2 \\[1.5pt]}
\newcommand{\cpp}{C\kern-0.04em\raisebox{0.06ex}{\small++}}
\newcommand{\Skills}{%
  \section{Skills}
  \begin{tabularx}{\linewidth}{@{}l@{\hspace{12pt}}>{\raggedright\arraybackslash}X@{}}
    \SkillRows
  \end{tabularx}}
